\documentclass[10pt,a4paper]{article}

% ----------------- Paquetes -------------------%
\usepackage[T1]{fontenc}
\usepackage[utf8]{inputenc}
\usepackage[a4paper,margin=2.5cm]{geometry}
\usepackage{mathptmx}             
\usepackage{changepage} 
\usepackage{setspace}
\usepackage[spanish]{babel}
\usepackage[labelfont=bf,labelsep=space]{caption}
\usepackage{amssymb}
\usepackage{amsmath}
\usepackage{ebgaramond}
\usepackage{placeins}
\usepackage{etoolbox}
\usepackage{setspace}
\usepackage{parskip} 
\usepackage{tcolorbox}
\usepackage{needspace}
\usepackage{xparse}
\usepackage{ocgx2}
\usepackage{hyperref}
\usepackage{hypcap}
\usepackage{soul}\sethlcolor{yellow}% resaltado amarillo: \hl{texto} (Panel TCEE, ⌃H)


%%%%%%%%%%%%%%%%%%%%%%%%%%%%%%%%%%%%%%%%%%%%%%%%%%%%%%%%%%%%%%%%
% --- Fija primero tus valores globales buenos ---
\setstretch{1.2}                 % Interlineado global
\setlength{\parskip}{0.7\baselineskip}
\setlength{\parindent}{0pt}
\newlength{\GlobalParskip}
\setlength{\GlobalParskip}{\parskip}

\newcounter{eqnum}[subsection]
\renewcommand{\theeqnum}{\arabic{eqnum}}
\newcounter{alphaitem}
\newcounter{numitem}

%% ------------------- Recuadros----------------------%%
\newcommand{\puntosclave}[1]{%
  \begin{tcolorbox}[
    colframe=black,
    title={Puntos clave},
    before upper={\setlength{\parindent}{0.5cm}\setlength{\parskip}{0.5\baselineskip}}
  ]
  #1
  \end{tcolorbox}
}

\newcommand{\modificaciones}[1]{
  \begin{tcolorbox}[
    colback=white,
    colframe=red,
    title={Modificaciones a realizar},
    before upper={\setlength{\parindent}{0.5cm}\setlength{\parskip}{0.5\baselineskip}}
  ]
  #1
  \end{tcolorbox}
}

%% ------------------- Listas----------------------%%
    % --- Lista alfabética ---
    \newenvironment{la}
      {\begin{list}{\alph{alphaitem})}{%
          \usecounter{alphaitem}%
          \setlength{\leftmargin}{1cm}%
          \setlength{\parskip}{3pt}% 
          \setlength{\itemsep}{0pt}% ← espacio entre ítems
          \setlength{\parsep}{0pt}%  ← párrafos dentro de un ítem
      }}
      {\end{list}}
      
    % --- Lista numérica ---
    \newenvironment{lnum}
      {\begin{list}{\arabic{numitem}.}{%
          \usecounter{numitem}%
          \setlength{\leftmargin}{1cm}%
          \setlength{\parskip}{3pt}% 
          \setlength{\itemsep}{0pt}% ← espacio entre ítems
          \setlength{\parsep}{0pt}%  ← párrafos dentro de un ítem
      }}
      {\end{list}}
      
% ------------------- Imágenes----------------------%
    \graphicspath{{Img/}}% Rutas por defecto 
    % --- Imagen adaptativa ---
    % Registros de dimensión definidos una sola vez
    \newdimen\imgcap
    \newdimen\imgmin
    \newdimen\imgmax
    \newdimen\imgremain
    \newdimen\imgh
    \newdimen\imgneed
    
    \newcommand{\imagenfit}[3]{%
      \addvspace{10pt}% separación antes
      \par\begingroup
        % Parámetros de composición (asignaciones, no \newdimen)
        \imgcap=1\baselineskip            % alto estimado de título+fuente
        \imgmin=0.15\textheight
        \imgmax=0.15\textheight
        \imgremain=\pagegoal
        \ifdim\pagegoal=\maxdimen \imgremain=0pt\fi
        \advance\imgremain by -\pagetotal
        \advance\imgremain by -\imgcap
        % Decisión de altura destino
        \ifdim\imgremain<\imgmin
          \newpage
          \imgh=0.20\textheight
        \else
          \imgh=\imgremain
          \ifdim\imgh>\imgmax \imgh=\imgmax \fi
        \fi
        % Reservar espacio para evitar cortes del bloque completo
        \imgneed=\imgh \advance\imgneed by \imgcap
        \Needspace{\imgneed}
        % Cuerpo
        \noindent\begin{minipage}{\textwidth}
          \centering
          \captionof{figure}{\textemdash\ #2}\par\vspace{0.3em}% título arriba
          \includegraphics[width=\linewidth,height=\imgh,keepaspectratio]{\detokenize{#1}}\par
          \vspace{0.3em}\small\textit{Fuente: }#3% fuente abajo
        \end{minipage}%
      \endgroup
      \par\addvspace{10pt}%
    }

%------------ Bloque de RECUADRO ESPECIAL -------------%
    \newenvironment{specialblock}[1]{%
      \begin{quote} % crea sangría a izquierda y derecha
      \small % texto más pequeño
      \noindent\textbf{#1}\\[-0.3em] % título en negrita
      \rule{\linewidth}{0.4pt}\\[0.6em]}{
      \end{quote}}
      

%-------------------------- Author Font -----------------------%
    \newcommand{\authorfont}[1]{{\fontfamily{EBGaramond-TLF}\selectfont #1}}

%-------------------------- Anexos -----------------------%
    \makeatletter
    \newcounter{anexo}
    \renewcommand{\theanexo}{\Roman{anexo}}
    \newcommand{\anexo}[1]{%
      \clearpage
      \phantomsection
      \refstepcounter{anexo}%
      \noindent\textbf{Anexo \theanexo.- }#1\par\medskip
      \addcontentsline{stc}{section}{Anexo \theanexo.- #1}%
    }
    \makeatother

    %-------------------------- Lista de figuras (personal) -----------------------%
\makeatletter
\newcommand{\MyListOfFigures}{%
  \clearpage
  \phantomsection
  {\centering\bfseries\Large Índice de figuras\par}%
  \medskip
  % Entrada en nuestro índice de contenido (stc)
  \addcontentsline{stc}{section}{Índice de figuras}%
  % Recuperación del listado (sin cabecera propia)
  \begingroup
    \parindent=0pt\parskip=2pt
    \@starttoc{lof}%
  \endgroup
}
\makeatother

%-------------------------- Bibliografía -----------------------%
\makeatletter
% Contador para las referencias
\newcounter{myref}
% Cabecera de la bibliografía, entra en el índice 'stc'
\newcommand{\MyBibliography}{%
  \clearpage
  \phantomsection
  {\centering\bfseries\Large Bibliografía y referencias\par}%
  \medskip
  \addcontentsline{stc}{section}{Bibliografía y referencias}%
  \setcounter{myref}{0}%
}
\newcommand{\MyBibItem}[4]{%
  \refstepcounter{myref}%
  \noindent[\themyref]\ #1.\ \emph{#2}. #3. #4\par
}
\makeatother

%--------- Establecer cuatro niveles de índice --------%
    \setcounter{tocdepth}{4}   
    \setcounter{secnumdepth}{4}% numera hasta \paragraph

%%%%%%%%%%%%%%%%%%%%%%%%%%%%%%%%

\makeatletter

    %--------- Interlineado Global --------%
    \edef\GlobalBaseLineStretch{\baselinestretch}
    
    %--------- Espaciado de seccinones --------%
    \renewcommand\section{\@startsection{section}
      {1}{\z@}
      {2.5ex \@plus 1ex \@minus .2ex} % espacio antes
      {1.5ex \@plus .2ex}             % espacio después
      {\normalfont\Large\bfseries}    % formato del título
    }
    \renewcommand\subsection{\@startsection{subsection}{2}{\z@}
      {3ex \@plus 0.8ex \@minus .2ex}
      {1.5ex \@plus .2ex}
      {\normalfont\large\bfseries}
    }
    \renewcommand\subsubsection{\@startsection{subsubsection}{3}{\z@}
      {3ex \@plus 0.5ex \@minus .2ex}
      {1ex \@plus .2ex}
      {\normalfont\normalsize\bfseries}
    }
    \renewcommand\paragraph{\@startsection{paragraph}{4}{\z@}%
    {-2ex \@plus -0.5ex \@minus -.2ex}% espacio antes
    {0.8ex \@plus .2ex}% espacio después
    {\normalfont\normalsize\itshape}}% ← cursiva en lugar de negrita

    %--------- Ecuaciones --------%   
    % Variables globales para almacenar títulos y notas
    \gdef\leq@current@section{}%
    \gdef\leq@current@subsection{}%
    \gdef\leq@last@printed@section{}%
    \gdef\leq@last@printed@subsection{}%
    
    % Comando para añadir nota (invisible en el cuerpo)
    \newcommand{\eqnote}[1]{%
      \addcontentsline{leq}{leqnote}{#1}%
    }
    
    % Comando para ecuaciones
    \newcommand{\eqblock}[2]{%
      % Verificar si necesitamos imprimir el título de sección
      \ifx\leq@current@section\leq@last@printed@section\else
        \addcontentsline{leq}{leqsection}{\leq@current@section}%
        \global\let\leq@last@printed@section\leq@current@section
        \gdef\leq@last@printed@subsection{}% Reset subsección al cambiar sección
      \fi
      % Verificar si necesitamos imprimir el título de subsección
      \ifx\leq@current@subsection\leq@last@printed@subsection\else
        \ifx\leq@current@subsection\@empty\else
          \addcontentsline{leq}{leqsubsection}{\leq@current@subsection}%
          \global\let\leq@last@printed@subsection\leq@current@subsection
        \fi
      \fi
      % Ecuación
      \refstepcounter{eqnum}%
      \begin{equation}
        #1\tag{E.\theeqnum}
      \end{equation}%
      \addcontentsline{leq}{leqt}{[E.\theeqnum] -- #2}%
      \addcontentsline{leq}{leqm}{#1}%
    }
    
    %%% ----- Índice de ecuaciones ----------%%%
    % Tamaño para []
    \newcommand{\leqIndexFont}{\normalsize}
    \newcommand{\leqTitleFont}{\tiny}
    \newcommand{\leqNoteFont}{\tiny}
    % Cabecera de sección 
    \newcommand\l@leqsection[2]{%
      \addvspace{25pt}% Mayor separación antes de nueva sección
      \noindent\leqIndexFont
      \makebox[\linewidth]{\rule{.38\linewidth}{0.4pt}\ \ \textbf{  \MakeUppercase{#1}}\ \ \rule{.38\linewidth}{0.4pt}}%
      \par\vspace{-10pt}
    }
    % Cabecera de subsección 
    \newcommand\l@leqsubsection[2]{%
      \par\addvspace{15pt}%
      \noindent\leqIndexFont #1
      \par\vspace{-7pt}
    }
    % Título de ecuación 
    \newcommand\l@leqt[2]{%
    \par\addvspace{10pt}%
      \par\noindent\leqTitleFont #1\par
    }
    % Miniatura de ecuación
    \newcommand\l@leqm[2]{%
      \vspace{0.3\baselineskip}%
      \par\scriptsize\noindent\hspace{1.5em}\(#1\)\par%
    }
    % Nota de ecuación 
    \newcommand\l@leqnote[2]{%
      \vspace{0.2\baselineskip}%
      \par\noindent\leqNoteFont\hspace{1.5em}\textbf{Nota.--} #1\par\vspace{0.5\baselineskip}%
    }
    % Generación del índice
    \newcommand{\mylistofequations}{%
      \clearpage
      \phantomsection
      % Título visual de la lista (ajústalo a tu gusto):
      {\centering\bfseries\Large Índice de ecuaciones\par}
      % Entrada en nuestro índice 'stc':
      \addcontentsline{stc}{section}{Índice de ecuaciones}%
      % Llamada a tu lista real (usa el nombre exacto de tu macro):
      \@starttoc{leq}%
    }
    
    %--------- Captura de títulos de sección --------%
    \let\leq@original@section\section
    \let\leq@original@subsection\subsection
    % Flag para saber si estamos en una sección asterisco
    \newif\ifLeqStarSection
    \LeqStarSectionfalse
    
    % Redefinir \section CON CONTROL DE ASTERISCO
    \renewcommand{\section}{%
      \@ifstar{\leq@section@star}{\@dblarg\leq@section@nostar}%
    }
    \def\leq@section@star#1{%
      \global\LeqStarSectiontrue
      \gdef\leq@current@section{#1}%
      \gdef\leq@current@subsection{}%
      \leq@original@section{#1}%
      \global\LeqStarSectionfalse
    }
    \def\leq@section@nostar[#1]#2{%
      \global\LeqStarSectionfalse
      \gdef\leq@current@section{#2}%
      \gdef\leq@current@subsection{}%
      \leq@original@section[#1]{#2}%
    }
    % Redefinir \subsection
    \renewcommand{\subsection}{%
      \@ifstar{\leq@subsection@star}{\@dblarg\leq@subsection@nostar}%
    }
    \def\leq@subsection@star#1{%
      \global\LeqStarSectiontrue
      \gdef\leq@current@subsection{#1}%
      \leq@original@subsection{#1}%
      \global\LeqStarSectionfalse
    }
    \def\leq@subsection@nostar[#1]#2{%
      \global\LeqStarSectionfalse
      \gdef\leq@current@subsection{#2}%
      \leq@original@subsection[#1]{#2}%
    }

    % ----- Restauración de valores Globales de estilo ----------%
    % Flags
    \newif\ifInFootnote
    \InFootnotefalse
    \pretocmd{\@footnotetext}{\InFootnotetrue}{}{}
    \apptocmd{\@footnotetext}{\InFootnotefalse}{}{}
    % Profundidad de listas (soporta anidadas)
    \newcount\MyListDepth \MyListDepth=0
    \AtBeginEnvironment{la}{\advance\MyListDepth by 1}
    \AfterEndEnvironment{la}{\advance\MyListDepth by -1}
    \AtBeginEnvironment{lnum}{\advance\MyListDepth by 1}
    \AfterEndEnvironment{lnum}{\advance\MyListDepth by -1}
    \AtBeginEnvironment{itemize}{\advance\MyListDepth by 1}
    \AfterEndEnvironment{itemize}{\advance\MyListDepth by -1}
    \AtBeginEnvironment{enumerate}{\advance\MyListDepth by 1}
    \AfterEndEnvironment{enumerate}{\advance\MyListDepth by -1}
    % Restauración diferida: hazlo al INICIO del próximo párrafo real
    \newif\if@pendingRestore
    \@pendingRestorefalse
    \newcommand{\RequestRestore}{\global\@pendingRestoretrue}
    \everypar{%
      \if@pendingRestore
        \global\@pendingRestorefalse
        \ifInFootnote\else
          \ifnum\MyListDepth=0
            \setlength{\parskip}{\GlobalParskip}%
            \setstretch{\GlobalBaseLineStretch}%
          \fi
        \fi
      \fi
    }
    % Al final de bloques:
    \AfterEndEnvironment{equation}{\RequestRestore}
    \AfterEndEnvironment{equation}{\RequestRestore}
    \AfterEndEnvironment{la}{\RequestRestore}
    \AfterEndEnvironment{lnum}{\RequestRestore}
    % Al final de macros:
    \apptocmd{\eqblock}{\RequestRestore}{}{}
    \apptocmd{\imagenfit}{\RequestRestore}{}{}
    
\makeatother

% ===================== ToC propio con 4 niveles (único bloque) =====================
\makeatletter

% --- Escritores: añadimos una línea al .stc en cada cabecera numerada ---
% Guardamos las versiones actuales para llamar al comportamiento normal
\let\MyTOC@orig@section\section
\let\MyTOC@orig@subsection\subsection
\let\MyTOC@orig@subsubsection\subsubsection
\let\MyTOC@orig@paragraph\paragraph

% SECTION
\renewcommand{\section}{%
  \@ifstar{\MyTOC@section@star}{\@dblarg\MyTOC@section@nostar}%
}
\def\MyTOC@section@star#1{%
  \MyTOC@orig@section{#1}% sin entrada al .stc
}
\def\MyTOC@section@nostar[#1]#2{%
  \MyTOC@orig@section[{#1}]{#2}% comportamiento normal (escribe al .toc)
  % Escribimos también nuestra línea al .stc (texto con \numberline)
  \addcontentsline{stc}{section}{\protect\numberline{\thesection}#2}%
}

% SUBSECTION
\renewcommand{\subsection}{%
  \@ifstar{\MyTOC@subsection@star}{\@dblarg\MyTOC@subsection@nostar}%
}
\def\MyTOC@subsection@star#1{%
  \MyTOC@orig@subsection{#1}%
}
\def\MyTOC@subsection@nostar[#1]#2{%
  \MyTOC@orig@subsection[{#1}]{#2}%
  \addcontentsline{stc}{subsection}{\protect\numberline{\thesubsection}#2}%
}

% SUBSUBSECTION
\renewcommand{\subsubsection}{%
  \@ifstar{\MyTOC@subsubsection@star}{\@dblarg\MyTOC@subsubsection@nostar}%
}
\def\MyTOC@subsubsection@star#1{%
  \MyTOC@orig@subsubsection{#1}%
}
\def\MyTOC@subsubsection@nostar[#1]#2{%
  \MyTOC@orig@subsubsection[{#1}]{#2}%
  \addcontentsline{stc}{subsubsection}{\protect\numberline{\thesubsubsection}#2}%
}

% PARAGRAPH
\renewcommand{\paragraph}{%
  \@ifstar{\MyTOC@paragraph@star}{\@dblarg\MyTOC@paragraph@nostar}%
}
\def\MyTOC@paragraph@star#1{%
  \MyTOC@orig@paragraph{#1}%
}
\def\MyTOC@paragraph@nostar[#1]#2{%
  \MyTOC@orig@paragraph[{#1}]{#2}%
  % Si no numeras los \paragraph, cambia \theparagraph por texto plano sin \numberline
  \addcontentsline{stc}{paragraph}{\protect\numberline{\theparagraph}#2}%
}

% --- Impresor del índice propio (lee el .stc) ---
\newcommand{\MyTOC}{%
  \par\bigskip
  {\centering\bfseries\Large Índice de contenido\par}%
  \medskip
  \begingroup
    \parindent=0pt\parskip=2pt
    \@starttoc{stc}%
  \endgroup
}

\makeatother
% ===================== fin ToC propio =====================


%%%%%%%%%%%%%%


% --- Datos del tema ---
\title{Política Presupuestaria (VI): Saldo presupuestario y deuda de las Administraciones Públicas en España\\
\large Análisis de su evolución reciente y comparación internacional. Política de financiación del Tesoro en la actualidad}
\author{Marcelo Ortega}
\date{Septiembre 2026}

\begin{document}
\maketitle
\newpage

% --- Página de resumen y modificaciones ---
\puntosclave{ %% Escribir puntos clave Apuntes CECO
     
    }
\vspace{1cm}

\modificaciones{
    \textbf{NOTA (04/10/2026):} Revisar (Rosell): https://www.bbvaresearch.com/publicaciones/espana-observatorio-fiscal-junio-2026/
}

\newpage
\MyTOC
\newpage

\mylistofequations
\clearpage

\MyListOfFigures
\clearpage


\pagenumbering{arabic}
\setcounter{page}{1}


\section{Introducción}



\subsection{Contextualización}


\subsubsection{Relevancia}




\subsubsection{Problemática}



\subsection{Estructura de la exposición}




\clearpage
\section{Saldo presupuestario y la deuda de las Administraciones Públicas en España: evolución}
\puntosclave{
    En 2025, España registró un déficit público de unos $37.000$ a $40.000$ millones de euros, según se incluyan o no los gastos de la DANA. Ese mismo año, el Estado pagó por sí solo unos $34.900$ millones de intereses. Es decir, casi todo el déficit de las AAPP se explica ya por la carga de la deuda heredada, y el saldo primario está cerca del equilibrio.    
}

En España, el déficit público es un fenómeno relativamente reciente. Hasta la Transición, el Estado mantuvo el equilibrio presupuestario e incluso pequeños superávits. Fue durante la Transición cuando el saldo empezó a mostrar déficits recurrentes, una situación que, con el breve paréntesis de 2005-2007, se ha mantenido hasta hoy. 

Podemos resumir el recorrido en tres puntos clave. El equilibrio previo a 1975 ocultaba el atraso del sistema tributario y solo era posible por el reducido tamaño del sector público. La creación del Estado de bienestar y del Estado autonómico trajo consigo un aumento estructural del gasto que no se acompañó de medidas equivalentes en los ingresos. Esos déficits han sido casi siempre estructuralmente negativos y, además, procíclicos, algo especialmente indeseable. La etapa de 1996 a 2007 fue muy favorable para el saneamiento y permitió los tres únicos superávits de la democracia. La recesión de 2008 provocó déficits superiores al $10\%$ del PIB y llevó la deuda al $105\%$ en 2014. Y la pandemia de 2020 frenó en seco la lenta consolidación posterior.

Este bloque reconstruye esa historia con una pregunta de fondo: ¿por qué España ha tenido déficit en casi todos los años de la democracia, incluidos muchos de expansión, y qué papel han desempeñado en su deuda los déficits, el crecimiento, los tipos de interés y las operaciones financieras? Para responderla, fragmentaremos la pregunta en cinco partes.

\textbf{¿Qué queremos decir?} En primer lugar, fijaremos el marco de análisis. ¿Qué se mide, cómo se leen en España los indicadores del saldo, cómo se relacionan saldo y deuda y cómo se pasa del déficit de las AAPP a las necesidades del Tesoro? Después, recorreremos la evolución en cuatro etapas: del equilibrio aparente al déficit estructural (1975-1995); la consolidación, el euro y los superávits (1996-2007); la Gran Recesión y la crisis de la deuda soberana (2008-2019); y la pandemia, la inflación y la situación actual (2020-2026). Como veremos, la historia fiscal de la democracia española combina un déficit estructural persistente, porque los ingresos nunca se han ajustado del todo a un Estado de bienestar construido en muy poco tiempo, con grandes perturbaciones que han llevado la deuda a saltos (1992-1993, 2008-2013, 2020). Los periodos de reducción de la deuda se han debido más al crecimiento nominal y a los tipos bajos que a superávits primarios sostenidos. Esa combinación explica que España llegue a 2026 con una ratio de deuda cercana al $100\%$ del PIB, un saldo primario próximo al equilibrio y una carga de intereses que vuelve a crecer.

\subsection{Marco de análisis: qué se mide y cómo se lee en España}
El título del tema habla de saldo presupuestario y de deuda de las Administraciones Públicas como si fueran magnitudes evidentes. No lo son. Cada una tiene varias definiciones, y la elección entre ellas cambia el diagnóstico. 

No repetiremos aquí los desarrollos conceptuales, pues corresponden a otras partes del temario, pero sí los presentaremos por su función operativa: proporcionarnos las claves para leer la evolución española y enlazarla con la política de financiación del Tesoro. Para ello, contestaremos a tres cuestiones: ¿Qué se mide? ¿Cómo se lee el saldo en España? ¿Cómo se relacionan saldo y deuda? Con esto, podremos ver cómo se relaciona el déficit de las AAPP con las necesidades del Tesoro.

Como resumen, podemos decir que el saldo presupuestario y la deuda de las AAPP se miden en contabilidad nacional, en términos de devengo y con un perímetro móvil. La deuda se valora por su nominal y se consolida entre subsectores, lo que convierte al Estado en el prestamista de las comunidades y de la Seguridad Social y explica que su deuda represente casi el $90\%$ del total. Con esto, veremos como en España, el saldo primario está hoy cerca del equilibrio, pero la carga de intereses vuelve a crecer. Además, el saldo estructural revela un déficit persistente y enseña, con los falsos superávits de 2005-2007, la fragilidad de las estimaciones en tiempo real. Sin embargo, la deuda no varía en la cuantía del déficit: su evolución depende también del crecimiento nominal, de los tipos de interés y de numerosas operaciones financieras. Por ello, los periodos españoles de reducción de la deuda se han debido más al crecimiento nominal y a los tipos bajos que a superávits primarios. Por último, entre el déficit de las AAPP y las necesidades del Tesoro median los préstamos del Estado a otros subsectores, las diferencias entre caja y devengo y los vencimientos, lo que explica que el Tesoro emita cada año mucho más de lo que dice el déficit. 


\subsubsection{Qué se mide: el saldo y la deuda en su definición oficial}
El saldo presupuestario se define como la diferencia entre los ingresos y los gastos totales de las AAPP durante el ejercicio, es decir, su capacidad o necesidad de financiación. Según el SEC 2010, corresponde al saldo de la cuenta de capital del sector institucional AAPP y es el indicador esencial para las comparaciones internacionales y para valorar el cumplimiento de las reglas fiscales europeas.\footnote{%
    En el SEC 2010, la capacidad ($+$) o necesidad ($-$) de financiación (B.9) es el saldo de la cuenta de capital del sector: el ahorro bruto más las transferencias de capital netas menos la formación bruta de capital y la adquisición neta de activos no financieros no producidos. Es el concepto de déficit público que utiliza el Protocolo de Déficit Excesivo (Reglamento (CE) 479/2009). ([\authorfont{Ver Tema 3.A.39}])
} 

Dos precisiones son relevantes para todo el tema. La primera: el saldo se mide en contabilidad nacional, con criterio de devengo, no en la contabilidad presupuestaria de caja con la que se ejecutan los presupuestos. De ahí que el déficit de las AAPP y las necesidades de financiación del Tesoro no coincidan.\footnote{%
    La distinción entre la contabilidad presupuestaria (caja y derechos u obligaciones reconocidos) y la contabilidad nacional (devengo), y los ajustes que la IGAE aplica para pasar de una a otra, se desarrollan en ([\authorfont{Ver Tema 4.B.21}])
} La segunda: el saldo se refiere al conjunto de las AAPP, pero se publica también por subsectores, lo que permite ver cómo se reparte el déficit entre la Administración central, las comunidades, las entidades locales y la Seguridad Social.

Por el contrario, la deuda se define como el saldo vivo, en un momento dado, de los pasivos contraídos para financiar los gastos que no pueden cubrirse con ingresos corrientes y de capital. Existen distintos conceptos según los instrumentos incluidos y el método de valoración, y el más relevante para este tema es el del Protocolo de Déficit Excesivo (PDE). Este criterio valora la deuda por su valor nominal o facial, lo que evita las grandes oscilaciones que produciría la valoración a precios de mercado.\footnote{%
    Con la valoración a precios de mercado, la deuda medida variaría con los tipos de interés: una subida de tipos reduciría el valor de mercado de la deuda viva, y con él la deuda contabilizada, sin que cambiaran en nada las obligaciones de pago del Estado, y una bajada la elevaría. Por eso, aunque las cuentas financieras del SEC 2010 registran los valores de deuda a precios de mercado, la deuda PDE se registra por el valor nominal, que es lo que el Estado debe reembolsar al vencimiento.
} La deuda PDE incluye el efectivo y los depósitos, los valores representativos de deuda (letras, bonos y obligaciones) y los préstamos, y excluye otros pasivos como los créditos comerciales, es decir, las facturas pendientes de pago, o los pasivos por pensiones. Por eso la literatura y los organismos internacionales complementan la deuda PDE con otras medidas: la deuda neta de activos financieros, el total de pasivos o los pasivos contingentes e implícitos, como los avales públicos o los compromisos de pensiones.\footnote{%
    Los pasivos contingentes (avales, colaboración público-privada) y los implícitos (compromisos de pensiones, recogidos en la tabla complementaria 29 del SEC 2010) se tratan en ([\authorfont{Ver Tema 4.B.24}]) (sobre la deuda implícita del sistema de pensiones) y, para las empresas públicas y los pasivos fuera del perímetro, ([\authorfont{Ver Tema 4.B.22}]).
}

El criterio de valoración nominal tiene una consecuencia práctica que conviene adelantar, porque afecta a la lectura de la deuda. Cuando el Tesoro emite un bono por encima de la par, lo que ocurrió mucho entre 2015 y 2021 al reabrir referencias con cupones superiores a los tipos de mercado, ingresa más de lo que computa como deuda PDE, porque solo computa el nominal. La diferencia, la prima de emisión, financia el déficit de caja sin aumentar la deuda PDE, a cambio de pagar cupones más altos en el futuro. Desde 2022 ocurre lo contrario: con la subida de tipos, muchas reaperturas se emiten por debajo de la par, y el Tesoro necesita emitir más nominal para obtener el mismo efectivo. Por tanto, constituye un ejemplo claro de por qué la deuda no varía exactamente en la cuantía del déficit.

\paragraph{El perímetro: quién forma parte de las AAPP}

La Administración central comprende las unidades con competencias en todo el territorio nacional: el Estado (Casa del Rey, Gobierno, Cortes Generales, Tribunal Constitucional, etc.), los organismos de la administración central y las empresas clasificadas como administración central. Las comunidades autónomas incluyen sus órganos de gobierno y sus organismos administrativos, sus universidades y las empresas clasificadas como administración regional. Las corporaciones locales agrupan ayuntamientos, diputaciones, cabildos y consejos insulares, mancomunidades y agrupaciones, las ciudades autónomas de Ceuta y Melilla y sus organismos dependientes. Las administraciones de la Seguridad Social incluyen las unidades que proveen prestaciones sociales.\footnote{
Además, con carácter general, forman parte de las AAPP las unidades institucionales controladas por ellas que sean productores no de mercado, definidos por el SEC-2010. Los criterios de delimitación (control público, carácter de productor de mercado o no de mercado y prueba del $50\%$ de los costes cubiertos con ventas), así como las reclasificaciones de unidades como la Sareb, el ADIF o las empresas públicas, se desarrollan en el ([\authorfont{Ver Tema 4.B.23}]).
}

El perímetro no es una cuestión técnica menor. En España, varias reclasificaciones han alterado de forma relevante la deuda contabilizada. La más importante fue la de la SAREB. En marzo de 2021, siguiendo el criterio de Eurostat, el INE, la IGAE y el BdE la incluyeron en el sector de las AAPP, lo que añadió unos $35.000$ millones de euros a la deuda de 2020 y contribuyó a que la ratio superara el $120\%$ del PIB ese año.\footnote{%
    La inclusión de la Sareb en las AAPP se basó en que el Estado controlaba la entidad a través del FROB y garantizaba su deuda sénior, y en que las pérdidas acumuladas, no previstas inicialmente, desvirtuaban su carácter de productor de mercado. Su vida limitada, con disolución prevista en 2027, había sido una de las condiciones para mantenerla fuera del perímetro. ([\authorfont{Ver Tema 4.B.23}])
} El episodio muestra que la frontera de las AAPP es móvil y que la deuda publicada depende de decisiones estadísticas, no solo de decisiones presupuestarias.

\paragraph{La consolidación entre subsectores}

Además, hay que tener en cuenta que la deuda de las AAPP se publica consolidada: se eliminan los pasivos de un subsector que son activos de otro. En España, esa eliminación es muy importante, porque una gran parte de la deuda de las comunidades autónomas y casi toda la de la Seguridad Social se deben a préstamos del Estado. 

Al cierre de 2025, la deuda consolidada de las AAPP ascendía a unos $1,698$ billones de euros ($100,7\%$ del PIB), mientras que la suma de los subsectores sin consolidar superaba esa cifra: Administración central, unos $1,56$ billones ($92,6\%$ del PIB); comunidades, $341.642$ millones ($20,2\%$); corporaciones locales, $20.729$ millones ($1,2\%$); y Seguridad Social, unos $136.200$ millones ($8,1\%$). 

En la práctica, como podemos ver, el Estado se ha convertido en el prestamista de los demás subsectores, y el Tesoro financia en el mercado no solo el déficit del Estado, sino también parte del de las comunidades y la Seguridad Social.


\subsubsection{Cómo se lee el saldo en España: total, primario y estructural}
Ahora que hemos repasado qué es el déficit y la deuda, conviene entender cómo se lee el saldo presupuestario, pues existen tres saldos que responden a tres preguntas distintas. Evidentemente, el \textbf{saldo total} responde a la pregunta de cuánto necesitan endeudarse las AAPP en un año. Pero para leer la política fiscal hacen falta otros dos indicadores. 

En primer lugar, el saldo primario, que excluye los intereses, responde a la pregunta de si los ingresos cubren el gasto corriente y de capital sin contar la carga de la deuda heredada. En segundo lugar, el saldo estructural, que corrige el efecto del ciclo y de las medidas excepcionales, responde a la pregunta de cuál sería el saldo con la economía en su nivel potencial.\footnote{%
    Saldo primario $=$ saldo total $+$ intereses. Saldo ajustado de ciclo $=$ saldo total $-$ $\varepsilon x_t$, donde $\varepsilon$ es la semielasticidad presupuestaria al ciclo (en torno a $0,5-0,6$ para España en la metodología común de la Unión) y la brecha de producción ($x_t$) es la diferencia porcentual entre el PIB observado y el potencial. Saldo estructural $=$ saldo ajustado de ciclo $-$ medidas excepcionales y temporales. Para los conceptos básicos [\authorfont{Ver Tema 3.A.39}]. Para la estimación del producto potencial y de la brecha de producción, y sus problemas en tiempo real, véanse [\authorfont{Ver Tema 3.A.42, Tema 3.A.43}].
} 

¿Qué dicen estos dos indicadores en el caso de España?

\paragraph{Saldo primario: la carga de la deuda heredada.}
En 2025, el déficit de las AAPP fue del $2,2\%$ del PIB, unos $36.800$ millones, o del $2,4\%$, unos $40.300$ millones, si se incluyen los gastos extraordinarios de la DANA de Valencia de octubre de 2024. 

El Estado pagó ese año unos $34.900$ millones en intereses, y el conjunto de las AAPP, algo más de $40.000$ millones. La consecuencia es que el saldo primario de las AAPP estuvo muy cerca del equilibrio en 2025, lo que no ocurría desde antes de la Gran Recesión. No obstante, esa mejora se produce justo cuando el gasto en intereses vuelve a subir. 

Hasta julio de 2026, los intereses pagados por el Estado crecían en torno a un $14\%$ interanual, y la AIReF (2025) y Fedea (2025) proyectan un aumento sostenido de la carga de intereses en la próxima década, a medida que se refinancia a tipos más altos la deuda emitida en la etapa de tipos cero.\footnote{%
    Según distintas proyecciones, la carga de intereses podría duplicarse en porcentaje del PIB hacia 2050, \textit{ceteris paribus}.
} Dicho de otro modo: el esfuerzo fiscal de los próximos años tendrá que hacerse para absorber el coste de la deuda heredada, no para corregir un exceso de gasto corriente.

\paragraph{Saldo estructural: el diagnóstico persistente y sus trampas}
No obstante, es la lectura del saldo estructural la que ofrece el diagnóstico más repetido. La AIReF, la Comisión, el Banco de España (BdE) y el FMI coinciden en que España tiene un déficit estructural persistente: incluso con la economía cerca de su nivel potencial, las AAPP registran un déficit relevante, que diversas estimaciones sitúan en los últimos años en torno a $2,5$ o $3$ puntos del PIB potencial 

Esta lectura presenta, también, un ejemplo histórico reciente. En 2005-2007, España registró superávits totales de hasta el $2\%$ del PIB, y las estimaciones de la época situaban también en superávit el saldo estructural. Con la información posterior, esas estimaciones se revisaron a fondo. 

Buena parte de los ingresos de aquellos años eran extraordinarios: procedían de la burbuja inmobiliaria (impuestos sobre transmisiones, IVA de la construcción, IRPF e Impuesto sobre Sociedades ligados al sector) y no desaparecían con el ciclo tal como lo medían las técnicas convencionales, sino con el pinchazo de la burbuja. Las estimaciones actuales sitúan el saldo estructural de 2007 en déficit. Es decir, España entró en la crisis con un déficit estructural oculto por los ingresos de la burbuja.\footnote{%
    El fenómeno lo documentaron, entre otros, el BdE y la Comisión tras la crisis, con análisis de los ingresos extraordinarios (revenue windfalls) ligados al precio de los activos y a la composición de la demanda. Los métodos estándar de ajuste cíclico, basados en la brecha de producción, no los detectan.
} 

La lección, que retomaremos más adelante, es que el saldo estructural estimado en tiempo real puede ser muy engañoso, y que los superávits de 2005-2007 no fueron tanto un éxito de la política fiscal como un espejismo de ingresos.

\begin{specialblock}{\textbf{Medidas excepcionales}: el caso de la DANA}
    El saldo estructural excluye también las medidas excepcionales y temporales, lo que plantea en cada caso la cuestión de qué se considera excepcional. 
    
    El Gobierno presentó el déficit de 2025 sin los gastos de la DANA ($2,2\%$), mientras que Eurostat y la IGAE publicaron la cifra total ($2,4\%$). La diferencia es una convención: los gastos de reconstrucción tras una catástrofe se consideran temporales en el marco fiscal europeo. Pero muestra que incluso el saldo total admite lecturas distintas según el propósito.    
\end{specialblock}

\subsubsection{Cómo se relacionan saldo y deuda: la contabilidad de la variación de la deuda en España}
Habida cuenta de todo lo expuesto, podemos ahora presentar la relación entre ambas. La variación de la ratio de deuda sobre el PIB puede descomponerse contablemente en tres componentes: (i) el saldo primario ($sp_t$); (ii) el llamado efecto bola de nieve ($\frac{i_t - g_t}{1 + g_t}\, b_{t-1}$), que recoge la diferencia entre el tipo de interés medio de la deuda y el crecimiento nominal de la economía; y, (iii) los ajustes entre déficit y deuda (ADD), que agrupan todas las operaciones que aumentan o reducen la deuda sin pasar por el déficit.\footnote{%
    Siendo $b$ la ratio de deuda sobre el PIB, $i$ el tipo de interés nominal implícito de la deuda, $g$ el crecimiento del PIB nominal, $sp$ el saldo primario en porcentaje del PIB y ADD los ajustes entre déficit y deuda: $\Delta b_t \approx \frac{i_t - g_t}{1 + g_t}\, b_{t-1} - sp_t + \text{ADD}_t$. El primer término es el efecto bola de nieve. Si $i < g$, la ratio puede reducirse aun con déficits primarios moderados. Para más información ([\authorfont{Ver Tema 3.A.39}]).
} 

Esta descomposición es puramente contable y es la herramienta que nos permitirá entender la evolución de las finanzas de las AAPP en España. Para verlo claramente y fijar los términos, vamos a presentar tres episodios (sobre los que incidiremos más adelante) que lo ilustran perfectamente.

En 2008-2014, la deuda pasó del $36\%\to105\%$ del PIB. Contribuyeron los tres componentes: grandes déficits primarios por la caída de los ingresos y los estabilizadores; un efecto bola de nieve desfavorable, con tipos de interés que subieron con la crisis de la deuda soberana mientras el PIB nominal caía; y ajustes entre déficit y deuda muy importantes, por la recapitalización bancaria con el préstamo del MEDE, la creación del FADE, el mecanismo de pago a proveedores y los préstamos a los países rescatados de la zona del euro.

Más adelante, entre 2015-2019, la deuda bajó lentamente, del $105\%\to98\%$ del PIB, a pesar de que hubo déficits en todos los años. La razón fue el efecto bola de nieve favorable: el crecimiento nominal superó el tipo de interés medio de la deuda, que se reducía gracias a la política del BCE. Finalmente, entre 2021-2025, la deuda bajó del máximo cercano al $120\%$ del PIB (2020) al $100,7\%$ de 2025, de nuevo con déficits todos los años. ¿Por qué? De nuevo, la razón fue el efecto denominador de un crecimiento nominal muy fuerte, impulsado por la recuperación y por la inflación de 2022-2023, con un coste medio de la deuda todavía bajo gracias a la vida media larga.

La conclusión es simple: en la historia reciente de España, los periodos de reducción de la deuda se deben al crecimiento nominal y a los tipos bajos más que a superávits primarios. Por tanto, no es de extrañar que sea una situación considerada \textit{frágil} por la literatura, ya que depende de que la diferencia entre el tipo de interés y el crecimiento siga siendo favorable.\footnote{
No es de extrañar, por tanto, que la reducción de la deuda proyectada por la AIReF se apoye principalmente en el crecimiento del PIB nominal, con una contribución muy notable del deflactor, y que, cuando termine ese impulso, las proyecciones dibujen una dinámica desfavorable a políticas constantes.

La literatura sobre el diferencial entre el tipo de interés y el crecimiento (BLANCHARD, 2019) sostiene que, con $r < g$, la deuda es menos costosa en términos de bienestar, pero advierte de que el diferencial puede invertirse bruscamente, como ocurrió desde 2022.
} 

\begin{specialblock}{\textbf{Resumen}: Principales ajustes entre déficit y deuda (ADD)}
    Para leer bien la historia conviene tener presentes los principales ajustes entre déficit y deuda, es decir, las operaciones que alteran la deuda sin pasar, o sin pasar íntegramente, por el saldo. En cada caso importan la dirección del ajuste (si aumenta o reduce la deuda) y el componente de la identidad por el que discurre (el déficit, la adquisición o venta de activos financieros, o un cambio de perímetro):
    \begin{lnum}
        \item \textbf{Recapitalización bancaria (FROB, 2009-2013)}. Aumenta la deuda. Las inyecciones de capital que no se esperaba recuperar se registraron como transferencias de capital y, por tanto, como déficit (en torno a $3,7$ puntos de PIB en 2012); el resto se contabilizó como adquisición de activos financieros (acciones y participaciones), que eleva la deuda bruta sin afectar al saldo. La parte financiada con el préstamo del MEDE de 2012 ($41.333$ millones) computó como deuda de las AAPP;
        \item \textbf{Asistencia a otros países de la zona del euro}. Aumenta la deuda sin déficit. Los préstamos bilaterales a Grecia, la parte de la deuda del FEEF imputada a España y el capital desembolsado en el MEDE son adquisiciones de activos financieros: el Estado se endeuda para prestar o para aportar capital, sin que haya gasto en el saldo;
        \item \textbf{Fondo de Amortización del Déficit Eléctrico (FADE)}. Aumenta la deuda sin déficit en el año. El FADE titulizó, con emisiones avaladas por el Estado, los derechos de cobro de las eléctricas por el déficit de tarifa acumulado. Al integrarse su deuda en la de las AAPP por decisión de Eurostat, pasó a computar como deuda pública, aunque el desequilibrio que la originó se había generado años antes y fuera del saldo;\footnote{%
            El FADE, creado en 2010 y 2011, emitió deuda con garantía del Estado para financiar los derechos de cobro de las eléctricas por el déficit de tarifa acumulado. El déficit de tarifa se atajó con la reforma del sector eléctrico de 2013 (Ley 24/2013).
        }
        \item \textbf{Sareb (reclasificación de 2021)}. Aumenta la deuda por un cambio de perímetro. No es una operación nueva, sino la incorporación a las AAPP de una entidad y de su deuda ya existente (unos $35.000$ millones), con efecto también sobre el déficit a través de sus pérdidas;
        \item \textbf{Mecanismos de pago a proveedores (2012)}. Aumentan la deuda PDE sin déficit. Convirtieron deuda comercial (facturas pendientes, que no computan como deuda PDE) en préstamos, que sí computan. Solo las facturas que no se habían registrado en su momento afloraron además como déficit;
        \item \textbf{Variación de los depósitos del Tesoro}. Aumenta la deuda cuando el Tesoro acumula liquidez y la reduce cuando la utiliza. Emitir para engrosar el colchón de tesorería eleva la deuda bruta, pero el efectivo obtenido es un activo, no un gasto;
        \item \textbf{Primas y descuentos de emisión}. Reducen la deuda PDE cuando se emite sobre la par y la aumentan cuando se emite bajo la par, porque la deuda computa por su nominal y no por el efectivo obtenido. El efecto se compensa a lo largo de la vida del bono a través de los cupones;
        \item \textbf{Privatizaciones}. Reducen la deuda sin mejorar el saldo. La venta de participaciones es una operación financiera (el Estado cambia un activo por liquidez) cuyo producto se destina a amortizar deuda o a reducir la emisión;
        \item \textbf{Préstamos del Estado a las CCAA (FLA, FFCCAA) y a la Seguridad Social}. No alteran la deuda consolidada de las AAPP, porque el pasivo del prestatario y el activo del Estado se eliminan al consolidar; sí aumentan la deuda del Estado y, por tanto, las necesidades del Tesoro.
    \end{lnum}
\end{specialblock}

\subsubsection{Del déficit de las AAPP a las necesidades de financiación del Tesoro}
Por último, y con todo sobre la mesa, debemos destacar que lo relevante para la política de financiación del Tesoro no será el saldo presupuestario de las AAPP, sino la \textbf{necesidad de endeudamiento del Estado}. 

El déficit de las AAPP es una magnitud de contabilidad nacional, en términos de devengo, para el conjunto de las AAPP. La necesidad de endeudamiento del Estado es una magnitud de caja y del Estado. Para pasar de una a otra hay que considerar cinco elementos:
\begin{la}
    \item \textbf{Emisión de otros subsectores}. La parte del déficit que corresponde a los otros subsectores, que el Estado no financia directamente salvo cuando les presta;
    \item \textbf{Criterio de caja}. Las diferencias entre caja y devengo en el propio Estado;\footnote{%
        La distinción entre la contabilidad presupuestaria (caja y derechos u obligaciones reconocidos) y la contabilidad nacional (devengo), y los ajustes que la IGAE aplica para pasar de una a otra, se desarrollan en ([\authorfont{Ver Tema 4.B.21}]).
    }
    \item \textbf{Préstamos intersectoriales}. Los préstamos del Estado a las comunidades autónomas, a través del FLA y la Facilidad Financiera, y a la Seguridad Social, que el Tesoro debe financiar en el mercado aunque se eliminen al consolidar;
    \item \textbf{Operaciones financieras}. Las adquisiciones netas de activos financieros por parte del Estado: aportaciones de capital a fondos y organismos (como el MEDE), préstamos a otros países o entidades y compras de participaciones, netas de los cobros por la amortización de préstamos concedidos o por la venta de activos. No son gasto en contabilidad nacional, porque el Estado cambia liquidez por otro activo, y por eso no aparecen en el déficit; pero exigen caja y, por tanto, emisión. A ellas se suma la variación de la tesorería del Estado: cuando el Tesoro decide aumentar su colchón de liquidez emite más de lo que necesita para pagar, y cuando lo reduce emite menos;
    \item \textbf{Vencimientos de deuda}. Los vencimientos de la deuda emitida en años anteriores, que no forman parte del déficit, pero sí de las necesidades brutas de financiación (emisión bruta, operaciones de refinanciación).
\end{la}

Con esto en mente podemos comprender por qué en 2025, el déficit de las AAPP fue de unos $37.000$ a $40.000$ millones de euros. Pero el Tesoro programó para 2025 y 2026 una emisión neta de unos $55.000$ millones y, para 2026, una emisión bruta de $285.677$ millones. 

La diferencia entre el déficit y la emisión neta se explica por los préstamos del Estado a los demás subsectores y por las operaciones financieras. Solo en 2025, la deuda de la Seguridad Social, casi toda con el Estado, aumentó en unos $10.000$ millones, y la de las comunidades, en unos $5.700$ millones. La diferencia entre la emisión neta y la bruta se explica por los vencimientos: el Tesoro debe refinanciar cada año en torno al $13\%$ del saldo vivo. 

He aquí, entonces, la clave para entender el Tema, la razón por la que el título une el saldo de las AAPP y la política de financiación del Tesoro: \textbf{el primero determina, a través de esta bisagra, la dimensión de la segunda}.


\subsection{Desarrollo histórico: ¿cómo ha evolucionado a lo largo de los años?}
Ahora que tenemos en mente lo necesario para comprender estas dos magnitudes y la relación entre ellas, pasamos a examinar la evolución histórica de éstas.

Remontémonos a la génesis misma de nuestra Nación, tal y como la concebimos hoy en día: al periodo que se expande desde 1975 hasta 1995. 

\subsubsection{Del equilibrio aparente al déficit estructural (1975-1995)}
Dividiremos el periodo en cuatro partes. El primero trata la herencia: por qué el equilibrio presupuestario de 1975 era un síntoma de atraso y no de solidez. El segundo cubre los años 1975-1985: la construcción simultánea del Estado democrático, del Estado de bienestar y del Estado autonómico, y la omnipresencia del déficit. El tercero cubre la integración en la CEE y una consolidación que fue más cíclica que estructural. Acabaremos con la crisis, el SME y la bola de nieve de la deuda que se prolonga a nuestros días.

Durante estos veinte años, la hacienda pública española pasó de ser una de las más pequeñas de Europa occidental a tener un tamaño homologable a la media comunitaria. Lo hizo en muy poco tiempo y en las peores condiciones macroeconómicas: dos crisis del petróleo, una reconversión industrial, una crisis bancaria sistémica, una transición política y una descentralización territorial sin precedentes. Todo a la vez y en todas partes. Como podemos suponer, entonces, el déficit de este periodo no será un accidente coyuntural. Es el precio de un catch-up fiscal: la construcción del Estado de bienestar se hizo antes de que el sistema tributario pudiera financiarla, y los mecanismos de financiación del déficit (primero la monetización, después la deuda a corto plazo y, al final, la deuda de mercado) fueron cambiando por el camino. La deuda pasa de menos del $10\%$ a más del $60\%$ del PIB, y su composición explica muchas decisiones posteriores del Tesoro.

Identificaremos dos lógicas distintas de acumulación de deuda. Hasta mediados de los ochenta la deuda crece por déficits primarios, pero la represión financiera y la inflación alta mantienen un diferencial tipo-crecimiento favorable, y el déficit se financia primero con el BdE y después con instrumentos casi monetarios como los pagarés. En los noventa, con la liberalización financiera, la defensa del SME y la recesión, el diferencial se vuelve fuertemente positivo: la bola de nieve explica la mayor parte del salto de la deuda de $(\sim41\%) \to (\sim62-65\%)$.\footnote{%
    Además, la autonomía del BdE (1994) pone fin a la monetización, y Maastricht introduce la restricción externa que disciplinará la siguiente década.
}

\paragraph{La herencia: un equilibrio sin Estado}
Como en cualquier proceso de cambio, el marco de ideas dominante condiciona las decisiones que se toman, y la política presupuestaria no es una excepción. Si repasamos brevemente la historia del pensamiento económico, hasta los años treinta del siglo XX predominó la ortodoxia clásica del presupuesto equilibrado. El \textit{laissez-faire} limitaba el papel del Estado a la provisión de los bienes públicos tradicionales (justicia, defensa, infraestructuras) y de las condiciones necesarias para el funcionamiento del mercado, lo que mantenía el gasto en niveles reducidos. La deuda se contemplaba con abierta desconfianza: RICARDO la consideraba una de las más terribles plagas jamás inventadas para afligir a una nación.\footnote{%
    La equivalencia entre la financiación mediante impuestos y mediante deuda, su formulación moderna por BARRO y el debate sobre la carga de la deuda se desarrollan en ([\authorfont{Ver Tema 3.A.39}]).
}

La Gran Depresión rompió ese consenso. Con KEYNES, la demanda agregada pasa a ocupar el centro del análisis y la política fiscal se convierte en el instrumento para reactivar una economía en crisis; el \textit{New Deal} (1933) es su ejemplo más conocido. La teoría justificó además la intervención estabilizadora por la inestabilidad intrínseca de las economías de mercado que se desprendía del modelo de HARROD-DOMAR, o por la necesidad de evitar el estancamiento secular que planteaba HANSEN (1938). A ello se sumó el arraigo del principio de equidad, que dio lugar al Estado de bienestar y a un aumento sin precedentes del peso del sector público en la economía. El resultado fue el abandono del equilibrio anual como norma y la aceptación de los déficits discrecionales. La reconstrucción de posguerra y el crecimiento de la etapa de Bretton Woods permitieron, sin embargo, que la mayoría de los países desarrollados mantuvieran cuentas relativamente equilibradas durante los cincuenta y los sesenta, pese a la expansión del gasto social.

Las crisis del petróleo de los setenta cambiaron el escenario. Con crecimiento bajo y desempleo elevado, el gasto comprometido en la etapa anterior no pudo recortarse y el déficit se convirtió en un rasgo estructural de las economías avanzadas. La estanflación puso además en entredicho la capacidad de las políticas de demanda para estabilizar el ciclo, tal como defendía la síntesis neoclásica, y ganaron terreno el monetarismo y la economía de la oferta. El giro antiinflacionista de la Reserva Federal desde el nombramiento de VOLCKER (1979) cerró la etapa de la monetización: el déficit pasó a financiarse fundamentalmente con deuda, los tipos reales se volvieron positivos y elevados, y la acumulación de deuda de los ochenta y noventa, junto con las crisis de deuda de muchos países en desarrollo, hizo visible el riesgo de insostenibilidad. De ahí la recuperación de la idea de equilibrio presupuestario y la introducción de reglas fiscales para anclar la credibilidad macroeconómica, con el Tratado de Maastricht y sus criterios de convergencia como hito europeo.\footnote{%
    Sobre la evolución de las teorías del equilibrio presupuestario ([\authorfont{Ver Tema 3.A.39}]); sobre las reglas fiscales de la Unión y su traslación española ([\authorfont{Ver Tema 3.B.45}]).
}

España recorrió esta secuencia con una particularidad decisiva: llegó tarde a cada una de sus fases y las recorrió comprimidas. Cuando los países de la OCDE construían sus Estados de bienestar con crecimiento alto y tipos reales bajos (1950-1973), España tenía un sector público mínimo. Cuando lo construyó (1977-1995), el contexto internacional ya era el contrario: crecimiento bajo, desempleo alto y, desde principios de los ochenta, tipos reales positivos y elevados. Esa desincronización es clave para entender por qué el déficit español aparece más tarde que en Europa pero crece más deprisa.

Las magnitudes de partida confirman ese retraso. En 1975, las cuentas públicas españolas presentaban equilibrio presupuestario o incluso pequeños superávits, con un gasto público en torno al $25\%$ del PIB, unos veinte puntos por debajo de la media comunitaria; un sistema tributario de escasa capacidad recaudatoria, ineficiente y poco redistributivo; y una deuda pública inferior al $10\%$ del PIB, muy por debajo de la de los países comunitarios, cuya media rondaba entre el $35$ y el $40\%$. Las cifras son coherentes con las series históricas de COMÍN y DÍAZ FUENTES (2005) y con la reconstrucción de LAGO PEÑAS (2024), que sitúa la ratio de deuda por debajo del $10\%$ en 1976 y 1977; otras series del BdE la dejan en torno al $7\%$ en 1975-1976. La presión fiscal, según la OCDE, rondaba el $18\%$ del PIB a mediados de los setenta, frente a una media de la OCDE cercana al $29\%$: un diferencial de diez puntos o más.


No obstante, bajo la aparente solidez de las cuentas públicas de 1975 se escondía un sistema tributario poco eficiente y un gasto público insuficiente por tres razones, esencialmente.

La primera es de economía política. COMÍN (1996) y COMÍN y DÍAZ FUENTES (2023) describen la hacienda franquista como insuficiente e injusta. Su equilibrio no respondía a una preferencia social por la austeridad. Era el resultado de un pacto implícito entre el régimen y sus bases de apoyo: impuestos bajos y poco progresivos a cambio de un gasto social mínimo. La reforma tributaria propuesta en 1973 (Libro Verde, impulsado por FUENTES QUINTANA) fue archivada precisamente porque aumentar la progresividad y la transparencia era políticamente inasumible para el régimen. Por tanto, vemos como España llega a la democracia con treinta años de retraso tributario acumulado.

La segunda dirección es financiera. El equilibrio del presupuesto del Estado convivía con mecanismos de financiación del sector público que no pasaban por el presupuesto. ESCARIO, SABATÉ y GADEA (2011) reconstruyen la sombra monetaria del déficit a lo largo de toda la era de la peseta. Tras la pignoración automática de la deuda (hasta el Plan de Estabilización de 1959), el canal principal en los sesenta y setenta fue el redescuento especial: crédito privilegiado del BdE a sectores prioritarios. A él se sumaron los coeficientes de inversión obligatoria, que desde 1971 obligaban a la banca a mantener una parte de su activo en fondos públicos y en crédito privilegiado.\footnote{%
    Caso de manual de la represión financiera de REINHART y SBRANCIA (2015): el sector público se financia a tipos reales bajos o negativos obligando a los intermediarios a absorber su deuda.
} La deuda explícita baja no significaba ausencia de financiación pública; significaba que ese coste lo soportaban de forma implícita los ahorradores, vía tipos de interés reprimidos e inflación.

La tercera dirección es perimetral. Una parte de la actividad pública (el INI, el crédito oficial y numerosos organismos autónomos) quedaba fuera del presupuesto del Estado y fuera de lo que entonces se consolidaba. Con los criterios actuales del SEC, parte de esa actividad se clasificaría dentro de las AAPP. No existe una reconstrucción homogénea que permita afirmar cuánto déficit oculto había. Pero el argumento cualitativo es sólido: el equilibrio de 1975 se mide con un perímetro más estrecho que el de hoy.

La lectura en términos de dinámica de la deuda es inmediata: $\Delta d_t \approx  (i - g)/(1 + g)\cdot d_{t-1} - sp_t + \text{ADD}_t$. En 1960-1975 la economía española creció a tasas nominales muy altas (real en torno al $7\%$ más una inflación creciente), y el tipo de interés implícito de la deuda estaba reprimido. El diferencial entre tipo de interés y crecimiento nominal era fuertemente negativo, de modo que incluso pequeños déficits primarios habrían sido compatibles con una ratio de deuda decreciente. 

Por eso la deuda cae a mínimos históricos justo cuando empieza la Transición: el punto de partida no reflejaba una disciplina fiscal extraordinaria, sino un entorno de crecimiento nominal alto y represión financiera que desaparecería en la década siguiente.

\paragraph{1975-1985: tres Estados construidos a la vez y la omnipresencia del déficit}
La Constitución de 1978 consagra el Estado de bienestar y el Estado autonómico, lo que trae un aumento del gasto público de carácter estructural. A ese componente se suma otro coyuntural: la crisis económica internacional de los setenta obligó a una reconversión industrial que absorbió grandes cantidades de recursos públicos. El gasto supera el $40\%$ del PIB en 1985.

La distinción entre componente estructural y coyuntural es pertinente, pero conviene precisar cuatro motores del gasto con su base normativa. Los tres primeros son estructurales: (i) \textbf{pensiones}, su gasto crece por maduración del sistema, extensión de la cobertura y revalorizaciones por encima de la inflación en los primeros años de la democracia ([\authorfont{Ver Tema 4.B.24}]); (ii) \textbf{protección por desempleo}, que la Ley 31/1984 la reordena en un contexto en que el paro supera el $20\%$ a mediados de los ochenta; y, (iii) \textbf{sanidad y educación}, cuya universalización culmina con la Ley 14/1986, General de Sanidad, y la LOGSE de 1990. El cuarto motor es coyuntural pero persistente: la reconversión industrial, con la Ley 27/1984 de reconversión y reindustrialización y las pérdidas de las empresas del INI.

Pero hay un quinto elemento especialmente relevante para la deuda: la crisis bancaria de 1977-1985 y la expropiación de Rumasa de 1983. SUDRIÀ (2014) documenta que afectó a $58$ bancos, el $52\%$ del total, que representaban el $27\%$ del sector en recursos ajenos, recursos propios o empleo. Las aportaciones públicas rondaron $1,3$ billones de pesetas hasta 1985, frente a unos ingresos presupuestados del Estado de $4,6$ billones ese año. Una parte se canalizó a través del Fondo de Garantía de Depósitos (creado en noviembre de 1977) y del propio BdE, y otra a través del presupuesto. Atendiendo a nuestra descomposición, por tanto, una fracción de ese coste aparecería como transferencia de capital en el déficit y otra como ajuste déficit-deuda.\footnote{%
    No existe una partición homogénea publicada, pero es el primer gran episodio de lo que conocemos como deuda que no es déficit. Por tanto, el paralelismo con 2012 (FROB, MEDE, Sareb) es claro y conviene explotarlo en el examen.
} 

A estos motores se suma la descentralización. La LOFCA (1980) y los traspasos de competencias iniciados a principios de los ochenta crean un nuevo nivel de gasto. En esta etapa el efecto no es tanto un aumento neto del gasto (las competencias se traspasan con su coste efectivo) como una pérdida de control centralizado sobre su dinámica futura, porque el gasto autonómico empieza a crecer con lógicas propias.\footnote{%
    Sobre la evolución del sistema de financiación autonómica y sus incentivos fiscales (restricción presupuestaria blanda, corresponsabilidad), ([\authorfont{Ver Tema 4.B.25}]).
}

La respuesta por el lado de los ingresos llegó con los Pactos de la Moncloa (1977), que acordaron una profunda reforma del sistema fiscal. Su traslación normativa fue rápida. La Ley 50/1977, de Medidas Urgentes de Reforma Fiscal, creó un Impuesto Extraordinario sobre el Patrimonio de las Personas Físicas, que, pese a su nombre, estuvo vigente hasta la Ley 19/1991; tipificó el delito fiscal y levantó el secreto bancario a efectos tributarios. La Ley 44/1978 creó el IRPF y la Ley 61/1978 reformó el Impuesto sobre Sociedades, que adquieren así su forma básica actual. Se eliminaron, además, figuras arcaicas: los impuestos de producto y el impuesto general sobre la renta se sustituyeron por un sistema de retenciones en la fuente sobre los rendimientos del trabajo y del capital y por un IRPF sintético y de tarifa progresiva ([\authorfont{Ver Tema 4.B.16}]).

Sin embargo, los ingresos tardaron mucho más en responder que el gasto, y de ahí la omnipresencia del déficit público. ¿Por qué tardó tanto la reforma en rendir? La literatura ofrece tres explicaciones complementarias. En primer lugar, la administración tributaria: faltaban medios, censos y cultura de cumplimiento, y un impuesto sintético y progresivo como el IRPF exige una administración que España no tenía (la AEAT no se crea hasta 1992). En segundo lugar, la progresividad en frío: con una inflación del $15$-$25\%$, la tarifa sin deflactar elevaba mecánicamente la recaudación del IRPF, lo que alivió en parte el problema, pero generó una fuerte presión política para deflactarla y erosionó la legitimidad del impuesto, un debate que ha reaparecido en 2022-2024. En tercer lugar, la composición del esfuerzo: según la OCDE, la presión fiscal pasó de en torno al $18\%\to27\%$ del PIB entre 1975 y 1985, un aumento de unos nueve puntos que en términos comparados es enorme; pero el gasto creció más de quince puntos en el mismo periodo. El déficit no fue consecuencia de que los ingresos no subieran, sino de que el gasto subió aún más deprisa.


El resultado de todo lo anterior fueron déficits primarios persistentes con deuda todavía moderada. La base histórica de deuda del FMI sitúa la ratio en el $40,9\%$ en 1985 ($16,1\%$ en 1980, $29,5\%$ en 1983). PRATS ALBENTOSA y ROCAMORA (2016) dan un $44\%$ de deuda y un $5,8\%$ de déficit. 

No obstante, lo interesante no es el nivel sino la mecánica de la acumulación. En diez años la deuda sube más de $30$ puntos de PIB, de $7\%\to 41\%$. Con un crecimiento nominal medio del orden del $15\%$ anual e intereses implícitos todavía reprimidos, el diferencial entre tipo de interés y crecimiento fue negativo y restó deuda. Por lo tanto, la totalidad del aumento (y algo más) se explica por los \textbf{déficits primarios acumulados} y por los \textbf{ajustes déficit-deuda} vinculados a la crisis bancaria y a la asunción de pasivos de empresas públicas. 

Y, ¿cómo se financió? ESCARIO, SABATÉ y GADEA (2011) precisan que, entre 1978 y 1983 el sector público alcanzó su máxima contribución histórica a la creación de base monetaria, fundamentalmente mediante el crédito del BdE al Tesoro a través de su cuenta corriente. Los autores sitúan en 1984 el fin efectivo de la financiación por señoreaje. A partir de ese año la contribución del sector público a la base monetaria se vuelve negativa, y el déficit se financia con deuda colocada fuera del banco central. El instrumento de transición fueron los Pagarés del Tesoro, creados en 1981: títulos a corto plazo con un régimen fiscal opaco (no sujetos a retención), que los hacía atractivos para capitales que buscaban anonimato. Además eran computables en los coeficientes obligatorios de la banca. Según los mismos autores, entre 1984 y 1989 los pagarés y las letras explican cerca de una cuarta parte del crecimiento de los activos líquidos en manos del público (ALP), y en 1986 más de la mitad. Ahora bien, estos activos eran casi dinero: no ampliaban la base monetaria pero sí la liquidez de la economía. Por tanto, dejar de monetizar no fue un salto a la ortodoxia, sino una transición gradual por instrumentos intermedios.\footnote{%
    La lectura teórica es la de la aritmética monetarista desagradable de SARGENT y WALLACE (1981). Si la autoridad fiscal es dominante y fija una senda de déficits que el público no absorbe como deuda de forma indefinida, la monetaria acabará teniendo que monetizarlos, y una política monetaria restrictiva hoy puede traducirse en más inflación mañana ([\authorfont{Ver Tema 3.A.39}]).

Mientras la autoridad fiscal fija la senda de déficits y la monetaria se acomoda, hay dominancia fiscal. Romperla exige que el banco central deje de financiar al Tesoro, y en España eso no se consolida legalmente hasta 1994. La reforma de 1983-1984, que acota la apelación del Tesoro al BdE, es el primer paso. En términos de credibilidad es incompleto, porque la separación dependía de acuerdos administrativos y no de una prohibición legal.
} 

\paragraph{1986-1990: la CEE, el IVA y una consolidación más cíclica que estructural}
El marco temporal que se expande de 1986 a 1990 fue una etapa de reducción del déficit gracias a cuatro factores: el crecimiento en torno al $5\%$ (que alivia los estabilizadores automáticos), la reducción de subvenciones exigida por la normativa comunitaria, la maduración de las reformas de los Pactos de la Moncloa y la introducción del IVA en 1986.\footnote{%
    El IVA entra en vigor el 1 de enero de 1986 (Ley 30/1985), como condición de la adhesión (Sexta Directiva, 77/388/CEE). Sustituye al Impuesto General sobre el Tráfico de Empresas, un impuesto en cascada, y a otras figuras indirectas. Sus efectos recaudatorios fueron positivos pero menores de lo que suele suponerse, porque sustituía a impuestos existentes ([\authorfont{Ver Tema 4.B.20}]). Por el lado del IRPF, la STC 45/1989 declara inconstitucional la tributación conjunta obligatoria de los matrimonios. Esto obligó a una reforma que culmina en la Ley 18/1991. La sentencia tuvo un coste recaudatorio y abrió un periodo de inestabilidad normativa ([\authorfont{Ver Tema 4.B.16}]).
} 

Conforme a la base histórica del FMI, la deuda se estabiliza en un $38,5-42\%$ durante el periodo: $42,1\%$ en 1986, $38,5\%$ en 1988 y $41,3\%$ en 1990. Por lo que respecta al déficit, su mínimo del periodo se alcanza en 1989, en torno al $3\%$ (o algo por debajo según la serie). No obstante, en 1990 el saldo ya se deteriora hasta el $4\%$ aproximadamente. La cuestión que resulta interesante dentro del periodo es, ¿por qué? ¿Por qué ese mínimo es tan esporádico? Aquí está el punto de mayor calado analítico. 

Entre 1986 y 1990 la economía española creció por encima de su potencial: un $4,5-5,5\%$ anual real con inflación en aumento. En una fase así, el saldo total mejora casi mecánicamente por los estabilizadores automáticos aunque la política fiscal discrecional sea expansiva. La lectura dominante en la literatura es que la política fiscal española fue procíclica en la fase alcista: el saldo estructural no mejoró o empeoró mientras mejoraba el total.\footnote{%
    Ver: OCDE, Economic Surveys de finales de los ochenta; el BdE en sus Informes Anuales de la época; y, con perspectiva, los análisis de ciclicidad de GALÍ y PEROTTI, 2003, para Europa.
} La prueba política fue la huelga general del 14 de diciembre de 1988. Tras ella el Gobierno aumentó el gasto social: mejora de las pensiones mínimas, extensión de la protección por desempleo y, poco después, la Ley 26/1990 de pensiones no contributivas. Evidentemente, son decisiones legítimas en términos de equidad, pero se toman en el pico del ciclo y comprometen gasto estructural justo antes de la recesión.\footnote{%
    ALESINA y DRAZEN (1991), con su modelo de guerra de desgaste, predicen que la estabilización se retrasa cuando hay conflicto distributivo sobre quién soporta el ajuste, lo que es plausible en el contexto de ruptura del diálogo social que culmina en el 14-D. Por su parte, ROUBINI y SACHS (1989) asocian los déficits a los gobiernos de coalición. Su predicción no encaja con España, que tuvo mayorías absolutas de un solo partido entre 1982 y 1993. La explicación española hay que buscarla más en la debilidad de las instituciones presupuestarias (VON HAGEN, 1992, sitúa a España en el grupo de procedimientos presupuestarios poco centralizados) y en la falta de reglas fiscales, que no existirán hasta Maastricht.
}

No obstante, el elemento estructural más duradero del periodo no fue fiscal sino institucional: la CEE, y desde 1989 el SME, introducen por primera vez un ancla externa. Con esto, España materializó lo que GIAVAZZI y PAGANO (1988) concibieron como \textit{atarse las manos}: un país con baja credibilidad antiinflacionista importa la del Bundesbank al fijar su tipo de cambio. 

No obstante, esto tuvo un problema: el ancla monetaria sin ancla fiscal genera tensiones. Con una política fiscal laxa (como la de la época), el BdE sostuvo tipos de interés muy altos para contener la inflación. Esos tipos atrajeron capitales y apreciaron la peseta, y tras la entrada en el SME en junio de 1989 la peseta se situó en la parte alta de la banda. Combinación (fiscal expansiva, monetaria restrictiva, tipo de cambio sobrevalorado) que se paga en 1992-1993.

Un patrón opuesto al que veremos a continuación.

\paragraph{1991-1995: la crisis, el SME y la bola de nieve}
La economía se desacelera desde 1990 y desemboca en la crisis de 1992-1993, en parte por las debilidades del SME. 

Varias economías europeas sufrieron crisis cambiarias en esos años. El Reino Unido e Italia abandonaron el mecanismo de cambios del SME en septiembre de 1992, y Finlandia, Suecia y Noruega, que mantenían vinculaciones unilaterales al ECU, las abandonaron también ese año. La peseta, que se había incorporado al SME en junio de 1989, se devaluó en cuatro ocasiones: un $5\%$ en septiembre de 1992, un $6\%$ en noviembre de 1992, un $8\%$ en mayo de 1993 y un $7\%$ en marzo de 1995. En agosto de 1993, las presiones persistentes obligaron a ampliar las bandas de fluctuación de las divisas que seguían en el mecanismo hasta el $\pm15\%$.

A la recesión se sumaron otros factores de gasto. Los años previos coincidieron con un fuerte esfuerzo inversor en infraestructuras y grandes eventos: la primera línea del AVE entre Madrid y Sevilla, los Juegos Olímpicos de Barcelona y la Expo de Sevilla, en 1992. El gasto en desempleo se disparó. La reforma laboral de 1984 (Ley 32/1984), que introdujo el contrato temporal de fomento del empleo, había hecho el empleo extremadamente sensible al ciclo (BENTOLILA y DOLADO, 1994), y las prestaciones se habían mejorado tras la huelga general de 1988; la tasa de paro rondó el $24\%$ en 1994. Y creció la carga de intereses, tanto por el aumento de la deuda como por los tipos más altos que exigían los inversores. Como resultado, el gasto público superó el $45\%$ del PIB en 1993, el déficit rondó el $7\%$ del PIB en 1993 y en 1995 (año del máximo registrado hasta entonces según ESTEVE y PRATS, 2024) y la deuda pasó del $41,3\%$ en 1990 al $61,6\%$ en 1995 y al $65,4\%$ en 1996.

¿Qué explica ese salto? Aquí la descomposición de la deuda es más reveladora que en ningún otro momento del periodo. Más de $10$ puntos del aumento se concentraron en 1993 ($44,1\%\to54,6\%$), y ese salto no se explica solo por el déficit. En orden de magnitud contribuyen tres factores: (i) el \textbf{déficit primario}, en torno a $2$ puntos de PIB en 1993, con un saldo total de alrededor del $7\%$ y unos intereses del $4,5$-$5\%$; (ii) el \textbf{efecto bola de nieve}, pues en 1993 el PIB real cayó en torno a un $1\%$ y el crecimiento nominal apenas alcanzó el $2$-$3\%$, mientras el tipo implícito de la deuda superaba el $10\%$, de modo que, con una ratio cercana al $45\%$, el diferencial aportó por sí solo $3$-$4$ puntos; y, (iii) los \textbf{ajustes déficit-deuda}, entre ellos la revalorización en pesetas de la deuda en divisas tras las devaluaciones, la asunción de deudas de entes públicos y la acumulación de activos líquidos del Tesoro.

La conclusión analítica es nítida. Entre 1975 y 1985 la deuda crece por déficits primarios con un diferencial favorable; entre 1991 y 1995 crece sobre todo por el diferencial, con déficits primarios moderados. Los intereses llegan a su máximo histórico en 1995-1996, en torno al $5,3\%$ del PIB, frente al $2,5\%$ de 2024-2025. El cambio de régimen tiene tres causas: la liberalización financiera, que acaba con la represión y eleva los tipos reales; la defensa de la peseta en el SME, que exige tipos muy altos; y la recesión, que hunde el crecimiento nominal. Este es el nexo con el tema del Tesoro. La estructura de la deuda de la época (predominio del corto plazo, vida media de pocos años) hacía que cada subida de tipos se trasladara rápidamente al coste medio. La vida media de la deuda del Estado rondaba entonces los $2-3$ años, frente a más de $8$ hoy. Alargar esa vida media será el objetivo central de la política de emisiones desde 1995-1996.

El marco institucional cambió en estos años de forma irreversible. El Tratado de Maastricht (1992) incluyó entre los requisitos para acceder a la UEM un déficit no superior al $3\%$ del PIB, lo que obligaba a iniciar la consolidación. Sin embargo, 1995 no fue el comienzo del ajuste, sino el año del déficit récord: los PGE para 1996 fueron rechazados en el Congreso en octubre de 1995, se prorrogaron los de 1995 y el ajuste efectivo comenzó en 1996, sobre esos presupuestos prorrogados.

Dos hitos del periodo son decisivos para el tema. El primero es el Programa de Convergencia de 1992, primer ejercicio de planificación fiscal plurianual con objetivos numéricos vinculados a una restricción externa. Sus objetivos se incumplieron en 1992-1993 por la recesión y hubo que revisarlo en 1994, pero es el precedente de los Programas de Estabilidad y de los actuales planes fiscales estructurales de medio plazo.\footnote{%
    Evolución de la planificación fiscal plurianual europea (Programas de Convergencia y de Estabilidad, Semestre Europeo y planes fiscales estructurales de medio plazo del marco de 2024) ([\authorfont{Ver Tema 3.B.45}]).
} El segundo es la Ley 13/1994, de Autonomía del BdE, cuyo artículo 13.2 prohíbe la autorización de descubiertos o la concesión de cualquier otro tipo de crédito por el BdE al Estado, en aplicación del artículo 104 del Tratado de Maastricht. El saldo deudor que el Tesoro mantenía con el BdE se reconvirtió en un préstamo a largo plazo. Es el cierre formal de la etapa de monetización: desde 1994, todo el déficit se financia en el mercado, premisa de la política de financiación del Tesoro que estudiaremos más adelante.

Por último, la crisis de 1992-1993 dejó un legado para la protección social. Obligó a recortar las prestaciones por desempleo (RDL 1/1992 y Ley 22/1992) y condujo al Pacto de Toledo (1995), que separa las fuentes de financiación de la Seguridad Social. En esos años se sitúa también el origen de los préstamos del Estado a la Seguridad Social para cubrir sus desequilibrios, que figuran entre los ajustes déficit-deuda.\footnote{%
    Los préstamos del Estado a la Seguridad Social no computan como déficit del Estado sino como activo financiero, con impacto nulo en la deuda consolidada ([\authorfont{Ver Tema 4.B.24}]).
}

En definitiva, la deuda de 1995 es alta, cara y corta, y reducir su coste y alargar su vida media será la tarea del Tesoro en los años siguientes.

\subsubsection{Consolidación de Maastricht y el espejismo del superávit (1996-2007)}
Entre 1996 y 2007 las AAPP españolas pasan de un déficit de más del $5\%$ del PIB a un superávit en torno al $2\%$, y la deuda pública cae a la mitad: de cerca del $65-70\%$ del PIB, según la serie, a en torno al $36\%$. En términos de saldo total es la mejora más larga y continuada de la historia fiscal española, y culmina con los tres únicos superávits de la democracia (2005, 2006 y 2007).

El periodo tiene, sin embargo, dos fases con lógicas distintas. La primera (1996-1998) es un ajuste deliberado con un objetivo externo: cumplir los criterios de Maastricht para entrar en la UEM desde el primer momento. La segunda (1999-2007) es una mejora del saldo impulsada sobre todo por dos fuerzas que no eran de política fiscal: la caída de la carga de intereses que trajo el euro, y una expansión de ingresos ligada al auge inmobiliario y crediticio, un auge transitorio. La mejora del saldo total, por tanto, ocultó un deterioro relativo del saldo estructural, que los instrumentos de medición y las reglas de la época no permitían ver. 

Por tanto, la conclusión a la que debemos llegar no es que la política fiscal de este periodo fuese irresponsable en nivel, porque la deuda baja fue un activo en 2008. Sino, más bien, que fue procíclica en composición y débil en estructura. Y que la lección institucional (reglas de gasto, desequilibrios macroeconómicos, instituciones fiscales independientes) se aprendió después de la crisis, no antes.

\paragraph{1996-1998: el ajuste de Maastricht}
Como vimos, España llega a 1995 con el déficit registrado más elevado del periodo democrático hasta entonces. En esa línea, los PGE para 1996 fueron rechazados en el Congreso en otoño de 1995, lo que obligó a prorrogar los de 1995 y precipitó las elecciones de marzo de 1996. 

El nuevo Gobierno empezó el ajuste hacia Maastricht sobre un presupuesto prorrogado, mediante acuerdos de no disponibilidad de créditos, y lo intensificó con los PGE de 1997. El objetivo era la primavera de 1998, cuando el examen de convergencia se evaluaría según los datos de 1997. Teniendo en cuenta que el criterio fiscal exigía un déficit no superior al $3\%$ del PIB y una deuda por debajo del $60\%$ o disminuyendo suficientemente hacia ese valor, la labor era complicada.

Entre 1995 y 1997 el déficit logró bajar más de cuatro puntos. Por el lado del gasto, contribuyeron la contención de la masa salarial pública (congelación de 1997 y contención de la oferta de empleo), el recorte de la inversión pública, la partida más fácil de ajustar a corto plazo y la que más había crecido en 1992-1993, y el inicio de la caída de la carga de intereses, a medida que los mercados descontaban la entrada en el euro y se estrechaban los diferenciales con Alemania. Además, la Ley 24/1997, de consolidación y racionalización del sistema de Seguridad Social, primera traslación legislativa del Pacto de Toledo, reformó el Sistema para contener el gasto, ampliando de $8$ a $15$ años el periodo de cálculo de la base reguladora, e iniciando la separación de fuentes de financiación.\footnote{%
    Contenido de la Ley 24/1997, del Pacto de Toledo y del Fondo de Reserva: ([\authorfont{Ver Tema 4.B.24}]).
}

Debemos tener en cuenta, también, que en 1996 se aprueba también el Programa de Modernización del Sector Público Empresarial del Estado: la privatización de hasta $36$ empresas en cuatro años, con unos ingresos de unos $30.000$ millones de euros, más del $6\%$ del PIB de 1996. Sabemos que en contabilidad nacional, la venta de acciones es una operación financiera: el Estado cambia un activo (participaciones) por otro (liquidez), que no reduce el déficit (no es ingreso no financiero). Sin embargo, las privatizaciones actuaron sobre la deuda a través del ajuste déficit-deuda: los ingresos destinados a amortizar deuda o a reducir la emisión neta.\footnote{%
    Además, también reducen indirectamente la carga futura de intereses, aunque a cambio se pierden los dividendos.
} Por tanto, los ingresos de privatización fueron un factor relevante en la reducción de la ratio de deuda desde el máximo de 1996. 

En definitiva, según el Informe de Convergencia del Instituto Monetario Europeo (marzo de 1998), España registró en 1997 un déficit del $2,6\%$ del PIB y una deuda del $68,8\%$, habiendo alcanzado su máximo en 1996 ($70,1\%$). Por tanto, la deuda seguía claramente por encima del $60\%$, pero se había reducido y podía considerarse en camino hacia el cumplimiento del criterio. De hecho, aunque el IME advirtió de que España necesitaba un sustancial progreso en la consolidación fiscal para situar la deuda por debajo del $60\%$, estimaba que, con presupuestos equilibrados desde 1999, ese nivel se alcanzaría hacia 2001.

Con todo lo expuesto, no debe extrañar que el caso español de 1996-1998 se haya citado a veces como ejemplo de consolidación fiscal compatible con un crecimiento elevado: el PIB crece en torno al $4\%$ en 1997-1998 mientras se reduce el déficit.\footnote{%
    La hipótesis de la consolidación fiscal expansiva (GIAVAZZI y PAGANO, 1990; ALESINA y PEROTTI, 1995; ALESINA y ARDAGNA, 2010) sostiene que un ajuste creíble, sobre todo por el lado del gasto, puede estimular la demanda privada a través de las expectativas: menos impuestos futuros y menores primas de riesgo. El argumento teórico, a priori, es sólido: en un modelo con agentes ricardianos o con primas de riesgo endógenas, una consolidación creíble puede elevar la renta permanente esperada y reducir los tipos, y el efecto riqueza puede superar el efecto demanda directo, dando un multiplicador negativo. Además, hay cierta evidencia empírica que la respalda, aunque sujeta a cierta controversia metodológica (acción de política frente a variación del saldo ajustado de ciclo; GUAJARDO, LEIGH y PESCATORI, 2014). Para más información ([\authorfont{Ver Tema 3.A.40}]).
} No obstante, esta lectura cuantitativa debe ser también matizada por los efectos de segunda ronda que conllevaba entrar en la UEM. En la línea de PEROTTI (2011), para los casos europeos de los ochenta y noventa, la mayoría de estas experiencias de consolidación fiscal y crecimiento se explican sobre todo por canales: (i) el desplome de los tipos de interés por la convergencia hacia la UEM, que en España fue extraordinario (el diferencial a diez años frente a Alemania pasa de unos $500$ puntos básicos a mediados de 1995 a prácticamente cero en 1998); (ii) la depreciación previa de la peseta en 1992-1995, que mejoró la competitividad; y, (iii) la moderación salarial.

La consolidación fue probablemente condición necesaria para que esa convergencia de tipos se materializase, porque sin ella España no habría entrado en el euro en 1999. Pero no es evidencia de que el ajuste fiscal fuese por sí mismo expansivo. 

\paragraph{1999-2007: el euro y el dividendo de los tipos}
Durante el periodo posterior a la entrada a la UEM hasta el estallido de la burbuja inmobiliaria, España experimenta un crecimiento económico sostenido en paralelo a una notable consolidación fiscal.

El déficit del periodo previo, que en 1999 se situaba algo por encima del $1\%$, se convirtió en superávit durante el periodo, en 2007 se publicó en tiempo real en el $2,23\%$ del PIB ($23.368$ millones), con una deuda del $36,2\%$ (a cierre del Ministerio). Paralelamente, el saldo primario mejora del orden de $3,5$ puntos, de un valor cercano al equilibrio en 1996 a un superávit en torno al $3,5\%$ en 2007. Y buena parte de esa mejora primaria se concentra en 2005-2007 (precisamente cuando los ingresos ligados al ciclo inmobiliario alcanzan su máximo). 

La pregunta relevante, por tanto, es cuánta mejora del saldo se debió a la política fiscal y cuánta a factores externos. En primer lugar, los intereses de las AAPP pasan de un máximo de alrededor del $5,3\%$ del PIB en 1995-1996 a en torno al $1,6\%$ en 2007: un ahorro de más de $3,5$ puntos de PIB. Si la mejora total del saldo entre 1996 y 2007 es de unos $7$ puntos: casi la mitad de la mejora del saldo total no fue, por tanto, mejora del saldo primario, sino un dividendo del euro en forma de menor carga de intereses.\footnote{%
    Esta descomposición tiene un paralelo comparado conocido. Los países de la periferia del euro recibieron el mismo dividendo de tipos, y la cuestión es qué hicieron con él. FERNÁNDEZ-VILLAVERDE, GARICANO y SANTOS (2013) sostienen que el acceso a financiación barata debilitó los incentivos a hacer reformas en la periferia. En su lectura, el crédito político del auge hizo menos urgente cualquier ajuste. España ahorró una parte del dividendo, a diferencia de Grecia o Portugal, pero no la suficiente para compensar el carácter transitorio de parte de sus ingresos.
}

Con la identidad $\Delta d \approx (i - g)/(1 + g)\cdot d_{t-1} - sp_t + \text{ADD}_t$, la caída de la deuda de casi $30$ puntos entre 1996 y 2007 tiene dos motores. Unos dos tercios de la reducción se explican por los superávits primarios acumulados, y aproximadamente un tercio por el efecto bola de nieve favorable: con un crecimiento nominal medio en torno al $7\%$ y un tipo implícito que baja del $6-7\%\to4\%$, el diferencial es negativo durante casi todo el periodo. Los ajustes déficit-deuda operan en ambos sentidos: los ingresos de privatización reducen la deuda, mientras que la acumulación de activos y ciertas asunciones de deuda la aumentan. Por tanto, es la imagen especular de 1991-1995. Entonces el diferencial empujaba la deuda hacia arriba con déficits primarios moderados. Ahora la empuja hacia abajo aunque el esfuerzo primario sea limitado.

¿Quién generaba el superávit? La composición por subsectores de los superávits se suele olvidar, y es muy reveladora. En el cierre de 2007: la Administración Central registró un superávit del $1,29\%$ del PIB; la Seguridad Social, uno del $1,25\%$; las CCAA, un déficit del $0,17\%$; y las corporaciones locales, uno del $0,14\%$.

El superávit de la Seguridad Social era estructural en el sentido demográfico del momento: cohortes numerosas en edad de trabajar, una inmigración masiva que elevó la afiliación (más de $19$ millones de afiliados en 2007) y unas pensiones aún moderadas. Una parte importante de ese superávit se canalizaba al Fondo de Reserva, que alcanzó $45.716$ millones de euros a finales de 2007 y $57.223$ millones en 2008. El Fondo invertía mayoritariamente en deuda pública española. Como esa deuda está en manos de otra administración, desaparece en la deuda PDE consolidada. El Fondo reducía así la deuda consolidada y, cuando se dispuso de él en 2012-2017, esa deuda reapareció en manos del mercado. Por tanto, parte del superávit español de 2005-2007 era ahorro de la Seguridad Social para pensiones futuras, no un margen de maniobra de libre disposición. 

La entrada al euro, además, cambió la política de financiación. La vida media de la deuda del Estado pasó de unos $3$ años a mediados de los noventa a más de $6$ años en 2007. Desaparece la deuda en divisas, se redenomina toda la deuda en euros en 1999 y se consolidan los programas de creadores de mercado y de segregación de principal y cupones (\textit{strips}, 1997). El volumen de emisión se reduce y, en algunos años, la emisión neta de medio y largo plazo es negativa o casi nula. Como lo veremos más adelante, aquí basta retener que el Tesoro llega a 2008 con una deuda larga y barata, lo que amortiguó el impacto inicial de la crisis. Pero también llega con un mercado de deuda con poca liquidez y una base inversora cada vez más concentrada en no residentes (en torno al $50\%$ de la deuda negociable en 2007-2008), factor que pesará en 2010-2012.

\begin{specialblock}{\textbf{Material complementario:} ¿Cómo estaba \textbf{realmente} España durante este periodo?}
    Conviene entrar un poco, aunque no sea durante la exposición, a los determinantes del crecimiento económico español del periodo y sus consecuencias sobre los datos que acabamos de arrojar. Tratamos, por tanto, algunas cuestiones que son fundamentales para comprender como las finanzas españolas pudieron deteriorse tan rápidamente durante un lapso de tiempo tan corto.

    \textbf{Espejismo del superávit: ingresos transitorios y saldo estructural en tiempo real}. Lo primero que debería llamarnos la atención es: ¿cómo pasó España de un superávit de más del $2\%$ del PIB a un déficit del $11\%$, en solo dos años? Evidentemente, la caída de la actividad y los estabilizadores automáticos explican una parte. Pero también es evidente que no puede explicarlo todo.
    
    La magnitud del deterioro solo se entiende si se acepta que el saldo de 2007 incluía un componente transitorio muy superior al que estimaban los métodos convencionales. Es el llamado \textbf{falso superávit}.

    DE CASTRO, ESTRADA, HERNÁNDEZ DE COS y MARTÍ (2008), estimaron que el auge inmobiliario generó un aumento acumulado de ingresos tributarios de alrededor de $2$ puntos de PIB entre 1998 y 2006. En ese mismo periodo, la presión fiscal total aumentó unos $4$ puntos: la mitad del aumento de la presión fiscal estaba ligado a la vivienda. Incorporando los efectos riqueza, su estimación del saldo estructural de 2007 bajaba del $1,78\%$ al $1,02\%$ del PIB. Ya advertían de que la desaceleración inmobiliaria había restado $0,6$ puntos de ingresos por IVA e ITP-AJD solo en 2007.\footnote{%
        Las figuras más sensibles fueron: el ITP-AJD, un tributo cedido, lo que explica parte del deterioro autonómico posterior ([\authorfont{Ver Tema 4.B.25}]); el IVA de la vivienda nueva; el Impuesto sobre Sociedades de los sectores de la construcción y financiero; el IRPF de las ganancias patrimoniales; y, los impuestos locales vinculados al urbanismo: ICIO, plusvalía municipal y aprovechamientos urbanísticos.
    } MARTÍ y PÉREZ (2015) cifran en torno al $2,5\%$ del PIB el peso de los impuestos sobre la propiedad y las transmisiones en el pico del ciclo. Documentan además que el gasto primario, excluido el desempleo, creció en términos reales alrededor de un $4\%$ anual, por encima del crecimiento tendencial.
    
    Además, tenemos que tener en cuenta que este hecho permaneció oculto porque la metodología de la Unión para estimar el saldo estructural corrige el saldo por la brecha de producción, pero no por los precios de los activos ni por la composición del crecimiento. En España esto produjo dos errores que se sumaron:
    \begin{lnum}
        \item \textbf{Brecha de producción}. El primero fue la infravaloración de la brecha de producción. ALBEROLA, ESTRADA y SANTABÁRBARA (2014) muestran que, con la metodología de la función de producción de la Comisión, la brecha de 2007 estimada con datos hasta 2007 era ligeramente negativa ($-0,4\%$), mientras que con información hasta 2011 pasó a ser positiva ($+1,5\%$): una revisión de casi dos puntos. Su método alternativo de crecimiento sostenible, que incorpora los desequilibrios financieros y exteriores, estimaba una brecha del $5,4\%$ en 2007, y apenas se revisaba con información nueva. Con una semielasticidad presupuestaria de $0,5$ aproximadamente, esa diferencia de brecha supone entre $1$ y $3$ puntos de saldo estructural. El superávit estructural que se veía en tiempo real se convierte, ex post, en equilibrio o déficit estructural según el método;
        \item \textbf{Falacia de composición}. El segundo es la omisión del componente de composición (DE CASTRO et al., 2008). Aunque la brecha de producción se hubiera estimado bien, el saldo ajustado de ciclo no captura los ingresos ligados a un crecimiento intensivo en construcción y en transacciones de activos: un crecimiento \textbf{nominal}. Estos generan mucha más recaudación por unidad de PIB que el consumo o la exportación. 
    \end{lnum}
    
    En definitiva, la conclusión que comparte hoy la literatura, del BdE a la AIReF y la Comisión, es que el saldo estructural español de 2007 estaba lejos del superávit que mostraba el saldo total. Las estimaciones ex post lo sitúan entre el equilibrio y un déficit de $1-2$ puntos, según el método.

    \textbf{El debate: ¿prudencia o prociclicidad?}. Dicho esto, lo que conviene determinar es si, dado que en realidad se experimentaba un déficit, el Gobierno actuaba con prudencia o con prociclicidad. Este debate conviene presentarlo desde dos posiciones, porque ambas tienen argumentos.
    
    La posición defensiva, que es la ortodoxia oficial de la época y la de una parte de la literatura posterior sobre la crisis del euro, sostiene cuatro cosas. En primer lugar, que España cumplió holgadamente el PEC, a diferencia de Alemania y Francia. En segundo lugar, que redujo su deuda a la mitad. En tercer lugar, que acumuló un Fondo de Reserva para la Seguridad Social. Y, en cuarto lugar, que la crisis española de 2008-2012 fue una crisis de balance privado (endeudamiento de hogares y empresas, burbuja inmobiliaria, cajas de ahorros) que acabó trasladándose al sector público, no una crisis de indisciplina fiscal. Esta linea es la lectura de DE GRAUWE (2011) y de buena parte de la literatura sobre la crisis del euro que distingue entre Grecia (crisis fiscal) e Irlanda y España (crisis bancaria y privada). Evidentemente, tiene un sólido apoyo empírico: con una deuda del $36\%$ en 2007, España tenía un margen que no tenían Italia o Grecia, y ese margen permitió absorber el shock inicial.

    No obstante, la posición crítica (BdE, MARTÍ y PÉREZ, HERNÁNDEZ DE COS, y la OCDE y el FMI en sus revisiones posteriores) sostiene tres cosas. En primer lugar, que en una economía que crecía muy por encima de su potencial sostenible, con un déficit exterior cercano al $10\%$ del PIB y una inflación diferencial persistente, la política fiscal debió haber sido mucho más restrictiva. En segundo lugar, que el superávit era insuficiente para contrarrestar el recalentamiento de la economía. Y, en tercer lugar, que la política fiscal fue procíclica en su composición, pues los ingresos extraordinarios se usaron para financiar gasto primario permanente y rebajas fiscales permanentes (bonificaciones/gasto fiscal). Entre estas se cuentan las reformas del IRPF de 1998 (Ley 40/1998), 2002 (Ley 46/2002) y 2006 (Ley 35/2006), y la rebaja del tipo del Impuesto sobre Sociedades del $35\%$ al $30\%$ en 2007-2008. A ello le añaden la ausencia de política macroprudencial en un país sin política monetaria propia: la política fiscal era el único instrumento nacional de estabilización, y no se usó con esa función.

    Evidentemente, las dos posiciones son compatibles. El nivel de deuda de 2007 era efectivamente bajo y fue un activo en la crisis; y el saldo estructural y la estructura de ingresos eran mucho más débiles de lo que parecía. El problema, por tanto, no fue el punto de partida en deuda, sino el agujero estructural de ingresos que quedó cuando desaparecieron los transitorios: la presión fiscal cayó más de $6$ puntos de PIB.

    Además, existían una serie de pasivos (públicos) que se estaban acumulando fuera del saldo. Como ya hemos visto, en esos años se gestan varios pasivos que emergerán después como deuda o como déficit, y que figuran en el inventario de ajustes déficit-deuda: el déficit de tarifa eléctrica, que empieza a acumularse a principios de los 2000 y acabará titulizado a través del FADE; las facturas no contabilizadas de las CCAA y las entidades locales, sobre todo en sanidad, que aflorarán con el plan de pago a proveedores de 2012; la deuda de entidades del sector público empresarial clasificadas fuera de las AAPP, como ADIF, creada en 2005; y, el pasivo contingente del sistema bancario, concentrado en las cajas de ahorros, cuya materialización en 2012 (FROB, MEDE, Sareb) es el mayor ajuste déficit-deuda de la historia fiscal española.
\end{specialblock}

En definitiva, la lectura más extendida de este periodo (España llegó a 2008 con las cuentas saneadas) es, en el mejor de los casos, incompleta. Los superávits de 2005-2007 se sostenían sobre unos ingresos transitorios ligados al ciclo inmobiliario y sobre un ahorro de intereses que no dependía de España, mientras el gasto primario crecía por encima de la tendencia de la economía. Este periodo, 1996-2007, no fue una historia de disciplina fiscal: de haberlo sido, su resultado no se habría evaporado en $18$ meses.

\subsubsection{Gran recesión, la crisis soberana y la consolidación incompleta (2008-2019)}
Como todos sabemos, el siguiente periodo es el de la Gran Recesión, la crisis de deuda soberana y el agónico letargo de la consolidación fiscal en 2019. Es la década más intensa de la historia fiscal española desde la Guerra Civil. 

La década tiene tres fases con lógicas distintas. La primer, 2008-2009, el desplome. Desaparecen los ingresos transitorios de la fase anterior, actúan los estabilizadores automáticos y se aplica un estímulo discrecional que acabó siendo en parte permanente. La segunda, 2010-2013, la crisis de deuda soberana. El ajuste fiscal se hace bajo presión de los mercados y de Bruselas, y el sector bancario se reestructura con apoyo europeo. La tercera, 2014-2019, la recuperación. El déficit total se reduce, sobre todo por el ciclo y por los tipos de interés, pero el déficit estructural queda prácticamente sin corregir.

Como ya sabemos, transforma la posición fiscal española: del $36\%$ al $98\%$ del PIB de deuda, con un déficit estructural cercano al $3\%$ al final.

La lectura habitual de la década atribuye la mejora del saldo desde 2015 a las reformas estructurales y al ajuste de 2010-2013. Sin embargo, la mayor parte de la reducción del déficit entre 2014 y 2019 fue cíclica y financiera: crecimiento, empleo y caída de los intereses gracias al BCE. España cerró 2019, en la fase alta del ciclo y con la brecha de producción ya positiva, con un déficit del $3,1\%$ del PIB, prácticamente el mismo déficit estructural con el que había salido de la recesión. Por el camino, el saldo pasó de un superávit del $1,9\%$ del PIB en 2007 a un déficit del $11,2\%$ en 2009, la deuda se triplicó (del $36\%$ en 2007 al $105\%$ en 2014)\footnote{%
    Las cifras de la década incorporan la reclasificación retroactiva de la Sareb en las AAPP (2021) y las revisiones posteriores del PIB: el déficit de 2012, publicado en su día en el $10,6\%$ del PIB, es hoy del $11,5\%$, y la deuda de 2014, publicada en 2015 en el $97,7\%$, es hoy del $105,1\%$.
} y, por primera vez desde la entrada en el euro, España perdió en la práctica el acceso normal a los mercados: en el verano de 2012 la prima de riesgo superó los $600$ puntos básicos y el Gobierno solicitó asistencia financiera europea para el sistema bancario.


\paragraph{2008-2009: el desplome}
El desplome de 2008-2009 fue, a priori, una crisis de ingresos: desaparecieron los ingresos transitorios del auge inmobiliario, y el estímulo discrecional, en parte con medidas permanentes, se aplicó sobre un diagnóstico erróneo de problema cíclico.

Durante la época, el saldo pasa de un superávit del $1,9\%$ en 2007 a un déficit del $4,6\%$ en 2008 y del $11,2\%$ en 2009, que no bajará del $9,5\%$ hasta 2012. Por su parte, la deuda sube casi $70$ puntos, del $36\%$ en 2007 al $105\%$ en 2014.

En términos generales, este deterioro se puede atribuir a cuatro factores: los estabilizadores automáticos (con un paro que llegó a superar el $25\%$ y una caída de la recaudación de más de $6$ puntos de PIB, sin paralelo en el entorno), las políticas expansivas (Plan E), las ayudas al sector bancario y la crisis de deuda soberana.

Estos cuatro factores pertenecen, sin embargo, a fases distintas: las ayudas bancarias y la crisis soberana corresponden a 2010-2013, mientras que en 2008-2009 dominan los dos primeros. La tasa de paro de la EPA alcanzó su máximo en el primer trimestre de 2013, con un $26,9\%$.

La variación del saldo entre 2007 y 2009 es de unos $13$ puntos de PIB. Más de la mitad procede de los ingresos. La presión fiscal, incluidas las cotizaciones, cae más de $6$ puntos en dos años: del $41\%$ al $35\%$ del PIB. No es un simple efecto de los estabilizadores automáticos, porque el PIB nominal cayó mucho menos: es la reversión de los ingresos transitorios y responde al esquema de DE CASTRO et al. (2008): el ITP-AJD y el IVA de la vivienda se desploman con las transacciones; el Impuesto sobre Sociedades pierde las bases del sector inmobiliario y financiero, y las pérdidas de ejercicios anteriores alimentan créditos fiscales; y el IRPF pierde las ganancias patrimoniales.

MARTÍ y PÉREZ (2015) documentan que el saldo empeoró $6,4$ puntos solo en 2008. Por el lado del gasto actúan dos fuerzas. La primera es el aumento del desempleo: las prestaciones pasan de $1,5\%\to3\%$ del PIB. La segunda es la inercia del gasto primario comprometido en la fase alcista, que siguió creciendo en términos nominales en 2008-2009 mientras el PIB se contraía. El efecto denominador amplifica todas las ratios.

La respuesta fiscal española se inscribió en el Plan Europeo de Recuperación Económica (Comisión, noviembre de 2008), que recomendaba un estímulo coordinado de en torno al $1,5\%$ del PIB de la Unión. Las principales medidas fueron: la deducción de $400$ euros en el IRPF (RDL 2/2008); la bonificación del $100\%$ del Impuesto sobre el Patrimonio, que supuso su supresión efectiva (Ley 4/2008); el Fondo Estatal de Inversión Local, dotado con $8.000$ millones de euros (RDL 9/2008) y, su continuación, el Fondo Estatal para el Empleo y la Sostenibilidad Local, con $5.000$ millones (RDL 13/2009). A ello se sumaron ayudas sectoriales, como el Plan 2000E de automoción, y medidas de liquidez. Todo ello se conoce como Plan E.

El debate sobre este estímulo tiene dos niveles. El primero es su eficacia. Los fondos locales de inversión se criticaron por su baja calidad y por un multiplicador probablemente reducido: eran proyectos pequeños, intensivos en obra menor y sin continuidad. Su defensa se basa en que actuaron rápido en el sector más castigado por el paro, la construcción, cuando el multiplicador en recesión y con política monetaria restringida debía ser alto ([\authorfont{Ver Tema 3.A.40}])

El segundo nivel, más importante para este tema, es el error de diagnóstico. El estímulo se diseñó como si el deterioro fuese cíclico. Pero, como vimos, buena parte de él era estructural: los ingresos perdidos no iban a volver. Varias medidas del estímulo, como la deducción de $400$ euros o la supresión del Impuesto sobre el Patrimonio, eran rebajas fiscales permanentes. En la práctica, España utilizó en 2008-2009 un margen fiscal que creía tener por el bajo nivel de deuda, cuando en realidad tenía un déficit estructural que se revelaría como tal en cuanto se estimó bien la brecha de producción.\footnote{%
    Esta es la lectura del BdE y de la OCDE ex post. Una lectura alternativa sostiene que el estímulo fue razonable en el contexto europeo de 2008-2009, porque el error de diagnóstico fue generalizado y era difícil detectarlo en tiempo real.
}

\paragraph{2010-2013: la crisis soberana, el ajuste y el rescate bancario}
La crisis soberana de 2010-2013 impuso un ajuste concentrado y probablemente excesivo en su calendario, en un contexto de multiplicadores altos y sin prestamista de última instancia. Se combinó con el círculo vicioso entre bancos y soberano, que acabó con la asistencia del MEDE ($41.333$ millones) y con un coste fiscal de la reestructuración bancaria que el Tribunal de Cuentas estima en unos $72.000$ millones. En esos años, los tres motores de la deuda (déficit primario, intereses y ajustes déficit-deuda) sumaron en la misma dirección.

El punto de inflexión se produjo en mayo de 2010. El Consejo había abierto a España el procedimiento de déficit excesivo (PDE) en abril de 2009, con plazo de corrección en 2012. Tras el primer programa griego y la creación del FEEF (9 de mayo de 2010), el RDL 8/2010 aprobó una reducción media del $5\%$ de las retribuciones del sector público, la congelación de las pensiones contributivas en 2011 (salvo las mínimas), la supresión del cheque bebé y un recorte de la inversión pública. Era la primera vez que el ajuste se imponía desde fuera: la presión de los socios europeos y de los mercados obligó a abandonar el estímulo. El objetivo de situar el déficit por debajo del $3\%$ no se alcanzaría hasta 2018, y la propia cronología del PDE da cuenta de las dificultades: el plazo se amplió a 2013 (diciembre de 2009), a 2014 (julio de 2012) y a 2016 (junio de 2013); en julio de 2016 el Consejo declaró que España no había adoptado medidas eficaces, aunque en agosto de ese año anuló la multa; y el procedimiento no se cerró hasta junio de 2019, con los datos de 2018. Las sucesivas ampliaciones son en sí mismas un dato: el ajuste fue más lento de lo previsto porque la recesión se prolongó y los multiplicadores resultaron mayores de lo supuesto.

El ajuste combinó recortes de gasto (retribuciones públicas, prestaciones por desempleo, inversión pública, educación, sanidad y ayuda oficial al desarrollo) con subidas de impuestos. El IVA pasó del $16\%\to18\%$ en julio de 2010 y al $21\%$ en septiembre de 2012, y el tipo reducido del $7\%\to8\%$ y después al $10\%$. El RDL 20/2011 introdujo un recargo temporal en el IRPF y en el IBI; el RDL 13/2011 restableció temporalmente el Impuesto sobre el Patrimonio; se elevó el gravamen de la base del ahorro y se limitaron las deducciones del Impuesto sobre Sociedades; y el RDL 20/2012 suprimió la paga extraordinaria de diciembre de los empleados públicos y redujo la prestación por desempleo a partir del día $181$ (del $60\%\to50\%$ de la base). En pensiones, las Leyes 27/2011 y 23/2013 introdujeron, respectivamente, el retraso de la edad de jubilación y el índice de revalorización y el factor de sostenibilidad ([\authorfont{Ver Tema 4.B.24}]). Se recurrió también a ingresos extraordinarios, como la regularización fiscal del RDL 12/2012, que recaudó en torno a $1.200$ millones (aproximadamente la mitad de lo previsto) y que la STC 73/2017 declaró inconstitucional por haberse aprobado mediante decreto-ley en una materia que afecta al deber de contribuir (art. 31 CE), sin afectar a las situaciones ya consolidadas.\footnote{%
    La salida a bolsa del $49\%$ de AENA en febrero de 2015, que suele citarse junto a estas medidas, es una privatización y, como tal, redujo la deuda a través del ajuste déficit-deuda, no el déficit. Del mismo modo, la caída del límite de gasto no financiero del Estado de $182.439$ millones (2010) a $117.353$ millones (2012) no mide un recorte homogéneo: una parte sustancial se debe al modelo de financiación autonómica de la Ley 22/2009, que amplió la cesión de tributos a las CCAA (IRPF al $50\%$, IVA al $50\%$, impuestos especiales al $58\%$) y redujo en contrapartida las transferencias del Estado que computan en su techo ([\authorfont{Ver Tema 4.B.25}]).
}

¿Fue un ajuste autodestructivo o inevitable? Es el debate central de la década. La posición crítica se apoya en dos hallazgos. BLANCHARD y LEIGH (2013) mostraron que las previsiones de crecimiento de 2010-2011 infravaloraron sistemáticamente el impacto de la consolidación: los multiplicadores efectivos estaban por encima de $1$, frente al $0,5$ supuesto. DE GRAUWE y JI (2013) mostraron que las primas de riesgo de la periferia en 2010-2012 no se explicaban por los fundamentales fiscales, sino por el pánico del mercado en una unión monetaria sin prestamista de última instancia: con peores ratios de deuda y déficit, el Reino Unido pagaba mucho menos porque tenía su propio banco central. De ahí la tesis de la \textbf{austeridad autodestructiva}: en recesión, con multiplicadores altos y sin política monetaria propia, la consolidación reduce el PIB más que el déficit y la ratio de deuda puede aumentar, como de hecho siguió aumentando en España en 2012-2014.

La posición defensiva, que es la de la Comisión, la del BdE en su momento y la de una parte de la literatura, sostiene que la alternativa no era un ajuste más lento sino la pérdida completa del acceso al mercado. Con necesidades de financiación del Estado en torno al $20\%$ del PIB anual y una prima de riesgo por encima de $600$ puntos básicos, un ajuste gradual no era financiable, y su credibilidad fue además condición política para que el BCE anunciase el OMT. ALESINA, FAVERO y GIAVAZZI (2019) añaden que los ajustes por el lado del gasto son mucho menos recesivos que los basados en impuestos; en España predominaron las subidas de impuestos, lo que en su marco explicaría parte del coste. La síntesis hoy más aceptada es que el ajuste era inevitable en su dirección, pero probablemente excesivo en su concentración temporal, y que el giro del BCE, más que el propio ajuste, fue el factor decisivo para romper el círculo vicioso.\footnote{%
    Sobre la hipótesis de la consolidación expansiva (canal de expectativas y de primas de riesgo) frente a la austeridad autodestructiva (multiplicador alto con tipos en su límite inferior e histéresis), y la evidencia de BLANCHARD y LEIGH (2013) y de GUAJARDO, LEIGH y PESCATORI (2014) ([\authorfont{Ver Tema 3.A.40}]). Sobre el OMT y la política monetaria no convencional ([\authorfont{Ver Tema 3.A.36, Tema 3.A.37}]).
}

Ese círculo vicioso explica por qué 2012 fue el momento crítico. Las dudas sobre la solvencia de los bancos, en particular de las cajas de ahorros expuestas al sector inmobiliario, elevaban el riesgo del soberano, que era su garante implícito; y, a su vez, las dudas sobre el soberano deterioraban el valor de la deuda pública que los bancos tenían en balance. Las operaciones de financiación a tres años del BCE de diciembre de 2011 y febrero de 2012 reforzaron el vínculo, porque los bancos españoles usaron esa liquidez para comprar deuda del Estado: la proporción de deuda en manos de no residentes cayó del $50\%$ a en torno al $33\%$ a mediados de 2012, y la deuda se renacionalizó en los balances bancarios. En julio de 2012 la prima de riesgo alcanzó los $638$ puntos básicos, con el bono a diez años por encima del $7,5\%$, y Moody's había rebajado la calificación en nueve escalones en menos de dos años (de Aaa a Baa3 en junio de 2012). En junio de 2012 España solicitó asistencia del MEDE para recapitalizar su sistema bancario. El Memorando de Entendimiento de julio abrió una línea de hasta $100.000$ millones, de la que se utilizaron $41.333$, con condicionalidad sobre las entidades, condicionalidad horizontal sobre el conjunto del sistema financiero y recordatorios de compromisos fiscales y estructurales ya adquiridos. España devolvió anticipadamente $17.612$ millones en 2014-2018, y el resto vence hasta 2027. Este diagnóstico del círculo vicioso es el que llevó a la Unión Bancaria.\footnote{%
    Sobre el FEEF, el MEDE y la Unión Bancaria (MUS y MUR) ([\authorfont{Ver Tema 3.B.45}]).
}

Sobre el coste fiscal, el Tribunal de Cuentas cifra el de la reestructuración bancaria en $71.833$ millones (datos a 1 de enero de 2022), de los que $50.622$ millones corresponden al FROB y $21.273$ millones a las propias entidades a través del Fondo de Garantía de Depósitos, sin que el coste sea aún definitivo. En contabilidad nacional, las aportaciones de capital que no se esperaba recuperar se registraron como transferencias de capital, con impacto en el déficit (en torno a $3,7$ puntos de PIB en 2012), y el resto como activos financieros, con impacto solo en la deuda.

A ello se sumaron los mecanismos de liquidez para las administraciones territoriales: el Fondo para la Financiación de los Pagos a Proveedores (RDL 4/2012 y RDL 7/2012) y el Fondo de Liquidez Autonómico (RDL 21/2012), después integrados en el Fondo de Financiación a CCAA (RDL 17/2014) ([\authorfont{Ver Tema 4.B.25}]). Su efecto sobre la deuda consolidada fue claro: el plan de pago a proveedores convirtió deuda comercial, que no computa como deuda PDE, en deuda financiera del Estado, que sí computa, y aumentó la deuda de las AAPP sin aumentar el déficit del año, salvo en la parte de las facturas no registradas en su momento. Ese afloramiento de facturas confirma que parte del déficit de los años anteriores no se había registrado. En paralelo se reformó el marco institucional con la reforma del artículo 135 de la Constitución (septiembre de 2011), la LOEPSF (2012) y la creación de la AIReF (2013).\footnote{%
    Sobre las reglas de la LOEPSF, su relación con el TECG y el \textit{six-pack}, y su suspensión en 2020-2023 ([\authorfont{Ver Tema 3.B.45}]); sobre su articulación con la LGP y el ciclo presupuestario ([\authorfont{Ver Tema 4.B.21}]).
}

La mecánica de la deuda de 2008-2014 resume el periodo. Según el BdE, del aumento acumulado de la ratio, unos $37$ puntos de PIB se debieron al déficit primario; unos $13$ puntos, a la carga de intereses, apenas compensados por un crecimiento nominal prácticamente nulo; y unos $6$ puntos netos, a los ajustes déficit-deuda (la parte financiera de las ayudas bancarias, la contribución a los programas europeos, el capital desembolsado en el MEDE y la acumulación de depósitos del Tesoro, con reversiones parciales). A diferencia de 1991-1995, dominado por el diferencial, y de 1975-1985, dominado por el déficit primario con un diferencial favorable, aquí los tres componentes sumaron en la misma dirección: es el único episodio de la serie en que eso ocurre.

\paragraph{2014-2019: la recuperación y la consolidación incompleta}
La recuperación de 2014-2019 redujo el déficit total gracias al ciclo y al BCE, no gracias a un esfuerzo estructural. La fragmentación política y las decisiones de gasto e ingreso de la fase alcista volvieron a hacer procíclica la política fiscal en la expansión, igual que en 1986-1990 y en 2005-2007. España llega a 2020 con una deuda del $98\%$ del PIB, un déficit estructural del orden del $3\%$ y una dependencia creciente del Eurosistema como tenedor de su deuda. Ese es el margen, muy estrecho, con el que afrontará la pandemia.

La lectura habitual atribuye a las reformas estructurales (la de pensiones de 2011, la laboral de 2012, la Ley de Garantía de la Unidad de Mercado y la del sector eléctrico, ambas de 2013) el retorno al crecimiento desde 2013 y la reducción del déficit, que se situó por debajo del $3\%$ en 2018 ($2,6\%$) tras más de diez años por encima, mientras la deuda bajaba del $105,1\%$ en 2014 al $98,2\%$ en 2019. Los datos son correctos, pero la atribución causal exige matices. ¿Qué redujo realmente el déficit? Entre 2014 y 2019 el déficit total bajó unos $3$ puntos de PIB, del $6\%\to3\%$, por tres vías. La primera fue el \textbf{ciclo}: la economía creció por encima del $2,5\%$ anual en 2015-2018, se crearon más de $2,5$ millones de empleos y la brecha de producción pasó de muy negativa a positiva hacia 2018-2019; la mejora cíclica explica la mayor parte de la reducción. La segunda fueron los \textbf{intereses}: tras el máximo de 2013-2014, en torno al $3,5\%$ del PIB, la carga de intereses bajó al $2,3\%$ en 2019 gracias a la política monetaria, pues el OMT eliminó el riesgo de redenominación y el programa de compras de deuda pública del BCE, desde marzo de 2015, llevó los tipos a mínimos; a finales de 2019 el Eurosistema tenía el $18,7\%$ de la deuda pública española.\footnote{%
    Sobre el programa de compras del BCE (APP/PSPP y PEPP), la regla de reparto por la clave de capital y sus efectos sobre las primas ([\authorfont{Ver Tema 3.A.36, Tema 3.A.37}]).
} La tercera, el \textbf{esfuerzo estructural primario}, fue prácticamente nula o negativa desde 2015: hubo rebajas del IRPF en 2015-2016 y 2018, revalorizaciones de las pensiones por encima del índice de revalorización en 2018-2019 y subidas salariales públicas, de modo que la Comisión y la AIReF sitúan el déficit estructural de 2019 en torno al $3\%$ del PIB, un nivel similar al de 2014-2015.

La conclusión es incómoda: el ajuste estructural de la década se concentró en 2010-2013, en plena recesión y cuando más costaba, y en la recuperación no se aprovechó para completarlo. La política fiscal volvió a ser procíclica en la fase expansiva, como en 1986-1990 y en 2005-2007, según el diagnóstico compartido por el FMI (consultas del artículo IV de 2018 y 2019), la Comisión (Semestre Europeo) y la AIReF.

¿Por qué no se completó el ajuste? La explicación más directa es política. Entre 2015 y 2019 España atravesó el periodo de mayor fragmentación parlamentaria de su democracia: dos elecciones generales en seis meses (2015-2016) y otras dos en 2019, una moción de censura (2018) y presupuestos prorrogados (los de 2018 rigieron en 2019 y 2020). Aquí la hipótesis de ROUBINI y SACHS (1989), que no encajaba con las mayorías absolutas de los ochenta, sí encuentra respaldo: la fragmentación dificulta el ajuste y favorece las concesiones de gasto a los socios de investidura. El caso más claro es la revalorización de las pensiones con el IPC en 2018-2019, que abandonó el índice de revalorización de 2013 y marcó el principio del fin del factor de sostenibilidad ([\authorfont{Ver Tema 4.B.24}]).

La mecánica de la deuda confirma que el crecimiento nominal fue el protagonista. Según el BdE, la reducción de la ratio en 2019 ($2,1$ puntos) se descompone en un déficit total que la aumentó en $2,8$ puntos ($0,5$ de déficit primario y $2,3$ de intereses), un crecimiento nominal que la redujo en $3,4$ puntos y unos ajustes déficit-deuda que la redujeron en $1,6$ puntos. Es el patrón de todo 2015-2019: la deuda bajó muy despacio, unos $7$ puntos en cinco años, y lo hizo porque el crecimiento nominal compensaba un déficit que todavía la aumentaba. Con un saldo primario casi equilibrado y un diferencial negativo entre tipo de interés y crecimiento la deuda baja, pero a ese ritmo habrían hecho falta décadas para volver al $60\%$.

Por último, el Tesoro llegó a 2020 con una deuda más larga y un inversor distinto. A finales de 2019 la vida media de la deuda era de unos $7,5$ años, el $93,5\%$ era a largo plazo y los no residentes tenían el $49\%$ (frente al $33\%$ de mediados de 2012). El dato relevante es que un inversor nuevo, el Eurosistema, se había convertido en el tenedor individual más importante, una dependencia cuyas consecuencias analizaremos al estudiar la política de financiación del Tesoro.


\subsubsection{Pandemia, inflación y nuevo marco fiscal (2020-2026)}
El último periodo que nos queda por recorrer son seis años en los que la hacienda española sufre tres shocks de naturaleza distinta. El primero, en 2020, es un shock de oferta y de demanda sin precedentes en tiempos de paz, que hunde el PIB, dispara el déficit por encima del $10\%$ y lleva la deuda a su máximo histórico. El segundo, en 2022-2023, es un shock inflacionario de origen energético. Cambia por completo la aritmética de la deuda: primero la alivia a través del denominador y después la encarece a través de los tipos de interés. El tercero es un conjunto de shocks que se superponen en 2024-2026: la DANA de octubre de 2024, las borrascas de principios de 2026 y la guerra en Oriente Medio. 

El relato habitual del periodo es que, tras la pandemia, España ha reducido su deuda en casi $20$ puntos. Como veremos, es cierto en la cifra y engañoso en la causa. Según la descomposición del BdE, entre 2021 y 2024 el déficit (primario más intereses) añadió unos $18$ puntos de PIB a la ratio de deuda y el crecimiento del PIB nominal restó unos $37$. La deuda bajó por el denominador, con inflación y crecimiento, no porque las administraciones generasen excedentes. No obstante, la hacienda española sale de este periodo con un saldo primario cercano al equilibrio y una deuda ligeramente por encima del $100\%$ del PIB: una posición claramente mejor que en 2020, pero gracias a un entorno de crecimiento nominal que no se va a repetir, y no ha corregido el componente estructural del déficit que arrastra desde 2008. 

En definitiva, que ambas lecturas son compatibles, tanto la positiva como la negativa. La posición de partida es mejor que en 2020 en el saldo primario y en la estructura de la deuda. Pero la sostenibilidad de la trayectoria depende de nuevo de que el diferencial entre el tipo de interés y el crecimiento siga siendo favorable, algo que la experiencia española de 1991-1995 y 2010-2013 enseña que no puede darse por descontado.

Además, 2026 es el año en que las condiciones que hicieron posible esa reducción empiezan a revertirse. El conflicto de Oriente Medio iniciado el 28 de febrero de 2026 ha devuelto la inflación al $4,3\%$ en agosto y ha llevado al BCE a subir los tipos en junio y septiembre, hasta el $2,50\%$ en la facilidad de depósito. La AIReF, por su parte, prevé un crecimiento del gasto neto del $6,4\%$ en 2026, frente al $3,5\%$ comprometido con Bruselas.




\paragraph{2020-2021: la pandemia y la respuesta fiscal}
El primer periodo que veremos será la pandemia. En 2020 el déficit volvió a situarse por encima del $10\%$ del PIB y la deuda superó el $120\%$: un máximo histórico.

La revisión de la contabilidad nacional publicada por el INE en septiembre de 2024 elevó el nivel del PIB nominal, y con ella la caída real de 2020 queda en torno al $10,9\%$. El BdE sitúa hoy la ratio de deuda de 2020 en el $119,3\%$, no por encima del $120\%$. El déficit de 2020 se publicó en marzo de 2021 en el $11\%$ del PIB, incluida la reclasificación de la Sareb, y con la serie vigente queda cerca del $10\%$.\footnote{%
    La diferencia entre las cifras de hoy y las que se publicaron en su momento es otro ejemplo de los desajustes: las cifras fiscales dependen del perímetro, en este caso la Sareb, y del denominador, en este caso la revisión del PIB.
} La lectura con la descomposición del BdE es especialmente clara para 2020. La ratio de deuda aumentó $21,6$ puntos: el déficit primario aportó $7,7$ puntos, los intereses $2,2$, los ajustes déficit-deuda $0,9$ y la caída del PIB nominal $10,8$. Por primera vez en la serie histórica, el efecto denominador explica prácticamente la mitad del aumento de la ratio en un solo año: la imagen inversa de lo que ocurriría en 2021-2024.

¿A qué se deben estas variaciones? Evidentemente a la respuesta a la Pandemia, que fue clara: mucha liquidez y un gasto directo relativamente contenido. Según la actualización del Programa de Estabilidad 2021-2024, se movilizaron unos $231.000$ millones de euros en 2020-2021, el $20,6\%$ del PIB de 2020. Algo más de $73.000$ millones fueron ayudas directas y unos $158.000$ millones, financiación movilizada mediante avales, garantías y moratorias.

La flexibilización de los expedientes de regulación temporal de empleo exoneró de cotizaciones a las empresas a cambio de mantener el empleo y se prolongó, con modificaciones, hasta marzo de 2022. Se crearon dos fondos de recapitalización: uno gestionado por la SEPI para empresas estratégicas, con $10.000$ millones de los que se utilizaron $3.255,8$, y otro gestionado por COFIDES para medianas empresas, con $1.000$ millones de los que se utilizaron $779$. Hubo además un fondo de ayudas directas de $7.000$ millones, gestionado por las CCAA, del que se utilizaron algo más de $6.000$, y se aprobó el Ingreso Mínimo Vital. Entre las medidas sin impacto inmediato en el gasto destacan las líneas de avales, por unos $143.000$ millones y en su mayoría gestionadas por el ICO, las moratorias de créditos y los aplazamientos tributarios. Según el ICO, a junio de 2022 se habían desplegado avales por $107.187$ millones, que movilizaron $140.737$ millones de financiación en cerca de $1,2$ millones de operaciones.

En definitiva, la respuesta fiscal española fue desproporcionadamente de liquidez y comparativamente contenida en gasto directo. El Fiscal Monitor del FMI de 2021 situaba el apoyo presupuestario directo (por encima de la línea) de España claramente por debajo del de Alemania, el Reino Unido o Estados Unidos, y el apoyo de liquidez (avales, préstamos, capital) en niveles comparables o superiores a los de sus pares. La razón es evidente, España entró en la pandemia con una deuda cercana al $98\%$ y un déficit estructural del $3\%$, y su espacio fiscal era menor. Por eso optó por instrumentos que no computan como déficit en el momento, sino como pasivos contingentes.\footnote{
Como sabemos ([\authorfont{Ver Tema 4.B.23}]), los avales no generan gasto salvo que se ejecuten, aunque se dotó una provisión por pérdida esperada con impacto en el déficit de 2020. La experiencia posterior, con las ejecuciones de avales y el marco de reestructuración de créditos avalados del RDL 5/2021, muestra que la pérdida efectiva fue moderada en relación con el volumen avalado. Esto avala \textit{ex post} el diseño.
}

No obstante, la diferencia decisiva frente a 2010-2012 fue la respuesta europea. En marzo de 2020 se activó la cláusula general de salvaguardia del Pacto de Estabilidad, que suspendía en la práctica los requisitos de ajuste y se mantuvo hasta finales de 2023. En España, el Congreso apreció la existencia de una situación de emergencia extraordinaria, lo que suspendía a su vez las reglas de la LOEPSF. El instrumento SURE financió los ERTE y mecanismos similares. España fue uno de sus principales beneficiarios, con unos $21.300$ millones en préstamos. En julio de 2020 se acordó el instrumento Next Generation EU.

Por otro lado, y en términos financieros, la pieza más importante fue la del BCE. Cuando la presidenta del BCE declaró el 12 de marzo de 2020 que su función no era cerrar diferenciales, la rentabilidad del bono español a diez años subió casi $100$ puntos básicos en una semana. El anuncio del programa de compras de emergencia frente a la pandemia (PEPP) el 18 de marzo y la flexibilización del colateral en abril devolvieron la rentabilidad por debajo de su nivel previo a la pandemia en julio de 2020. Lo que refleja la absoluta dependencia de los tipos de la deuda respecto de las decisiones de política monetaria. Lección que la crisis de 2010-2012 había enseñado a un coste muy alto y que en 2020 se aplicó desde el primer momento.

Dicho en los términos de DE GRAUWE, en 2020 los países de la zona del euro sí tuvieron un prestamista de última instancia de facto. Sin él, la respuesta fiscal española habría sido mucho más limitada. Es también una advertencia para el presente: la dependencia se ha vuelto a manifestar, en sentido contrario, desde que el BCE dejó de comprar deuda.

¿Había que haber hecho más? Hubo un debate genuino sobre si la respuesta española fue suficiente. Una posición, sostenida por parte de la literatura y por organismos como el FMI en 2020, defendía que en una crisis de esta naturaleza el coste de hacer demasiado poco (quiebras de empresas viables, histéresis en el empleo) era mayor que el de hacer demasiado, y que España debió aprobar más ayudas directas a empresas y antes, sobre todo en el turismo y la hostelería: el fondo de ayudas directas de $7.000$ millones no llegó hasta marzo de 2021 (RDL 5/2021), un año después del inicio de la crisis.

La posición contraria recordaba que la recuperación del empleo fue relativamente rápida gracias a los ERTE, que el tejido productivo no sufrió la ola de quiebras temida y que, dada la posición de partida, un gasto directo mayor habría elevado todavía más la deuda. El balance que hoy parece más razonable es que el diseño de los ERTE fue un acierto, porque evitó la destrucción masiva de empleo de 2008-2013, mientras que el apoyo directo a empresas llegó tarde y se ejecutó por debajo de lo dotado, como muestra la infrautilización de los fondos de la SEPI y de COFIDES.

\paragraph{2022-2024: inflación, Plan de Recuperación y reducción de la deuda por el denominador}
El segundo periodo viene caracterizado por la inflación. Sin embargo, antes de entrar a valorar su impacto, veamos las principales medidas de política fiscal y monetaria que se adoptaron.

Las medidas que se adoptaron como respuesta a la crisis energética de 2022-2023 fueron diversas, y su coste conjunto en 2022-2023 se sitúa en torno a los $45.000$ millones de euros según distintas estimaciones.\footnote{%
    Incluyó la bonificación de $20$ céntimos por litro de combustible en 2022, las rebajas del IVA y del impuesto especial de la electricidad y del gas, las rebajas del IVA de los alimentos básicos, las ayudas al transporte público y el llamado mecanismo ibérico de limitación del precio del gas para la generación eléctrica, este último sin coste presupuestario directo.

La literatura, con la OCDE, el FMI y el BdE a la cabeza, criticó que la mayoría de estas medidas no estaban focalizadas: beneficiaban a todos los consumidores, con independencia de su renta, y mantenían la señal de precios de la energía en lugar de corregirla. Una transferencia de renta a los hogares vulnerables habría tenido menor coste y habría preservado el incentivo al ahorro energético. La defensa de las medidas generales se basaba en la rapidez de aplicación y en su efecto directo sobre el IPC, que a su vez moderaba la indexación de pensiones y salarios.
} Además, por el lado de los ingresos, la Ley 38/2022 creó gravámenes temporales sobre las empresas energéticas y las entidades de crédito, así como el Impuesto Temporal de Solidaridad de las Grandes Fortunas. El gravamen bancario se transformó a finales de 2024 en un impuesto sobre el margen de intereses y comisiones de las entidades financieras (Ley 7/2024), mientras que el energético decayó. 

Esa recaudación merece ser comentada. Los ingresos tributarios del Estado superaron los $325.000$ millones de euros en 2025, con un crecimiento del $10,4\%$. La presión fiscal alcanzó máximos históricos, en torno al $37-38\%$ del PIB en 2024-2025. Es la contracara del efecto denominador: la inflación y la progresividad en frío del IRPF, que no se ha deflactado con carácter general en el ámbito estatal, elevaron la recaudación por encima del crecimiento de las bases. Una parte de este aumento es estructural (el empleo récord, más de $21$ millones de afiliados, y los cambios normativos), pero otra parte es transitoria o depende del nivel de inflación. Estimar cuál es cuál es de nuevo el problema central para leer el saldo estructural.

Por otro lado, el Mecanismo de Recuperación y Resiliencia, núcleo de Next Generation EU, asignó a España el mayor volumen de subvenciones tras Italia. Tras la adenda de 2023, el Plan de Recuperación español alcanzaba unos $163.000$ millones de euros, aproximadamente la mitad en transferencias no reembolsables y la otra mitad en préstamos. El Gobierno solicitó finalmente solo una fracción de los préstamos disponibles, en torno a $21.500$ millones, al considerar que el coste de financiarse directamente en el mercado era similar o inferior. A 31 de agosto de 2026, España había recibido unos $78.000$ millones en seis desembolsos. El sexto, de $6.234$ millones, llegó el 11 de agosto de 2026.\footnote{%
    Quedaban pendientes unos $25.862$ millones del último tramo, condicionados al cumplimiento de 148 hitos y objetivos, con plazo para solicitar el desembolso hasta el 30 de septiembre de 2026.
} Por lo que a nuestro tema respecta, las subvenciones no generan deuda española. La deuda la emite la Comisión en nombre de la Unión y se amortizará con cargo al presupuesto europeo a partir de 2028. En contabilidad nacional, las subvenciones se registran como ingreso cuando se realiza el gasto que financian, de modo que su efecto sobre el déficit español es en principio neutro: financian gasto adicional sin aumentar el déficit. Los préstamos, en cambio, computan como deuda de las AAPP españolas igual que cualquier otra financiación del Tesoro.

Esto tiene dos implicaciones. La primera es que el Plan permitió mantener niveles de inversión pública que de otro modo habrían exigido más déficit. La segunda es que el fin del Plan en 2026 plantea el llamado precipicio fiscal: si la inversión financiada con fondos europeos debe mantenerse, tendrá que pasar a financiarse con recursos propios, con impacto en el déficit a partir de 2027.

Por otro lado, entre julio de 2022 y septiembre de 2023 el BCE elevó la facilidad de depósito del $-0,5\%$ al $4\%$, el ciclo de subidas más rápido de su historia. Dejó además de reinvertir los vencimientos del programa de compra de activos (APP) desde julio de 2023 y del PEPP desde enero de 2025. Para el Tesoro, el efecto fue doble: el coste de las nuevas emisiones subió de forma abrupta, y el Eurosistema, que había sido el principal comprador neto de deuda española en 2015-2022, pasó a reducir sus tenencias. El Tesoro tuvo que colocar en el mercado no solo sus necesidades netas, sino también los vencimientos que antes absorbía el banco central: la ratio de no residentes subió al $44,9\%$ en 2024. La prima de riesgo española se mantuvo en niveles moderados, e incluso por debajo de la francesa en algunos momentos de 2024-2025, lo que sugiere que los mercados valoraban positivamente la evolución del saldo primario y del crecimiento.\footnote{%
    El rápido aumento de los tipos desde 2022 puso en entredicho la hipótesis, cada vez más aceptada antes de la pandemia, de que cambios estructurales en las economías avanzadas (demografía, exceso de ahorro, caída de la productividad) mantendrían los tipos reales bajos o negativos de forma persistente, lo que haría sostenibles niveles de deuda más altos. Sobre esta hipótesis, véase GALESI, NUÑO y THOMAS (2017).

La formulación más influyente de esa tesis es la de BLANCHARD (2019): con $r < g$, la deuda pública puede tener un coste fiscal nulo y un coste en bienestar menor del supuesto. El propio BLANCHARD advertía de que la condición podía invertirse, y la subida de tipos de 2022-2023 la puso a prueba. La evidencia del periodo es matizada. Los tipos nominales subieron, pero la inflación subió más, de modo que el diferencial entre tipo de interés y crecimiento siguió siendo favorable en España durante 2022-2024. Lo que sí se ha demostrado es que la condición puede cambiar rápidamente y que depende de la política monetaria. El debate sigue abierto sobre si el tipo de interés natural ha vuelto a los niveles previos a 2008 o si se mantiene bajo, con una inflación transitoria que ha alterado temporalmente los tipos observados.
}

Con todo lo expuesto, sorprenderá la evolución de la deuda. La ratio de deuda pasa del $119,3\%$ en 2020 al $115,7\%$ en 2021, al $109,5\%$ en 2022, al $105,1\%$ en 2023 y al $101,8\%$ en 2024. En esos cuatro años se redujo en unos $17,5$ puntos. La suma de los déficits primarios aportó unos $8,6$ puntos al alza ($4,5$ en 2021, $2,3$ en 2022, $1,1$ en 2023 y $0,7$ en 2024) y los intereses, unos $9,2$ (entre $2,1$ y $2,4$ puntos cada año). Los ajustes déficit-deuda añadieron en conjunto algo menos de $2$ puntos. ¿Cómo es posible, entonces, que experimentáramos semejante descenso?

Como es de esperar, el crecimiento del PIB nominal restó unos $37$ puntos: $10,3$ en 2021, $11,6$ en 2022, $9,1$ en 2023 y $6,2$ en 2024. La conclusión es inequívoca. Sin el crecimiento nominal, la deuda habría aumentado unos $20$ puntos en ese periodo, porque las administraciones siguieron teniendo déficit todos los años. La mejora del saldo primario es real, de un déficit del $7,7\%$ a uno del $0,7\%$ en cuatro años. Pero un déficit primario casi nulo solo estabiliza la deuda si el crecimiento nominal supera al tipo de interés, que es precisamente la condición que empieza a debilitarse en 2026.

Ahora bien, la inflación afecta a las finanzas públicas de forma asimétrica en el tiempo. A corto plazo, la inflación favorece las cuentas públicas por tres vías. La primera es que los ingresos se ajustan automáticamente a los precios: el IVA recauda sobre bases nominales y el IRPF, que no se deflacta, sufre la llamada progresividad en frío, porque los contribuyentes saltan de tramo sin haber ganado poder adquisitivo. La segunda es que buena parte del gasto se ajusta con retraso. La tercera es que el tipo de interés implícito de la deuda apenas se mueve, porque la vida media es larga (unos $8$ años). El resultado es que la ratio de deuda baja con fuerza por el denominador mientras la carga de intereses sigue siendo baja.

No obstante, a medio plazo el efecto se invierte cuando el gasto indexado alcanza a la inflación: las pensiones se revalorizaron un $8,5\%$ en 2023 por aplicación de la Ley 21/2021, y los salarios públicos y la contratación pública se actualizan. Al mismo tiempo, la deuda se va refinanciando a los nuevos tipos. Por tanto, la reducción de la deuda se sustentaba principalmente en el crecimiento del PIB nominal, con una contribución muy notable del deflactor. Una vez agotado ese impulso, las proyecciones de la AIReF dibujaban una dinámica desfavorable a políticas constantes, agravada por la subida de tipos. 

\paragraph{2025-2026: la situación actual}
Como anticipamos, España cerró 2025 con un déficit del $2,2\%$ del PIB sin los gastos extraordinarios de la DANA de Valencia, o del $2,4\%$ incluyéndolos, y con una deuda del $100,7\%$, la ratio más baja desde el inicio de la pandemia. 

El Estado pagó unos $34.900$ millones en intereses y el conjunto de las AAPP algo más de $40.000$, en torno al $2,4\%$ del PIB. El saldo primario estuvo, por tanto, muy cerca del equilibrio, la primera vez desde 2007 que se da esta situación.

El año 2025 es también el primero de aplicación del nuevo marco de gobernanza económica de la Unión (Reglamento (UE) 2024/1263), que sustituye los objetivos de saldo estructural por una senda plurianual de gasto primario neto negociada con cada Estado miembro.\footnote{%
    Contenido del nuevo marco de gobernanza económica (Reglamento (UE) 2024/1263 y Reglamento (UE) 2024/1264 de la vertiente correctiva, con el análisis de sostenibilidad de la deuda, la trayectoria técnica, la cuenta de control y las cláusulas de salvaguardia general y nacional) ([\authorfont{Ver Tema 3.B.45}]).
}El Plan Fiscal y Estructural de Medio Plazo de España, aprobado por el Consejo el 21 de enero de 2025, fija un periodo de ajuste de siete años (2025-2031), ampliado en contrapartida por compromisos de reformas e inversiones. Establece un crecimiento máximo del gasto neto del $3,7\%$ en 2025, del $3,5\%$ en 2026, del $3,2\%$ en 2027 y del $3\%$ en 2028. Según el Plan, la deuda bajaría del $102,5\%$ del PIB en 2024 al $90,6\%$ en 2031 y al $76,8\%$ en 2041.

La traslación al ordenamiento interno está pendiente. La LOEPSF no se ha adaptado al nuevo marco, y la regla de gasto nacional, el objetivo de estabilidad estructural y el límite de deuda de la ley orgánica siguen formalmente vigentes pero conviven con una regla europea de naturaleza distinta. A ello se añade una anomalía institucional sin precedentes en la democracia: España funciona desde 2024 con los Presupuestos Generales del Estado de 2023 prorrogados, porque ni en 2024, ni en 2025, ni en 2026 se aprobaron unos nuevos. Las sendas de estabilidad presentadas en 2024 y 2025 no obtuvieron el respaldo parlamentario necesario. La prórroga presupuestaria limita el margen de actuación discrecional, lo que paradójicamente ha contribuido a contener el gasto en algunos capítulos. Pero impide a la vez aplicar una estrategia fiscal de medio plazo y obliga a canalizar las nuevas medidas por decreto-ley y créditos extraordinarios.

El retorno de los shocks fue 2026. El 28 de febrero de 2026 comenzó la operación militar conjunta de Estados Unidos e Israel contra el régimen iraní, seguida de la respuesta de Irán y de un bloqueo del estrecho de Ormuz que disparó los precios del petróleo y del gas. El Gobierno aprobó el 20 de marzo de 2026 un Plan de Respuesta que movilizaba más de $5.000$ millones de euros, y la inflación española alcanzó el $4,3\%$ en agosto de 2026, la tasa más alta en más de tres años.\footnote{%
    Incluía rebajas escalonadas del impuesto sobre hidrocarburos, descuentos en el bono social eléctrico, ayudas a los transportistas y al sector primario y la supresión progresiva del impuesto sobre la producción de energía eléctrica. El 29 de septiembre de 2026, el Real Decreto-ley 25/2026 prorrogó buena parte de estas medidas hasta el 31 de diciembre, incluida una rebaja de $20$ céntimos por litro en el impuesto sobre hidrocarburos en octubre, con reducciones condicionadas en noviembre y diciembre.
} 

No obstante, la consecuencia monetaria es la más relevante. El BCE, que había llevado la facilidad de depósito al $2\%$ en junio de 2025, la subió al $2,25\%$ el 11 de junio de 2026 (la primera subida en tres años) y al $2,50\%$ el 10 de septiembre de 2026. Justificó la decisión porque el conflicto seguía generando presiones inflacionistas y la inflación se mantendría claramente por encima del objetivo durante un periodo prolongado. Sus proyecciones sitúan la inflación en el $3\%$ en 2026 y en el $2,5\%$ en 2027. Además, las carteras del APP y del PEPP siguen reduciéndose al no reinvertirse los vencimientos. Para España esto significa que \textbf{el coste marginal de la deuda vuelve a subir justo cuando se refinancia una parte creciente de la deuda emitida en la etapa de tipos cero}.

En su informe de julio de 2026, la AIReF mantenía su previsión de déficit para 2026 en el $2,6\%$ del PIB. Precisaba que, sin las medidas de emergencia (DANA, borrascas y conflicto de Oriente Medio), el déficit se situaría en el $1,8\%$. Rebajaba su previsión de deuda al $99,3\%$ y elevaba la de crecimiento al $2,5\%$. Pero advertía de que el crecimiento del gasto neto previsto para 2026 subía del $5,8\%$ al $6,4\%$, muy por encima del $3,5\%$ comprometido en el Plan Fiscal, y reclamaba un ajuste del $0,6\%$ del PIB para no incumplir las reglas nacionales. Una parte de esa desviación puede quedar justificada por el carácter excepcional de las medidas de emergencia. La valoración final dependerá de cómo las trate la Comisión en el Semestre Europeo y de la cuenta de control del Plan.

A las presiones inmediatas se suman otras de medio plazo que se tratan en otros temas pero que condicionan la lectura de la situación actual. Una es el gasto en pensiones, cuyo crecimiento por el envejecimiento se analiza, junto con la cláusula de salvaguarda y el informe de la AIReF de mayo de 2026, en otro tema ([\authorfont{Ver Tema 4.B.24}]). Otra es el compromiso de gasto en defensa, que España elevó al $2\%$ del PIB en 2025 mediante un plan de más de $10.000$ millones de euros sin recurrir a la cláusula nacional de salvaguardia prevista en el marco europeo. La tercera es la condonación de deuda autonómica: el proyecto de ley que prevé que el Estado asuma $83.252$ millones de deuda de las CCAA, enviado al Congreso en diciembre de 2025, superó el debate de totalidad en julio de 2026 y reanudó su tramitación el 29 de septiembre.\footnote{%
    El contenido de la condonación, su distribución por comunidades y el análisis de la AIReF de junio de 2025 ([\authorfont{Ver Tema 4.B.25}]).
} Su efecto sobre la deuda consolidada de las AAPP es nulo, porque se trata de una transferencia entre subsectores. Pero sí altera la distribución de la deuda entre el Estado y las CCAA y los incentivos de estas.





\clearpage
\section{España en perspectiva comparada}
\puntosclave{
    
}

El primer bloque ha reconstruido la trayectoria del saldo y de la deuda de las AAPP españolas desde 1975. Cerrado el recorrido, conviene presentar las dos lecturas de la situación actual: ¿son hoy sanas las cuentas públicas españolas? La lectura favorable subraya tres cosas: la reducción de la deuda en unos $19$ puntos de PIB desde 2020; un saldo primario cercano al equilibrio, que no se veía desde 2007; y un crecimiento económico por encima de la media de la zona del euro, que ha permitido a España financiarse con una prima de riesgo moderada, en ocasiones inferior a la de Francia. A ello añade una estructura de deuda larga (unos $8$ años de vida media), que amortigua el impacto de la subida de tipos, y una base inversora diversificada. La lectura crítica subraya que el nivel de deuda sigue siendo uno de los más altos de la Unión. Añade que su reducción se ha debido casi por completo al crecimiento nominal, que el déficit estructural se mantiene en torno al $2,5$-$3\%$ del PIB según la Comisión y la AIReF, que el gasto neto crece muy por encima de lo comprometido y que España lleva tres años sin presupuestos aprobados. Sostiene que, con el BCE subiendo tipos, una nueva perturbación energética y presiones estructurales de gasto (pensiones, defensa, fin del Plan de Recuperación), la ventana favorable de 2021-2025 se ha cerrado sin haber completado la consolidación estructural. Es la misma crítica que se hizo a 1986-1990, a 2005-2007 y a 2015-2019: no se aprovechó la fase expansiva para crear margen.

Decantarse por una u otra valoración, es una cuestión que se reduce al sistema de creencias de cada individuo. Como sabemos (o deberíamos), las ciencias sociales parten de supuestos (implícitos u explícitos) sin los que no se podría construir un sistema de relaciones tan complejo como el que tenemos. A diferencia de las ciencias puras: dan argumentos, pero no ofrecen verdades.

Por ello puede ser interesante alejarse de la disyuntiva y optar por una vía más empírica, comparada. Al fin y al cabo, el título de deuda de un Gobierno es, por definición, el activo más seguro que se ofrece en el mercado nacional; por tanto, parece razonable que la situación de éste se extraiga por comparación por pares. 

Es frecuente que este ejercicio comparativo se realice sobre todo por niveles de déficit y deuda, acabando por definir un ranking que dice poco sobre la situación real del activo de España. Precisamente, la lección más importante que ofrece la comparación histórica es que el nivel de deuda pública es un mal predictor de la vulnerabilidad fiscal. En 2007, los dos países de la zona del euro con menos deuda eran Irlanda ($24\%$) y España ($36\%$), y fueron los que sufrieron los mayores deterioros en 2008-2012. 

Por ello, optaremos por restringir nuestro análisis y preguntarnos si la trayectoria presentada es peculiar de España o común a otras economías avanzadas, y en qué se diferencia de las de los países con los que se compara habitualmente: los grandes países de la zona del euro, las economías de la periferia europea y el conjunto de la OCDE.\footnote{
Hay que advertir de un problema de método que condiciona la lectura de todos los datos del bloque: la comparabilidad de las estadísticas fiscales entre países. Los países de la Unión publican el déficit y la deuda según el SEC 2010 y el Protocolo de Déficit Excesivo, con la deuda valorada por su nominal y consolidada dentro del sector AAPP. Las bases del FMI y de la OCDE, que son las que permiten comparar con Estados Unidos, Japón o Canadá, siguen el Manual de Estadísticas de Finanzas Públicas (GFSM 2014). Este manual incluye en la deuda bruta instrumentos que la deuda PDE excluye, como los créditos comerciales, y en algunos países los pasivos de los planes de pensiones de los empleados públicos. Por eso la deuda de Estados Unidos que publica el FMI es sensiblemente superior a la que resultaría con el criterio europeo.

A ello se suma la distinción entre deuda bruta y deuda neta de activos financieros. Para países con grandes fondos públicos, como Noruega, Japón o los nórdicos, la deuda neta puede ser muy inferior a la bruta, e incluso negativa. En el caso de Japón, la deuda bruta supera el $230\%$ del PIB, pero la neta ronda la mitad, porque el sector público japonés tiene muchos activos financieros, entre ellos los del fondo de pensiones público y las tenencias cruzadas con el banco central. En España la diferencia es moderada, porque las AAPP tienen relativamente pocos activos financieros líquidos. Estas cautelas obligan a comparar siempre series homogéneas de la misma fuente y a no mezclar datos de Eurostat con datos del FMI en una misma frase.
}

Como veremos, España llega tarde al ciclo de expansión del sector público de las economías avanzadas y lo recorre comprimido en el tiempo. Por eso su deuda es inferior a la media hasta 2010, pero su déficit estructural es persistentemente mayor. Las comparaciones muestran también que la vulnerabilidad fiscal de un país depende menos de su nivel de deuda que de la estructura de sus ingresos, de sus desequilibrios privados y exteriores, de sus instituciones presupuestarias y de su pertenencia a una unión monetaria sin prestamista de última instancia. En todas esas dimensiones España ha sido, hasta fechas recientes, más vulnerable de lo que sugería su nivel de deuda.

\subsection{Trayectoria histórica comparada (1970-2007)}
Empecemos por comparar la trayectoria española con la de las economías avanzadas en las cuatro décadas que van de las crisis del petróleo a la víspera de la Gran Recesión. 

La comparación persigue tres objetivos. El primero es situar el punto de partida español, cuya anomalía (equilibrio con un sector público pequeño) solo se aprecia bien frente a los demás. El segundo es identificar los mecanismos que explican por qué, entre 1975 y 1995, algunos países duplicaron o triplicaron su deuda y otros la contuvieron, lo que permite leer el caso español a la luz de la literatura comparada de economía política. El tercero es mostrar cómo la convergencia de Maastricht y los primeros años del euro produjeron la paradoja con la que se abre 2008: los países con menos deuda no eran los menos vulnerables.

Como decíamos, la comparación muestra a España como un caso de llegada tardía y trayectoria comprimida. En 1975 tenía una de las deudas más bajas de las economías avanzadas, pero con un sector público veinte puntos de PIB menor que la media comunitaria, lo que hacía de su equilibrio un síntoma de atraso y no de virtud. En los ochenta, mientras la subida de los tipos reales provocaba la gran divergencia de las deudas (Bélgica, Italia, Irlanda, Grecia), España multiplicaba su ratio por dos veces y media, pero desde niveles tan bajos que en 1990 seguía por debajo de Alemania o Francia. Su déficit, sin embargo, era ya superior a la media de la OCDE. La crisis de principios de los noventa tuvo epicentros más graves que España (los nórdicos, Italia, Bélgica, Canadá y, a partir de entonces, Japón). Las consolidaciones más duraderas que siguieron, como las de Suecia, Canadá y Bélgica, combinaron ajuste y reforma institucional, sobre todo reglas de gasto. Maastricht ancló la convergencia de los déficits, con una contabilidad creativa cuya huella estadística documentó la literatura. Los primeros años del euro produjeron la gran paradoja con la que se abre el 2008: Irlanda y España, los países con menos deuda y con superávit, eran los más vulnerables, porque su posición fiscal dependía de ingresos transitorios y de un endeudamiento privado y exterior insostenible.

\subsubsection{Los años setenta: la ruptura del equilibrio en las economías avanzadas y la anomalía española}
Las economías avanzadas llegan a los años setenta con ratios de deuda pública históricamente bajas. La deuda heredada de la Segunda Guerra Mundial, que en el Reino Unido superaba el $200\%$ del PIB en 1946 y en Estados Unidos rozaba el $120\%$, se redujo durante los treinta gloriosos (1945-1973) sin grandes superávits primarios. 

La literatura (ABBAS et al., 2010; REINHART y SBRANCIA, 2015) atribuye esa reducción sobre todo a un diferencial muy negativo entre el tipo de interés y el crecimiento. El crecimiento real fue alto, la inflación moderada pero persistente, y en muchos países hubo represión financiera, con controles de capital, coeficientes obligatorios y techos a los tipos de interés, que mantenía los tipos reales bajos o negativos.\footnote{%
    La misma mecánica que vimos para la España franquista. La diferencia es que en las demás economías avanzadas esa mecánica convivía con una rápida expansión del Estado de bienestar, y en España no.
} 

TANZI y SCHUKNECHT (2000) documentan que el gasto público de los países industrializados pasó de un promedio de en torno al $28\%$ del PIB en 1960 a más del $40\%$ en 1980. Fue la etapa de mayor crecimiento del sector público en tiempos de paz. 

\paragraph{Crisis del petróleo y la aparición del déficit estructural}
Las crisis del petróleo convierten el déficit en un rasgo estructural de las economías avanzadas. En los setenta el déficit medio de la OCDE siguió siendo bajo, entre el $1\%$ y el $2\%$ del PIB, pero ya por encima del casi equilibrio español.

No obstante, lo relevante de los setenta no es tanto el nivel medio del déficit como el cambio de régimen. Los shocks de oferta de 1973 y 1979 redujeron el crecimiento potencial, mientras que los compromisos de gasto social adquiridos en la fase de crecimiento alto se mantuvieron. Los gobiernos intentaron además amortiguar el shock con políticas de demanda y el resultado fue la aparición de déficits que ya no se corregían en la fase alcista.\footnote{%
    Por eso la literatura habla del fin del presupuesto equilibrado a lo largo del ciclo como norma implícita de la posguerra (BUCHANAN y WAGNER, 1977).
}

La inflación de los setenta, sin embargo, siguió conteniendo la deuda. En 1975, la deuda era del $23,7\%$ del PIB en Alemania, del $20,9\%$ en Japón, del $44,4\%$ en Estados Unidos, del $49,6\%$ en el Reino Unido y del $49,9\%$ en Canadá. Los casos más altos de Europa continental eran Bélgica ($59,5\%$) e Italia ($57,9\%$), y en el extremo opuesto estaban Dinamarca ($11,2\%$) y Finlandia ($6,8\%$). Con los datos, la media de la CEE de los nueve, ponderada por el PIB, se situaba en torno al $35$-$40\%$, porque el mayor peso de Alemania y Francia (esta con una deuda en torno al $20\%$ a finales de los setenta) compensaba los niveles altos de Bélgica, Italia, Irlanda y el Reino Unido. 

España, con una deuda inferior al $10\%$ del PIB, se situaba en 1975 en el grupo de menor endeudamiento, junto a Finlandia y Dinamarca. Pero su caso era muy distinto del de estos dos países. Finlandia y Dinamarca tenían deuda baja con un sector público ya grande y un sistema tributario moderno. España la tenía con un sector público pequeño y un sistema tributario atrasado. De nuevo, el equilibrio español no era una virtud fiscal sino un síntoma de atraso.


\subsubsection{Los años ochenta: la gran divergencia de las deudas}
Los años ochenta son el periodo en que la deuda de las economías avanzadas diverge como nunca lo había hecho en tiempos de paz. El detonante común es el cambio de política monetaria: el giro antiinflacionista de la Reserva Federal desde 1979, seguido por los demás bancos centrales, elevó los tipos de interés reales a niveles positivos y altos. 

Al mismo tiempo, la liberalización financiera fue desmontando la represión financiera. El resultado fue que el diferencial entre tipo de interés y crecimiento, que había sido negativo durante tres décadas, se volvió positivo en la mayoría de las economías avanzadas: $\Delta d \approx [(i - g)/(1 + g)]\cdot d_{t-1} - sp_t + \text{ADD}_t$, el paso de $i < g$ a $i > g$, convierte la deuda heredada en un motor autónomo de endeudamiento, tanto mayor cuanto mayor sea $d$. Por tanto, con la liberalización apareció el efecto bola de nieve: un país con déficits primarios moderados y una deuda ya relevante podía ver crecer su ratio de forma casi automática.

El déficit medio de la OCDE de los ochenta se situaba en torno al $3\%$ del PIB y el de España algo por encima del $4\%$. A diferencia de los setenta, por tanto, España ya no estaba por debajo de la media sino por encima, y esa diferencia persistirá hasta 2005.

\paragraph{Gran Divergencia}
Entre 1975 y 1990 la deuda belga pasó del $59,5\%$ al $130,3\%$ del PIB, y la italiana del $57,9\%$ al $95,3\%$. Irlanda llegó al $108,3\%$ en 1987, desde el $64,6\%$ de 1980. Grecia pasó del $22,8\%$ en 1980 al $74,2\%$ en 1990, y los Países Bajos del $36\%$ en 1975 al $75\%$ en 1990. Canadá pasó del $44,6\%$ en 1980 al $73,7\%$ en 1990. Estados Unidos también subió, del $41,3\%$ en 1980 al $62,2\%$ en 1990, con la combinación de rebajas de impuestos y aumento del gasto en defensa de la primera administración Reagan. 

En el extremo opuesto, el Reino Unido redujo su deuda del $43,7\%$ en 1980 al $28,3\%$ en 1990, con superávits a finales de la década e ingresos de privatizaciones. Alemania la contuvo en torno al $40\%$. Japón la subió del $47,8\%$ en 1980 al $68,3\%$ en 1985 y la estabilizó después con la burbuja de finales de la década.

En este contexto, España pasa del $16,1\%$ en 1980 al $40,9\%$ en 1985 y al $41,3\%$ en 1990. Es uno de los mayores aumentos relativos del periodo (la deuda se multiplica por más de dos veces y media en cinco años), pero desde un nivel tan bajo que en 1990 la ratio española seguía por debajo de la de Alemania, Francia o Estados Unidos.

La pregunta que le surgiría a cualquiera es por qué. ¿Por qué divergieron? La explicación económica más simple (distintos shocks, distinto crecimiento) no basta, porque los países que más aumentaron su deuda no fueron necesariamente los que sufrieron los peores shocks. Por eso, al buscar en otros ámbitos, la divergencia de los ochenta dio lugar a una de las literaturas más fecundas de la economía política de la hacienda pública.

ROUBINI y SACHS (1989) mostraron que los déficits eran mayores y más persistentes en los países con gobiernos de coalición amplios, frágiles y de corta duración, como Bélgica, Italia, los Países Bajos, Irlanda o Dinamarca en ciertos periodos. En estos gobiernos cada socio tiene capacidad de veto sobre los recortes que afectan a su electorado y el ajuste se retrasa. GRILLI, MASCIANDARO y TABELLINI (1991) ampliaron el análisis a la estructura institucional. Encontraron que los países con sistemas electorales proporcionales y gobiernos inestables acumularon más deuda que los de sistemas mayoritarios, y que la independencia del banco central estaba asociada a menor inflación sin coste en crecimiento. 

ALESINA y DRAZEN (1991) formalizaron el mecanismo con su modelo de guerra de desgaste: cuando los grupos sociales discrepan sobre quién debe soportar el coste del ajuste, cada uno espera que el otro ceda, y la estabilización se retrasa hasta que el coste de no hacer nada se vuelve insoportable. PERSSON y SVENSSON (1989) y ALESINA y TABELLINI (1990) añadieron el argumento de la deuda estratégica: un gobierno que prevé perder el poder puede endeudarse para limitar las opciones de gasto de su sucesor.

Por último, la línea iniciada por VON HAGEN (1992) mostró la importancia de las instituciones presupuestarias. Los países con procedimientos centralizados, un ministro de Hacienda con poder de agenda frente a los ministerios de gasto y restricciones en la fase parlamentaria tenían déficits significativamente menores que los de procedimientos fragmentados.\footnote{%
    El caso español encaja parcialmente en este cuadro. Evidentemente, no encaja en la hipótesis de las coaliciones, porque entre 1982 y 1993 España tuvo gobiernos de mayoría absoluta de un solo partido. Sin embargo, encaja en la de las instituciones presupuestarias débiles, porque VON HAGEN situaba a España en el grupo de procedimientos poco centralizados, y en la de la guerra de desgaste, a la luz de la ruptura del diálogo social que culminó en la huelga general de 1988. Encaja también, y sobre todo, en una explicación que la literatura comparada menciona menos: el catch-up del Estado de bienestar. España estaba construyendo en los ochenta lo que los demás países habían construido en los sesenta y setenta. Su déficit no era de un Estado de bienestar maduro que no quería ajustarse, sino de uno en construcción que no había terminado de ajustar sus ingresos.
}

\paragraph{Consolidaciones: ¿las consolidaciones expansivas?}
El contrapunto de la divergencia son los países que consiguieron corregir la trayectoria en los ochenta, porque de su experiencia nace una tesis que se invocará en todas las consolidaciones posteriores, incluida la española de 2010-2013. 

Dinamarca (1983-1986) e Irlanda (1987-1989) redujeron fuertemente sus déficits y, al mismo tiempo, crecieron con fuerza. GIAVAZZI y PAGANO (1990) interpretaron estos casos como consolidaciones fiscales expansivas. Un ajuste creíble y centrado en el gasto habría mejorado las expectativas de los hogares sobre sus impuestos futuros y reducido las primas de riesgo, estimulando la demanda privada lo bastante para compensar el efecto contractivo directo.

La relectura posterior, cuyo exponente más citado es PEROTTI (2011), ha mostrado que en esos casos pesaron mucho otros factores. En Irlanda, la devaluación de 1986 y el pacto social de moderación salarial de 1987. En Dinamarca, la fijación del tipo de cambio al marco alemán y la caída de los tipos de interés. En ambos, también, la recuperación de la demanda de sus respectivos socios comerciales. 

Con todo, lo que debemos retener es que el canal de las expectativas existe, pero rara vez es suficiente por sí solo, y requiere circunstancias (política monetaria acomodaticia, tipo de cambio o salarios flexibles) que no se dan en una unión monetaria en recesión. 

\subsubsection{1990-1997: la crisis de principios de los noventa y la convergencia de Maastricht}
La crisis de 1991-1993 desencadenó un periodo de fuerte aumento generalizado del déficit pero con epicentros muy distintos.

El caso más grave de deterioro fiscal fue el de los países nórdicos, que sufrieron crisis bancarias sistémicas tras la liberalización financiera de los ochenta. La deuda de Suecia pasó del $39\%$ del PIB en 1990 al $66,1\%$ en 1993, y la de Finlandia del $13,9\%$ al $54,2\%$ en solo tres años. El déficit público sueco superó el $10\%$ del PIB en 1993. En Italia, la crisis del SME de septiembre de 1992 coincidió con un sistema político en descomposición, y la deuda superó el $115\%$ en 1993. Bélgica alcanzó su máximo en 1993 ($138,9\%$) y Grecia superó el $100\%$. Fuera de Europa, Canadá llegó al $100\%$ en 1995, y Japón inició, tras el estallido de su burbuja en 1990-1991, un aumento de la deuda que no se ha detenido desde entonces: del $63,2\%$ en 1990 al $92,5\%$ en 1995, al $135,6\%$ en 2000 y al $173\%$ en 2007. El déficit medio de la OCDE fue del $4,9\%$.

En ese cuadro, España (del $41,3\%$ en 1990 al $61,6\%$ en 1995) muestra un deterioro intenso pero no extremo. La comparación con Suecia y Finlandia es la más instructiva: en los tres casos una crisis de balance privado, bancaria en los nórdicos y de competitividad y tipos de cambio en España, se tradujo en una crisis fiscal. Un anticipo de lo que ocurriría en 2008-2012.

\paragraph{Consolidaciones de los noventa: el papel de las instituciones}
La segunda mitad de los noventa es la etapa de las grandes consolidaciones, cuya comparación ilustra cuán importante es el diseño institucional.

Suecia reformó a fondo su marco presupuestario tras la crisis: techos de gasto nominales plurianuales aprobados por el Parlamento, un procedimiento presupuestario de arriba abajo y un objetivo de superávit a lo largo del ciclo (1997). Redujo su deuda del $68,7\%$ en 1995 al $39,2\%$ en 2007. 

Estados Unidos, con la Budget Enforcement Act de 1990 y sus reglas de compensación (PAYGO), el crecimiento de finales de los noventa y los ingresos de la burbuja bursátil, registró superávits federales entre 1998 y 2001. Canadá aplicó en 1995 su \textit{Program Review}, una revisión sistemática de todos los programas de gasto federal, con recortes muy significativos y una regla de previsiones prudentes. Su deuda bajó del $100\%$ al $67,2\%$ en 2007. 

Bélgica, el caso más extremo de Europa, mantuvo durante más de una década superávits primarios de en torno al $5$-$6\%$ del PIB, lo que le permitió bajar del $138,9\%$ en 1993 al $87,3\%$ en 2007. 

Como vemos, las consolidaciones más duraderas combinaron un ajuste sustancial con reglas de gasto y procedimientos presupuestarios reforzados, no solo con objetivos de saldo. 

\paragraph{Maastricht como ancla de la convergencia}
Además, dentro de este periodo destacan en la literatura comparada los efectos que tuvo Maastricht entre sus integrantes. Como sabemos, Maastricht fijó unos criterios de déficit ($3\%$ del PIB) y de deuda ($60\%$, o disminuyendo suficientemente) uniformes para acceder a la moneda única.\footnote{%
    Los criterios de Maastricht fueron objeto de críticas desde su origen. BUITER, CORSETTI y ROUBINI (1993) cuestionaron tres aspectos. En primer lugar, la arbitrariedad de las cifras de referencia (el $60\%$ era aproximadamente la media de la época, y el $3\%$ es el déficit que estabiliza la deuda en el $60\%$ con un crecimiento nominal del $5\%$). En segundo lugar, su carácter procíclico, porque exigía ajustes en recesión. Y, por último, la ausencia de fundamento en una teoría de la sostenibilidad. 

La ratio de deuda se estabiliza cuando el déficit total es $d_t\cdot g_t/(1 + g_t)$. Con $d = 60\%$ y $g = 5\%$ nominal ($3\%$ real y $2\%$ de inflación), el déficit estabilizador es $\sim 2,9\%$, de donde procede la coherencia aritmética entre el $3\%$ y el $60\%$.
} Los once países que formaron la primera oleada de la UEM pasaron de un déficit medio del $5,7\%$ en 1993 a uno inferior al $2\%$ a finales de los noventa.

En la evaluación de 1998, con datos de 1997, todos los candidatos salvo Grecia cumplieron el criterio de déficit. En cuanto a la deuda, se aplicó una interpretación flexible: Bélgica e Italia, con ratios superiores al $115\%$, fueron admitidas por la tendencia decreciente. España, con un $68,8\%$, se situaba en una posición intermedia, lejos de los casos más problemáticos. Grecia se incorporó en 2001.

Ahora bien, la otra cara de la convergencia de Maastricht fue el recurso y proliferación de operaciones que mejoraban el saldo del año de referencia sin mejorar la posición fiscal subyacente (con efectos perversos sobre el periodo subsiguiente). Los ejemplos más conocidos son los de Italia y Francia.

Italia aprobó en 1997 una eurotasa, un impuesto extraordinario que luego se devolvió parcialmente. Francia contabilizó en 1997 como ingreso el pago de France Télécom al Estado a cambio de que este asumiera sus compromisos de pensiones, de alrededor de medio punto de PIB, lo que Eurostat aceptó. KOEN y VAN DEN NOORD (2005), para la OCDE, documentaron el uso sistemático de medidas puntuales en el periodo.

VON HAGEN y WOLFF (2006) aportaron la evidencia más sólida. Los ajustes déficit-deuda, es decir, la diferencia entre la variación de la deuda y el déficit, eran significativamente mayores en los años en que el déficit de un país se acercaba al límite del $3\%$. Es la huella estadística de la contabilidad creativa: se reducía el déficit declarado desplazando operaciones al ajuste déficit-deuda. MILESI-FERRETTI (2004) había modelizado ese incentivo: una regla basada en un indicador manipulable induce la manipulación del indicador, no el ajuste.

El caso extremo, conocido después, fue el de Grecia. La revisión de sus estadísticas fiscales, iniciada en 2004 tras el cambio de gobierno, mostró que sus déficits habían superado el $3\%$ en todos los años desde 1997, incluido el de referencia para su entrada en el euro.\footnote{%
    En el caso de España no hay evidencia de manipulación comparable en la convergencia. Pero sí se acumularon pasivos fuera del saldo, lo que pertenece a la misma familia de problemas: la distancia entre el indicador vigilado y la posición fiscal real.
}

\subsubsection{1999-2007: los primeros años del euro y la paradoja de los países con menos deuda}
La década previa a la crisis fue muy dispar. 

Empezando por el ámbito extra europeo, Estados Unidos cerró los superávits de 1998-2001 con las rebajas fiscales de 2001-2003, el gasto militar posterior al 11 de septiembre y la recesión de 2001. Su deuda subió del $53,2\%$ en 2000 al $64,7\%$ en 2007. Japón, por otro lado, representa el caso extremo de deuda creciente sin crisis de financiación: su ratio superó el $170\%$ en 2007. La explicación habitual es que se trata de una deuda en moneda propia, en manos de residentes, con un banco central que puede comprarla y con un exceso de ahorro privado que mantiene los tipos en mínimos.\footnote{%
    Confirma, por contrafactual, la tesis de DE GRAUWE: lo que hace vulnerable a un país no es solo cuánto debe, sino en qué moneda, a quién y con qué respaldo del banco central.
} 

En el ámbito europeo, con la llegada de la moneda única, la evolución del déficit fue muy dispar. España e Irlanda continuaron la consolidación, mientras que Francia, Alemania, Italia y Portugal tuvieron grandes dificultades para cumplir el límite del $3\%$ del PEC. Fuera del euro, Estados Unidos registró superávits en 1999-2000, y Japón transformó sus superávits de principios de los noventa en déficits persistentes.

Las cifras de deuda confirman y completan esa descripción. Entre 2000 y 2007, la deuda de Irlanda bajó del $36,4\%$ al $24\%$ y la de España del $57,8\%$ al $35,7\%$, y fueron los únicos grandes descensos de la zona del euro en ese periodo. Bélgica siguió reduciendo su deuda (del $109,6\%$ al $87,3\%$), pero desde niveles muy altos. En cambio, Alemania subió del $59,2\%$ al $63,7\%$ y Francia del $59,7\%$ al $65,5\%$, ambas por encima del $60\%$. Italia se estancó en torno al $104\%$, Portugal pasó del $54,2\%$ al $72,7\%$ y Grecia se mantuvo por encima del $100\%$. Alemania y Francia incumplieron el límite de déficit a partir de 2002-2003, y la decisión del Ecofin de noviembre de 2003 de no aplicarles el procedimiento llevó a la sentencia del TJUE de 2004 y a la reforma del PEC de 2005.

Alemania, por tanto, es el contraejemplo del relato habitual. En los primeros años del euro no fue un modelo de disciplina fiscal: soportaba todavía el coste de la reunificación, crecía muy poco y era considerada el enfermo de Europa. Su consolidación fiscal y la mejora de su competitividad, con las reformas del mercado laboral de 2003-2005 y la moderación salarial, son posteriores. En 2007 tenía una deuda superior en casi treinta puntos a la española.

\paragraph{Paradoja: }
Con todo lo expuesto es prácticamente imposible dar respuesta a una pregunta. ¿Por qué los países con menos deuda fueron los más vulnerables? Aquí está la lección central de la comparación histórica. 

En 2007, Irlanda y España eran, con diferencia, los grandes países de la zona del euro con menos deuda pública y los únicos con superávit presupuestario significativo. Sin embargo, entre 2008 y 2012 fueron dos de los países con mayor deterioro fiscal de toda la OCDE, y ambos necesitaron asistencia financiera europea: Irlanda, un programa completo en 2010, y España, la asistencia para su sector bancario en 2012.

La literatura ha explicado la paradoja con tres argumentos complementarios. El primero, como habíamos visto para el caso español, es que sus saldos de 2005-2007 contenían un gran componente de ingresos transitorios ligados a burbujas inmobiliarias. El saldo estructural real era mucho peor que el publicado, y ese problema de medición afectaba especialmente a los países con auges de activos: \textbf{grandes economías financieras}.

El segundo argumento es que la vulnerabilidad procedía del sector privado y del sector exterior. LANE (2012) y BALDWIN y GIAVAZZI (2015) han mostrado que la crisis de la zona del euro fue en su origen una crisis de balanza de pagos. Los países de la periferia acumularon grandes déficits por cuenta corriente, financiados con entradas de capital privado del centro de la zona, que se detuvieron bruscamente en 2008-2010. En 2007 el déficit por cuenta corriente español rondaba el $9$-$10\%$ del PIB, y la deuda privada de hogares y empresas superaba el $200\%$. Cuando las entradas se detuvieron, la deuda privada, sobre todo la bancaria, se convirtió en deuda pública por los rescates bancarios y la caída de ingresos.

El tercer argumento es la pertenencia a una unión monetaria sin prestamista de última instancia, que DE GRAUWE (2011) formuló como la fragilidad de la zona del euro: un país que se endeuda en una moneda que no controla puede sufrir una crisis de liquidez autocumplida aunque sus fundamentales sean sólidos.

La consecuencia para el marco de análisis es de primer orden. La comparación de los niveles de deuda de 2007 habría situado a España entre los países más sólidos de la OCDE. Una comparación que hubiera incluido el saldo estructural corregido por los ingresos transitorios, la posición de inversión internacional neta y la deuda privada la habría situado entre los más vulnerables. Evidentemente, esta lección está detrás del Procedimiento de Desequilibrios Macroeconómicos de 2011 y de la inclusión del análisis de sostenibilidad de la deuda en el nuevo marco fiscal de 2024, que se examinarán más adelante.\footnote{%
    Procedimiento de Desequilibrios Macroeconómicos (Reglamento (UE) 1176/2011) y su cuadro de indicadores, y el análisis de sostenibilidad de la deuda del Reglamento (UE) 2024/1263 ([\authorfont{Ver Tema 3.B.45}]).
}


\subsection{Las crisis de 2008-2023 en perspectiva comparada}
Entre 2008 y 2023 las economías avanzadas atraviesan tres crisis que ponen a prueba sus finanzas públicas de formas distintas. La primera, la Gran Recesión de 2008-2009, es un shock financiero global que deteriora los saldos en todas partes, pero con intensidades muy diferentes. La segunda, la crisis de la deuda soberana de 2010-2013, es un fenómeno esencialmente europeo, y más precisamente de la zona del euro, que separa a los países de la moneda única en dos grupos: los que pierden el acceso a los mercados o se acercan a perderlo, y los que actúan como refugio. La tercera, la pandemia de 2020 seguida del shock inflacionario de 2022-2023, vuelve a ser global y sincronizada, pero llega con una respuesta institucional europea radicalmente distinta.

La comparación de las crisis de 2008-2023 muestra tres rasgos del caso español. El primero es la \textbf{elevada sensibilidad de su saldo al ciclo}. En 2008-2009, con una caída del PIB menor que la alemana o la italiana, España registró un deterioro del saldo del doble de la media de la zona del euro, comparable solo al de Irlanda y Grecia. La causa fue la desaparición de los ingresos del auge inmobiliario y la intensidad de los estabilizadores del desempleo, propia de su mercado de trabajo dual.

El segundo rasgo es su \textbf{posición intermedia en la crisis soberana}. España no recibió un programa completo como Grecia, Irlanda y Portugal, pero necesitó asistencia europea para su banca, y sufrió primas de riesgo que la comparación con el Reino Unido, Estados Unidos o Japón muestra que respondían más al diseño de la zona del euro que a sus fundamentales. El coste de su crisis bancaria fue mucho menor que el irlandés en porcentaje del PIB, pero el agujero de ingresos produjo un aumento de la deuda comparable.

El tercer rasgo, y el más revelador, es su \textbf{incapacidad para generar superávits primarios en las recuperaciones}. En 2014-2019 y, sobre todo, en 2021-2023, países que habían pasado por programas de ajuste, como Portugal, Grecia, Chipre e Irlanda, redujeron su deuda por debajo del nivel previo a la pandemia gracias a superávits primarios. España, con el mismo viento de cola del crecimiento nominal, salió de 2023 con más deuda que en 2019. 

El resultado de los quince años es un cambio de posición completo: de ser uno de los países con menos deuda de la zona del euro en 2007 ($36\%$ frente a una media del $66\%$) a estar entre los cinco con más. 

\subsubsection{La Gran Recesión (2008-2009): un deterioro general con intensidades muy distintas}
Tras la crisis financiera internacional prácticamente todos los países de la OCDE incurrieron en déficit o lo agravaron, y el déficit medio pasó de poco más del $1\%$ del PIB en 2007 a más del $8\%$ en 2009, con grandes diferencias entre países.

Las cifras de Eurostat de la época permiten concretar esas diferencias. La zona del euro pasó de un déficit del $0,7\%$ en 2007 al $6,4\%$ en 2009, un deterioro de unos seis puntos. España pasó de un superávit del $1,9\%$ a un déficit del $11,2\%$, un deterioro de unos $13$ puntos, el doble que la media. Solo Irlanda (del $0,1\%$ de superávit al $14,2\%$ de déficit) y Grecia (del $6,5\%$ al $15,8\%$ de déficit) registraron deterioros comparables. Fuera de la zona del euro, el Reino Unido pasó del $2,7\%$ al $11,5\%$ de déficit. En el otro extremo, Alemania pasó de un ligero superávit a un déficit del $3,2\%$ e Italia del $1,6\%$ al $5,4\%$.

\paragraph{Por qué España se deterioró el doble que la media: los ingresos}
La comparación permite descartar una explicación simple. La caída del PIB español en 2009 (en torno al $-3,8\%$) fue menor que la de Alemania ($-5,7\%$), Italia ($-5,5\%$) o el Reino Unido ($-4,5\%$). Si el deterioro del saldo dependiera solo del tamaño de la recesión, el de España debería haber sido menor que el alemán, y fue cuatro veces mayor.

La explicación está en la composición del deterioro. La recaudación española cayó más de seis puntos de PIB, una caída que no se dio en ningún otro país de nuestro entorno. Es el efecto, ya analizado, de la desaparición de los ingresos transitorios del auge inmobiliario, que tenía su paralelo en Irlanda y, en menor medida, en el Reino Unido y en Estados Unidos. Los tres países de la OCDE con mayor deterioro de los ingresos en 2008-2009 fueron precisamente los que habían vivido los mayores auges inmobiliarios.

Un segundo factor comparado es la mayor intensidad de los estabilizadores automáticos del gasto en España. DOLLS, FUEST y PEICHL (2012) estimaron que, en Europa, los estabilizadores automáticos absorbían una proporción del shock de renta de los hogares bastante mayor que en Estados Unidos. Dentro de Europa, España destacaba por la reacción del gasto en desempleo, porque su mercado laboral dual, con una elevada proporción de temporales, destruye empleo con mucha más intensidad que el de otros países ante una misma caída del PIB. Entre 2007 y 2013 la tasa de paro española pasó del $8\%$ al $26,9\%$, mientras que la alemana simplemente bajó. Esa elasticidad del empleo al ciclo, que es un rasgo estructural del mercado de trabajo español, se traduce en una elasticidad del saldo al ciclo mayor que la de sus pares.\footnote{%
    Son cuestiones importantes. La elasticidad presupuestaria al ciclo en la metodología común de la Unión y su mayor valor para España por la reacción del desempleo. Sobre la dualidad del mercado de trabajo español y su efecto en la destrucción de empleo, véanse los temas dedicados al mercado de trabajo.
}

La respuesta discrecional fue coordinada a escala del G20 (cumbres de Washington y Londres, 2008-2009) y, en la Unión, a través del Plan Europeo de Recuperación Económica, que recomendaba un estímulo de en torno al $1,5\%$ del PIB. Los mayores estímulos discrecionales correspondieron a Estados Unidos (la American Recovery and Reinvestment Act de 2009, en torno al $5,5\%$ del PIB repartido en varios años) y a China, con un paquete de cuatro billones de yuanes (en torno al $12\%$ de su PIB) concentrado en inversión en infraestructuras. En Europa, Alemania aprobó dos paquetes coyunturales de tamaño relevante con la economía más saneada. El de España, con el Plan E, las rebajas del IRPF y la supresión efectiva del Impuesto sobre el Patrimonio, fue de tamaño similar en porcentaje del PIB. La diferencia es que España lo aplicó sobre un saldo estructural que, como vimos, era mucho peor de lo que parecía.

\subsubsection{La crisis de la deuda soberana (2010-2013): la fractura de la zona del euro}
Otro rasgo particular que muestra la comparación es que la crisis de 2010-2013 no fue una crisis de las economías con más deuda o más déficit. Fue una crisis de los países de la zona del euro con desequilibrios exteriores y problemas bancarios, como ya sabemos. 

El caso más claro es el del Reino Unido. En 2010 tenía un déficit del $10,3\%$ del PIB, superior al español ($9,3\%$), y Estados Unidos y Japón tenían déficits y deudas mayores. Sin embargo, sus rentabilidades a diez años bajaron durante la crisis mientras las de la periferia del euro se disparaban. Japón, con una deuda que superaba el $220\%$ del PIB en 2012, se financiaba a tipos inferiores al $1\%$.\footnote{
De nuevo, la base empírica de la tesis de DE GRAUWE (2011). Los países que se endeudan en su propia moneda y cuentan con un banco central que puede actuar como prestamista de última instancia no sufren crisis de liquidez autocumplidas. Los que se endeudan en una moneda que no controlan, como los miembros de la zona del euro, sí pueden sufrirlas. DE GRAUWE y JI (2013) mostraron que el aumento de las primas de la periferia en 2010-2012 no se explicaba por sus fundamentales fiscales. La fuerte caída que siguió al anuncio del OMT en 2012, sin cambios en esos fundamentales, confirma que buena parte de la prima era un componente de pánico y de riesgo de redenominación.
}

Ante el sucesivo deterioro de las finanzas públicas, la zona del euro respondió con una escala de programas de intensidad creciente. Grecia recibió un primer programa en mayo de 2010, en forma de préstamos bilaterales de los países de la zona del euro y del FMI por $110.000$ millones de euros. El segundo (2012) incluyó la mayor reestructuración de deuda soberana de la historia: el canje con participación del sector privado (PSI) de marzo de 2012 impuso a los acreedores privados una quita nominal del $53,5\%$. Hubo un tercero en 2015, ya con el MEDE. Irlanda recibió un programa de $85.000$ millones en noviembre de 2010, y Portugal uno de $78.000$ millones en mayo de 2011. Chipre recibió en 2013 un programa de $10.000$ millones condicionado a un bail-in de los depositantes de sus dos principales bancos, un precedente de la resolución bancaria europea. España solicitó en junio de 2012 una línea para su sector bancario de hasta $100.000$ millones, de la que utilizó $41.333$ millones. Italia no pidió asistencia, pero sufrió tensiones comparables a las españolas en 2011-2012, y el BCE compró deuda de ambos países con el programa SMP a partir de agosto de 2011.

En esta escala, España ocupa una posición intermedia que se suele malinterpretar. No fue rescatada en el sentido de los programas completos de Grecia, Irlanda y Portugal, que incluían financiación del Estado y condicionalidad macroeconómica, fiscal y estructural sujeta a la supervisión de la troika. Pero tampoco quedó al margen. La asistencia fue sectorial, con condicionalidad financiera, y en términos de deuda pública tuvo el mismo efecto que un préstamo al Estado, porque el FROB era la entidad receptora y el Estado el garante.

La comparación más instructiva en este contexto es con Irlanda, porque ambos países partían en 2007 de una deuda baja y un auge inmobiliario financiado por el crédito bancario. La decisión irlandesa de garantizar en septiembre de 2008 prácticamente todos los pasivos de sus bancos llevó a un déficit del $31,3\%$ del PIB en 2010, del que unos veinte puntos correspondían a la recapitalización bancaria. La deuda irlandesa pasó del $24\%$ en 2007 al $118,7\%$ en 2012. Según la base de crisis bancarias de LAEVEN y VALENCIA (2018), el coste fiscal bruto directo de la crisis bancaria irlandesa se situó en torno a un tercio del PIB. El español, en torno al $4$-$6\%$ según las estimaciones ($5,5\%$ del PIB en el cómputo de fondos utilizados). Aun así, el coste de la reestructuración bancaria española, que el Tribunal de Cuentas cifra en $71.833$ millones, es uno de los mayores de Europa en términos absolutos, solo por detrás de Alemania y del Reino Unido.

La diferencia de escala se explica por el tamaño relativo del sistema bancario (en Irlanda superaba varias veces su PIB), por la garantía general de 2008 y por el hecho de que en España el problema se concentró en las cajas de ahorros, mientras los grandes bancos resistían. Pero la consecuencia fiscal agregada fue parecida por otra vía. En España pesó mucho menos el rescate bancario y mucho más la caída estructural de los ingresos. Ambos países llegaron a ratios de deuda cercanas al $100\%$-$120\%$: Irlanda por el rescate bancario y España por el agujero de los ingresos.

\paragraph{Consolidación: Ajuste comparado y el debate sobre sus costes}

Los países de la periferia aplicaron entre 2010 y 2013 consolidaciones fiscales de gran tamaño, en buena parte simultáneas y en plena recesión. La consolidación discrecional acumulada, medida por la mejora del saldo estructural primario, fue del orden de 15-20 puntos de PIB en Grecia, de 7-9 en Portugal y en Irlanda y de 6-8 en España en ese periodo, según las estimaciones de la Comisión y del FMI. Es una magnitud sin precedentes en economías avanzadas en tiempos de paz.

La literatura comparada ha evaluado el coste de este ajuste con gran intensidad. BLANCHARD y LEIGH (2013), como vimos, mostraron que los multiplicadores eran mayores de lo supuesto. HOUSE, PROEBSTING y TESAR (2020) estiman que la austeridad explica una parte muy importante de la caída del PIB de la periferia europea entre 2010 y 2014, y que en algunos casos la ratio de deuda habría aumentado menos sin ella. FATÁS y SUMMERS (2018) añaden el argumento de la histéresis: si la recesión reduce de forma permanente el producto potencial, por la pérdida de capital humano y la caída de la inversión, la consolidación en recesión puede empeorar la sostenibilidad a largo plazo.

La defensa del ajuste se basa en que sin él no habría habido financiación, ni europea ni de mercado. La comparación internacional ofrece un matiz relevante. El Reino Unido, fuera del euro, también aplicó un programa de consolidación desde 2010, pero con su propio banco central comprando deuda y con una libra que se había depreciado en torno a un $25\%$ desde 2007. El coste en crecimiento fue menor. 

La diferencia no está tanto en el ajuste fiscal como en la ausencia de amortiguadores monetarios y cambiarios en la zona del euro hasta 2012.

\subsubsection{La recuperación (2014-2019): dos modelos de salida}
Entre 2014 y 2019, las trayectorias de deuda de la zona del euro divergen de nuevo. Alemania, con superávits presupuestarios desde 2014 y su regla constitucional de freno de la deuda (Schuldenbremse, 2009), redujo su deuda del $79,8\%$ del PIB en 2012 al $58,7\%$ en 2019, por debajo del umbral de Maastricht. Irlanda la redujo del $118,7\%$ en 2012 al $55,9\%$ en 2019. Portugal, del $132,5\%$ en 2014 al $116,1\%$, y los Países Bajos, del $67,2\%$ al $47,6\%$. En el grupo contrario, la de Francia subió del $96,2\%$ al $98,2\%$ y la de Italia se estancó en torno al $134\%$. España pasó del $104,4\%$ al $97,6\%$, un descenso de unos siete puntos en cinco años de crecimiento superior a la media.\footnote{
\textbf{NOTA.-} La reducción irlandesa está exagerada por un problema de denominador. En 2015, la relocalización de activos intangibles de multinacionales elevó el PIB irlandés un $25\%$ en un solo año, lo que redujo mecánicamente su ratio de deuda. Las autoridades irlandesas utilizan desde entonces un indicador alternativo, la renta nacional bruta modificada (GNI), sobre el que su deuda es muy superior a la que resulta sobre el PIB. 
}

La comparación más reveladora de este periodo, y la que anticipa el contraste de 2019-2023, es la de España con Portugal. Ambos países crecieron en 2015-2019 a tasas parecidas, y ambos se beneficiaron del programa de compras del BCE. Pero Portugal, que salió de su programa en 2014 con compromisos de consolidación supervisados, mantuvo superávits primarios de en torno al $3\%$ del PIB a partir de 2016, y cerró 2019 con un superávit total, el primero de su democracia. España mantuvo un déficit primario hasta 2018 y un saldo total cercano al $-3\%$ al final del periodo.

La explicación que ofrece la economía política comparada combina dos factores. Por un lado, los programas de ajuste dejaron en los países rescatados un anclaje institucional más fuerte: consejos fiscales reforzados, supervisión posprograma y una memoria social del coste de perder el acceso a los mercados. Por otro, España vivió en 2015-2019 un periodo de fragmentación política sin precedentes, con presupuestos prorrogados, que la literatura (ROUBINI y SACHS, 1989) asocia a una menor capacidad de ajuste. 

En 2019 España tenía el cuarto déficit público más elevado de las economías avanzadas, solo superado por Estados Unidos, Australia e Israel, e igualado por Francia. La conclusión es robusta: en plena fase alta del ciclo y con tipos de interés en mínimos históricos, España estaba entre las pocas economías avanzadas con un déficit en torno al $3\%$ del PIB. 

Es la misma posición que en 2007 con respecto a su saldo estructural oculto. La diferencia es que ahora el déficit era visible, porque ya no había ingresos transitorios que lo enmascararan.

\subsubsection{La pandemia y la inflación (2020-2023): una crisis sincronizada con respuestas desiguales}
Según los datos del FMI, entre 2019 y 2023 la ratio de deuda cambió de forma muy desigual. Portugal la redujo del $116,1\%$ al $97,7\%$ ($-18,4$ puntos), Grecia del $183,7\%$ al $165,2\%$ ($-18,5$), Irlanda del $55,9\%$ al $43,3\%$ ($-12,6$) y Chipre del $92,3\%$ al $73,7\%$ ($-18,6$). En el otro extremo, España la aumentó del $97,6\%$ al $105\%$ (+$7,4$), Francia del $98,2\%$ al $109,8\%$ (+$11,6$), Estados Unidos del $108,2\%$ al $119\%$ (+$10,8$) y el Reino Unido del $85,7\%$ al $100,4\%$ (+$14,7$). Alemania (+$4,2$) e Italia (+$0,7$) se situaron en posiciones intermedias.

Si se comparan 2019 y 2022, España fue la quinta economía avanzada con mayor aumento de su ratio (unos $15$ puntos), solo superada por San Marino, Japón, Nueva Zelanda y Malta. Su ratio, del $113,1\%$ en 2022, era la séptima más alta de las economías avanzadas y la cuarta de la zona del euro, unos $20$ puntos por encima de la media.

El contraste no se explica por las condiciones macroeconómicas, que fueron parecidas. Todos los países de la periferia se beneficiaron del crecimiento nominal de 2021-2023, de la inflación y de los fondos europeos, y todos soportaron la subida de tipos de 2022-2023. Se explica por el saldo primario. Portugal, Grecia y Chipre volvieron a superávits primarios en 2022-2023: Grecia del orden del $2\%$ del PIB, Portugal por encima del $2\%$ y Chipre cerca del $3\%$. España mantuvo déficits primarios en todos los años hasta 2024, y su saldo primario no se acercó al equilibrio hasta 2025. Con el mismo viento de cola del denominador, los países que generaban superávits primarios redujeron su deuda por debajo del nivel de 2019, y España no llegó a recuperarlo. En 2025 su ratio ($100,7\%$) sigue por encima de la de 2019.

La comparación con Francia merece un comentario, porque es el país con el que España ha convergido. Francia tenía en 2007 una deuda casi treinta puntos superior a la española y en 2023 era solo cinco puntos superior. Ambos países comparten un déficit estructural persistente, un gasto público elevado en relación con sus ingresos y dificultades políticas para aprobar presupuestos. Francia entró en procedimiento de déficit excesivo en 2024, y España no, porque su déficit de 2023 era inferior y se preveía por debajo del $3\%$ en 2024. En 2024-2025, la prima de riesgo española se situó en algunos momentos por debajo de la francesa, lo que indica que los mercados valoraban mejor la dinámica española que el nivel. 

Analicemos más en concreto los shocks del periodo para ver que nos dicen los datos.


\paragraph{Pandemia}
Evidentemente, el cataclismo del periodo fue la pandemia de la COVID-19, donde la capacidad de respuesta fiscal en 2020 dependió del acceso a la financiación, y ese acceso dependió a su vez del respaldo del banco central. 

En las economías avanzadas, el déficit medio pasó de menos del $3\%$ del PIB en 2019 a casi el $11\%$ en 2020, y dos años después había vuelto a niveles próximos a los de 2019. En las economías emergentes (excluida China) pasó de casi el $4\%$ a casi el $8\%$, y había vuelto a niveles cercanos al $4\%$ en 2022. China, que ya tenía un déficit cercano al $6\%$ en 2019, llegó a casi el $10\%$ en 2020 y a un $9\%$ en 2022, por la prolongación de su política de covid cero. Las economías de bajo ingreso pasaron de casi el $4\%$ a casi el $5\%$. La deuda pública mundial aumentó unos $15$ puntos de PIB en 2020, el mayor aumento anual desde la Segunda Guerra Mundial, hasta alcanzar el $99\%$ del PIB. Más del $90\%$ del aumento se concentró en las economías avanzadas y en China, que pudieron apoyarse en sus bancos centrales y en mercados financieros más desarrollados. En las economías avanzadas la deuda subió $19$ puntos, hasta el $124\%$ del PIB, con Estados Unidos cerca del $135\%$. En las emergentes y de bajo ingreso subió unos $7$ puntos, sobre todo por la caída del PIB.

España fue la séptima economía avanzada con mayor aumento de su deuda en 2020, con más de $21$ puntos de PIB. El mayor aumento fue el de Canadá, con más de $30$ puntos, y Taiwán incluso la redujo ligeramente. La deuda española pasó del $97,6\%$ en 2019 al $119,2\%$ en 2020. 

La comparación con la zona del euro muestra que el aumento español se explica en gran medida por el denominador. La caída del PIB español en 2020 (en torno al $-11\%$) fue la mayor de la zona del euro, casi el doble de la media (en torno al $-6\%$), por el peso del turismo, de la hostelería y de los servicios de contacto personal en su estructura productiva. Como vimos, el efecto denominador explicó en torno a la mitad del aumento de la ratio española en 2020. La respuesta discrecional española fue, en comparación, contenida en gasto directo y amplia en liquidez. La base de datos de medidas fiscales del FMI sitúa el apoyo directo (por encima de la línea) de España por debajo del de Alemania, el Reino Unido, Estados Unidos o Japón, y su apoyo de liquidez (avales, préstamos y capital) en niveles comparables a los de Francia, aunque por debajo de Italia y Alemania. 

\paragraph{Crisis energética de 2022: el mismo patrón}
La respuesta a la crisis energética de 2022 muestra de nuevo una diferencia de espacio fiscal. Según la base de datos de Bruegel, los países europeos destinaron casi $800.000$ millones de euros a proteger a hogares y empresas frente a la subida de los precios de la energía entre septiembre de 2021 y enero de 2023. 

Alemania encabezaba la lista con unos $270.000$ millones, seguida del Reino Unido, Italia y Francia, con menos de $150.000$ millones cada uno. España movilizó un volumen menor en porcentaje del PIB, en parte gracias al mecanismo ibérico, que redujo el precio de la electricidad sin coste presupuestario directo. La crítica de Bruegel, extensible a España, es que la mayor parte del apoyo europeo no estaba focalizado y actuaba como una subvención a los combustibles fósiles.

\subsection{Posición actual de España en perspectiva comparada (2024-2026)}
La verdad, con todo lo expuesto, es que la foto comparada de 2025-2026 favorece a España más de lo que lo explicado haría esperar, y eso encierra una trampa. 

Con datos de Eurostat de abril de 2026, España tuvo en 2025 un déficit ($2,4\%$) inferior a la media de la zona del euro ($2,9\%$) y muy inferior al de Francia ($5,1\%$) o Bélgica ($5,2\%$). Su prima de riesgo cotiza por debajo de $60$ puntos básicos, el diferencial de Francia frente a España está en máximos históricos y S\&P le da la misma calificación que a Francia. Pero España sigue teniendo la quinta deuda más alta de la Unión, con un déficit estructural que no ha corregido y unas necesidades brutas de financiación entre las mayores de Europa. 

Hasta ahora, hemos reconstruido la trayectoria comparada histórica. Sin embargo, para ver el por qué de la afirmación anterior debemos analizar su posición actual con tres grupos de indicadores. Los de nivel, que comparan el déficit, la deuda, el saldo primario y el tamaño del sector público. Los de dinámica y sostenibilidad, proyecciones de deuda, necesidades de financiación, envejecimiento y los indicadores de vulnerabilidad macrofinanciera. Y, por último, los de percepción recogen las primas de riesgo y las calificaciones crediticias.

Como veremos, España ocupa hoy una posición comparada mejor en dinámica y peor en nivel. 

Por nivel de déficit, España está por debajo de la media de la zona del euro ($2,4\%$ frente a $2,9\%$) y muy por debajo de Francia, Bélgica o Austria. Por nivel de deuda, tiene la quinta ratio más alta de la Unión ($100,7\%$), unos $13$ puntos por encima de la media, detrás de Grecia, Italia, Francia y Bélgica, y ha sido rebasada a la baja por Portugal. Su saldo primario está cerca del equilibrio, pero lejos de los superávits de la antigua periferia; sus intereses pesan más que la media en el gasto público, y sus ingresos siguen varios puntos por debajo de la media europea. Los indicadores prospectivos muestran un riesgo alto a medio plazo según la Comisión y unas necesidades brutas de financiación entre las más altas de la Unión. En cambio, los indicadores de vulnerabilidad macrofinanciera que fallaron en 2007 (cuenta corriente, posición exterior, deuda privada, solvencia bancaria) muestran hoy una posición mucho más sólida: la configuración de 2007 se ha invertido, con deuda pública alta y deuda privada moderada. Los mercados y las agencias premian la dinámica, pero las primas son hoy una señal menos fiable de disciplina por el respaldo del BCE. España ha mejorado su posición relativa gracias al crecimiento y al deterioro de otros, pero no ha resuelto el problema estructural que la comparación con Portugal pone de relieve.


\subsubsection{Nivel: déficit, deuda y tamaño del sector público en 2025}
Según la notificación de Eurostat de abril de 2026, el déficit de la zona del euro se situó en el $2,9\%$ del PIB en 2025, frente al $3\%$ de 2024, y el de la Unión se mantuvo en el $3,1\%$. En ese cuadro, España ($2,4\%$, incluidos los gastos de la DANA) está claramente por debajo de la media y muy por debajo de los grandes países con déficit elevado: Francia ($5,1\%$), Bélgica ($5,2\%$), Austria ($4,2\%$), Finlandia ($3,4\%$) e Italia ($3,1\%$). También está por debajo de Alemania, cuyo déficit ($2,7\%$) empieza a crecer con la reforma de su freno de la deuda para financiar defensa e infraestructuras. Los mayores déficits de la Unión fueron los de Rumanía ($7,9\%$) y Polonia ($7,3\%$). 

Solo cinco Estados miembros registraron superávit: Chipre ($3,4\%$), Dinamarca ($2,9\%$), Irlanda ($1,8\%$), Grecia ($1,7\%$) y Portugal ($0,7\%$).

La comparación muestra dos grupos en la zona del euro. En uno están los países de la antigua periferia (Grecia, Portugal, Chipre e Irlanda), con superávits y deudas en rápida reducción. En el otro, varios países del centro y del norte (Francia, Bélgica, Austria, Finlandia) con déficits superiores al $3\%$ y procedimientos de déficit excesivo abiertos o en perspectiva. España se sitúa entre ambos: sin superávit, pero con un déficit por debajo del $3\%$ y decreciente. Una inversión casi completa del mapa de 2010-2012, que conviene señalar.

Por otro lado, al cierre de 2025, la deuda de la zona del euro era del $87,8\%$ del PIB y la de la Unión del $81,7\%$. España ($100,7\%$) tiene la quinta ratio más alta de la Unión, detrás de Grecia ($146,1\%$), Italia ($137,1\%$), Francia ($115,6\%$) y Bélgica ($107,9\%$). Supera en unos $13$ puntos la media de la zona del euro y en más de 10 a Portugal ($89,7\%$), que en 2014 tenía casi treinta puntos más que España. Queda muy por encima de Alemania ($63,5\%$), de los Países Bajos ($44,4\%$) y de Irlanda ($32,9\%$). Las ratios más bajas de la Unión son las de Estonia ($24,1\%$), Luxemburgo ($26,5\%$) y Dinamarca ($27,9\%$).

Frente a 2022, la distancia con la media de la zona del euro se ha reducido de unos 20 a unos $13$ puntos, y España ha pasado del cuarto al quinto lugar de la zona del euro, al superarla Bélgica. Pero la reducción de esa distancia se debe sobre todo a que la ratio española ha caído por el denominador, no a un cambio en su posición estructural. Además, varios países que estaban por debajo de España, como Portugal o Chipre, la han rebasado a la baja en estos años.

Por último, estaría la comparación del saldo primario, más informativa que la del saldo total. En 2025 España tiene un saldo primario cercano al equilibrio. Es mejor que el de Francia o Bélgica, que tienen déficits primarios de varios puntos, pero peor que el de Italia, que viene registrando superávits primarios, y claramente peor que los de Grecia, Portugal o Chipre, con superávits primarios del $2$-$4\%$. 

Eurostat añade un indicador de coste de la deuda. En 2025 los intereses pagados por las AAPP representaron el $5,3\%$ del gasto público en España, frente al $3,8\%$ de media en la zona del euro. Solo Hungría ($8\%$), Italia ($7,6\%$), Rumanía ($6,6\%$) y Grecia ($6,5\%$) tuvieron una proporción mayor. En porcentaje del PIB, los intereses españoles (en torno al $2,4\%$) superan la media de la zona del euro. Son una carga que otros países con menos deuda no soportan, y que en 2026 vuelve a crecer.

Por último, el tamaño del sector público de la zona del euro en 2025 fue del $49,8\%$ del PIB, por gastos, y del $46,9\%$, por ingresos. España se sitúa varios puntos por debajo de la media en ambas magnitudes. Su gasto está en torno al $45$-$46\%$ del PIB y sus ingresos en torno al $42$-$43\%$, siendo la brecha de ingresos frente a la media es mayor: una constante desde la Transición. España ha homologado el tamaño de su Estado de bienestar con Europa más que su capacidad recaudatoria, y su déficit estructural refleja en buena parte esa brecha de ingresos.\footnote{%
    La brecha fiscal de España frente a la Unión (tipos efectivos, bases imponibles, beneficios fiscales y economía sumergida) y el debate sobre la suficiencia del sistema tributario corresponden a los temas de sistema fiscal del temario. Ver también el Libro Blanco sobre la Reforma Tributaria (2022) del Comité de personas expertas.

La literatura española discute si la brecha procede de tipos más bajos o de bases más estrechas (economía sumergida, beneficios fiscales, fraude); la mayoría de los análisis de la AIReF, del BdE y de la Comisión apuntan más a lo segundo.
}

\subsubsection{Sostenibilidad y vulnerabilidad: lo que miden los indicadores prospectivos}
Los indicadores de nivel describen el punto de partida; los prospectivos, la trayectoria. El nuevo marco fiscal europeo ha convertido el análisis de sostenibilidad de la deuda (DSA, por sus siglas en inglés) en el eje de la supervisión: las sendas de gasto neto se derivan de una trayectoria que debe situar la deuda en una senda descendente y plausible a medio plazo.\footnote{%
    Sobre el análisis de sostenibilidad de la deuda en el Reglamento (UE) 2024/1263 (trayectoria técnica y salvaguardias de resiliencia del déficit y de sostenibilidad de la deuda) ([\authorfont{Ver Tema 3.B.45}]); sobre los indicadores de sostenibilidad S0, S1 y S2 de la Comisión ([\authorfont{Ver Tema 3.A.39}]).
} La Comisión publica anualmente su \textit{Debt Sustainability Monitor}, que clasifica a los Estados miembros por su riesgo a corto, medio y largo plazo.

En la edición de 2024, la Comisión situó a España entre los nueve Estados miembros con riesgo alto de sostenibilidad a medio plazo. En un escenario sin cambios de política proyectaba una deuda en torno al $109\%$ del PIB en 2030 y al $118\%$ en 2034, con un ajuste anual necesario de entre el $0,4$ y el $0,6\%$ del PIB, e identificaba unas necesidades brutas de financiación cercanas al $20\%$ del PIB, entre las más altas de la Unión junto con Italia, Francia y Bélgica. La edición de 2025, publicada el 12 de febrero de 2026, eleva a doce el número de países con riesgo alto a medio plazo, grupo en el que España seguiría figurando. En la dimensión a corto plazo, en cambio, la Comisión no identifica a España entre los países con necesidades de financiación preocupantes, que en esta edición son Bélgica, Francia, Italia y Finlandia.

La diferencia entre la senda del Plan Fiscal y Estructural español, con la deuda bajando al $90,6\%$ en 2031, y la proyección de la Comisión sin cambios de política, con la deuda subiendo, mide precisamente el esfuerzo que España se ha comprometido a hacer. Ese esfuerzo exige respetar una senda de gasto neto que, como vimos, la AIReF considera en riesgo en 2026. A largo plazo, la Comisión clasifica el riesgo vinculado al envejecimiento como medio o alto según el escenario de pensiones que se considere.\footnote{%
    Sobre las proyecciones del gasto en pensiones (\textit{Ageing Report} 2024), la reforma de 2021-2023, la cláusula de salvaguarda y el informe de la AIReF de mayo de 2026 ([\authorfont{Ver Tema 4.B.24}]).
}

Pero la lección central de la comparación histórica era que el nivel de deuda pública, por sí solo, no mide la vulnerabilidad fiscal: en 2007 España e Irlanda tenían deuda baja, pero enormes desequilibrios privados y exteriores. Aplicada a la foto de 2025, esa lección arroja un resultado simétrico. Los indicadores que en 2007 señalaban una vulnerabilidad extrema muestran hoy una posición mucho más sólida. España tiene superávit por cuenta corriente desde 2013, que en 2024-2025 rondaba el $3\%$ del PIB; su posición de inversión internacional neta ha mejorado desde en torno al $-98\%$ del PIB en 2014 hasta en torno al $-40\%$ en 2025; la deuda de hogares y empresas se ha reducido desde más del $200\%$ del PIB en 2010 hasta en torno al $125\%$, por debajo de la media de la zona del euro; y el sistema bancario tiene más capital y menos exposición inmobiliaria, bajo la supervisión única del BCE.

En otras palabras, la configuración de 2007 se ha invertido. España tenía entonces una \textbf{deuda pública baja con una deuda privada y exterior muy alta}, y tiene ahora una \textbf{deuda pública alta con una deuda privada y exterior moderada}. La segunda configuración es menos vulnerable a una crisis de balanza de pagos como la de 2010-2012, porque ya no hay entradas de capital privado cuya interrupción pueda provocarla. Pero es más vulnerable a una subida de tipos y a un shock que exija gasto público, porque el espacio fiscal para absorberlo es menor. Por eso la mejor valoración de los mercados es compatible con la advertencia de la Comisión: miden riesgos distintos.\footnote{%
    Sobre los indicadores del Procedimiento de Desequilibrios Macroeconómicos (cuenta corriente, posición de inversión internacional neta, deuda privada, crédito, precios de la vivienda) y la salida de España de la categoría de desequilibrios ([\authorfont{Ver Tema 3.B.45}]).
}

La estructura de la deuda actúa, además, como amortiguador comparativo. Según el BdE, la vida media de la deuda española era de unos $8$ años en 2024, frente a $8,9$ años en la zona del euro; el $13,9\%$ de la deuda vencía en menos de un año y los no residentes tenían el $44,9\%$. Frente a Italia o Francia, España no tiene una estructura de vencimientos especialmente desfavorable, aunque su vida media, algo menor que la media de la zona del euro, hace que las subidas de tipos se trasladen algo más deprisa al coste medio, una cuestión que retomaremos al analizar la política del Tesoro.

\subsubsection{Percepción de los mercados y de las agencias de calificación}
¿Cómo leen los mercados esta posición? En septiembre de 2026, la prima de riesgo española frente al bono alemán a diez años ha bajado de $60$ puntos básicos, un nivel que no se veía desde antes de la crisis financiera, y el diferencial entre la deuda francesa y la española se sitúa en máximos históricos: Francia paga más que España a diez años, con su bono cerca del $4\%$ en el verano de 2026. Italia, que vivió en septiembre la mayor caída de primas de la zona del euro, cotiza también con primas históricamente bajas. Es un cambio radical respecto a julio de 2012, cuando la prima española alcanzó los $638$ puntos básicos.

La literatura y los analistas ofrecen tres interpretaciones complementarias. La primera es que los mercados valoran la dinámica más que el nivel: el crecimiento español superior a la media, un déficit decreciente y la estabilidad del saldo primario, frente a la parálisis presupuestaria y política de Francia. La segunda es el respaldo institucional del BCE: desde julio de 2022 existe el Instrumento para la Protección de la Transmisión (TPI), que permite al BCE comprar deuda de un país que sufra un aumento de su prima no justificado por sus fundamentales, siempre que cumpla el marco fiscal europeo, lo que reduce el componente de pánico que DE GRAUWE identificó en 2010-2012. La tercera es la escasez relativa de activos seguros en euros y la diversificación de los inversores, que favorecen a los emisores de la zona del euro con fundamentales en mejora.

La crítica a esta lectura sostiene que los mercados pueden estar siendo complacientes. Si el TPI actúa como un seguro implícito, las primas han dejado de ser una señal fiable de disciplina fiscal: es el problema de riesgo moral que se discutió al diseñar el instrumento. Además, las primas de riesgo cayeron también en 2005-2007, cuando la deuda española bajaba pero su vulnerabilidad aumentaba.\footnote{%
    Sobre el Instrumento para la Protección de la Transmisión (julio de 2022), sus criterios de elegibilidad y el debate sobre el riesgo moral ([\authorfont{Ver Tema 3.A.36, Tema 3.A.37}]).
}

Las agencias de calificación confirman la mejora. Tras las subidas de 2025, S\&P asigna a España una calificación A+ con perspectiva estable; Fitch, A; Moody's, A3, ambas estables; Morningstar DBRS, A (\textit{high}); y Scope la elevó a A+ con perspectiva estable el 25 de septiembre de 2026. En el caso de S\&P, la calificación de España coincide hoy con la de Francia, que perdió notas con Fitch y S\&P en septiembre y octubre de 2025, en ambos casos hasta A+. Italia sigue en la zona BBB+/Baa, y España e Italia son los dos grandes países cuyas notas han mejorado de forma más continuada desde 2023. Las agencias atribuyen las mejoras de España a su crecimiento, a la corrección de los desequilibrios exteriores y a la reducción del déficit, y condicionan nuevas subidas, que llevarían a la banda AA, a una reducción clara de la ratio de deuda, que sigue siendo su principal debilidad.

La historia de las calificaciones españolas ilustra el ciclo completo. España tuvo la máxima calificación (AAA/Aaa) hasta 2009-2010; Moody's la rebajó nueve escalones en menos de dos años, hasta Baa3 en junio de 2012, a un escalón del grado especulativo; y desde entonces ha recuperado posiciones de forma gradual, pero sigue lejos de la máxima calificación que tenía con una deuda del $36\%$. La distancia entre ambas notas resume la lección que ya hemos visto: las agencias, como los mercados, no detectaron la vulnerabilidad de 2007 y la penalizaron después con dureza.

\subsubsection{Una valoración comparada}
Cerrado el recorrido, conviene integrar las dimensiones anteriores en una valoración. La lectura favorable sostiene que, en términos relativos, España es hoy uno de los grandes países de la zona del euro con mejor trayectoria fiscal: su déficit está por debajo de la media y bajando, su saldo primario está cerca del equilibrio, su crecimiento es de los más altos, sus desequilibrios privados y exteriores se han corregido, su prima de riesgo está en mínimos de casi dos décadas y su calificación mejora. Frente a Francia, Bélgica o Austria, su posición ha mejorado de forma nítida, y la comparación con 2012 muestra un cambio de régimen.

La lectura crítica sostiene que esa mejora es relativa y en parte coyuntural. España tiene la quinta deuda más alta de la Unión y una de las mayores cargas de intereses en proporción al gasto; la Comisión la clasifica con riesgo alto a medio plazo, y sus necesidades brutas de financiación están entre las más elevadas de Europa; mantiene un déficit estructural, una brecha de ingresos frente a Europa y un gasto neto que crece por encima de lo comprometido; y, a diferencia de Portugal, Grecia o Chipre, no ha generado superávits primarios en la fase favorable del ciclo. Frente a la antigua periferia, su posición ha empeorado de forma igualmente nítida.

Ambas lecturas describen realidades distintas. La \textbf{posición relativa} de España ha mejorado porque otros países han empeorado (Francia, Bélgica) y porque su crecimiento ha sido alto; su \textbf{posición absoluta} sigue siendo vulnerable porque no ha corregido su déficit estructural. La comparación más exigente, y la más útil para la política fiscal, no es con Francia sino con Portugal: con una historia fiscal parecida y un shock mucho mayor en 2011-2014, Portugal ha reducido su deuda por debajo de la española generando superávits primarios. Esa es la vara de medir que aplicará el nuevo marco europeo, basado en la dinámica de la deuda.


\clearpage
\section{Política de financiación del Tesoro Público}
\puntosclave{
    La política de financiación del Tesoro no es un catálogo de instrumentos y procedimientos: responde a un problema concreto, financiar las necesidades del Estado al menor coste posible a medio y largo plazo con un riesgo prudente, y se evalúa con indicadores de coste y de riesgo. Desde 2012 el Tesoro ha alargado la vida media de la deuda, ha estandarizado sus procedimientos y ha diversificado su base inversora; el reto actual es sostener esa estrategia sin el BCE como comprador y con tipos que vuelven a subir.
}

Los dos primeros bloques han mostrado que el déficit es un rasgo estructural de las cuentas públicas españolas y que la deuda de las AAPP ronda hoy el $100\%$ del PIB. La trayectoria fiscal española es una versión tardía y comprimida del ciclo de las economías avanzadas: las crisis de 2008-2023 han llevado a España de estar entre los países con menos deuda de la zona del euro a estar entre los cinco con más, y el rasgo que la distingue de los que han reducido su deuda es la ausencia de superávits primarios en las recuperaciones. Con este diagnóstico, cabe preguntarse cómo se financia una deuda de esta dimensión.

El estudio de la forma de financiación se justifica por dos razones. La primera son sus efectos económicos, que remiten al vínculo entre la política de deuda, la política monetaria y la estabilidad financiera: la deuda pública es el activo sin riesgo de referencia de la economía nacional, su curva de rentabilidades sirve de base para fijar el precio de la deuda privada, y el volumen y los plazos de su emisión afectan a las condiciones financieras de toda la economía. La segunda es la eficiencia: una política de financiación adecuada reduce la carga de intereses. En 2025 las AAPP pagaron algo más de $40.000$ millones de euros en intereses, una cifra comparable a la recaudación del Impuesto sobre Sociedades, y cada punto básico de ahorro en el coste medio de una cartera de más de $1,5$ billones de euros equivale a unos $150$ millones anuales una vez renovada toda la cartera a ese coste.

Recordemos, además, la bisagra que establecimos al principio: el Tesoro no financia el déficit de las AAPP en contabilidad nacional, sino las \textbf{necesidades de endeudamiento del Estado} en términos de caja, que incluyen su propio déficit, los préstamos a las comunidades autónomas y a la Seguridad Social, las operaciones financieras y, sobre todo, los vencimientos de la deuda emitida en años anteriores.

Como veremos, desde la crisis de 2010-2012 el Tesoro ha seguido una estrategia coherente y reconocida internacionalmente: ha alargado la vida media de la deuda para reducir el riesgo de refinanciación, ha estandarizado y hecho predecibles sus procedimientos y ha diversificado su base inversora con instrumentos nuevos. Esa estrategia se desarrolló, sin embargo, en un entorno excepcional, en el que el BCE absorbía la mayor parte de la emisión neta y los tipos estaban en mínimos. El reto actual es mantenerla sin ese respaldo, con tipos que vuelven a subir en 2026 y con necesidades brutas cercanas al $16\%$ del PIB, de modo que el mercado privado absorba lo que antes absorbía el banco central sin encarecer en exceso la financiación. Para analizarlo, estudiaremos primero el marco institucional y los objetivos de la política de financiación; después, la evolución de las necesidades de financiación y sus determinantes; a continuación, los instrumentos y los procedimientos de emisión; y, por último, la relación con el BCE y la estrategia de 2026.

\subsection{Marco institucional y objetivos de la política de financiación}
Antes de estudiar las necesidades, los instrumentos y los procedimientos, conviene responder a tres preguntas previas: cuál es el marco jurídico que autoriza y limita el endeudamiento del Estado, qué órgano gestiona la deuda y con qué grado de autonomía, y qué objetivos persigue esa gestión y cómo se concilian. A ellas se añade una cuarta, más aplicada: con qué indicadores se mide el cumplimiento de esos objetivos.

Como veremos, la política de financiación del Tesoro se apoya en un marco jurídico de cuatro niveles, resultado de una transición histórica desde la monetización de los setenta y principios de los ochenta hasta una financiación plenamente de mercado. La gestión corresponde a la Secretaría General del Tesoro y Financiación Internacional, un modelo de unidad ministerial con alta autonomía técnica que la literatura comparada no considera inferior al de agencia si hay mandato claro, transparencia y capacidad técnica. Su objetivo, el menor coste a medio y largo plazo con un riesgo prudente, encierra una disyuntiva entre coste y riesgo cuyo fundamento teórico está en la nivelación impositiva y en las crisis de refinanciación autocumplidas. Los indicadores muestran una cartera larga, con poca deuda a corto plazo y un coste medio aún bajo pero en aumento.

\subsubsection{Marco jurídico: de la monetización a la financiación de mercado}
Hasta 1982, la financiación del déficit recayó en buena medida en el BdE: ese año, más del $80\%$ del déficit del Estado se financió con su crédito. El fuerte crecimiento del déficit complicaba una política monetaria orientada a reducir una inflación que superaba el $14\%$ en 1981. Por eso, desde 1982 la financiación del Estado se reorientó hacia mecanismos de mercado, lo que obligó al Tesoro a desarrollar una política de deuda: un conjunto de instrumentos que le permitieran financiar el déficit sin monetización y al menor coste posible.

Esa transición exigió, además, crear demanda donde no la había, porque los mercados de valores estaban poco desarrollados y apenas existían inversores institucionales. En 1984 se estableció un coeficiente de inversión obligatoria en pagarés del Tesoro, del $12\%$ de los pasivos computables de bancos y cajas, que se redujo desde 1989 y se eliminó en 1992. Durante los primeros años, el sistema bancario sustituyó al BdE como principal financiador del déficit. En 1990 se restringió legalmente el acceso del Tesoro al crédito del BdE y en 1994 se prohibió. Como vimos, la monetización alcanzó su máximo en 1978-1983 y terminó en torno a 1984 (ESCARIO, SABATÉ y GADEA, 2011); la transición pasó por instrumentos casi monetarios, como los pagarés, y por la represión financiera, y la financiación de mercado nació apoyada en una demanda obligatoria. Solo en los noventa, con la liberalización financiera, los incentivos a los no residentes y la creación de los creadores de mercado, pasó a ser una financiación verdaderamente de mercado.

En paralelo se construyó la infraestructura del mercado. En 1987 se crearon el Mercado de Deuda Pública en Anotaciones y la Central de Anotaciones del BdE, que eliminaron los títulos físicos; en 1988 se creó la figura de los creadores de mercado; en 1995 empezó a anunciarse el rango objetivo de colocación de las subastas; y en 2003, con la Ley 44/2002, el registro de la deuda anotada pasó a Iberclear.

El marco jurídico vigente se articula en cuatro niveles. El \textbf{primer nivel} es el constitucional y europeo. El artículo 149.1.14.ª de la Constitución atribuye al Estado la competencia exclusiva sobre la Hacienda general y la deuda del Estado. El artículo 135, tras la reforma de 2011, exige que la emisión de deuda esté autorizada por ley y establece que los créditos para pagar sus intereses y su capital se entenderán siempre incluidos en el estado de gastos y que su pago gozará de prioridad absoluta. En el ámbito europeo, el artículo 123 del TFUE prohíbe al BCE y a los bancos centrales nacionales conceder crédito a las administraciones públicas o adquirir directamente su deuda en el mercado primario; el artículo 124 prohíbe el acceso privilegiado de las administraciones a las entidades financieras, lo que impide reproducir coeficientes obligatorios como los de los ochenta; y el artículo 125 recoge la cláusula de no corresponsabilidad financiera (\textit{no bail-out}).\footnote{%
    Sobre la prohibición de financiación monetaria, su interpretación por el TJUE en las sentencias \textit{Gauweiler} (C-62/14) y \textit{Weiss} (C-493/17) y la cláusula de no corresponsabilidad ([\authorfont{Ver Tema 3.A.36, Tema 3.A.37}]) y ([\authorfont{Ver Tema 3.B.45}]).
}

El \textbf{segundo nivel} es la Ley 47/2003, General Presupuestaria, que regula la deuda del Estado en su título IV. Sus artículos 94 y 98 atribuyen al Gobierno, y por delegación al ministro competente en economía, la creación de deuda dentro del límite fijado en la Ley de Presupuestos, y definen las modalidades de emisión; la misma ley regula la gestión de la tesorería del Estado.\footnote{%
    Sobre la LGP, el ciclo presupuestario y la prórroga de los presupuestos (art. 134.4 CE y art. 38 LGP) ([\authorfont{Ver Tema 4.B.21}]).
}

El \textbf{tercer nivel} es la Ley de Presupuestos Generales del Estado, que fija cada año el límite al incremento del saldo vivo de la deuda del Estado. Como España funciona con los PGE de 2023 prorrogados, el límite vigente en 2024, 2025 y 2026 es el que fijó el artículo 46 de la Ley 31/2022, de PGE para 2023: un incremento del saldo vivo de la deuda del Estado de $96.021,98$ millones de euros. En la práctica el límite no ha sido restrictivo, porque la emisión neta de 2024-2026 se ha situado muy por debajo; pero ilustra que la autorización parlamentaria del endeudamiento lleva cuatro ejercicios fijada por una ley pensada para 2023.

El \textbf{cuarto nivel} es la normativa de desarrollo. Cada año, una orden ministerial dispone la creación de deuda del Estado y fija los instrumentos y procedimientos autorizados. Para 2026 es la Orden ECM/2/2026, de 9 de enero, que faculta a la Secretaría General del Tesoro y Financiación Internacional para emitir durante 2026 y enero de 2027. Autoriza letras de hasta $24$ meses, bonos de $2$ a $5$ años, obligaciones de más de $5$ años y sus variantes indexadas, y contempla como procedimientos la subasta, la sindicación, las operaciones con creadores de mercado y otras técnicas que se consideren adecuadas. A la orden se suman las resoluciones de la Secretaría General, que publican el calendario de subastas, regulan las condiciones de los creadores de mercado y presentan la Estrategia de Financiación anual.

\subsubsection{Organización: ¿ministerio o agencia de deuda?}
Desde la reestructuración ministerial de noviembre de 2023, las competencias de economía corresponden al Ministerio de Economía, Comercio y Empresa, cuya estructura orgánica desarrolla el Real Decreto 410/2024. La gestión de la deuda corresponde a la Secretaría General del Tesoro y Financiación Internacional, dependiente de la Secretaría de Estado de Economía y Apoyo a la Empresa, creada a finales de 2011 como Secretaría General del Tesoro y Política Financiera y con su denominación actual desde 2018. Dentro de ella, la Dirección General del Tesoro y Política Financiera gestiona la deuda y la tesorería del Estado, junto con la Dirección General de Financiación Internacional. Un detalle reciente apunta hacia un perfil más técnico: el Real Decreto 644/2026, de 28 de julio, permite que la titularidad de la Dirección General del Tesoro y Política Financiera recaiga en empleados públicos que no pertenezcan al subgrupo A1 de funcionarios de carrera, por la elevada especialización técnica y proyección estratégica del puesto.

Este modelo contrasta con el de otros países, que desde finales de los ochenta trasladaron la gestión de la deuda a agencias especializadas: Irlanda creó la \textit{National Treasury Management Agency} en 1990; el Reino Unido, la \textit{Debt Management Office} en 1998, tras conceder la independencia al Banco de Inglaterra, que hasta entonces gestionaba la deuda; Francia, la \textit{Agence France Trésor} en 2001, como servicio con competencia nacional dentro del Ministerio de Economía; y Alemania, la \textit{Finanzagentur} en 2000-2001, como sociedad pública de responsabilidad limitada. Portugal tiene el IGCP y Suecia la \textit{Riksgälden}, con orígenes en el siglo XVIII. Italia y España, en cambio, mantienen la gestión dentro de una unidad ministerial.

¿Qué modelo es preferible? Los argumentos a favor de la agencia son cuatro: (i) la especialización técnica y la capacidad de contratar perfiles de mercado con retribuciones competitivas; (ii) la separación entre las decisiones de gestión de la deuda y las presiones políticas de corto plazo, como la tentación de emitir a corto plazo para abaratar el coste inmediato a costa de más riesgo; (iii) un mandato explícito y medible que facilita la rendición de cuentas, en la línea de la independencia de los bancos centrales; y, (iv) una separación nítida respecto del banco central. A favor de la gestión ministerial se aducen la mejor coordinación con la política presupuestaria y con la gestión de la tesorería, la responsabilidad política directa ante el Parlamento y el menor coste administrativo. La conclusión de la literatura comparada (CURRIE, DETHIER y TOGO, 2003; FMI y BANCO MUNDIAL, 2014) es que no hay un modelo institucional superior: lo decisivo es que existan un mandato claro, objetivos explícitos de coste y riesgo, una separación operativa respecto de la política monetaria, transparencia y capacidad técnica. El Tesoro español es un caso habitualmente citado de unidad ministerial con elevada autonomía técnica y reputación en los mercados, que opera guiado por criterios financieros.

Hay, sin embargo, un rasgo específico del caso español. Desde 2018, con la separación entre el Ministerio de Hacienda (presupuesto, ingresos y financiación territorial) y el de Economía (Tesoro), la gestión de la deuda y la gestión presupuestaria dependen de ministerios distintos. Esto aproxima de hecho el modelo español al de agencia en su separación de la política presupuestaria, aunque sin su estatuto jurídico, y exige coordinación en tres frentes: las previsiones de ingresos y pagos de caja, que determinan las necesidades; la financiación de los mecanismos autonómicos, gestionados por Hacienda pero financiados por el Tesoro; y los préstamos a la Seguridad Social.

El principio de separación entre gestión de la deuda y política monetaria, uno de los pilares de las directrices internacionales, también está en revisión. El gestor de la deuda decide plazos y volúmenes con criterios de coste y riesgo, y el banco central fija los tipos con criterios de estabilidad de precios. Pero la expansión cuantitativa de 2009-2022 ha cuestionado ese principio en la práctica: cuando el banco central compra deuda a largo plazo para reducir la prima por plazo y el Tesoro alarga la vida media para aprovechar los tipos bajos, ambas políticas actúan en direcciones opuestas sobre la duración de la deuda en manos privadas. GREENWOOD, HANSON, RUDOLPH y SUMMERS (2014) documentaron ese conflicto en Estados Unidos, y BLOMMESTEIN y TURNER (2012) advirtieron del riesgo de que una política se subordine a la otra en contextos de dominancia fiscal. En la zona del euro, con un solo banco central y veinte gestores de deuda, el problema tiene una dimensión adicional: el Tesoro español alargó la vida media mientras el Eurosistema absorbía más de la mitad de su emisión neta, una cuestión que retomaremos al final del bloque.\footnote{%
    Sobre la interacción entre la política monetaria no convencional y la gestión de la deuda (efecto de las compras de activos sobre la prima por plazo y canal de reequilibrio de carteras) ([\authorfont{Ver Tema 3.A.36, Tema 3.A.37}]).
}

\subsubsection{Objetivos de la política de financiación y su fundamento teórico}
La formulación de referencia internacional del objetivo de la gestión de la deuda es la de las directrices del FMI y el Banco Mundial (2014): asegurar que las necesidades de financiación del Estado y sus obligaciones de pago se cubran \textbf{al menor coste posible a medio y largo plazo, con un grado de riesgo prudente}. El Tesoro español la hace suya en sus estrategias, a las que añade la contribución al buen funcionamiento del mercado de deuda.

La formulación contiene tres elementos. El primero es el \textbf{horizonte}: el coste relevante no es el del año, sino el de medio y largo plazo. Emitir todo a muy corto plazo abarata normalmente el coste inmediato, porque la curva de tipos suele tener pendiente positiva, pero expone al Estado a refinanciar toda su deuda a los tipos que haya en el futuro. El segundo es el \textbf{riesgo}, en varias dimensiones: el de refinanciación (no poder renovar los vencimientos, o hacerlo en condiciones muy desfavorables); el de tipo de interés (la sensibilidad del coste de la cartera a las variaciones de los tipos); el de tipo de cambio, hoy prácticamente inexistente porque toda la deuda está en euros, pero relevante en los noventa; el de inflación, que afecta a los bonos indexados; el de liquidez (la dificultad de colocar deuda en un mercado poco profundo); y el operativo. El tercero es la \textbf{disyuntiva entre coste y riesgo}: alargar la vida media reduce el riesgo de refinanciación y el de tipo de interés, pero normalmente eleva el coste esperado, porque los inversores exigen una prima por plazo. El gestor elige un punto de esa frontera en función de su aversión al riesgo, y ese punto es el núcleo de su estrategia; el FMI y el Banco Mundial lo han formalizado en su marco de estrategia de gestión de la deuda a medio plazo (MTDS), que compara carteras alternativas por su coste esperado y su coste en escenarios adversos.

¿Por qué los gestores prefieren la deuda larga y por qué emiten bonos indexados? Detrás del objetivo canónico hay una literatura teórica que lo explica.\footnote{%
    Sobre la teoría de la gestión óptima de la deuda (nivelación impositiva, deuda contingente, estructura de plazos e indexación) y la de las crisis de deuda autocumplidas ([\authorfont{Ver Tema 3.A.39}]). Con impuestos distorsionantes de coste convexo, el Gobierno minimiza la pérdida de eficiencia manteniendo constante el tipo impositivo esperado, lo que convierte a la deuda en el amortiguador de las perturbaciones presupuestarias.
} El punto de partida es la \textbf{nivelación impositiva} de BARRO (1979): si los impuestos generan distorsiones crecientes con el tipo, el Gobierno debe mantener tipos estables en el tiempo y absorber las perturbaciones presupuestarias con la deuda. De ahí se deriva la pregunta de qué estructura de deuda protege mejor al presupuesto. LUCAS y STOKEY (1983) mostraron que, idealmente, el Gobierno emitiría deuda contingente, cuyo rendimiento bajaría cuando el presupuesto sufre una perturbación adversa y subiría en caso contrario. Como esa deuda no existe, la literatura ha buscado aproximaciones: ANGELETOS (2002) y BUERA y NICOLINI (2004) mostraron que una combinación adecuada de plazos puede replicar esa contingencia, porque el valor de la deuda larga cae cuando suben los tipos, y BARRO (2003) concluyó que la mejor aproximación es la deuda larga indexada a la inflación. Es el fundamento del programa de bonos indexados del Tesoro, cuyo atractivo se basa en que los ingresos públicos están fuertemente ligados a la inflación, aunque la cobertura se debilita cuando también se indexa el gasto, como ocurre con las pensiones desde 2021.

Una segunda línea destaca la \textbf{credibilidad}. CALVO y GUIDOTTI (1990) y MISSALE y BLANCHARD (1994) mostraron que un Gobierno con mucha deuda nominal a largo plazo tiene incentivos a generar inflación para reducir su valor real, de modo que, con baja credibilidad antiinflacionista, la deuda corta o indexada es una señal de compromiso. En la zona del euro, sin política monetaria nacional, el argumento pierde fuerza: España no puede inflar su deuda. Una tercera línea, la más relevante para la experiencia española de 2010-2012, es la de las \textbf{crisis de refinanciación autocumplidas}. COLE y KEHOE (2000) mostraron que un país con mucha deuda a corto plazo es vulnerable a una crisis en la que los inversores, temiendo que no pueda refinanciarse, se niegan a hacerlo y provocan así el impago que temían. Alargar la vida media reduce la proporción de deuda que hay que refinanciar cada año y, con ella, esa vulnerabilidad. ARELLANO y RAMANARAYANAN (2012) formalizaron la disyuntiva: la deuda corta disciplina al Gobierno y es más barata, y la larga lo protege frente a los pánicos. Esta literatura es el fundamento de la estrategia seguida por el Tesoro desde 2012: alargar la vida media hasta los $8$ años.

Al objetivo principal se añaden otros que el Tesoro formula en sus estrategias. El primero es la \textbf{liquidez} y el buen funcionamiento del mercado: el Tesoro concentra la emisión en pocas referencias de gran volumen, en torno a $25.000$ millones para bonos y obligaciones, porque una deuda más líquida es más atractiva y más barata, y porque su curva es la referencia de la deuda privada. El segundo es la \textbf{transparencia y la previsibilidad}: el Tesoro publica en enero su estrategia anual, con el calendario de subastas y los rangos objetivo de colocación, porque la previsibilidad reduce la incertidumbre de los inversores y, con ella, la prima que exigen. El tercero es la \textbf{diversificación de la base inversora}, porque una base amplia y variada reduce el riesgo de que la retirada de un grupo de inversores desestabilice la financiación, como ocurrió en 2011-2012, cuando la salida de los no residentes dejó a la banca nacional como principal comprador. El cuarto, más reciente, es la \textbf{sostenibilidad ambiental}, que el programa de bonos verdes vincula a gastos elegibles del presupuesto.

Estos objetivos no siempre son compatibles. La liquidez exige pocas referencias grandes, mientras que la diversificación empuja a crear instrumentos nuevos, como los indexados y los verdes, que fragmentan la oferta y reducen la liquidez de cada segmento. La previsibilidad exige atenerse al calendario, mientras que la minimización del coste invitaría a aprovechar las ventanas de mercado favorables. Y el alargamiento de la vida media reduce el riesgo pero, con una curva positiva, eleva el coste esperado. La estrategia de cada año es, en esencia, la resolución de estas tensiones.

\subsubsection{Indicadores de coste y riesgo de la cartera española}
Los objetivos anteriores se traducen en un conjunto de indicadores que el Tesoro publica y que permiten evaluar su gestión. Los indicadores de \textbf{coste} son dos. El coste medio de emisión mide el tipo al que se ha emitido la deuda en el año: en 2021 fue negativo por primera vez en la historia del Tesoro ($-0,04\%$), y con la subida de tipos pasó al $3,44\%$ en 2023, al $3,16\%$ en 2024 y al $2,70\%$ en 2025. El coste medio de la deuda en circulación mide el tipo medio del conjunto de la cartera: alcanzó su mínimo histórico en 2021 ($1,64\%$) y desde entonces sube de forma gradual, hasta el $2,31\%$ en 2025, a medida que la deuda antigua y barata se refinancia a tipos más altos. La diferencia entre ambos indicadores es la clave de la dinámica de la carga de intereses: mientras el coste de emisión supere al coste medio de la cartera, este seguirá subiendo aunque los tipos de mercado se estabilicen.

Los indicadores de \textbf{riesgo} son varios. La vida media de la deuda del Estado alcanzó por primera vez los $8$ años en 2021 y se ha mantenido en torno a ese nivel ($7,93$ años a finales de 2025, frente a $8,9$ en la zona del euro según el BdE). La proporción de letras, que llegó al $62\%$ de la deuda en 1990, se sitúa hoy en torno al $7\%$, y el $13,9\%$ de la deuda de las AAPP vencía en menos de un año en 2024. Completan el cuadro el perfil de vencimientos, es decir, la concentración de amortizaciones en determinados años, y la composición de la base inversora: el Eurosistema tenía en torno a una cuarta parte de la deuda en 2025 y los no residentes el $47,8\%$. A estos indicadores la literatura añade el plazo medio hasta la revisión del tipo (ATR, por sus siglas en inglés), que mide cuánto tarda la cartera en trasladar un cambio de tipos al coste. En España coincide prácticamente con la vida media, porque casi toda la deuda es a tipo fijo, con la excepción de los bonos indexados. Es el dato que explica por qué la subida de tipos de 2022-2023 tardó en trasladarse a la carga de intereses y por qué seguirá trasladándose en 2026-2030, amplificada por las subidas del BCE de junio y septiembre de 2026.

\subsection{Necesidades de financiación del Tesoro: evolución y determinantes}
En los últimos quince años, las necesidades de financiación del Tesoro han aumentado mucho. En 2007 la necesidad bruta fue de unos $50.000$ millones de euros y la emisión neta, negativa en casi $10.000$ millones; desde entonces, la emisión neta ha sido positiva sin interrupción y superó los $100.000$ millones en 2009 y en 2020. Para entender esta evolución conviene distinguir tres magnitudes. Las necesidades \textbf{netas} equivalen a la variación del saldo de deuda del Estado y cubren su déficit de caja, los préstamos a las comunidades y a la Seguridad Social, las operaciones financieras y la variación de la tesorería, menos la financiación europea. Las \textbf{brutas} añaden los vencimientos, que dependen del saldo y de la vida media. Y la \textbf{oferta neta al sector privado} añade a las netas los vencimientos que el Eurosistema ya no reinvierte.

Como veremos, la evolución muestra una explosión en 2008-2013, una normalización en 2014-2019 con el BCE absorbiendo más que la emisión neta, un máximo de $277.000$ millones de emisión bruta en 2020 y, desde 2021, una emisión neta que baja hasta estabilizarse en $55.000$ millones mientras la bruta sube hasta unos $285.700$ millones en 2026, en torno al $16\%$ del PIB. El rasgo estructural del caso español es que el Estado actúa como prestamista de las demás administraciones, sobre todo de la Seguridad Social; y el cambio más relevante desde 2022 es que la oferta neta al sector privado ha pasado de ser negativa a rondar probablemente los $90.000$-$100.000$ millones anuales, absorbida sobre todo por inversores no residentes.

\subsubsection{Conceptos: necesidades netas, necesidades brutas y su medición}
Como vimos, el Tesoro financia las necesidades de endeudamiento del Estado, una magnitud de caja referida solo al Estado. La emisión neta, es decir, el aumento del saldo vivo de la deuda del Estado en el año, cubre la suma de cuatro componentes: (i) el \textbf{déficit de caja del Estado}, que difiere del déficit del Estado en contabilidad nacional por los ajustes de devengo y por la periodificación de los pagos; (ii) los \textbf{préstamos netos del Estado a otras administraciones}, principalmente a las comunidades autónomas, a través del Fondo de Financiación a CCAA, y a la Seguridad Social; (iii) las demás \textbf{operaciones financieras} del Estado (aportaciones a fondos y organismos, préstamos a otros países o instituciones, adquisición de participaciones y, en sentido contrario, los cobros por amortización de préstamos concedidos o por venta de activos); y, (iv) la \textbf{variación del saldo de tesorería}, el colchón de liquidez del Estado en el BdE y en el mercado. A esa suma hay que restar la financiación que el Estado obtiene fuera del mercado, que en los últimos años ha procedido sobre todo de la Unión: los préstamos del instrumento SURE y los del Mecanismo de Recuperación. Las subvenciones del Mecanismo de Recuperación no son financiación, sino ingresos del Estado, pero su efecto de caja es el mismo: reducen lo que el Tesoro necesita emitir.

Las necesidades \textbf{brutas} añaden a las netas los vencimientos de la deuda emitida en años anteriores. En la estrategia de 2026, la emisión neta prevista es de $55.000$ millones: $50.000$ en bonos y obligaciones y $5.000$ en letras. La emisión bruta de bonos y obligaciones asciende a $176.935$ millones, resultado de sumar a la neta unos $126.935$ millones de amortizaciones, y la de letras ronda los $108.700$ millones, porque cada año se renueva todo el saldo de letras. La emisión bruta total asciende así a $285.677$ millones. Los vencimientos dependen del saldo vivo y de su vida media: con una vida media de unos $8$ años, cada año vence en promedio alrededor de una octava parte de la deuda a medio y largo plazo, lo que concuerda con el $13\%$ del saldo que el Tesoro refinancia anualmente. Por eso las necesidades brutas crecen de forma casi automática con el saldo de deuda aunque la emisión neta sea constante, como muestra el paso de una emisión bruta de $232.570$ millones en 2022 a unos $285.700$ en 2026 con una emisión neta decreciente.

La emisión bruta del Estado prevista para 2026 equivale a en torno al $16\%$ del PIB nominal, una proporción menor que la de hace unos años porque el PIB nominal ha crecido más deprisa que los vencimientos. Conviene distinguirla, sin embargo, de las necesidades brutas de financiación de las AAPP, que incluyen además el déficit y los vencimientos de las comunidades y de las entidades locales que se financian en el mercado, y que la Comisión sitúa cerca del $20\%$ del PIB, entre las más altas de la Unión. La cifra importa por su relación con el riesgo de refinanciación: el marco de análisis de sostenibilidad del FMI para países con acceso a los mercados utiliza como referencia para las economías avanzadas unas necesidades brutas del $15\%$ del PIB, por encima de las cuales la vulnerabilidad a un cierre de los mercados aumenta de forma significativa. España supera ese umbral en términos de AAPP y lo roza con el Estado solo, y es uno de los motivos por los que la Comisión y el FMI insisten en mantener una vida media larga: cada año de vida media adicional reduce los vencimientos anuales y, con ellos, las necesidades brutas.\footnote{%
    Las necesidades brutas de financiación equivalen al déficit más las amortizaciones de la deuda a medio y largo plazo y el saldo de la deuda a corto plazo. Sobre los indicadores de sostenibilidad y de riesgo de refinanciación ([\authorfont{Ver Tema 3.A.39}]).
}

\subsubsection{Determinantes: por qué el Tesoro emite más de lo que dice el déficit}
El rasgo más característico de las necesidades del Tesoro español es que financia no solo el déficit del Estado, sino una parte del de las comunidades autónomas y de la Seguridad Social. Desde 2012, la mayoría de las comunidades se financian a través de los mecanismos estatales (FFPP, FLA y, desde 2015, FFCCAA), que movilizaron $126.895$ millones entre 2012 y 2024, en lugar de acudir al mercado; y la Seguridad Social cubre su déficit desde 2017 con préstamos del Estado, además de con transferencias. En 2025, la deuda de la Seguridad Social, casi toda con el Estado, aumentó en unos $10.000$ millones, y la de las comunidades en unos $5.700$ millones.

Esto tiene tres consecuencias. La primera es contable: los préstamos del Estado a otras administraciones no alteran la deuda consolidada de las AAPP, pero sí aumentan la deuda del Estado y lo que el Tesoro emite; por eso la deuda del Estado representa casi el $90\%$ del total. La segunda es financiera: las comunidades y la Seguridad Social se financian al coste del Tesoro, inferior al que obtendrían en el mercado, lo que reduce el coste agregado. La tercera, la más debatida, es de incentivos: los mecanismos debilitan la restricción presupuestaria de las comunidades, que ya no afrontan el juicio del mercado, y generan el riesgo moral asociado a la restricción presupuestaria blanda.\footnote{%
    Sobre los mecanismos de financiación autonómica, la restricción presupuestaria blanda y la condonación ([\authorfont{Ver Tema 4.B.25}]); sobre los préstamos del Estado a la Seguridad Social ([\authorfont{Ver Tema 4.B.24}]).
} La condonación de $83.252$ millones de deuda autonómica en tramitación no aumentará la emisión del Tesoro, porque la mayor parte de la deuda condonada es deuda con el propio Estado a través del FFCCAA: convertirá un derecho de cobro del Estado en nada. Su efecto sobre las necesidades será diferido, porque el Estado dejará de recibir las amortizaciones que las comunidades habrían devuelto, lo que aumentará marginalmente sus necesidades netas en los próximos años.

El segundo determinante son los fondos europeos, que han actuado como financiación fuera del mercado. Los préstamos del SURE ($21.326$ millones en 2020-2021) y las transferencias del programa \textit{Next Generation EU} ($31.036$ millones en 2021-2022) contuvieron las necesidades de emisión, y a ellos se han sumado los préstamos del Mecanismo de Recuperación: unos $16.000$ millones movilizados en 2025 y $6.500$ millones solicitados para 2026. Conviene hacer dos precisiones. La primera es que los préstamos europeos son deuda del Estado español a todos los efectos: reducen la emisión en el mercado, pero no la deuda. Su ventaja es un coste cercano al de la Unión, que se financia a tipos inferiores a los españoles, y unos plazos largos; por eso, cuando el diferencial entre España y la Unión se estrechó, se solicitó solo una parte de los préstamos disponibles, porque el ahorro dejó de compensar la condicionalidad y la gestión asociadas. La segunda es que la financiación europea desaparece como amortiguador a partir de 2027, cuando los préstamos del SURE y del Mecanismo empiecen a amortizarse y dejen de llegar subvenciones. A ello se suma la amortización del préstamo del MEDE de 2012, que termina de devolverse en 2027 (unos $3.643$ millones en 2026).

El tercer determinante, que suele omitirse, es la variación del saldo de tesorería del Estado. El Tesoro mantiene un colchón de liquidez para cubrir los desfases entre cobros y pagos a lo largo del año, afrontar los grandes vencimientos y no depender de una sola subasta en un momento de tensión. En 2020 lo aumentó sustancialmente como precaución, lo que elevó la emisión neta por encima de lo que exigía el déficit, y en años posteriores lo ha reducido, lo que le ha permitido emitir menos. Las reglas del BCE influyen en esa gestión: desde diciembre de 2024, los depósitos de las administraciones en los bancos centrales se remuneran como máximo al ESTR menos $20$ puntos básicos, lo que incentiva a los Tesoros a colocar sus excedentes en el mercado. Estas variaciones de tesorería aparecen en la deuda como ajustes déficit-deuda: aumentan la deuda bruta sin aumentar el déficit.

Con estos elementos puede ensayarse, en orden de magnitud, una descomposición de la emisión neta de 2026 que ilustra la bisagra entre las dos partes del tema. De los $55.000$ millones, en torno a la mitad correspondería al déficit de caja del Estado propiamente dicho; una parte importante, del orden de $10.000$ a $15.000$ millones, a los préstamos a la Seguridad Social; otra, menor y variable, a la financiación neta de las comunidades a través del FFCCAA, una vez descontadas sus amortizaciones; y el resto, a operaciones financieras, mientras los préstamos europeos y la amortización del MEDE actúan en sentido contrario. Lo relevante es la estructura: la Seguridad Social se ha convertido en uno de los principales destinos de la emisión neta del Tesoro, que es la manifestación financiera del problema de las pensiones.

\subsubsection{Evolución, 2007-2026}
En 2007, con superávit, el Estado amortizaba deuda: la necesidad bruta fue de unos $50.000$ millones y la emisión neta, negativa en casi $10.000$ millones. Con la Gran Recesión, la emisión neta superó los $100.000$ millones en 2009, y el Tesoro, con un acceso difícil a los mercados, aumentó las emisiones a corto plazo, más baratas y más fáciles de colocar en condiciones de tensión: ese año cubrió casi el $30\%$ de sus necesidades netas con letras, que llegaron a representar el $18\%$ de la deuda del Estado, y emitió incluso un bono a tipo variable (Euribor a doce meses menos $10$ puntos básicos) para diversificar la base inversora, una modalidad que hoy no utiliza. Entre 2010 y 2013 las necesidades se mantuvieron muy elevadas por tres motivos: los déficits siguieron por encima del $9\%$ del PIB; desde 2012 el Estado empezó a prestar a las comunidades a través de los mecanismos de liquidez; y hubo que refinanciar vencimientos crecientes en el peor momento de la crisis soberana. Parte de la financiación no pasó por el Tesoro: los $41.333$ millones del MEDE para la recapitalización bancaria se canalizaron a través del FROB, en forma de bonos del propio MEDE, y computaron como deuda de las AAPP. En 2012 la demanda se sostuvo gracias a la banca nacional, que absorbió deuda con la liquidez de las operaciones a tres años del BCE.

En 2014-2019, la emisión neta se redujo de forma gradual con el déficit, mientras la bruta se mantenía elevada por los vencimientos. El aumento de las necesidades coincidió con una caída sin precedentes del coste de emisión, favorecida por el cambio del patrón de crecimiento hacia uno más equilibrado, que mejoró la calificación y redujo la prima de riesgo, y por los estímulos del BCE, cuyas compras netas de deuda española superaron los $500.000$ millones entre 2015 y 2022, más de la tercera parte de la deuda de las AAPP a finales de 2022. En varios años, esas compras superaron la emisión neta del Tesoro, de modo que la deuda en manos del sector privado disminuyó: descontada la deuda en manos del BCE, la deuda de las AAPP a finales de 2022 equivalía a la de principios de 2013.

En enero de 2020, antes de la pandemia, el Tesoro esperaba emitir en bruto menos de $200.000$ millones por primera vez en casi diez años. La caída de los ingresos y el aumento del gasto obligaron a emitir $277.000$ millones, la cifra más alta de su historia hasta entonces, con una emisión neta superior a los $100.000$ millones. A diferencia de 2009, el PEPP absorbió una parte sustancial de la emisión neta de 2020-2021, el coste de emisión siguió cayendo y el Tesoro mantuvo incluso su estrategia de reducir las letras.

Desde 2021 la emisión neta se ha reducido de forma escalonada: $70.063$ millones en 2022, unos $65.000$ en 2023 y $55.000$ en 2024, 2025 (el plan inicial de $60.000$ se recortó en septiembre de 2025 por la mejora de los ingresos) y 2026. En cambio, la emisión bruta ha crecido: $232.570$ millones en 2022, unos $252.000$ en 2023, unos $257.600$ en 2024, unos $274.200$ en 2025 y $285.677$ en 2026. La divergencia se explica por los vencimientos, porque el saldo crece cada año en la cuantía de la emisión neta y vencen las abultadas emisiones de 2016-2021. En la presentación de la Estrategia de 2026, el Tesoro destacó que había mantenido la vida media en torno a los $8$ años por quinto año consecutivo ($7,93$ años a finales de 2025); que el coste medio de la deuda en circulación era del $2,31\%$ en 2025, $10$ puntos básicos más que en 2024; que el coste medio de la nueva emisión, del $2,70\%$, era $74$ puntos básicos inferior al de 2023; y que los no residentes tenían el $47,8\%$ de la deuda, el segundo registro más alto de la serie. La brecha entre el coste de emisión y el de la cartera implica que el coste medio seguirá subiendo, más aún tras las subidas del BCE de 2026.

\subsubsection{Oferta neta al sector privado y las presiones de los próximos años}
La emisión neta mide cuánto aumenta el saldo de deuda del Estado, pero no cuánta deuda nueva tiene que absorber el mercado privado. Cuando el BCE compra deuda en el mercado secundario por encima de la emisión neta, como en 2015-2021, la oferta neta al sector privado es negativa: los inversores privados reducen sus tenencias. Cuando deja de reinvertir los vencimientos de sus carteras, la deuda que vence en su poder tiene que refinanciarse en el mercado, y la oferta neta al sector privado se aproxima a la emisión neta más los vencimientos no reinvertidos del Eurosistema. Con el Eurosistema en torno a una cuarta parte de la deuda española en 2025 y sin reinversiones del APP desde julio de 2023 ni del PEPP desde finales de 2024, esa oferta neta rondaría los $90.000$-$100.000$ millones anuales en 2025-2026, frente a cifras negativas en 2015-2021.

Este cambio explica que los inversores internacionales se hayan consolidado como los principales compradores de bonos y obligaciones del Estado y que la proporción de no residentes haya subido al $47,8\%$. La emisión sindicada a diez años de enero de 2026, de $15.000$ millones, recibió una demanda récord de $145.000$ millones; los no residentes tomaron el $90\%$ y los bancos centrales e instituciones oficiales fueron el primer grupo inversor ($26,5\%$). La demanda ha respondido, pero a cambio de una base inversora más sensible a las condiciones financieras globales y a la comparación con otros emisores.

Hacia delante, las necesidades brutas se mantendrán elevadas por cinco factores: (i) los \textbf{vencimientos}, que seguirán creciendo con el saldo y con la llegada a vencimiento de las grandes emisiones de la etapa de tipos cero, que se refinanciarán a tipos más altos; (ii) el \textbf{fin de la financiación europea}, pues las subvenciones del Mecanismo de Recuperación terminan en 2026 y desde 2027-2028 empezarán a amortizarse los préstamos del SURE y del Mecanismo; (iii) los \textbf{préstamos a la Seguridad Social}, que tenderán a crecer con el gasto en pensiones mientras no se corrija su desequilibrio; (iv) la \textbf{financiación autonómica}, cuyos flujos alterarán la condonación y la eventual reforma del modelo; y, (v) los \textbf{nuevos gastos}, en particular las medidas frente a la crisis de Oriente Medio y el compromiso de defensa, que, sin presupuestos aprobados, se canalizan por créditos extraordinarios y aumentan el déficit de caja del Estado. En sentido contrario, el crecimiento nominal y la buena evolución de la recaudación han permitido mantener estable la emisión neta. Si se cumple la previsión de la AIReF de un déficit de las AAPP del $2,6\%$ en 2026, la emisión neta podría mantenerse en niveles parecidos; el riesgo está en la combinación de un déficit que no baja y unos tipos que suben, que convertiría el coste de la deuda en el principal motor de las necesidades futuras.

\subsection{Instrumentos de financiación del Tesoro}
La gama de instrumentos del Tesoro responde a una idea de diseño: atender las preferencias de los inversores, pero evitando la segmentación del mercado que produciría un número excesivo de instrumentos, porque esa segmentación reduciría la liquidez de la deuda y, con ella, su atractivo. Como veremos, la gama se construyó desde 1982 sustituyendo los instrumentos con ventajas fiscales o demanda obligatoria por instrumentos de mercado, y el resultado es una cartera dominada por los bonos y obligaciones nominales a tipo fijo, que superan el $90\%$ del total, con una vida media de unos $8$ años frente a los menos de $2$ de 1990. Las letras cumplen funciones de gestión de tesorería, de amortiguador en las crisis y de activo monetario; los programas de diversificación (indexados y verdes) son deliberadamente pequeños; la deuda en divisas ha desaparecido con el euro; y los préstamos oficiales europeos han amortiguado las necesidades de mercado en las crisis. La diversificación relevante de los últimos años se ha producido, en suma, en la base inversora y no en los instrumentos.

\subsubsection{La construcción de la gama: liquidez frente a segmentación}
Cuando en 1982 comenzó el giro hacia la financiación de mercado, el Tesoro disponía solo de dos instrumentos: la deuda desgravable del Estado, un instrumento a medio plazo dirigido a los contribuyentes del IRPF cuya demanda se basaba sobre todo en su ventaja fiscal; y los pagarés del Tesoro, activos a corto plazo (a seis meses en un primer momento y a dieciocho después) creados por el BdE como instrumento de regulación de la liquidez que, al no tener restricciones de acceso, servían también para financiar al Estado. A finales de 1982 se crearon los bonos del Estado, de $2$ a $5$ años, y en 1983 las obligaciones del Estado, a plazos más largos, ya sin ventajas fiscales; en 1987 se crearon las letras del Tesoro, que sustituyeron a los pagarés. La secuencia refleja la transición del marco jurídico: los primeros instrumentos dependían de ventajas fiscales o de demanda obligatoria, y los que aparecen desde 1982 se basan en su rentabilidad de mercado.

La gama actual responde a cuatro principios. El primero es la \textbf{concentración en pocos instrumentos estandarizados}: un mercado con pocas referencias grandes es más líquido que uno con muchas pequeñas, y la liquidez reduce la prima que exigen los inversores. Por eso el Tesoro aplica desde finales de los ochenta el sistema de \textit{agregación de emisiones}, que reabre sucesivamente cada referencia hasta que alcanza un volumen que garantiza su liquidez, en torno a $25.000$ millones para bonos y obligaciones. El segundo es la \textbf{curva completa de referencias}: el Tesoro mantiene referencias a $3$, $5$, $10$, $15$, $20$, $30$ y $50$ años, y cuando emite una nueva referencia a un plazo esta se convierte en el nuevo \textit{benchmark}, mientras la anterior pasa a ser una referencia \textit{off-the-run}. Una curva completa y líquida no solo abarata la financiación del Tesoro: sirve de referencia para valorar la deuda de las comunidades, de las empresas y de los bancos, lo que convierte la política de emisión en un bien público para el sistema financiero. El tercero es la \textbf{adaptación a la demanda sin fragmentar la oferta}, cuyo mejor ejemplo es la segregación: todos los bonos y obligaciones son segregables desde 1997, de modo que los creadores de mercado pueden separar cada cupón y el principal en bonos de cupón cero (\textit{strips}) y reconstituir después el bono original. Los \textit{strips} atienden a inversores que necesitan flujos únicos en fechas concretas, como las aseguradoras y los fondos de pensiones que cubren pasivos, o que buscan su ventaja fiscal (sus rendimientos no están sujetos a retención en el Impuesto sobre Sociedades), sin que el Tesoro tenga que emitir instrumentos nuevos; su creación respondió a la proximidad de la UEM, cuando la deuda española tenía que competir con la de los socios europeos. El cuarto es la \textbf{diversificación de la base inversora} mediante programas específicos, los indexados (2014) y los verdes (2021), que entran en tensión con el primer principio y por eso tienen un tamaño limitado.

El resultado es una transformación radical de la composición de la cartera. En 1990, las letras llegaron a representar el $62\%$ de la deuda en circulación y los bonos y obligaciones el $30\%$; hoy los bonos y obligaciones superan el $90\%$ y las letras rondan el $7\%$. Los demás instrumentos (deuda en divisas y préstamos), que a finales de los noventa llegaron a superar el $15\%$, son hoy menos del $4\%$. Con ello, y con la emisión a plazos cada vez más largos, la vida media de la deuda del Estado pasó de menos de $2$ años en 1990 a superar los $8$ en 2021, un alargamiento que las agencias de calificación citaron entre los argumentos de sus subidas de 2014-2019.

\subsubsection{Letras del Tesoro: la flexibilidad a corto plazo}
Las letras del Tesoro son valores de renta fija a corto plazo, emitidos al descuento y sin cupón, cuya rentabilidad es implícita: la diferencia entre el precio de compra y el valor nominal. En los años de tipos negativos llegaron a emitirse a premio, con rentabilidad negativa. Se introdujeron en 1987 a doce meses; cuando en 1990 se restringió el acceso del Tesoro al crédito del BdE, se crearon letras a tres y seis meses para gestionar los desfases transitorios de tesorería. Hoy se emiten a $3$, $6$, $9$ y $12$ meses, aunque para maximizar su liquidez cada mes se emite una nueva letra a doce meses, por un importe inicial elevado, que se reabre cuando su vida residual es de nueve, seis y tres meses: en la práctica solo hay letras a doce meses, y en todo momento circulan doce referencias. La inversión mínima es de $1.000$ euros, y la Orden de creación de deuda de 2026 autoriza letras de hasta $24$ meses, un margen que no se ha utilizado en los últimos años pero que da flexibilidad ante una situación de tensión.

Las letras cumplen tres funciones que justifican mantener una proporción de deuda a corto plazo, aunque aumente el riesgo de refinanciación: la \textbf{gestión de tesorería}, porque permiten ajustar la emisión a los desfases de caja sin alterar el programa de bonos y obligaciones; el papel de \textbf{amortiguador en las crisis}, porque cuando el acceso al mercado se complica pueden aumentarse con más facilidad, como en 2009 y en 2020; y la de \textbf{activo del mercado monetario}, útil para fondos monetarios, tesorerías de empresas y bancos. A ello se añade un argumento de coste: GREENWOOD, HANSON y STEIN (2015) sostienen que la deuda pública a corto plazo tiene un \textbf{rendimiento de conveniencia}, porque los inversores la valoran como un sustituto del dinero y aceptan por ella una rentabilidad inferior a la que justificaría su riesgo, de modo que al Estado le convendría emitir deuda corta hasta que el ahorro marginal iguale el coste marginal del mayor riesgo de refinanciación. Es un contrapeso teórico al argumento de COLE y KEHOE a favor de la deuda larga, y la proporción actual de letras, en torno al $7\%$, refleja el peso que el Tesoro da al riesgo de refinanciación tras la experiencia de 2010-2012.

La historia de las letras ilustra ese riesgo. Gracias a los elevados tipos a los que se emitieron al principio (un $15,5\%$ anual), demostraron una gran capacidad de financiación, y en 1990 llegaron a suponer el $62\%$ de la deuda, lo que obligaba a renovar en menos de un año la mayor parte de la cartera en condiciones desconocidas. Para corregirlo se promovió la participación de los inversores institucionales, que prefieren plazos más largos: en 1990 se introdujeron incentivos fiscales para los no residentes, que elevaron su proporción de bonos y obligaciones del $10\%$ a más del $30\%$ en apenas un año, y se fomentó la demanda minorista con campañas de publicidad y con los FONDTESORO, fondos de inversión especializados en deuda del Tesoro con comisiones limitadas por convenio. Las letras bajaron por debajo del $30\%$ de la deuda en 1996. Esa cartera tan corta explica la vulnerabilidad de la deuda española en 1992-1995: la subida de tipos para defender la peseta en el SME se trasladó casi de inmediato al coste de la deuda.

Un fenómeno reciente merece atención: el aumento súbito de la inversión minorista en 2023. Con la subida de los tipos del BCE, la rentabilidad de las letras pasó de negativa en 2022 a superar el $3,5$-$4\%$ en 2023, mientras los depósitos bancarios apenas se remuneraban. La deuda del Estado en manos de personas físicas, que había tocado mínimo en abril de 2022 ($898$ millones, el $0,08\%$ del total, frente al $3,8\%$ de media de la Unión, el $7,9\%$ de Italia o el $13,2\%$ de Portugal a finales de ese año), se disparó: la inversión de los particulares en letras superó los $28.000$ millones en 2023, y la proporción de letras en manos de hogares e instituciones sin fines de lucro pasó del $0,1\%$ en septiembre de 2022 a cerca del $30\%$ un año después. El episodio admite tres lecturas. Para el Tesoro supuso una diversificación de la base inversora en el momento en que el BCE dejaba de comprar; para el sistema financiero, una presión competitiva que obligó a los bancos a mejorar la remuneración de los depósitos, confirmando la transmisión de la política monetaria; y para la gestión del riesgo, una advertencia, porque una base minorista concentrada en el corto plazo es sensible a la rentabilidad relativa y puede retirarse con rapidez cuando los tipos bajan.

\subsubsection{Los bonos y obligaciones del Estado: el núcleo de la financiación}
Los bonos y obligaciones son valores de renta fija a más de dos años que solo se diferencian en el plazo: de $2$ a $5$ años los bonos y de más de $5$ años las obligaciones, hasta $50$ años en la actualidad. Pagan normalmente un cupón anual fijo, aunque en la etapa de tipos negativos llegaron a emitirse con cupón cero y al descuento. La obligación a diez años es la referencia para calcular la prima de riesgo frente a Alemania. Se emiten habitualmente por subasta, salvo el primer tramo de las nuevas obligaciones a plazos largos, que suele colocarse por sindicación, y de forma puntual se han realizado colocaciones privadas.

Las reaperturas de una referencia pueden emitirse a la par, sobre la par o bajo la par, según la relación entre su cupón y el tipo de mercado. Cuando los tipos caen, como en 2015-2021, la reapertura de una referencia con cupón alto se emite sobre la par: el Tesoro ingresa más que el nominal, y la diferencia (la prima de emisión) financia el déficit de caja sin aumentar la deuda PDE, a cambio de pagar cupones altos hasta el vencimiento. Desde 2022 ocurre lo contrario: muchas reaperturas se emiten bajo la par y el Tesoro necesita emitir más nominal para obtener el mismo efectivo. La elección entre reabrir una referencia antigua, con un cupón muy distinto del tipo de mercado, o crear una nueva, con un cupón próximo a él, no es, por tanto, neutral para la deuda publicada: aunque en valor presente el coste sea equivalente, la reapertura sobre la par reduce la deuda PDE del momento y eleva los intereses futuros en contabilidad nacional. La literatura sobre contabilidad creativa ha señalado este tipo de decisiones como un margen de maniobra de los gestores (VON HAGEN y WOLFF, 2006), aunque el criterio declarado del Tesoro es la liquidez de las referencias.

La emisión a plazos muy largos plantea un debate propio. La primera obligación a $50$ años se emitió en 2016, y le siguieron otras en la etapa de tipos mínimos, con la justificación de asegurar financiación barata durante medio siglo y atender la demanda de aseguradoras y fondos de pensiones, que necesitan activos muy largos para cubrir sus pasivos. A favor está la lógica del seguro: el Estado se protege frente a subidas futuras de tipos, como la de 2022-2026. En contra, el coste esperado, porque con una curva de pendiente positiva la deuda a $50$ años es más cara en promedio que la renovación sucesiva de deuda más corta, y su menor liquidez. La evaluación \textit{ex post} depende del periodo elegido: las emisiones de 2016-2021 se hicieron a tipos muy inferiores a los actuales y han resultado baratas frente a la refinanciación de 2023-2026, aunque su volumen es pequeño en relación con la cartera.

Una característica jurídica afecta a todos los bonos y obligaciones emitidos desde 2013. En aplicación del artículo 12.3 del Tratado del MEDE, todos los valores de deuda pública de la zona del euro a más de un año emitidos desde el 1 de enero de 2013 incluyen \textbf{cláusulas de acción colectiva} estandarizadas, que permiten que una mayoría cualificada de tenedores acuerde una reestructuración vinculante para todos, lo que evita que una minoría la bloquee. La reforma del Tratado del MEDE de 2021 prevé sustituirlas por cláusulas de agregación única (\textit{single-limb}), en las que la mayoría se calcula para todas las emisiones a la vez, pero no ha entrado en vigor al no haberla ratificado Italia. Su relevancia práctica para España es pequeña, pero muestran que la deuda soberana de la zona del euro no se considera jurídicamente libre de riesgo de reestructuración, una consecuencia de la experiencia griega de 2012.\footnote{%
    Sobre las reestructuraciones de deuda soberana, las cláusulas de acción colectiva y la experiencia griega de 2012 ([\authorfont{Ver Tema 3.A.38}]); sobre la reforma del Tratado del MEDE ([\authorfont{Ver Tema 3.B.45}]).
}

\subsubsection{Los programas de diversificación: bonos indexados y bonos verdes}
El programa de \textbf{bonos indexados a la inflación}, lanzado en 2014, permite al Tesoro emitir bonos y obligaciones cuya rentabilidad está indexada al IPC armonizado de la zona del euro sin tabaco, publicado por Eurostat; si la inflación acumulada fuera negativa, se devolvería al menos el nominal. El programa ha diversificado la base inversora, con acceso a inversores interesados en protegerse de la inflación, y su saldo vivo equivale a algo más del $6\%$ de la deuda del Estado. Su fundamento teórico es doble: BARRO (2003) concluyó que, en ausencia de deuda contingente, la deuda larga indexada es la que mejor protege al presupuesto, y CAMPBELL y SHILLER (1996) defendieron su emisión por el ahorro de la \textbf{prima de riesgo de inflación}, la compensación que los inversores exigen a la deuda nominal por la incertidumbre inflacionista y que el Estado se ahorra si asume él ese riesgo.

La evaluación del programa español tiene, sin embargo, cuatro matices. El primero, y el más relevante para España, es la cobertura: funciona si los ingresos públicos crecen con la inflación más que los gastos, pero desde la Ley 21/2021, que revaloriza las pensiones con el IPC, una parte muy importante del gasto también está indexada, y el Estado soporta el riesgo de inflación por los dos lados de su balance. El segundo es el coste \textit{ex post} de 2022-2023: con una inflación de la zona del euro de más del $8\%$ en 2022, el principal de los indexados se revalorizó en la misma cuantía, con un coste superior al de la deuda nominal. No es un argumento contra el programa, porque se trata de un seguro y ese coste es la contrapartida de la protección que ofrece a los inversores, pero sí a favor de mantenerlo pequeño mientras el gasto esté indexado; el Reino Unido, con más del $20\%$ de su deuda indexada, sufrió por este motivo un fuerte aumento de su carga de intereses en 2022-2023. El tercero es el índice de referencia: el Tesoro indexa a la inflación de la zona del euro, que es la que tiene un mercado líquido de inversores y coberturas, pero los ingresos españoles dependen de la inflación española, de modo que si esta se desvía de la media la cobertura se debilita. El cuarto es la liquidez: el segmento indexado es menos líquido que el nominal, y la prima de liquidez que exigen los inversores compensa en parte el ahorro de la prima de inflación. Por todo ello, el Tesoro mantiene los indexados en torno al $6\%$ de la cartera como instrumento de diversificación y no como apuesta de coste.

El programa de \textbf{bonos verdes soberanos}, lanzado en 2021, persigue un doble objetivo: diversificar la base inversora y reforzar el compromiso de España con la transición ecológica. Los fondos deben destinarse a gastos elegibles (mitigación y adaptación al cambio climático, uso sostenible y protección de los recursos hídricos y marinos, economía circular, prevención de la contaminación y protección de la biodiversidad), lo que limita el volumen que puede emitirse. La emisión inaugural, en septiembre de 2021, fue una obligación a $20$ años con vencimiento en 2042 por $5.000$ millones, reabierta después en varias ocasiones, y la elevada demanda, unida a su escasez relativa, permitió un ahorro frente a un bono convencional (el \textit{greenium}) de $2$ puntos básicos. El 9 de septiembre de 2026, el Tesoro colocó por sindicación una nueva referencia verde a $20$ años (vencimiento en 2047) por $4.000$ millones, a una rentabilidad en torno al $4,2\%$ y con una demanda de $68.500$ millones, con lo que el saldo verde se acerca a los $22.000$ millones, en torno al $1,5\%$ de la deuda del Estado.

El programa suscita tres debates. El primero es la \textbf{adicionalidad}: el dinero es fungible, el Estado habría realizado los gastos verdes elegibles con o sin bono verde y la emisión no altera la composición del presupuesto. Sus defensores responden que su valor está en la transparencia, porque los informes anuales de asignación y de impacto obligan a identificar y medir los gastos verdes, y en la señal de compromiso; sus críticos lo ven como una etiqueta sin efecto real. El segundo es la \textbf{magnitud del \textit{greenium}}: la evidencia sobre bonos verdes soberanos (DORONZO, SIRACUSA y ANTONELLI, 2021) encuentra ahorros positivos pero pequeños, de pocos puntos básicos y variables en el tiempo, y una parte de la diferencia puede reflejar la menor liquidez de estos bonos y no una preferencia de los inversores. El tercero es la \textbf{competencia europea}: la Comisión Europea, que financia una parte del \textit{Next Generation EU} con bonos verdes, se ha convertido en el mayor emisor verde del mundo, lo que reduce la escasez relativa que sostenía el \textit{greenium} de los emisores nacionales.

\begin{specialblock}{\textbf{Material complementario:} un instrumento que el Tesoro no emite, la deuda vinculada al PIB}
    La teoría de la gestión óptima de la deuda recomienda un instrumento que ningún gran emisor utiliza. Siguiendo la lógica de LUCAS y STOKEY (1983), la deuda cuyo pago se vincula al crecimiento del PIB nominal protegería al presupuesto frente a las recesiones: se pagaría menos cuando la economía cae y más cuando crece. SHILLER la propuso hace décadas, BORENSZTEIN y MAURO (2004) la defendieron, y en 2016-2017 el Banco de Inglaterra y un grupo de trabajo del G20 elaboraron un modelo de contrato (el \textit{London Term Sheet}).

    Los obstáculos que han impedido su adopción son varios: el riesgo de revisiones y de manipulación del PIB, la falta de liquidez de un instrumento nuevo, la prima de novedad que exigirían los inversores y la ausencia de un primer emisor de referencia. La experiencia más próxima son las cláusulas de deuda resistente a catástrofes climáticas, que suspenden los pagos tras un desastre natural y que han adoptado algunos países en sus préstamos a países en desarrollo. Tras la DANA de 2024, la idea tendría cierta lógica para España, aunque no figura en la agenda del Tesoro.\footnote{%
        Sobre la deuda contingente al estado de la economía como solución teórica a la nivelación impositiva ([\authorfont{Ver Tema 3.A.39}]).
    }
\end{specialblock}

\subsubsection{Otros instrumentos: deuda en divisas y préstamos}
El crecimiento de las necesidades de financiación ha llevado a primar las letras y los bonos y obligaciones sobre otros instrumentos, que hoy suponen menos del $4\%$ de la deuda del Estado. La \textbf{deuda en divisas} tuvo un papel destacado para atraer inversores internacionales en los ochenta y noventa, y llegó a superar el $10\%$ de la deuda a finales de los noventa, pero con el euro dejó de ser necesaria para acceder a los no residentes. La última emisión en divisas tuvo lugar en 2014, y quedan emisiones residuales en yenes, dólares y libras, de menos del $0,1\%$ de la deuda, cuyo riesgo de tipo de cambio el Tesoro cubría con permutas financieras (\textit{swaps}). Su historia ilustra el riesgo cambiario: las devaluaciones de 1992-1995 elevaron el valor en pesetas de la deuda en divisas no cubierta y contribuyeron al salto de la ratio de 1993. Con el euro, ese riesgo ha desaparecido para la deuda del Estado, pero la renuncia a la moneda propia tiene la contrapartida que analizamos en la comparación internacional: la deuda está denominada en una moneda que el país no controla.

Los \textbf{préstamos} han permitido al Tesoro reducir sus emisiones en momentos concretos, y todos se obtuvieron a tipos inferiores a los que el Tesoro pagaba en cada momento. Los más importantes han sido el del MEDE de 2012 para la recapitalización bancaria ($41.333$ millones), los del SURE en 2020-2021 ($21.326$ millones) y los del Mecanismo de Recuperación y Resiliencia (en torno a $21.500$ millones solicitados en total), a los que se añaden varios préstamos del Banco Europeo de Inversiones y algunos préstamos \textit{Schuldschein}, sujetos a la legislación alemana y adaptados a las necesidades de cada inversor, utilizados por última vez en 2014. Los del MEDE y del SURE no fueron préstamos al Tesoro en sentido estricto: el primero se desembolsó al FROB en bonos del propio MEDE, y el segundo se financió con emisiones de la Unión garantizadas por los Estados miembros. Pero ambos computan como deuda de las AAPP y redujeron lo que el Tesoro necesitaba emitir en el mercado. La financiación oficial europea ha actuado así como amortiguador de las necesidades de mercado en 2012 y en 2020-2026: en el primer caso, con condicionalidad financiera y en una crisis de acceso; en el segundo, sin estigma y en el marco de una respuesta común. Desaparece, sin embargo, a partir de 2027, lo que devolverá al Tesoro a una dependencia completa del mercado.

\subsection{Procedimientos de emisión y el funcionamiento del mercado}
La subasta sigue siendo el procedimiento central del Tesoro, con más del $80\%$ de la emisión bruta anual. El formato español es híbrido: las pujas por encima del precio medio ponderado se adjudican a ese precio, y las demás, al precio ofrecido. Con frecuencia se presenta como más barato que la subasta holandesa, pero, con pujas estratégicas, la teoría no establece una ordenación general de los formatos y la evidencia empírica es mixta; su justificación es que corrige parcialmente la rebaja de las pujas altas y la reducción estratégica de la demanda. Las peticiones no competitivas han ganado importancia con la entrada de los minoristas desde 2023. La previsibilidad (calendario anual, $48$ subastas en 2026, rangos objetivo desde 1995) aplica a la gestión de la deuda la lógica de las reglas frente a la discrecionalidad. Las sindicaciones, algo más del $10\%$ de la emisión en años normales y casi el $20\%$ en 2020, permiten lanzar nuevas referencias largas con volumen y distribución controlada, a cambio de una prima de nueva emisión y de una dependencia del sindicato; las de 2026 han registrado demandas récord y una base inversora predominantemente internacional. Los canjes, las recompras y las operaciones de tesorería gestionan la estructura de la cartera y la liquidez, y no la financiación. Por último, el sistema de creadores de mercado remunera de forma implícita sus compromisos con privilegios como la segunda vuelta, que equivale a una opción gratuita, y afronta el reto de absorber una oferta neta creciente con los balances bancarios sujetos a restricciones regulatorias, ahora que el BCE se retira.

El punto de partida es histórico. El paso a la financiación de mercado desde principios de los ochenta exigía procedimientos que asegurasen el abastecimiento regular del mercado y la financiación al menor coste posible. En 1983 se adoptó la subasta competitiva para las obligaciones del Estado, que después se extendió a los bonos y a las letras, y a ella se fueron añadiendo la sindicación, las colocaciones privadas, los canjes y los préstamos, con todo el sistema apoyado en los creadores de mercado.

\subsubsection{La subasta: el formato español y el debate sobre el diseño}
Los Tesoros utilizan dos formatos de subasta puros. En la \textbf{subasta holandesa} o de precio único, todas las peticiones aceptadas se adjudican al precio marginal, es decir, al mínimo aceptado. En la \textbf{subasta convencional} o de precios múltiples, cada petición aceptada se adjudica al precio que ofreció. El sistema español combina ambos. El Tesoro fija primero el precio mínimo aceptado y calcula después el precio medio ponderado de las peticiones aceptadas; las peticiones a un precio superior al medio ponderado se adjudican a ese precio medio, y las demás, al precio que ofrecieron. El Tesoro decide el precio mínimo en la propia subasta, a la vista de las peticiones recibidas, del importe que desea emitir (anunciado antes como un rango orientativo), del número de referencias subastadas y del saldo en circulación de cada una. En las subastas de letras se ofrecen siempre dos referencias, y en las de bonos y obligaciones, normalmente tres o cuatro.

Junto a las peticiones competitivas, los inversores pueden presentar \textbf{peticiones no competitivas}, en las que solo indican la cantidad y que se adjudican al precio medio ponderado. Salvo excepciones, no pueden superar los $5$ millones de euros y, con carácter general, no se adjudican si ese precio implica una rentabilidad negativa. En todo el proceso, el BdE actúa como agente financiero del Estado: recibe las peticiones, gestiona la resolución técnica de la subasta, liquida las operaciones y publica los resultados, mientras que las decisiones de fondo, el importe y el precio mínimo, corresponden al Tesoro.

¿Qué formato minimiza el coste? El argumento habitual a favor del formato español es sencillo: como las pujas altas se adjudican al precio medio ponderado, que es superior al marginal, el Tesoro ingresaría más que con la subasta holandesa, en la que todo se adjudicaría al marginal. El razonamiento sería correcto si los inversores pujasen igual en ambos formatos, pero no lo hacen: el comportamiento de los pujadores depende de las reglas de adjudicación.\footnote{%
    Con valoraciones privadas independientes y demanda unitaria, el teorema de equivalencia de ingresos (VICKREY, 1961; MYERSON, 1981) implica que los formatos de precio único y de precios múltiples generan el mismo ingreso esperado. Con demanda de varias unidades y valoraciones con un componente común, como en la deuda pública, la equivalencia se rompe. En la subasta de precios múltiples, cada inversor paga lo que puja, y tiene incentivo a pujar por debajo de su valoración (\textit{bid shading}) para no pagar de más y para protegerse de la \textbf{maldición del ganador}: si gana, es probablemente porque ha sobrevalorado el título respecto a los demás. En la de precio único, el inversor paga un precio marginal determinado por las pujas ajenas, lo que le anima a pujar cerca de su valoración; FRIEDMAN (1960) la defendió por esa razón. Pero AUSUBEL y CRAMTON (2002) mostraron que, cuando cada inversor demanda varias unidades, surge otro incentivo, la \textbf{reducción estratégica de la demanda}: rebajar las pujas en las unidades marginales para bajar el precio que pagará por todas las demás, de modo que el formato de precio único puede generar menos ingresos. El efecto neto es ambiguo y depende de la distribución de valoraciones, del número de pujadores y de su información ([\authorfont{Ver Tema 3.A.14}]).
} En los contextos de unidades múltiples y valoraciones interdependientes, que son los de la deuda pública, no se cumple el teorema de equivalencia de ingresos, y ningún formato domina teóricamente a los demás en ingreso esperado.

La evidencia empírica es igualmente mixta. El Tesoro de Estados Unidos, que entre 1992 y 1998 utilizó el formato de precio único en parte de sus subastas y lo generalizó en 1998, encontró ahorros modestos y difíciles de atribuir con seguridad al formato (MALVEY y ARCHIBALD, 1998). UMLAUF (1993), con el caso de México, encontró ingresos mayores con el precio único, y HORTAÇSU y MCADAMS (2010), con datos de Turquía, estimaron diferencias pequeñas entre formatos. HORTAÇSU, KASTL y ZHANG (2018) añadieron un elemento que la teoría clásica no recogía: el comportamiento estratégico de los creadores de mercado, que conocen el flujo de órdenes de sus clientes, influye de forma importante en el resultado.

Sobre el formato español existe una literatura específica, aunque reducida. ÁLVAREZ y MAZÓN (2007) lo compararon teóricamente con el de precios múltiples y encontraron que puede reducir el \textit{bid shading} sin eliminar la protección frente a pujas excesivamente altas; ABBINK, BRANDTS y PEZANIS-CHRISTOU (2006), en un experimento de laboratorio con los tres formatos, encontraron que el español no genera sistemáticamente ingresos menores que los otros dos, pero tampoco mayores en todas las condiciones. La conclusión defendible es, por tanto, que el formato español no es más barato por construcción. Su justificación es que corrige parcialmente los dos problemas de los formatos puros: reduce el incentivo a rebajar las pujas altas, porque las pujas por encima de la media no se pagan a su precio, y limita la reducción estratégica de la demanda, porque las pujas por debajo de la media sí se pagan a su precio. A ello se suman su antigüedad y el conocimiento que tienen de él los inversores, que reducen la incertidumbre.

Las peticiones no competitivas han adquirido un papel relevante desde 2023, con el aumento súbito de la inversión minorista en letras que vimos al analizar los instrumentos. Permiten a los pequeños ahorradores participar sin necesidad de estimar el precio y les garantizan el precio medio, que no es el peor de la subasta. Para el Tesoro, sus importes son relativamente previsibles y constituyen una demanda estable, aunque sensible a la rentabilidad relativa, y la regla que excluye su adjudicación con rentabilidades negativas protegía a los minoristas en la etapa de tipos negativos. Desde 2025, el Tesoro ofrece además a los particulares una plataforma propia para invertir directamente en letras, bonos y obligaciones, como parte de su estrategia de diversificación de la base inversora.

\subsubsection{El calendario, los rangos y la transparencia}
Para asegurar la estabilidad de las emisiones y el acceso continuado a la financiación, esencial desde que en 1994 se prohibió la financiación del BdE, cada año se publica en el BOE el calendario de subastas. Cada mes hay dos subastas de letras, una a $6$ y $12$ meses y otra a $3$ y $9$ meses, y dos de bonos y obligaciones. El Tesoro anuncia el viernes anterior las referencias concretas que subastará y, desde 1995, anuncia el lunes de la semana de la subasta el \textbf{rango objetivo de colocación}, una práctica que se introdujo para evitar las grandes oscilaciones que había hasta entonces en los volúmenes emitidos. La Estrategia de Financiación de 2026 prevé $48$ subastas ordinarias.

El calendario y los rangos son la aplicación del objetivo de previsibilidad y transparencia que vimos entre los objetivos de la política de financiación, y su fundamento es doble. Por un lado, reducen la incertidumbre de los inversores sobre la oferta futura, lo que disminuye la prima que exigen por esa incertidumbre. Por otro, reducen el margen de discrecionalidad del Tesoro. Un emisor que decidiera cada semana cuánto emite según las condiciones del momento podría aprovechar las ventanas favorables, pero los inversores anticiparían ese comportamiento, esperarían que emitiera más cuando los precios le son favorables y menos cuando no, y exigirían una prima mayor por la volatilidad resultante. Es la aplicación a la gestión de la deuda de la idea de \textbf{reglas frente a discrecionalidad} que KYDLAND y PRESCOTT (1977) formularon para la política económica.

El rango objetivo tiene, sin embargo, una contrapartida. Al comprometerse con un intervalo, el Tesoro renuncia a ajustar mucho el importe cuando la demanda es débil. En las subastas de 2011-2012, cuando la demanda caía, colocar en la parte baja del rango a tipos altos era preferible a no colocar, porque no hacerlo habría enviado una señal de pérdida de acceso al mercado: la regla, que en tiempos normales abarata la financiación, obliga en las crisis a aceptar precios desfavorables para preservar la credibilidad del propio compromiso.

\subsubsection{Las sindicaciones: colocación frente a coste}
Las sindicaciones son el segundo procedimiento en importancia y se usan, con carácter general, para emitir el primer tramo de las nuevas obligaciones. En una sindicación, el Tesoro contrata a un grupo de bancos (el sindicato), fija el precio de emisión y los inversores indican la cantidad que desean a ese precio. Sus ventajas son tres. La primera es que permiten emitir de una vez un volumen mayor que una subasta, lo que da liquidez al nuevo título desde el inicio: en plena pandemia, el Tesoro captó $15.000$ millones de una nueva obligación en una sola operación. La segunda es que dan al Tesoro el control de la distribución: como la demanda suele superar con mucho la oferta, puede asignar los títulos por categorías de inversores y dar un porcentaje mayor a los que previsiblemente los mantendrán más tiempo, como los bancos centrales y las instituciones oficiales, o a los de regiones con menor inversión acumulada en deuda española. La tercera es la notoriedad que da al emisor. Su inconveniente principal es un mayor coste: el Tesoro no llega a conocer el precio marginal al que los inversores estarían dispuestos a comprar, como sí ocurre en la subasta, debe pagar comisiones al sindicato y depende de este, que tiene más información de primera mano sobre los inversores finales. Normalmente, el Tesoro obtiene algo más del $10\%$ de su emisión bruta anual mediante sindicaciones, y en 2020 llegó a casi el $20\%$.

Las sindicaciones recientes muestran un apetito por la deuda española muy superior al de hace una década. En enero de 2026, el Tesoro emitió $15.000$ millones de una nueva obligación a diez años con una demanda de $145.000$ millones, un récord histórico que supera los $139.000$ millones de 2025. El cupón era del $3,30\%$, y la rentabilidad del $3,323\%$ quedaba $6$ puntos básicos por encima de la referencia a diez años vigente; los no residentes tomaron el $90\%$ de la emisión, y los bancos centrales e instituciones oficiales fueron el primer grupo inversor ($26,5\%$). En mayo de 2026 colocó otra obligación a diez años por $13.000$ millones, con una demanda diez veces superior, y en septiembre, como vimos, $4.000$ millones de un nuevo bono verde a veinte años, con una demanda de $68.500$ millones.

¿Cuánto cuesta la sindicación? La literatura distingue dos componentes. El primero es la comisión pagada al sindicato, explícita y pequeña en relación con el importe. El segundo, más importante, es la \textbf{prima de nueva emisión}, la diferencia entre el precio de emisión y el precio al que cotiza después el título en el mercado secundario. El diferencial de $6$ puntos básicos de la emisión de enero de 2026 sobre la curva es un ejemplo de esa prima: el Tesoro paga algo más para asegurar una colocación grande y bien distribuida. La cuestión es si ese coste se compensa con una base inversora más estable y con la liquidez inmediata del título. Los análisis de la OCDE y la práctica de los gestores de deuda europeos sugieren que sí para las nuevas referencias largas, en las que la subasta es arriesgada porque no hay precio de referencia en el secundario y el volumen necesario para garantizar la liquidez es grande. Por eso los Tesoros europeos, incluido el español, reservan la sindicación para el lanzamiento de nuevas referencias a plazos largos, para los programas especiales (verdes e indexados) y para las situaciones de tensión, y reabren después esas referencias por subasta. La experiencia de 2020 muestra que la sindicación es también un instrumento de gestión de crisis, porque permite asegurar grandes volúmenes en un solo día cuando la volatilidad hace arriesgada la subasta.

La crítica principal a la sindicación es la \textbf{asimetría de información}. Los bancos del sindicato conocen la demanda de sus clientes antes que el Tesoro y pueden fijar el precio de forma favorable para ellos o para sus clientes. El Tesoro la mitiga con varios mecanismos: rota a los bancos del sindicato según su desempeño como creadores de mercado, publica los criterios de asignación y compara el precio con el de la curva secundaria.

\subsubsection{Otros procedimientos: colocaciones privadas, canjes y recompras}
El Tesoro utiliza además otros procedimientos de importancia residual. Las \textbf{colocaciones privadas} se realizan a propuesta de los inversores y se ejecutan si contribuyen a diversificar la base inversora, reducen la carga de intereses y encajan en la estrategia de financiación. Los \textbf{canjes} se utilizan de forma excepcional para aumentar la liquidez y suavizar el perfil de vencimientos, retirando referencias poco líquidas o próximas al vencimiento a cambio de otras. Los préstamos, como vimos, tienen por objetivo reducir la carga de intereses y alargar la vida media.

Los canjes y las \textbf{recompras}, en las que el Tesoro compra en el mercado referencias propias antes de su vencimiento, son instrumentos de gestión de pasivos, no de financiación: su objetivo no es captar recursos, sino modificar la estructura de la cartera. Sirven para suavizar los picos de vencimientos, que concentran el riesgo de refinanciación en fechas concretas, y para retirar referencias poco líquidas que fragmentan el mercado. Su uso ha sido limitado en España y suele concentrarse en los años previos a grandes amortizaciones. En un contexto de vencimientos crecientes como el actual, podrían ganar importancia, como ya ocurre en otros Tesoros europeos que los utilizan con más frecuencia.

A estos procedimientos se añaden las \textbf{operaciones de gestión de tesorería}. Desde que el BCE limitó la remuneración de los depósitos de las administraciones públicas en el banco central, el Tesoro coloca sus excedentes de liquidez en el mercado mediante subastas de depósitos y adquisiciones temporales de activos con las entidades, lo que le permite remunerarlos sin mantenerlos en el banco central. Estas operaciones no aumentan la deuda, pero forman parte de la gestión integrada de la deuda y la tesorería que recomiendan las directrices internacionales.

\subsubsection{Los creadores de mercado y el mercado secundario}
Los \textbf{creadores de mercado} son entidades financieras que se comprometen con el Tesoro a favorecer la liquidez del mercado secundario y a colaborar en la difusión de la deuda del Estado. Hay dos grupos, uno de letras y otro de bonos y obligaciones, aunque casi todas las entidades pertenecen a ambos: en 2022 eran $19$ y $18$, respectivamente. Para acceder a la condición deben cumplir requisitos técnicos y humanos y demostrar su compromiso con el mercado español, y adquieren obligaciones de participación en las subastas y de provisión de liquidez en el secundario. A cambio, disfrutan de condiciones más favorables en las subastas (solo ellos pueden presentar peticiones en los treinta minutos previos al cierre), forman parte de los sindicatos, acceden a las segundas vueltas y pueden segregar y reconstituir valores. El Tesoro evalúa mensualmente su desempeño con criterios objetivos y públicos, y los resultados determinan el acceso a las segundas vueltas y la selección de los sindicatos. En 2022, los mejores creadores de bonos y obligaciones fueron Deutsche Bank, BBVA y JP Morgan, y los de letras, Crédit Agricole, Banco Santander y BBVA.

El privilegio más relevante desde el punto de vista económico es la \textbf{segunda vuelta}, introducida en 1991. Permite a los creadores de mercado solicitar, hasta un máximo del $24\%$ del nominal que adquirieron en la subasta, títulos adicionales al precio medio ponderado redondeado; el porcentaje de cada creador depende de su participación y de su desempeño. Como se celebra dos o tres días después de la subasta, funciona como una \textbf{opción de compra gratuita}: el creador la ejercerá si en ese intervalo los tipos han bajado y el precio ha subido. Para el Tesoro tiene, por tanto, un coste esperado, el valor de la opción que regala, que es mayor cuanto más volátil es el mercado. Ese coste es la remuneración implícita de los compromisos de los creadores de mercado, participar en las subastas aunque sean difíciles y dar precios en el secundario: una forma de pagar por un servicio sin hacerlo de forma explícita. La literatura sobre sistemas de creadores de mercado discute si es preferible remunerar esos servicios con privilegios de este tipo o con comisiones explícitas, y la mayoría de los Tesoros europeos utiliza combinaciones de privilegios parecidas a la española.\footnote{%
    Sobre el diseño de los sistemas de creadores de mercado de deuda pública y las alternativas de remuneración de sus compromisos, véase ARNONE e IDEN (2003), elaborado para el FMI.
}

La deuda del Estado se negocia en el Mercado de Deuda Pública en Anotaciones y se registra en Iberclear. La mayor parte de la negociación entre profesionales se realiza en plataformas electrónicas (MTS España, BrokerTec y plataformas con clientes como Tradeweb o Bloomberg) y en operaciones bilaterales. A ello se suma un mercado de repos muy activo, que permite financiar las posiciones en deuda pública y convierte a esta en un colateral esencial para el sistema financiero. La liquidez del secundario, medida por los diferenciales entre precios de compra y de venta y por el volumen negociado, es uno de los objetivos del Tesoro, y se ha mantenido elevada en los últimos años, incluso durante la retirada del BCE.

\begin{specialblock}{\textbf{Material complementario:} el futuro del modelo de creadores de mercado}
    El modelo de creadores de mercado afronta dos retos que la literatura reciente ha señalado. El primero es regulatorio. Las reformas bancarias posteriores a 2008, en particular la ratio de apalancamiento de Basilea III y los requisitos de capital y liquidez, han encarecido que los bancos mantengan grandes carteras de deuda pública y actúen como intermediarios. DUFFIE (2020) mostró cómo esta restricción contribuyó a las tensiones del mercado de deuda estadounidense en marzo de 2020, cuando los intermediarios no pudieron absorber las ventas masivas. En la zona del euro, el problema se amortiguó por la presencia del BCE; pero con la retirada del Eurosistema y el aumento de la oferta neta al sector privado, la capacidad de los creadores de mercado para absorber la emisión vuelve a ser una cuestión relevante.

    El segundo es estructural. La electronificación y la entrada de intermediarios no bancarios, como los fondos de alta frecuencia y los creadores de mercado electrónicos, han cambiado la forma en que se provee la liquidez: una liquidez abundante en tiempos normales, pero que puede evaporarse en los episodios de tensión, precisamente cuando más se necesita. Algunos Tesoros europeos han revisado sus sistemas de creadores de mercado para adaptarlos a este entorno, mientras que el español se ha mantenido estable, con ajustes en los criterios de evaluación.
\end{specialblock}

\subsection{El BCE y la estrategia de financiación del Tesoro en 2026}
La estrategia de financiación del Tesoro no puede entenderse sin su contexto decisivo, la relación con el BCE. Entre 2015 y 2022, el Eurosistema compró más de $500.000$ millones de deuda española con el APP y el PEPP, más de un tercio de la deuda de las AAPP, lo que permitió al Tesoro financiarse a coste negativo en 2021 y alargar la vida media por encima de $8$ años, incluso renunciando al ahorro de las letras durante la pandemia. Pero esa protección era completa para la cartera del Tesoro y no para el sector público en su conjunto: desde el balance consolidado, la expansión cuantitativa acortó la duración efectiva de la deuda, porque una parte se financió con reservas bancarias remuneradas a un día, y la subida de tipos trasladó ese coste al BdE, con pérdidas de $8.389$ millones en política monetaria en 2024 y ninguna transferencia de beneficios al Tesoro en 2023-2024. Es un coste fiscal que no aparece en la cartera del Tesoro y que ha alimentado el debate sobre la remuneración de las reservas. La retirada del BCE, con el fin de las reinversiones del APP en julio de 2023 y del PEPP en enero de 2025, ha multiplicado la oferta neta de deuda que absorbe el sector privado; la han cubierto sobre todo los no residentes, sin aumento de la prima de riesgo, gracias a los fundamentales, al TPI y a la gradualidad de la retirada. La estrategia de 2026 es de continuidad ($55.000$ millones netos, unos $285.677$ millones brutos, una vida media de unos $8$ años y $48$ subastas) en un entorno que ha dejado atrás el dinero barato, con las subidas del BCE de junio y septiembre de 2026, y que afronta riesgos de tipos, fiscales, de contagio y de concentración de la base inversora.\footnote{%
    Sobre la política monetaria no convencional (expansión cuantitativa, canales de transmisión, efectos sobre la prima por plazo) y el Instrumento para la Protección de la Transmisión ([\authorfont{Ver Tema 3.A.36, Tema 3.A.37}]).
}

\subsubsection{2015-2022: el BCE como comprador}
A finales de 2014, las medidas convencionales no bastaban para devolver la inflación a su objetivo, y el BCE anunció un paquete de medidas no convencionales que incluía las operaciones de financiación a largo plazo con objetivo específico (TLTRO) y el \textbf{programa de compra de activos} (APP), que abarcaba deuda pública (PSPP), bonos corporativos (CSPP), titulizaciones (ABSPP) y bonos garantizados (CBPP3). En marzo de 2020 se añadió el \textbf{programa de compras de emergencia frente a la pandemia} (PEPP). En ambos casos, las compras de deuda pública se repartían según la clave de capital de cada banco central nacional, con cierta flexibilidad temporal. Las compras netas del APP alcanzaron su máximo en 2016-2017, con $80.000$ millones mensuales, y las del PEPP llegaron a unos $120.000$ millones en junio de 2020. Entre 2015 y 2022, las compras netas de deuda española superaron los $500.000$ millones, más de la tercera parte de la deuda de las AAPP a finales de 2022; descontada esa deuda, la deuda de las AAPP en manos de otros inversores era a finales de 2022 equivalente a la de principios de 2013.

Dos precisiones de diseño son relevantes para el Tesoro. La primera es que el PSPP limitaba las tenencias del Eurosistema al $33\%$ de cada emisión y de cada emisor. El límite se estableció para no bloquear las cláusulas de acción colectiva que vimos al analizar los bonos y obligaciones y para evitar que el BCE se convirtiera en el acreedor dominante de un Estado. La segunda es que el PEPP flexibilizó esos límites y la clave de capital, lo que permitió comprar más deuda de los países más afectados en 2020; esa flexibilidad fue decisiva para contener las primas en la primavera de 2020.

Estos programas redujeron de forma significativa el coste de emisión, porque absorbieron la mayor parte de la emisión neta desde 2015. En 2021, el coste medio de emisión fue negativo por primera vez en la historia del Tesoro ($-0,04\%$), y el coste de la deuda en circulación alcanzó un mínimo histórico del $1,64\%$. El Tesoro aprovechó esas condiciones para alargar la vida media de la deuda, que en 2021 superó por primera vez los $8$ años, y mantuvo la estrategia de reducir las letras incluso durante la pandemia, a diferencia de muchos países de su entorno. Esta última decisión ilustra la disyuntiva entre coste y riesgo que vimos entre los objetivos de la política de financiación. En 2020-2021 la curva era muy plana y los tipos a corto eran negativos, de modo que emitir letras habría sido más barato a corto plazo; el Tesoro eligió comprar un seguro frente a una subida de tipos futura, renunciando al ahorro inmediato de la deuda corta, y la subida de 2022-2023 mostró \textit{ex post} que la decisión fue acertada para su cartera.

\begin{specialblock}{\textbf{Material complementario:} el debate jurídico y económico sobre la compra de deuda pública por el BCE}
    Los programas del BCE generaron un debate jurídico y económico que afecta directamente al Tesoro. En el plano jurídico, el TJUE avaló el programa OMT en la sentencia \textit{Gauweiler} (2015) y el PSPP en la sentencia \textit{Weiss} (2018), al considerar que no vulneraban la prohibición de financiación monetaria del artículo 123 del TFUE siempre que existieran garantías como los límites por emisión, la compra exclusivamente en el mercado secundario y un periodo de espera tras la emisión. En mayo de 2020, el Tribunal Constitucional alemán calificó la sentencia \textit{Weiss} de \textit{ultra vires} y exigió al BCE justificar la proporcionalidad del programa, un episodio de tensión institucional sin consecuencias prácticas duraderas.

    En el plano económico, el debate gira en torno a la \textbf{dominancia fiscal}. Los críticos sostienen que, al absorber durante años buena parte de la emisión neta de los Estados, el BCE redujo la disciplina de mercado y facilitó que los gobiernos retrasaran la consolidación; España, que no corrigió su déficit estructural en 2015-2019, sería un ejemplo. Los defensores responden que los programas se diseñaron con un objetivo de estabilidad de precios, en un contexto de inflación persistentemente baja, que las compras según la clave de capital no favorecían a ningún país y que su efecto sobre las primas fue una consecuencia de la reducción del riesgo de redenominación, no una financiación selectiva.\footnote{%
        Sobre la dominancia fiscal y la interacción entre política monetaria y fiscal (SARGENT y WALLACE, 1981; teoría fiscal del nivel de precios) ([\authorfont{Ver Tema 3.A.39}]); sobre las sentencias \textit{Gauweiler} (C-62/14), \textit{Weiss} (C-493/17) y la del Tribunal Constitucional alemán de 5 de mayo de 2020 ([\authorfont{Ver Tema 3.A.36, Tema 3.A.37, Tema 3.B.45}]).
    }
\end{specialblock}

\subsubsection{La perspectiva del balance consolidado: el coste que no se ve en el Tesoro}
Cuando el Eurosistema compra deuda del Estado, sustituye en manos del sector privado un bono del Tesoro, largo y a tipo fijo, por una reserva bancaria en el banco central, a la vista y remunerada al tipo de la facilidad de depósito. Si se consolida el balance del Estado con el de su banco central nacional, que es de su propiedad y le transfiere sus beneficios, el efecto neto de la expansión cuantitativa es un \textbf{acortamiento de la duración de la deuda del sector público consolidado}: una parte de la deuda larga que el Tesoro emitió pasa a estar financiada con pasivos a un día. GREENWOOD, HANSON, RUDOLPH y SUMMERS (2014) documentaron este efecto en Estados Unidos, donde, mientras el Tesoro alargaba la vida media de su deuda, la Reserva Federal la acortaba en manos privadas. En la zona del euro, el razonamiento se aplica al BdE, que tenía en su balance la mayor parte de la deuda española comprada por el Eurosistema, porque las compras de deuda soberana se hacían sobre todo a través de los bancos centrales nacionales y sin compartir riesgos.\footnote{%
    Sobre el reparto de riesgos en el Eurosistema (las compras de deuda soberana del PSPP se hicieron en un $90\%$ sin mutualización de riesgos, a través de los bancos centrales nacionales) y el tratamiento de las pérdidas de los bancos centrales ([\authorfont{Ver Tema 3.A.36, Tema 3.A.37}]).
}

La consecuencia se hizo visible con la subida de tipos de 2022-2023. Los bancos centrales del Eurosistema tenían en el activo bonos comprados a tipos muy bajos y en el pasivo reservas bancarias que llegaron a remunerarse al $4\%$ en 2023, y su margen se volvió negativo. El BdE registró pérdidas en sus operaciones de política monetaria, que en 2024 ascendieron a $8.389$ millones; pudo absorberlas con las provisiones acumuladas en años anteriores y cerró 2023 y 2024 con un resultado nulo, sin transferencia al Tesoro. En 2025, gracias a las bajadas de la facilidad de depósito entre junio de 2024 y junio de 2025, volvió a beneficios ($234$ millones, íntegramente transferidos al Tesoro), tras liberar $541$ millones de provisiones y con unas pérdidas por política monetaria reducidas a $2.071$ millones. En los años previos a 2023, en cambio, el BdE transfería al Tesoro beneficios del orden de miles de millones anuales.

En términos fiscales, por tanto, la subida de tipos de 2022-2024 no solo elevó el coste de la nueva emisión del Tesoro: eliminó además durante dos años un ingreso público relevante, los beneficios del banco central. Esa pérdida de ingreso es la contrapartida del acortamiento de la deuda consolidada que produjo la expansión cuantitativa. Con las subidas de junio y septiembre de 2026, el margen del BdE volverá a estrecharse, y sus beneficios de 2026 previsiblemente se reducirán, aunque el efecto depende de la evolución de su balance, que se contrae a medida que vencen los bonos.

La cuestión ha dado lugar a un debate de política con implicaciones fiscales directas. DE GRAUWE y JI (2023) calcularon que, con unas reservas de $4,3$ billones en el Eurosistema, remunerarlas al tipo de la facilidad de depósito suponía transferir a los bancos unos $129.000$ millones anuales, y propusieron un sistema de dos niveles, con una parte de las reservas sujeta a un coeficiente obligatorio no remunerado, que habría reducido esa transferencia a unos $32.000$ millones anuales con un tipo del $3\%$. El BCE adoptó una medida en esa línea, aunque limitada: en julio de 2023 fijó en el $0\%$ la remuneración de las reservas mínimas obligatorias, cuyo volumen es pequeño, pero no amplió el coeficiente. Los críticos de la propuesta advierten de que un coeficiente amplio no remunerado actuaría como un impuesto sobre los bancos de la zona del euro, podría desplazar depósitos fuera del sistema bancario y dificultaría la transmisión de la política monetaria.

La conclusión es que la estrategia del Tesoro y la política del BCE no pueden evaluarse por separado. El alargamiento de la vida media protegió la cartera del Tesoro, pero la expansión cuantitativa trasladó al banco central una parte del riesgo de tipo de interés que el Tesoro había cubierto, y ese riesgo se materializó en la cuenta de resultados pública por otra vía. Presentar el alargamiento de la vida media como una protección completa frente a la subida de tipos exige, por tanto, un matiz: lo fue para el Tesoro, no para el sector público consolidado.

\subsubsection{2022-2026: la retirada del BCE y la vuelta de las subidas de tipos}
La retirada del BCE fue gradual y distinta en cada programa. Las compras netas del PEPP terminaron en marzo de 2022 y las del APP en junio de 2022. En el APP, la reinversión de los vencimientos se redujo en $15.000$ millones mensuales entre marzo y junio de 2023 y se suprimió por completo desde julio de 2023, con lo que comenzó el \textbf{endurecimiento cuantitativo} (QT). En el PEPP, los vencimientos se reinvirtieron íntegramente hasta junio de 2024, se redujeron en unos $7.500$ millones mensuales en el segundo semestre de 2024 y dejaron de reinvertirse desde enero de 2025. Desde entonces, el Eurosistema no reinvierte nada y sus carteras se reducen al ritmo de los vencimientos, algo que el BCE reiteró en su decisión de septiembre de 2026. En ese contexto, la estrategia del Tesoro de 2023 preveía estabilizar la vida media, cubrir toda la financiación neta ($70.000$ millones) con instrumentos a medio y largo plazo y reducir las letras en $5.000$ millones.

En cuanto a los tipos, el BCE subió la facilidad de depósito del $-0,50\%$ al $4\%$ entre julio de 2022 y septiembre de 2023, y la bajó después, desde junio de 2024, hasta el $2\%$ en junio de 2025. Tras la crisis de Oriente Medio, la subió al $2,25\%$ en junio de 2026 y al $2,50\%$ en septiembre de 2026. En marzo de 2024, el BCE revisó su marco operativo, redujo a $15$ puntos básicos el diferencial entre el tipo de las operaciones principales de financiación y la facilidad de depósito y confirmó que esta seguirá siendo el tipo que orienta la política monetaria; por eso, tras la subida de septiembre de 2026, las operaciones principales se remuneran al $2,65\%$.

La retirada del BCE ha tenido tres efectos sobre el Tesoro, que ya hemos ido viendo y que aquí se reúnen. El primero es el aumento de la \textbf{oferta neta al sector privado}, que ha pasado de ser negativa a rondar los $90.000$-$100.000$ millones anuales. El segundo es el cambio en la \textbf{base inversora}: los no residentes han pasado a ser los principales compradores de bonos y obligaciones, con el $47,8\%$ de la deuda en 2025, y los minoristas volvieron a las letras en 2023, tras tocar mínimo en abril de 2022, una demanda concentrada en el corto plazo que justificaba un ligero aumento de la emisión de letras. El tercero es el \textbf{coste}: el coste medio de emisión subió al $3,44\%$ en 2023 y bajó al $3,16\%$ en 2024 y al $2,70\%$ en 2025 con las bajadas del BCE, mientras que el de la cartera subió del $1,64\%$ de 2021 al $2,31\%$ de 2025.

Lo que no ha ocurrido es un aumento de la prima de riesgo. Contra lo que temían muchos analistas en 2022, la retirada del BCE ha coincidido con una caída de la prima española por debajo de $60$ puntos básicos, y ya vimos las tres explicaciones complementarias: la mejora de los fundamentales relativos de España, la existencia del TPI, que actúa como red de seguridad aunque nunca se haya activado, y la gradualidad de la retirada, que el BCE ha descrito como un proceso en segundo plano. El riesgo es que esa calma dependa de que el TPI siga siendo creíble y de que España cumpla sus condiciones de elegibilidad, entre ellas el respeto del marco fiscal europeo, que es precisamente lo que la AIReF pone en duda para 2026.

\subsubsection{La estrategia de financiación de 2026}
La Estrategia de Financiación de 2026, presentada en diciembre de 2025, prevé una emisión neta de $55.000$ millones, igual que en 2025: $50.000$ en bonos y obligaciones y $5.000$ en letras. La emisión bruta será de unos $285.677$ millones, un $4\%$ más que en 2025: $176.935$ millones de bonos y obligaciones, para cubrir unas amortizaciones de unos $126.935$ millones, y unos $108.742$ millones de letras. El programa se ejecuta mediante $48$ subastas ordinarias y varias sindicaciones, para nuevas referencias y para el programa verde, y se completa con $6.500$ millones de préstamos del Mecanismo de Recuperación y Resiliencia, mientras que en 2026 se amortizan unos $3.643$ millones del préstamo del MEDE. El Tesoro se fija como objetivo preservar la vida media de la cartera, que fue de $7,93$ años en 2025, en torno a su máximo histórico por quinto año consecutivo, y limitar los riesgos de refinanciación; mantiene los bonos verdes como elemento estructural del programa y parte de una base inversora en la que los internacionales se han consolidado como principales compradores.

La ejecución hasta septiembre de 2026 ha confirmado la solidez de la demanda, con las sindicaciones de enero, mayo y septiembre que ya vimos. Pero también muestra el cambio del entorno: la obligación a diez años de enero se emitió con un cupón del $3,30\%$, y la verde a veinte años de septiembre, a una rentabilidad cercana al $4,2\%$, ambas muy por encima del $2,31\%$ de coste medio de la cartera, lo que garantiza que el coste medio seguirá subiendo en 2026 y 2027.

Las decisiones de la estrategia pueden leerse a la luz de los objetivos de la política de financiación. La \textbf{continuidad} (misma emisión neta, misma vida media, mismo calendario) responde al objetivo de previsibilidad: en un año de shocks (la guerra en Oriente Medio, las subidas del BCE, la ausencia de presupuestos), el Tesoro evita añadir incertidumbre sobre la oferta. El \textbf{mantenimiento de la vida media} en torno a $8$ años, con tipos subiendo, responde a la disyuntiva entre coste y riesgo: acortar la deuda abarataría la emisión inmediata si la curva mantiene su pendiente, pero aumentaría la exposición a nuevas subidas de tipos y el volumen a refinanciar cada año, que ya ronda el $16\%$ del PIB; con un BCE que sube tipos por un shock de oferta de duración incierta, el Tesoro opta por el seguro frente al ahorro, como en 2020-2021. El \textbf{ligero aumento neto de las letras} ($5.000$ millones), tras años de reducción, admite dos lecturas compatibles: recoge la demanda minorista y de fondos monetarios a corto plazo, y da un pequeño margen de flexibilidad de tesorería en un año con necesidades inciertas por las medidas de emergencia.

La estrategia afronta cinco riesgos. El primero son los \textbf{tipos de interés}: si la inflación energética se prolonga y el BCE sigue subiendo, aumentará el coste de la nueva emisión y, con él, la carga de intereses proyectada, en un contexto global en el que la rentabilidad del bono estadounidense a diez años se acercaba al $5\%$ en septiembre de 2026. El segundo es el \textbf{riesgo fiscal}: la AIReF prevé un crecimiento del gasto neto del $6,4\%$ frente al $3,5\%$ comprometido, y España lleva tres años con los Presupuestos Generales del Estado de 2023 prorrogados, de modo que un incumplimiento prolongado del marco europeo podría afectar a la percepción de los mercados y, en el límite, a la elegibilidad para el TPI. El tercero es el \textbf{contagio}: la inestabilidad fiscal y política de Francia, cuya prima frente a España está en máximos, puede trasladarse al conjunto de la zona del euro en un episodio de aversión al riesgo. El cuarto es la \textbf{concentración de la base inversora} en no residentes ($47,8\%$), más sensibles a las condiciones financieras globales. El quinto es el \textbf{fin de la financiación europea} desde 2027. En sentido contrario, hay factores favorables: las calificaciones siguen mejorando, con la subida de Scope a A+ el 25 de septiembre de 2026; el crecimiento está por encima de la media, el saldo primario cerca del equilibrio y la vida media es larga. Todo ello da al Tesoro un margen que no tenía en 2011-2012.

\subsection{Una valoración}
La política de financiación del Tesoro español es, con criterios internacionales, un ejemplo de buena práctica. Tiene un marco jurídico claro, un mandato de coste y riesgo, procedimientos previsibles y transparentes, una cartera larga, instrumentos estandarizados y líquidos y una base inversora diversificada, y su gestión en 2020-2021, al alargar la vida media en lugar de aprovechar los tipos cortos negativos, resultó acertada con la subida de 2022-2023.

Pero la gestión de la deuda tiene un límite. El Tesoro gestiona la composición de la deuda, su coste y su riesgo, no su volumen, que depende de la política fiscal, y la mejor estrategia de financiación no puede compensar un déficit estructural que no se corrige. Como vimos, la reducción de la deuda en 2021-2025 se debió al crecimiento nominal y no a superávits primarios, y España no ha generado superávits primarios en las fases favorables del ciclo, a diferencia de Portugal. La calidad de la gestión del Tesoro reduce el coste y el riesgo de esa deuda, pero no sustituye a la consolidación fiscal. El éxito de la política de financiación dependerá, en último término, de la capacidad de la economía para crecer y de acometer una nueva consolidación fiscal, porque unos niveles muy elevados de deuda son fuente de vulnerabilidad y pueden afectar al crecimiento, aunque la literatura no ha encontrado un umbral universal a partir del cual lo hagan.\footnote{%
    Sobre el umbral del $90\%$ de REINHART y ROGOFF (2010), la crítica de HERNDON, ASH y POLLIN (2014), que cuestionan el umbral y sugieren que la causalidad puede ir del bajo crecimiento a la deuda alta, y la literatura posterior sobre la relación entre deuda y crecimiento (causalidad inversa, heterogeneidad entre países, ausencia de umbral universal) ([\authorfont{Ver Tema 3.A.38}]).
}


\clearpage
\section{Conclusión}

\subsection{Recapitulación}


\subsection{Extensiones y relación con otras partes del Temario}



\subsection{Opinión}

\subsubsection{Idea final (Salida o cierra)}


\subsection{Cuestiones pendientes de verificar}
\begin{lnum}
    \item El tratamiento de los gastos de la DANA en el cómputo del déficit de 2025 y en la evaluación de la Comisión.
    \item El saldo estructural de 2025 y 2026 según las últimas estimaciones de la Comisión y de la AIReF, y la valoración del cumplimiento de la senda de gasto neto en 2025.
    \item El gasto en intereses del conjunto de las AAPP en 2025 y 2026 en porcentaje del PIB.
    \item El saldo estructural de 2007 y el de 2019 en la última estimación disponible (AMECO, Comisión y AIReF).
    \item La fecha de reclasificación del Fondo de Titulización del Déficit del Sistema Eléctrico (FADE) y su saldo vivo.
    \item La serie homogénea del déficit de las AAPP en 1985-1999 (SEC 2010 retropolado), en particular el máximo de 1993 o 1995, el mínimo de 1986-1990 y el déficit de 1999.
    \item Las ratios de presión fiscal de España y de la media de la OCDE en 1975, 1985 y 1990, y la de 2007, 2009, 2024 y 2025 en la serie SEC 2010.
    \item La media de deuda pública de la CEE en 1975.
    \item El saldo primario de 1993 y la descomposición de la variación de la deuda de ese año.
    \item Los términos de la conversión del saldo deudor del Tesoro con el BdE en 1994 (norma, importe y plazo).
    \item La vida media de la deuda del Estado en 1990, 1995 y 2007, y la proporción de no residentes en 2007-2008.
    \item La fecha de supresión definitiva del coeficiente de inversión obligatoria en deuda pública.
    \item El mecanismo y el porcentaje del requisito de 1984 en pagarés del Tesoro, y la norma de la restricción del crédito del BdE al Tesoro en 1990.
    \item El gasto público en porcentaje del PIB en 1985 y en 1993.
    \item Los importes del Programa de Modernización del Sector Público Empresarial de 1996-2001.
    \item La serie de intereses de las AAPP en porcentaje del PIB en 1995-2007.
    \item Los importes de la asunción de deuda de RENFE (2004) y RTVE (2006-2007) y su impacto en el saldo.
    \item Los umbrales de crecimiento de la reforma de la Ley General de Estabilidad Presupuestaria de 2006.
    \item La serie vigente de déficit y deuda de 2008-2024, en particular el efecto de la Sareb sobre 2012-2014 y el de la revisión del PIB de 2024.
    \item El impacto en el déficit de 2012 de las ayudas al sector bancario.
    \item La recaudación final de la regularización fiscal especial de 2012.
    \item El saldo vivo del préstamo del MEDE en 2026 y las amortizaciones anticipadas posteriores a 2018.
    \item El límite de gasto no financiero del Estado de 2011 y 2012 en términos homogéneos.
    \item El número máximo de trabajadores en ERTE en abril de 2020.
    \item Las pérdidas definitivas por la ejecución de los avales del ICO.
    \item El coste presupuestario de las medidas frente a la crisis energética de 2022-2023.
    \item El volumen del Plan de Recuperación tras la adenda, los préstamos recibidos y la solicitud del último desembolso.
    \item El coste final del Plan de Respuesta a la crisis de Oriente Medio y su tratamiento como medida excepcional por la Comisión.
    \item La fecha en que España recuperó el nivel de PIB real previo a la pandemia.
    \item La tramitación del proyecto de ley de condonación de deuda autonómica.
    \item El déficit medio de la OCDE en los años setenta, en los ochenta y en 1993.
    \item El déficit de Suecia en 1993 y el de Canadá en 1992-1993.
    \item La cuantía de la eurotasa italiana y del pago de France Télécom de 1997.
    \item El alcance de la revisión de las estadísticas griegas de 2004.
    \item El déficit por cuenta corriente y la deuda privada de España en 2007.
    \item La diferencia entre la deuda de Estados Unidos según el GFSM y según el criterio de Maastricht.
    \item Las caídas del PIB de 2009 de España, Alemania, Italia y el Reino Unido en las series actuales.
    \item La consolidación estructural primaria acumulada en 2010-2013 en Grecia, Portugal, Irlanda y España.
    \item El coste fiscal bruto de las crisis bancarias de Irlanda y España según LAEVEN y VALENCIA (2018).
    \item Los saldos primarios de Portugal, Grecia, Italia y Chipre en 2022-2025.
    \item El ranking de déficits de las economías avanzadas en 2019 con los datos actuales del FMI.
    \item Las cifras de apoyo directo y de liquidez frente a la pandemia de España y sus pares en la base de medidas del FMI.
    \item El apoyo energético de España en la base de datos de Bruegel.
    \item La deuda irlandesa sobre la renta nacional bruta modificada en 2019 y 2023.
    \item La media de deuda de la zona del euro en 2014 y 2023.
    \item La clasificación de riesgo de España en el \textit{Debt Sustainability Monitor} 2025 y sus indicadores S1 y S2.
    \item El gasto y los ingresos públicos de España en 2025 en porcentaje del PIB.
    \item La posición de inversión internacional neta y la deuda privada consolidada de España en 2025.
    \item La prima de riesgo española y la rentabilidad del bono a diez años a finales de septiembre de 2026.
    \item Las calificaciones de Francia e Italia a septiembre de 2026.
    \item La deuda y el déficit de Estados Unidos, Japón y el Reino Unido en el \textit{Fiscal Monitor} de abril de 2026, y la deuda neta de Japón.
    \item El déficit alemán de 2025 y su trayectoria tras la reforma del freno de la deuda.
    \item El coste medio y la vida media de la deuda en circulación a finales de 2025.
    \item Las tenencias de deuda española del Eurosistema y de los no residentes a mediados de 2026, y el ritmo anual de vencimientos no reinvertidos.
    \item La serie anual de emisión neta y bruta del Tesoro de 2008 a 2025.
    \item La descomposición de la emisión neta de 2026 por destinos.
    \item El umbral de necesidades brutas de financiación del marco de sostenibilidad del FMI para las economías avanzadas.
    \item El calendario de amortización de los préstamos del SURE y del Mecanismo de Recuperación y Resiliencia.
    \item La evolución del colchón de tesorería del Estado en 2020-2025.
    \item El peso de las letras en la deuda del Estado a mediados de 2026.
    \item El saldo y el peso de los bonos indexados, y la última referencia indexada emitida.
    \item El saldo de bonos verdes tras la emisión de septiembre de 2026 y el último informe de asignación e impacto.
    \item Las fechas, importes y rentabilidades de las emisiones a $50$ años.
    \item El saldo de inversión minorista en letras en 2024-2026.
    \item El estado de ratificación de la reforma del Tratado del MEDE.
    \item La proporción de deuda indexada del Reino Unido y su coste en 2022-2023.
    \item El reparto de la emisión bruta entre subastas y sindicaciones en 2024-2025, y el de las $48$ subastas de 2026 por tipo de instrumento.
    \item El número de subastas resueltas por debajo del rango objetivo.
    \item El número actual de creadores de mercado y la resolución vigente que regula su actuación.
    \item Las operaciones de canje y recompra del Tesoro en 2023-2026.
    \item El volumen de las operaciones de colocación de tesorería.
    \item La fecha y el funcionamiento de la plataforma del Tesoro para inversores particulares.
    \item Las pérdidas del BdE por política monetaria en 2023 y las transferencias de beneficios al Tesoro en 2019-2022.
    \item El calendario de reducción de las reinversiones del PEPP en el segundo semestre de 2024.
    \item El coste medio de emisión en 2026 y la ejecución del programa de emisión hasta septiembre.
    \item La rentabilidad del bono estadounidense a diez años en septiembre de 2026.
    \item Una eventual revisión de la emisión neta de 2026 por las medidas de emergencia.
\end{lnum}

\clearpage
\pagenumbering{gobble}
\MyBibliography

{\bfseries\large Legislación\par}
\medskip

\noindent\hangindent=1.5em\hangafter=1 Tratado de la Unión Europea, firmado en Maastricht el 7 de febrero de 1992. \textit{Diario Oficial de las Comunidades Europeas}, C 191, de 29 de julio de 1992, pp. 1--112. \href{http://data.europa.eu/eli/treaty/teu/sign}{\nolinkurl{http://data.europa.eu/eli/treaty/teu/sign}}

\noindent\hangindent=1.5em\hangafter=1 Reglamento (CE) n.º 479/2009 del Consejo, de 25 de mayo de 2009, relativo a la aplicación del Protocolo sobre el procedimiento aplicable en caso de déficit excesivo, anejo al Tratado constitutivo de la Comunidad Europea (versión codificada). \textit{Diario Oficial de la Unión Europea}, L 145, de 10 de junio de 2009, pp. 1--9. \href{http://data.europa.eu/eli/reg/2009/479/oj}{\nolinkurl{http://data.europa.eu/eli/reg/2009/479/oj}}

\noindent\hangindent=1.5em\hangafter=1 Reglamento (UE) n.º 1176/2011 del Parlamento Europeo y del Consejo, de 16 de noviembre de 2011, relativo a la prevención y corrección de los desequilibrios macroeconómicos. \textit{Diario Oficial de la Unión Europea}, L 306, de 23 de noviembre de 2011, pp. 25--32. \href{http://data.europa.eu/eli/reg/2011/1176/oj}{\nolinkurl{http://data.europa.eu/eli/reg/2011/1176/oj}}

\noindent\hangindent=1.5em\hangafter=1 Tratado Constitutivo del Mecanismo Europeo de Estabilidad, hecho en Bruselas el 2 de febrero de 2012. \textit{Boletín Oficial del Estado}, 239, de 4 de octubre de 2012. \href{https://www.boe.es/buscar/doc.php?id=BOE-A-2012-12378}{\nolinkurl{https://www.boe.es/buscar/doc.php?id=BOE-A-2012-12378}}

\noindent\hangindent=1.5em\hangafter=1 Reglamento (UE) n.º 549/2013 del Parlamento Europeo y del Consejo, de 21 de mayo de 2013, relativo al Sistema Europeo de Cuentas Nacionales y Regionales de la Unión Europea. \textit{Diario Oficial de la Unión Europea}, L 174, de 26 de junio de 2013, pp. 1--727. \href{http://data.europa.eu/eli/reg/2013/549/oj}{\nolinkurl{http://data.europa.eu/eli/reg/2013/549/oj}}

\noindent\hangindent=1.5em\hangafter=1 Versión consolidada del Tratado de Funcionamiento de la Unión Europea. \textit{Diario Oficial de la Unión Europea}, C 202, de 7 de junio de 2016, pp. 47--390. \href{http://data.europa.eu/eli/treaty/tfeu_2016/oj}{\nolinkurl{http://data.europa.eu/eli/treaty/tfeu_2016/oj}}

\noindent\hangindent=1.5em\hangafter=1 Reglamento (UE) 2024/1263 del Parlamento Europeo y del Consejo, de 29 de abril de 2024, relativo a la coordinación eficaz de las políticas económicas y a la supervisión presupuestaria multilateral y por el que se deroga el Reglamento (CE) n.º 1466/97 del Consejo. \textit{Diario Oficial de la Unión Europea}, L, 2024/1263, de 30 de abril de 2024. \href{http://data.europa.eu/eli/reg/2024/1263/oj}{\nolinkurl{http://data.europa.eu/eli/reg/2024/1263/oj}}

\noindent\hangindent=1.5em\hangafter=1 Reglamento (UE) 2024/1264 del Consejo, de 29 de abril de 2024, por el que se modifica el Reglamento (CE) n.º 1467/97, relativo a la aceleración y clarificación del procedimiento de déficit excesivo. \textit{Diario Oficial de la Unión Europea}, L, 2024/1264, de 30 de abril de 2024. \href{http://data.europa.eu/eli/reg/2024/1264/oj}{\nolinkurl{http://data.europa.eu/eli/reg/2024/1264/oj}}

\noindent\hangindent=1.5em\hangafter=1 Recomendación del Consejo, de 21 de enero de 2025, por la que se aprueba el plan fiscal-estructural nacional a medio plazo de España. \textit{Diario Oficial de la Unión Europea}, C/2025/643, de 10 de febrero de 2025. \href{http://data.europa.eu/eli/C/2025/643/oj}{\nolinkurl{http://data.europa.eu/eli/C/2025/643/oj}}

\noindent\hangindent=1.5em\hangafter=1 Sentencia del Tribunal de Justicia de 16 de junio de 2015, \textit{Gauweiler y otros}, C-62/14, ECLI:EU:C:2015:400. \href{https://eur-lex.europa.eu/legal-content/ES/TXT/?uri=CELEX:62014CJ0062}{\nolinkurl{https://eur-lex.europa.eu/legal-content/ES/TXT/?uri=CELEX:62014CJ0062}}

\noindent\hangindent=1.5em\hangafter=1 Sentencia del Tribunal de Justicia de 11 de diciembre de 2018, \textit{Weiss y otros}, C-493/17, ECLI:EU:C:2018:1000. \href{https://eur-lex.europa.eu/legal-content/ES/TXT/?uri=CELEX:62017CJ0493}{\nolinkurl{https://eur-lex.europa.eu/legal-content/ES/TXT/?uri=CELEX:62017CJ0493}}

\noindent\hangindent=1.5em\hangafter=1 Bundesverfassungsgericht. Sentencia de la Sala Segunda de 5 de mayo de 2020, 2 BvR 859/15, 2 BvR 1651/15, 2 BvR 2006/15, 2 BvR 980/16. \href{https://www.bundesverfassungsgericht.de/SharedDocs/Entscheidungen/EN/2020/05/rs20200505_2bvr085915en.html}{\nolinkurl{https://www.bundesverfassungsgericht.de/SharedDocs/Entscheidungen/EN/2020/05/rs20200505_2bvr085915en.html}}

\noindent\hangindent=1.5em\hangafter=1 Ley 50/1977, de 14 de noviembre, sobre medidas urgentes de reforma fiscal. \textit{Boletín Oficial del Estado}, de 16 de noviembre de 1977. \href{https://www.boe.es/buscar/doc.php?id=BOE-A-1977-27150}{\nolinkurl{https://www.boe.es/buscar/doc.php?id=BOE-A-1977-27150}}

\noindent\hangindent=1.5em\hangafter=1 Ley 44/1978, de 8 de septiembre, del Impuesto sobre la Renta de las Personas Físicas. \textit{Boletín Oficial del Estado}, 217, de 11 de septiembre de 1978. \href{https://www.boe.es/buscar/doc.php?id=BOE-A-1978-23326}{\nolinkurl{https://www.boe.es/buscar/doc.php?id=BOE-A-1978-23326}}

\noindent\hangindent=1.5em\hangafter=1 Ley 61/1978, de 27 de diciembre, del Impuesto sobre Sociedades. \textit{Boletín Oficial del Estado}, 312, de 30 de diciembre de 1978. \href{https://www.boe.es/eli/es/l/1978/12/27/61}{\nolinkurl{https://www.boe.es/eli/es/l/1978/12/27/61}}

\noindent\hangindent=1.5em\hangafter=1 Ley Orgánica 8/1980, de 22 de septiembre, de Financiación de las Comunidades Autónomas. \textit{Boletín Oficial del Estado}, 236, de 1 de octubre de 1980. \href{https://www.boe.es/eli/es/lo/1980/09/22/8}{\nolinkurl{https://www.boe.es/eli/es/lo/1980/09/22/8}}

\noindent\hangindent=1.5em\hangafter=1 Ley 27/1984, de 26 de julio, sobre reconversión y reindustrialización. \textit{Boletín Oficial del Estado}, 180, de 28 de julio de 1984. \href{https://www.boe.es/eli/es/l/1984/07/26/27}{\nolinkurl{https://www.boe.es/eli/es/l/1984/07/26/27}}

\noindent\hangindent=1.5em\hangafter=1 Ley 31/1984, de 2 de agosto, de protección por Desempleo, por la que se modifica el título II de la Ley 51/1980, de 8 de octubre. \textit{Boletín Oficial del Estado}, 186, de 4 de agosto de 1984. \href{https://www.boe.es/eli/es/l/1984/08/02/31}{\nolinkurl{https://www.boe.es/eli/es/l/1984/08/02/31}}

\noindent\hangindent=1.5em\hangafter=1 Ley 32/1984, de 2 de agosto, sobre modificación de determinados artículos de la Ley 8/1980, de 10 de marzo, del Estatuto de los Trabajadores. \textit{Boletín Oficial del Estado}, 186, de 4 de agosto de 1984. \href{https://www.boe.es/eli/es/l/1984/08/02/32}{\nolinkurl{https://www.boe.es/eli/es/l/1984/08/02/32}}

\noindent\hangindent=1.5em\hangafter=1 Ley 13/1985, de 25 de mayo, de coeficientes de inversión, recursos propios y obligaciones de información de los intermediarios financieros. \textit{Boletín Oficial del Estado}, 127, de 28 de mayo de 1985. \href{https://www.boe.es/eli/es/l/1985/05/25/13}{\nolinkurl{https://www.boe.es/eli/es/l/1985/05/25/13}}

\noindent\hangindent=1.5em\hangafter=1 Ley 30/1985, de 2 de agosto, del Impuesto sobre el Valor Añadido. \textit{Boletín Oficial del Estado}, 190, de 9 de agosto de 1985. \href{https://www.boe.es/eli/es/l/1985/08/02/30}{\nolinkurl{https://www.boe.es/eli/es/l/1985/08/02/30}}

\noindent\hangindent=1.5em\hangafter=1 Ley 14/1986, de 25 de abril, General de Sanidad. \textit{Boletín Oficial del Estado}, 102, de 29 de abril de 1986. \href{https://www.boe.es/eli/es/l/1986/04/25/14}{\nolinkurl{https://www.boe.es/eli/es/l/1986/04/25/14}}

\noindent\hangindent=1.5em\hangafter=1 Ley 26/1990, de 20 de diciembre, por la que se establecen en la Seguridad Social prestaciones no contributivas. \textit{Boletín Oficial del Estado}, 306, de 22 de diciembre de 1990. \href{https://www.boe.es/eli/es/l/1990/12/20/26}{\nolinkurl{https://www.boe.es/eli/es/l/1990/12/20/26}}

\noindent\hangindent=1.5em\hangafter=1 Ley 18/1991, de 6 de junio, del Impuesto sobre la Renta de las Personas Físicas. \textit{Boletín Oficial del Estado}, 136, de 7 de junio de 1991. \href{https://www.boe.es/eli/es/l/1991/06/06/18}{\nolinkurl{https://www.boe.es/eli/es/l/1991/06/06/18}}

\noindent\hangindent=1.5em\hangafter=1 Ley 19/1991, de 6 de junio, del Impuesto sobre el Patrimonio. \textit{Boletín Oficial del Estado}, 136, de 7 de junio de 1991. \href{https://www.boe.es/eli/es/l/1991/06/06/19}{\nolinkurl{https://www.boe.es/eli/es/l/1991/06/06/19}}

\noindent\hangindent=1.5em\hangafter=1 Real Decreto-ley 1/1992, de 3 de abril, de medidas urgentes sobre fomento del empleo y protección por desempleo. \textit{Boletín Oficial del Estado}, 84, de 7 de abril de 1992. \href{https://www.boe.es/eli/es/rdl/1992/04/03/1}{\nolinkurl{https://www.boe.es/eli/es/rdl/1992/04/03/1}}

\noindent\hangindent=1.5em\hangafter=1 Ley 22/1992, de 30 de julio, de medidas urgentes sobre fomento del empleo y protección por desempleo. \textit{Boletín Oficial del Estado}, 186, de 4 de agosto de 1992. \href{https://www.boe.es/eli/es/l/1992/07/30/22}{\nolinkurl{https://www.boe.es/eli/es/l/1992/07/30/22}}

\noindent\hangindent=1.5em\hangafter=1 Ley 13/1994, de 1 de junio, de Autonomía del Banco de España. \textit{Boletín Oficial del Estado}, 131, de 2 de junio de 1994. \href{https://www.boe.es/eli/es/l/1994/06/01/13}{\nolinkurl{https://www.boe.es/eli/es/l/1994/06/01/13}}

\noindent\hangindent=1.5em\hangafter=1 Ley 24/1997, de 15 de julio, de Consolidación y Racionalización del Sistema de Seguridad Social. \textit{Boletín Oficial del Estado}, 169, de 16 de julio de 1997. \href{https://www.boe.es/eli/es/l/1997/07/15/24}{\nolinkurl{https://www.boe.es/eli/es/l/1997/07/15/24}}

\noindent\hangindent=1.5em\hangafter=1 Ley 40/1998, de 9 de diciembre, del Impuesto sobre la Renta de las Personas Físicas y otras Normas Tributarias. \textit{Boletín Oficial del Estado}, 295, de 10 de diciembre de 1998. \href{https://www.boe.es/eli/es/l/1998/12/09/40}{\nolinkurl{https://www.boe.es/eli/es/l/1998/12/09/40}}

\noindent\hangindent=1.5em\hangafter=1 Ley 18/2001, de 12 de diciembre, General de Estabilidad Presupuestaria. \textit{Boletín Oficial del Estado}, 298, de 13 de diciembre de 2001. \href{https://www.boe.es/eli/es/l/2001/12/12/18}{\nolinkurl{https://www.boe.es/eli/es/l/2001/12/12/18}}

\noindent\hangindent=1.5em\hangafter=1 Ley 44/2002, de 22 de noviembre, de Medidas de Reforma del Sistema Financiero. \textit{Boletín Oficial del Estado}, 281, de 23 de noviembre de 2002. \href{https://www.boe.es/eli/es/l/2002/11/22/44}{\nolinkurl{https://www.boe.es/eli/es/l/2002/11/22/44}}

\noindent\hangindent=1.5em\hangafter=1 Ley 46/2002, de 18 de diciembre, de reforma parcial del Impuesto sobre la Renta de las Personas Físicas y por la que se modifican las Leyes de los Impuestos sobre Sociedades y sobre la Renta de no Residentes. \textit{Boletín Oficial del Estado}, 303, de 19 de diciembre de 2002. \href{https://www.boe.es/eli/es/l/2002/12/18/46}{\nolinkurl{https://www.boe.es/eli/es/l/2002/12/18/46}}

\noindent\hangindent=1.5em\hangafter=1 Ley 47/2003, de 26 de noviembre, General Presupuestaria. \textit{Boletín Oficial del Estado}, 284, de 27 de noviembre de 2003. \href{https://www.boe.es/eli/es/l/2003/11/26/47}{\nolinkurl{https://www.boe.es/eli/es/l/2003/11/26/47}}

\noindent\hangindent=1.5em\hangafter=1 Ley 35/2006, de 28 de noviembre, del Impuesto sobre la Renta de las Personas Físicas y de modificación parcial de las leyes de los Impuestos sobre Sociedades, sobre la Renta de no Residentes y sobre el Patrimonio. \textit{Boletín Oficial del Estado}, 285, de 29 de noviembre de 2006. \href{https://www.boe.es/eli/es/l/2006/11/28/35}{\nolinkurl{https://www.boe.es/eli/es/l/2006/11/28/35}}

\noindent\hangindent=1.5em\hangafter=1 Real Decreto-ley 2/2008, de 21 de abril, de medidas de impulso a la actividad económica. \textit{Boletín Oficial del Estado}, 97, de 22 de abril de 2008. \href{https://www.boe.es/eli/es/rdl/2008/04/21/2}{\nolinkurl{https://www.boe.es/eli/es/rdl/2008/04/21/2}}

\noindent\hangindent=1.5em\hangafter=1 Real Decreto-ley 9/2008, de 28 de noviembre, por el que se crean un Fondo Estatal de Inversión Local y un Fondo Especial del Estado para la Dinamización de la Economía y el Empleo y se aprueban créditos extraordinarios para atender a su financiación. \textit{Boletín Oficial del Estado}, 290, de 2 de diciembre de 2008. \href{https://www.boe.es/eli/es/rdl/2008/11/28/9}{\nolinkurl{https://www.boe.es/eli/es/rdl/2008/11/28/9}}

\noindent\hangindent=1.5em\hangafter=1 Ley 4/2008, de 23 de diciembre, por la que se suprime el gravamen del Impuesto sobre el Patrimonio, se generaliza el sistema de devolución mensual en el Impuesto sobre el Valor Añadido, y se introducen otras modificaciones en la normativa tributaria. \textit{Boletín Oficial del Estado}, 310, de 25 de diciembre de 2008. \href{https://www.boe.es/eli/es/l/2008/12/23/4}{\nolinkurl{https://www.boe.es/eli/es/l/2008/12/23/4}}

\noindent\hangindent=1.5em\hangafter=1 Real Decreto-ley 13/2009, de 26 de octubre, por el que se crea el Fondo Estatal para el Empleo y la Sostenibilidad Local. \textit{Boletín Oficial del Estado}, 259, de 27 de octubre de 2009. \href{https://www.boe.es/eli/es/rdl/2009/10/26/13}{\nolinkurl{https://www.boe.es/eli/es/rdl/2009/10/26/13}}

\noindent\hangindent=1.5em\hangafter=1 Ley 22/2009, de 18 de diciembre, por la que se regula el sistema de financiación de las Comunidades Autónomas de régimen común y Ciudades con Estatuto de Autonomía y se modifican determinadas normas tributarias. \textit{Boletín Oficial del Estado}, 305, de 19 de diciembre de 2009. \href{https://www.boe.es/eli/es/l/2009/12/18/22}{\nolinkurl{https://www.boe.es/eli/es/l/2009/12/18/22}}

\noindent\hangindent=1.5em\hangafter=1 Real Decreto-ley 8/2010, de 20 de mayo, por el que se adoptan medidas extraordinarias para la reducción del déficit público. \textit{Boletín Oficial del Estado}, 126, de 24 de mayo de 2010. \href{https://www.boe.es/eli/es/rdl/2010/05/20/8}{\nolinkurl{https://www.boe.es/eli/es/rdl/2010/05/20/8}}

\noindent\hangindent=1.5em\hangafter=1 Real Decreto-ley 13/2011, de 16 de septiembre, por el que se restablece el Impuesto sobre el Patrimonio, con carácter temporal. \textit{Boletín Oficial del Estado}, 224, de 17 de septiembre de 2011. \href{https://www.boe.es/eli/es/rdl/2011/09/16/13}{\nolinkurl{https://www.boe.es/eli/es/rdl/2011/09/16/13}}

\noindent\hangindent=1.5em\hangafter=1 Reforma del artículo 135 de la Constitución Española, de 27 de septiembre de 2011. \textit{Boletín Oficial del Estado}, 233, de 27 de septiembre de 2011. \href{https://www.boe.es/buscar/doc.php?id=BOE-A-2011-15210}{\nolinkurl{https://www.boe.es/buscar/doc.php?id=BOE-A-2011-15210}}

\noindent\hangindent=1.5em\hangafter=1 Real Decreto-ley 20/2011, de 30 de diciembre, de medidas urgentes en materia presupuestaria, tributaria y financiera para la corrección del déficit público. \textit{Boletín Oficial del Estado}, 315, de 31 de diciembre de 2011. \href{https://www.boe.es/eli/es/rdl/2011/12/30/20}{\nolinkurl{https://www.boe.es/eli/es/rdl/2011/12/30/20}}

\noindent\hangindent=1.5em\hangafter=1 Real Decreto-ley 4/2012, de 24 de febrero, por el que se determinan obligaciones de información y procedimientos necesarios para establecer un mecanismo de financiación para el pago a los proveedores de las entidades locales. \textit{Boletín Oficial del Estado}, 48, de 25 de febrero de 2012. \href{https://www.boe.es/eli/es/rdl/2012/02/24/4}{\nolinkurl{https://www.boe.es/eli/es/rdl/2012/02/24/4}}

\noindent\hangindent=1.5em\hangafter=1 Real Decreto-ley 7/2012, de 9 de marzo, por el que se crea el Fondo para la financiación de los pagos a proveedores. \textit{Boletín Oficial del Estado}, 60, de 10 de marzo de 2012. \href{https://www.boe.es/eli/es/rdl/2012/03/09/7}{\nolinkurl{https://www.boe.es/eli/es/rdl/2012/03/09/7}}

\noindent\hangindent=1.5em\hangafter=1 Real Decreto-ley 12/2012, de 30 de marzo, por el que se introducen diversas medidas tributarias y administrativas dirigidas a la reducción del déficit público. \textit{Boletín Oficial del Estado}, 78, de 31 de marzo de 2012. \href{https://www.boe.es/eli/es/rdl/2012/03/30/12}{\nolinkurl{https://www.boe.es/eli/es/rdl/2012/03/30/12}}

\noindent\hangindent=1.5em\hangafter=1 Ley Orgánica 2/2012, de 27 de abril, de Estabilidad Presupuestaria y Sostenibilidad Financiera. \textit{Boletín Oficial del Estado}, 103, de 30 de abril de 2012. \href{https://www.boe.es/eli/es/lo/2012/04/27/2}{\nolinkurl{https://www.boe.es/eli/es/lo/2012/04/27/2}}

\noindent\hangindent=1.5em\hangafter=1 Real Decreto-ley 20/2012, de 13 de julio, de medidas para garantizar la estabilidad presupuestaria y de fomento de la competitividad. \textit{Boletín Oficial del Estado}, 168, de 14 de julio de 2012. \href{https://www.boe.es/eli/es/rdl/2012/07/13/20}{\nolinkurl{https://www.boe.es/eli/es/rdl/2012/07/13/20}}

\noindent\hangindent=1.5em\hangafter=1 Real Decreto-ley 21/2012, de 13 de julio, de medidas de liquidez de las Administraciones públicas y en el ámbito financiero. \textit{Boletín Oficial del Estado}, 168, de 14 de julio de 2012. \href{https://www.boe.es/eli/es/rdl/2012/07/13/21}{\nolinkurl{https://www.boe.es/eli/es/rdl/2012/07/13/21}}

\noindent\hangindent=1.5em\hangafter=1 Ley 23/2013, de 23 de diciembre, reguladora del Factor de Sostenibilidad y del Índice de Revalorización del Sistema de Pensiones de la Seguridad Social. \textit{Boletín Oficial del Estado}, 309, de 26 de diciembre de 2013. \href{https://www.boe.es/eli/es/l/2013/12/23/23}{\nolinkurl{https://www.boe.es/eli/es/l/2013/12/23/23}}

\noindent\hangindent=1.5em\hangafter=1 Ley 24/2013, de 26 de diciembre, del Sector Eléctrico. \textit{Boletín Oficial del Estado}, 310, de 27 de diciembre de 2013. \href{https://www.boe.es/eli/es/l/2013/12/26/24}{\nolinkurl{https://www.boe.es/eli/es/l/2013/12/26/24}}

\noindent\hangindent=1.5em\hangafter=1 Real Decreto-ley 17/2014, de 26 de diciembre, de medidas de sostenibilidad financiera de las comunidades autónomas y entidades locales y otras de carácter económico. \textit{Boletín Oficial del Estado}, 315, de 30 de diciembre de 2014. \href{https://www.boe.es/eli/es/rdl/2014/12/26/17}{\nolinkurl{https://www.boe.es/eli/es/rdl/2014/12/26/17}}

\noindent\hangindent=1.5em\hangafter=1 Real Decreto-ley 5/2021, de 12 de marzo, de medidas extraordinarias de apoyo a la solvencia empresarial en respuesta a la pandemia de la COVID-19. \textit{Boletín Oficial del Estado}, 62, de 13 de marzo de 2021. \href{https://www.boe.es/eli/es/rdl/2021/03/12/5}{\nolinkurl{https://www.boe.es/eli/es/rdl/2021/03/12/5}}

\noindent\hangindent=1.5em\hangafter=1 Ley 21/2021, de 28 de diciembre, de garantía del poder adquisitivo de las pensiones y de otras medidas de refuerzo de la sostenibilidad financiera y social del sistema público de pensiones. \textit{Boletín Oficial del Estado}, 312, de 29 de diciembre de 2021. \href{https://www.boe.es/eli/es/l/2021/12/28/21}{\nolinkurl{https://www.boe.es/eli/es/l/2021/12/28/21}}

\noindent\hangindent=1.5em\hangafter=1 Ley 31/2022, de 23 de diciembre, de Presupuestos Generales del Estado para el año 2023. \textit{Boletín Oficial del Estado}, 308, de 24 de diciembre de 2022. \href{https://www.boe.es/eli/es/l/2022/12/23/31}{\nolinkurl{https://www.boe.es/eli/es/l/2022/12/23/31}}

\noindent\hangindent=1.5em\hangafter=1 Ley 38/2022, de 27 de diciembre, para el establecimiento de gravámenes temporales energético y de entidades de crédito y establecimientos financieros de crédito y por la que se crea el impuesto temporal de solidaridad de las grandes fortunas, y se modifican determinadas normas tributarias. \textit{Boletín Oficial del Estado}, 311, de 28 de diciembre de 2022. \href{https://www.boe.es/eli/es/l/2022/12/27/38}{\nolinkurl{https://www.boe.es/eli/es/l/2022/12/27/38}}

\noindent\hangindent=1.5em\hangafter=1 Real Decreto 410/2024, de 23 de abril, por el que se desarrolla la estructura orgánica básica del Ministerio de Economía, Comercio y Empresa. \textit{Boletín Oficial del Estado}, 100, de 24 de abril de 2024. \href{https://www.boe.es/eli/es/rd/2024/04/23/410}{\nolinkurl{https://www.boe.es/eli/es/rd/2024/04/23/410}}

\noindent\hangindent=1.5em\hangafter=1 Ley 7/2024, de 20 de diciembre, por la que se establecen un Impuesto Complementario para garantizar un nivel mínimo global de imposición para los grupos multinacionales y los grupos nacionales de gran magnitud, un Impuesto sobre el margen de intereses y comisiones de determinadas entidades financieras y un Impuesto sobre los líquidos para cigarrillos electrónicos y otros productos relacionados con el tabaco, y se modifican otras normas tributarias. \textit{Boletín Oficial del Estado}, 307, de 21 de diciembre de 2024. \href{https://www.boe.es/eli/es/l/2024/12/20/7}{\nolinkurl{https://www.boe.es/eli/es/l/2024/12/20/7}}

\noindent\hangindent=1.5em\hangafter=1 Orden ECM/2/2026, de 9 de enero, por la que se dispone la creación de Deuda del Estado durante el año 2026 y enero de 2027. \textit{Boletín Oficial del Estado}, 11, de 13 de enero de 2026. \href{https://www.boe.es/eli/es/o/2026/01/09/ecm2}{\nolinkurl{https://www.boe.es/eli/es/o/2026/01/09/ecm2}}

\noindent\hangindent=1.5em\hangafter=1 Real Decreto 644/2026, de 28 de julio, por el que se modifica el Real Decreto 1009/2023, de 5 de diciembre, por el que se establece la estructura orgánica básica de los departamentos ministeriales. \textit{Boletín Oficial del Estado}, 184, de 29 de julio de 2026. \href{https://www.boe.es/diario_boe/txt.php?id=BOE-A-2026-16441}{\nolinkurl{https://www.boe.es/diario_boe/txt.php?id=BOE-A-2026-16441}}

\noindent\hangindent=1.5em\hangafter=1 Real Decreto-ley 25/2026, de 29 de septiembre, por el que se adoptan y se prorrogan determinadas medidas en el marco del Plan Integral de Respuesta a la Crisis en Oriente Medio. \textit{Boletín Oficial del Estado}, 241, de 30 de septiembre de 2026. \href{https://www.boe.es/diario_boe/txt.php?id=BOE-A-2026-20265}{\nolinkurl{https://www.boe.es/diario_boe/txt.php?id=BOE-A-2026-20265}}

\bigskip
{\bfseries\large Investigación\par}
\medskip

\noindent\hangindent=1.5em\hangafter=1 Abbas, S. M. A., Belhocine, N., ElGanainy, A. y Horton, M. (2010). \textit{A historical public debt database} (IMF Working Paper n.º 10/245). Fondo Monetario Internacional. \href{https://doi.org/10.5089/9781455209453.001}{\nolinkurl{https://doi.org/10.5089/9781455209453.001}}

\noindent\hangindent=1.5em\hangafter=1 Abbink, K., Brandts, J. y Pezanis-Christou, P. (2006). Auctions for government securities: A laboratory comparison of uniform, discriminatory and Spanish designs. \textit{Journal of Economic Behavior \& Organization, 61}(2), 284--303. \href{https://doi.org/10.1016/j.jebo.2004.12.007}{\nolinkurl{https://doi.org/10.1016/j.jebo.2004.12.007}}

\noindent\hangindent=1.5em\hangafter=1 AIReF (2025). \textit{Informe sobre la ejecución presupuestaria, deuda pública y regla de gasto 2025}. Autoridad Independiente de Responsabilidad Fiscal. \href{https://www.airef.es/es/centro-documental/informes/}{\nolinkurl{https://www.airef.es/es/centro-documental/informes/}}

\noindent\hangindent=1.5em\hangafter=1 AIReF (2026a). \textit{Informe de seguimiento 2026 del Plan Fiscal y Estructural de Medio Plazo 2025-2028}. Autoridad Independiente de Responsabilidad Fiscal. \href{https://www.airef.es/es/centro-documental/informes/informe-de-seguimiento-2026-del-plan-fiscal-y-estructural-de-medio-plazo-2025-2028/}{\nolinkurl{https://www.airef.es/es/centro-documental/informes/informe-de-seguimiento-2026-del-plan-fiscal-y-estructural-de-medio-plazo-2025-2028/}}

\noindent\hangindent=1.5em\hangafter=1 AIReF (2026b). \textit{Informe sobre la ejecución presupuestaria, deuda pública y regla de gasto 2026} (julio). Autoridad Independiente de Responsabilidad Fiscal. \href{https://www.airef.es/es/centro-documental/informes/}{\nolinkurl{https://www.airef.es/es/centro-documental/informes/}}

\noindent\hangindent=1.5em\hangafter=1 Alberola, E., Estrada, Á. y Santabárbara, D. (2014). Growth and imbalances in Spain: A reassessment of the output gap. \textit{SERIEs, 5}(2--3), 333--356. \href{https://doi.org/10.1007/s13209-014-0112-z}{\nolinkurl{https://doi.org/10.1007/s13209-014-0112-z}}

\noindent\hangindent=1.5em\hangafter=1 Alesina, A. y Ardagna, S. (2010). Large changes in fiscal policy: Taxes versus spending. \textit{Tax Policy and the Economy, 24}, 35--68. \href{https://doi.org/10.1086/649828}{\nolinkurl{https://doi.org/10.1086/649828}}

\noindent\hangindent=1.5em\hangafter=1 Alesina, A. y Drazen, A. (1991). Why are stabilizations delayed? \textit{The American Economic Review, 81}(5), 1170--1188. \href{https://econpapers.repec.org/RePEc:aea:aecrev:v:81:y:1991:i:5:p:1170-88}{\nolinkurl{https://econpapers.repec.org/RePEc:aea:aecrev:v:81:y:1991:i:5:p:1170-88}}

\noindent\hangindent=1.5em\hangafter=1 Alesina, A. y Perotti, R. (1995). Fiscal expansions and adjustments in OECD countries. \textit{Economic Policy, 10}(21), 205--248. \href{https://doi.org/10.2307/1344590}{\nolinkurl{https://doi.org/10.2307/1344590}}

\noindent\hangindent=1.5em\hangafter=1 Alesina, A. y Tabellini, G. (1990). A positive theory of fiscal deficits and government debt. \textit{The Review of Economic Studies, 57}(3), 403--414. \href{https://doi.org/10.2307/2298021}{\nolinkurl{https://doi.org/10.2307/2298021}}

\noindent\hangindent=1.5em\hangafter=1 Alesina, A., Favero, C. y Giavazzi, F. (2019). \textit{Austerity: When it works and when it doesn't}. Princeton University Press. \href{https://doi.org/10.2307/j.ctvc77f4b}{\nolinkurl{https://doi.org/10.2307/j.ctvc77f4b}}

\noindent\hangindent=1.5em\hangafter=1 Álvarez, F. y Mazón, C. (2007). Comparing the Spanish and the discriminatory auction formats: A discrete model with private information. \textit{European Journal of Operational Research, 179}(1), 253--266. \href{https://ideas.repec.org/a/eee/ejores/v179y2007i1p253-266.html}{\nolinkurl{https://ideas.repec.org/a/eee/ejores/v179y2007i1p253-266.html}}

\noindent\hangindent=1.5em\hangafter=1 Angeletos, G.-M. (2002). Fiscal policy with noncontingent debt and the optimal maturity structure. \textit{The Quarterly Journal of Economics, 117}(3), 1105--1131. \href{https://doi.org/10.1162/003355302760193977}{\nolinkurl{https://doi.org/10.1162/003355302760193977}}

\noindent\hangindent=1.5em\hangafter=1 Arellano, C. y Ramanarayanan, A. (2012). Default and the maturity structure in sovereign bonds. \textit{Journal of Political Economy, 120}(2), 187--232. \href{https://doi.org/10.1086/666589}{\nolinkurl{https://doi.org/10.1086/666589}}

\noindent\hangindent=1.5em\hangafter=1 Arnone, M. e Iden, G. (2003). \textit{Primary dealers in government securities: Policy issues and selected countries' experience} (IMF Working Paper n.º 03/45). Fondo Monetario Internacional. \href{https://doi.org/10.5089/9781451846492.001}{\nolinkurl{https://doi.org/10.5089/9781451846492.001}}

\noindent\hangindent=1.5em\hangafter=1 Ausubel, L. M. y Cramton, P. (2002). \textit{Demand reduction and inefficiency in multi-unit auctions} (Working Paper n.º 96-07, revisado). University of Maryland. \href{https://drum.lib.umd.edu/items/4caeafaa-7735-488d-97bf-8c4a6bf35909}{\nolinkurl{https://drum.lib.umd.edu/items/4caeafaa-7735-488d-97bf-8c4a6bf35909}}

\noindent\hangindent=1.5em\hangafter=1 Baldwin, R. y Giavazzi, F. (Eds.). (2015). \textit{The Eurozone crisis: A consensus view of the causes and a few possible solutions}. CEPR Press. \href{https://cepr.org/voxeu/columns/towards-consensus-causes-ez-crisis}{\nolinkurl{https://cepr.org/voxeu/columns/towards-consensus-causes-ez-crisis}}

\noindent\hangindent=1.5em\hangafter=1 Banco de España (2015). La evolución de la deuda pública en España en 2014. \textit{Boletín Económico}, julio. \href{https://www.bde.es/f/webbde/SES/Secciones/Publicaciones/InformesBoletinesRevistas/BoletinEconomico/15/Jul/Fich/be1507-art4.pdf}{\nolinkurl{https://www.bde.es/f/webbde/SES/Secciones/Publicaciones/InformesBoletinesRevistas/BoletinEconomico/15/Jul/Fich/be1507-art4.pdf}}

\noindent\hangindent=1.5em\hangafter=1 Banco de España (2020). La evolución de la deuda pública en España en 2019. \textit{Boletín Económico}, 3/2020, Notas Económicas. \href{https://www.bde.es/f/webbde/SES/Secciones/Publicaciones/InformesBoletinesRevistas/NotasEconomicas/20/T3/descargar/Fich/be2003-ne05.pdf}{\nolinkurl{https://www.bde.es/f/webbde/SES/Secciones/Publicaciones/InformesBoletinesRevistas/NotasEconomicas/20/T3/descargar/Fich/be2003-ne05.pdf}}

\noindent\hangindent=1.5em\hangafter=1 Banco de España (2025). La evolución de la deuda pública en España en 2024. \textit{Boletín Económico}. \href{https://www.bde.es/f/webbe/SES/Secciones/Publicaciones/PublicacionesSeriadas/DocumentosOcasionales/25/Fich/be2503-art04.pdf}{\nolinkurl{https://www.bde.es/f/webbe/SES/Secciones/Publicaciones/PublicacionesSeriadas/DocumentosOcasionales/25/Fich/be2503-art04.pdf}}

\noindent\hangindent=1.5em\hangafter=1 Banco de España (2026). \textit{Cuentas anuales del Banco de España 2025}. \href{https://www.bde.es/f/webbe/SES/Secciones/Publicaciones/PublicacionesAnuales/Cuentas_anuales_del_BdE/2025/fich/BE_CuentasAnuales_2025.pdf}{\nolinkurl{https://www.bde.es/f/webbe/SES/Secciones/Publicaciones/PublicacionesAnuales/Cuentas_anuales_del_BdE/2025/fich/BE_CuentasAnuales_2025.pdf}}

\noindent\hangindent=1.5em\hangafter=1 Barro, R. J. (1979). On the determination of the public debt. \textit{Journal of Political Economy, 87}(5, Parte 1), 940--971. \href{https://doi.org/10.1086/260807}{\nolinkurl{https://doi.org/10.1086/260807}}

\noindent\hangindent=1.5em\hangafter=1 Barro, R. J. (2003). Optimal management of indexed and nominal debt. \textit{Annals of Economics and Finance, 4}(1), 1--15. \href{https://ideas.repec.org/a/cuf/journl/y2003v4i1p1-15.html}{\nolinkurl{https://ideas.repec.org/a/cuf/journl/y2003v4i1p1-15.html}}

\noindent\hangindent=1.5em\hangafter=1 Bentolila, S. y Dolado, J. J. (1994). Labour flexibility and wages: Lessons from Spain. \textit{Economic Policy, 9}(18), 53--99. \href{https://doi.org/10.2307/1344458}{\nolinkurl{https://doi.org/10.2307/1344458}}

\noindent\hangindent=1.5em\hangafter=1 Blanchard, O. (2019). Public debt and low interest rates. \textit{American Economic Review, 109}(4), 1197--1229. \href{https://doi.org/10.1257/aer.109.4.1197}{\nolinkurl{https://doi.org/10.1257/aer.109.4.1197}}

\noindent\hangindent=1.5em\hangafter=1 Blanchard, O. J. y Leigh, D. (2013). Growth forecast errors and fiscal multipliers. \textit{American Economic Review, 103}(3), 117--120. \href{https://doi.org/10.1257/aer.103.3.117}{\nolinkurl{https://doi.org/10.1257/aer.103.3.117}}

\noindent\hangindent=1.5em\hangafter=1 Blommestein, H. J. y Turner, P. (2012). \textit{Interactions between sovereign debt management and monetary policy under fiscal dominance and financial instability} (OECD Working Papers on Sovereign Borrowing and Public Debt Management n.º 3). OECD Publishing. \href{https://doi.org/10.1787/5k9fdwrnd1g3-en}{\nolinkurl{https://doi.org/10.1787/5k9fdwrnd1g3-en}}

\noindent\hangindent=1.5em\hangafter=1 Borensztein, E. y Mauro, P. (2004). The case for GDP-indexed bonds. \textit{Economic Policy, 19}(38), 166--216. \href{https://doi.org/10.1111/j.1468-0327.2004.00121.x}{\nolinkurl{https://doi.org/10.1111/j.1468-0327.2004.00121.x}}

\noindent\hangindent=1.5em\hangafter=1 Buchanan, J. M. y Wagner, R. E. (1977). \textit{Democracy in deficit: The political legacy of Lord Keynes}. Academic Press.

\noindent\hangindent=1.5em\hangafter=1 Buera, F. y Nicolini, J. P. (2004). Optimal maturity of government debt without state contingent bonds. \textit{Journal of Monetary Economics, 51}(3), 531--554. \href{https://doi.org/10.1016/j.jmoneco.2003.06.002}{\nolinkurl{https://doi.org/10.1016/j.jmoneco.2003.06.002}}

\noindent\hangindent=1.5em\hangafter=1 Buiter, W., Corsetti, G. y Roubini, N. (1993). Excessive deficits: Sense and nonsense in the Treaty of Maastricht. \textit{Economic Policy, 8}(16), 57--100. \href{https://doi.org/10.2307/1344568}{\nolinkurl{https://doi.org/10.2307/1344568}}

\noindent\hangindent=1.5em\hangafter=1 CaixaBank Research (2025). \textit{La Estrategia del Tesoro 2025 en un contexto de reducción del déficit público}. \href{https://www.caixabankresearch.com/en/economics-markets/public-sector/2025-treasury-strategy-context-reduction-spains-public-deficit}{\nolinkurl{https://www.caixabankresearch.com/en/economics-markets/public-sector/2025-treasury-strategy-context-reduction-spains-public-deficit}}

\noindent\hangindent=1.5em\hangafter=1 Calvo, G. A. y Guidotti, P. E. (1990). Indexation and maturity of government bonds: An exploratory model. En R. Dornbusch y M. Draghi (Eds.), \textit{Public debt management: Theory and history} (pp. 52--82). Cambridge University Press. \href{https://doi.org/10.1017/CBO9780511628528.006}{\nolinkurl{https://doi.org/10.1017/CBO9780511628528.006}}

\noindent\hangindent=1.5em\hangafter=1 Campbell, J. Y. y Shiller, R. J. (1996). A scorecard for indexed government debt. \textit{NBER Macroeconomics Annual, 11}, 155--208. \href{https://www.nber.org/chapters/c11029}{\nolinkurl{https://www.nber.org/chapters/c11029}}

\noindent\hangindent=1.5em\hangafter=1 Cole, H. L. y Kehoe, T. J. (2000). Self-fulfilling debt crises. \textit{The Review of Economic Studies, 67}(1), 91--116. \href{https://doi.org/10.1111/1467-937X.00123}{\nolinkurl{https://doi.org/10.1111/1467-937X.00123}}

\noindent\hangindent=1.5em\hangafter=1 Comín, F. (1996). \textit{Historia de la Hacienda pública, II: España (1808-1995)}. Crítica.

\noindent\hangindent=1.5em\hangafter=1 Comín, F. y Díaz Fuentes, D. (2005). Sector público administrativo y Estado del bienestar. En A. Carreras y X. Tafunell (Coords.), \textit{Estadísticas históricas de España: siglos XIX-XX} (2.ª ed., vol. 2). Fundación BBVA. \href{https://www.fbbva.es/wp-content/uploads/2017/05/dat/DE_2006_estadisticas_historicas.pdf}{\nolinkurl{https://www.fbbva.es/wp-content/uploads/2017/05/dat/DE_2006_estadisticas_historicas.pdf}}

\noindent\hangindent=1.5em\hangafter=1 Comín, F. y Díaz Fuentes, D. (2023). Las transformaciones fiscales de la Hacienda y las Administraciones públicas españolas, 1898-2023. \textit{ICE, Revista de Economía}, (933). \href{https://doi.org/10.32796/ice.2023.933.7685}{\nolinkurl{https://doi.org/10.32796/ice.2023.933.7685}}

\noindent\hangindent=1.5em\hangafter=1 Comisión Europea (2024a). \textit{Debt Sustainability Monitor 2023} (Institutional Paper). Oficina de Publicaciones de la Unión Europea. \href{https://economy-finance.ec.europa.eu/publications/debt-sustainability-monitor-2023_en}{\nolinkurl{https://economy-finance.ec.europa.eu/publications/debt-sustainability-monitor-2023_en}}

\noindent\hangindent=1.5em\hangafter=1 Comisión Europea (2024b). \textit{2024 Ageing Report: Economic and budgetary projections for the EU Member States (2022-2070)} (Institutional Paper). Oficina de Publicaciones de la Unión Europea. \href{https://economy-finance.ec.europa.eu/publications/2024-ageing-report-economic-and-budgetary-projections-eu-member-states-2022-2070_en}{\nolinkurl{https://economy-finance.ec.europa.eu/publications/2024-ageing-report-economic-and-budgetary-projections-eu-member-states-2022-2070_en}}

\noindent\hangindent=1.5em\hangafter=1 Comisión Europea (2026). \textit{Debt Sustainability Monitor 2025} (Institutional Paper). Oficina de Publicaciones de la Unión Europea. \href{https://economy-finance.ec.europa.eu/publications/debt-sustainability-monitor-2025_en}{\nolinkurl{https://economy-finance.ec.europa.eu/publications/debt-sustainability-monitor-2025_en}}

\noindent\hangindent=1.5em\hangafter=1 Comité de personas expertas (2022). \textit{Libro Blanco sobre la reforma tributaria}. Instituto de Estudios Fiscales. \href{https://www.ief.es}{\nolinkurl{https://www.ief.es}}

\noindent\hangindent=1.5em\hangafter=1 Currie, E., Dethier, J.-J. y Togo, E. (2003). \textit{Institutional arrangements for public debt management} (Policy Research Working Paper n.º 3021). Banco Mundial. \href{http://hdl.handle.net/10986/18226}{\nolinkurl{http://hdl.handle.net/10986/18226}}

\noindent\hangindent=1.5em\hangafter=1 De Castro, F., Estrada, Á., Hernández de Cos, P. y Martí, F. (2008). Una aproximación al componente transitorio del saldo público en España. \textit{Boletín Económico} (Banco de España), junio, 71--81. \href{https://www.bde.es/f/webbde/SES/Secciones/Publicaciones/InformesBoletinesRevistas/BoletinEconomico/08/Jun/Fich/art4.pdf}{\nolinkurl{https://www.bde.es/f/webbde/SES/Secciones/Publicaciones/InformesBoletinesRevistas/BoletinEconomico/08/Jun/Fich/art4.pdf}}

\noindent\hangindent=1.5em\hangafter=1 De Grauwe, P. (2011). \textit{The governance of a fragile Eurozone} (CEPS Working Document n.º 346). Centre for European Policy Studies. \href{https://www.ceps.eu/wp-content/uploads/2011/05/WD\%20346\%20De\%20Grauwe\%20on\%20Eurozone\%20Governance.pdf}{\nolinkurl{https://www.ceps.eu/wp-content/uploads/2011/05/WD\%20346\%20De\%20Grauwe\%20on\%20Eurozone\%20Governance.pdf}}

\noindent\hangindent=1.5em\hangafter=1 De Grauwe, P. y Ji, Y. (2013). Self-fulfilling crises in the Eurozone: An empirical test. \textit{Journal of International Money and Finance, 34}, 15--36. \href{https://doi.org/10.1016/j.jimonfin.2012.11.003}{\nolinkurl{https://doi.org/10.1016/j.jimonfin.2012.11.003}}

\noindent\hangindent=1.5em\hangafter=1 Dolls, M., Fuest, C. y Peichl, A. (2012). Automatic stabilizers and economic crisis: US vs. Europe. \textit{Journal of Public Economics, 96}(3--4), 279--294. \href{https://doi.org/10.1016/j.jpubeco.2011.11.001}{\nolinkurl{https://doi.org/10.1016/j.jpubeco.2011.11.001}}

\noindent\hangindent=1.5em\hangafter=1 Doronzo, R., Siracusa, V. y Antonelli, S. (2021). \textit{Green bonds: The sovereign issuers' perspective} (Markets, Infrastructures, Payment Systems n.º 3). Banca d'Italia. \href{https://www.bancaditalia.it/pubblicazioni/mercati-infrastrutture-e-sistemi-di-pagamento/approfondimenti/2021-003/N.3-MISP.pdf}{\nolinkurl{https://www.bancaditalia.it/pubblicazioni/mercati-infrastrutture-e-sistemi-di-pagamento/approfondimenti/2021-003/N.3-MISP.pdf}}

\noindent\hangindent=1.5em\hangafter=1 Duffie, D. (2020). \textit{Still the world's safe haven? Redesigning the U.S. Treasury market after the COVID-19 crisis} (Hutchins Center Working Paper n.º 62). Brookings Institution. \href{https://www.brookings.edu/articles/still-the-worlds-safe-haven/}{\nolinkurl{https://www.brookings.edu/articles/still-the-worlds-safe-haven/}}

\noindent\hangindent=1.5em\hangafter=1 Fatás, A. y Summers, L. H. (2018). The permanent effects of fiscal consolidations. \textit{Journal of International Economics, 112}, 238--250. \href{https://doi.org/10.1016/j.jinteco.2017.11.007}{\nolinkurl{https://doi.org/10.1016/j.jinteco.2017.11.007}}

\noindent\hangindent=1.5em\hangafter=1 Fedea (2025). \textit{Estimación del gasto futuro en intereses de la deuda pública} (Apuntes 2025/28). Fundación de Estudios de Economía Aplicada. \href{https://documentos.fedea.net/pubs/ap/2025/ap2025-28.pdf}{\nolinkurl{https://documentos.fedea.net/pubs/ap/2025/ap2025-28.pdf}}

\noindent\hangindent=1.5em\hangafter=1 Fernández-Villaverde, J., Garicano, L. y Santos, T. (2013). Political credit cycles: The case of the Eurozone. \textit{Journal of Economic Perspectives, 27}(3), 145--166. \href{https://doi.org/10.1257/jep.27.3.145}{\nolinkurl{https://doi.org/10.1257/jep.27.3.145}}

\noindent\hangindent=1.5em\hangafter=1 Fondo Monetario Internacional (2014). \textit{Government Finance Statistics Manual 2014}. \href{https://www.imf.org/external/Pubs/FT/GFS/Manual/2014/gfsfinal.pdf}{\nolinkurl{https://www.imf.org/external/Pubs/FT/GFS/Manual/2014/gfsfinal.pdf}}

\noindent\hangindent=1.5em\hangafter=1 Fondo Monetario Internacional (2021). \textit{Database of country fiscal measures in response to the COVID-19 pandemic} (Fiscal Monitor, octubre). \href{https://www.imf.org/-/media/files/topics/covid/fm-database/october-2021/oct-2021-country-fiscal-measures-database-publication.pdf}{\nolinkurl{https://www.imf.org/-/media/files/topics/covid/fm-database/october-2021/oct-2021-country-fiscal-measures-database-publication.pdf}}

\noindent\hangindent=1.5em\hangafter=1 Fondo Monetario Internacional (2022). \textit{Staff guidance note on the sovereign risk and debt sustainability framework for market access countries}. \href{https://www.imf.org/en/Publications/Policy-Papers/Issues/2022/08/08/Staff-Guidance-Note-on-the-Sovereign-Risk-and-Debt-Sustainability-Framework-for-Market-521884}{\nolinkurl{https://www.imf.org/en/Publications/Policy-Papers/Issues/2022/08/08/Staff-Guidance-Note-on-the-Sovereign-Risk-and-Debt-Sustainability-Framework-for-Market-521884}}

\noindent\hangindent=1.5em\hangafter=1 Fondo Monetario Internacional y Banco Mundial (2009). \textit{Developing a medium-term debt management strategy (MTDS): Guidance note for country authorities}. \href{https://www.imf.org}{\nolinkurl{https://www.imf.org}}

\noindent\hangindent=1.5em\hangafter=1 Fondo Monetario Internacional y Banco Mundial (2014). \textit{Revised guidelines for public debt management}. \href{https://www.imf.org/external/np/pp/eng/2014/040114.pdf}{\nolinkurl{https://www.imf.org/external/np/pp/eng/2014/040114.pdf}}

\noindent\hangindent=1.5em\hangafter=1 Friedman, M. (1960). \textit{A program for monetary stability}. Fordham University Press.

\noindent\hangindent=1.5em\hangafter=1 Galesi, A., Nuño, G. y Thomas, C. (2017). El tipo de interés natural: concepto, determinantes e implicaciones para la política monetaria. \textit{Boletín Económico} (Banco de España), 1/2017, Artículos analíticos. \href{http://www.bde.es/f/webbde/SES/Secciones/Publicaciones/InformesBoletinesRevistas/ArticulosAnaliticos/2017/T1/fich/beaa1701-art7.pdf}{\nolinkurl{http://www.bde.es/f/webbde/SES/Secciones/Publicaciones/InformesBoletinesRevistas/ArticulosAnaliticos/2017/T1/fich/beaa1701-art7.pdf}}

\noindent\hangindent=1.5em\hangafter=1 Galí, J. y Perotti, R. (2003). Fiscal policy and monetary integration in Europe. \textit{Economic Policy, 18}(37), 533--572. \href{https://doi.org/10.1111/1468-0327.00115}{\nolinkurl{https://doi.org/10.1111/1468-0327.00115}}

\noindent\hangindent=1.5em\hangafter=1 Giavazzi, F. y Pagano, M. (1988). The advantage of tying one's hands: EMS discipline and central bank credibility. \textit{European Economic Review, 32}(5), 1055--1075. \href{https://doi.org/10.1016/0014-2921(88)90065-7}{\nolinkurl{https://doi.org/10.1016/0014-2921(88)90065-7}}

\noindent\hangindent=1.5em\hangafter=1 Giavazzi, F. y Pagano, M. (1990). Can severe fiscal contractions be expansionary? Tales of two small European countries. \textit{NBER Macroeconomics Annual, 5}, 75--111. \href{https://doi.org/10.1086/654131}{\nolinkurl{https://doi.org/10.1086/654131}}

\noindent\hangindent=1.5em\hangafter=1 Greenwood, R., Hanson, S. G. y Stein, J. C. (2015). A comparative-advantage approach to government debt maturity. \textit{The Journal of Finance, 70}(4), 1669--1719. \href{https://doi.org/10.1111/jofi.12253}{\nolinkurl{https://doi.org/10.1111/jofi.12253}}

\noindent\hangindent=1.5em\hangafter=1 Greenwood, R., Hanson, S. G., Rudolph, J. S. y Summers, L. H. (2014). \textit{Government debt management at the zero lower bound} (Hutchins Center Working Paper n.º 5). Brookings Institution. \href{https://www.brookings.edu/articles/government-debt-management-at-the-zero-lower-bound/}{\nolinkurl{https://www.brookings.edu/articles/government-debt-management-at-the-zero-lower-bound/}}

\noindent\hangindent=1.5em\hangafter=1 Grilli, V., Masciandaro, D. y Tabellini, G. (1991). Political and monetary institutions and public financial policies in the industrial countries. \textit{Economic Policy, 6}(13), 341--392. \href{https://doi.org/10.2307/1344630}{\nolinkurl{https://doi.org/10.2307/1344630}}

\noindent\hangindent=1.5em\hangafter=1 Guajardo, J., Leigh, D. y Pescatori, A. (2014). Expansionary austerity? International evidence. \textit{Journal of the European Economic Association, 12}(4), 949--968. \href{https://doi.org/10.1111/jeea.12083}{\nolinkurl{https://doi.org/10.1111/jeea.12083}}

\noindent\hangindent=1.5em\hangafter=1 Hansen, A. H. (1938). \textit{Full recovery or stagnation?} W. W. Norton.

\noindent\hangindent=1.5em\hangafter=1 Herndon, T., Ash, M. y Pollin, R. (2014). Does high public debt consistently stifle economic growth? A critique of Reinhart and Rogoff. \textit{Cambridge Journal of Economics, 38}(2), 257--279. \href{https://doi.org/10.1093/cje/bet075}{\nolinkurl{https://doi.org/10.1093/cje/bet075}}

\noindent\hangindent=1.5em\hangafter=1 Hortaçsu, A. y McAdams, D. (2010). Mechanism choice and strategic bidding in divisible good auctions: An empirical analysis of the Turkish Treasury auction market. \textit{Journal of Political Economy, 118}(5), 833--865. \href{https://doi.org/10.1086/657948}{\nolinkurl{https://doi.org/10.1086/657948}}

\noindent\hangindent=1.5em\hangafter=1 Hortaçsu, A., Kastl, J. y Zhang, A. (2018). Bid shading and bidder surplus in the US Treasury auction system. \textit{American Economic Review, 108}(1), 147--169. \href{https://doi.org/10.1257/aer.20160675}{\nolinkurl{https://doi.org/10.1257/aer.20160675}}

\noindent\hangindent=1.5em\hangafter=1 House, C. L., Proebsting, C. y Tesar, L. L. (2020). Austerity in the aftermath of the Great Recession. \textit{Journal of Monetary Economics, 115}, 37--63. \href{https://doi.org/10.1016/j.jmoneco.2019.05.004}{\nolinkurl{https://doi.org/10.1016/j.jmoneco.2019.05.004}}

\noindent\hangindent=1.5em\hangafter=1 Instituto Monetario Europeo (1998). \textit{Informe de Convergencia} (marzo). \href{https://www.ecb.europa.eu/pub/pdf/conrep/cr1998es.pdf}{\nolinkurl{https://www.ecb.europa.eu/pub/pdf/conrep/cr1998es.pdf}}

\noindent\hangindent=1.5em\hangafter=1 Koen, V. y van den Noord, P. (2005). \textit{Fiscal gimmickry in Europe: One-off measures and creative accounting} (OECD Economics Department Working Papers n.º 417). OECD Publishing. \href{https://doi.org/10.1787/237714513517}{\nolinkurl{https://doi.org/10.1787/237714513517}}

\noindent\hangindent=1.5em\hangafter=1 Kydland, F. E. y Prescott, E. C. (1977). Rules rather than discretion: The inconsistency of optimal plans. \textit{Journal of Political Economy, 85}(3), 473--491. \href{https://doi.org/10.1086/260580}{\nolinkurl{https://doi.org/10.1086/260580}}

\noindent\hangindent=1.5em\hangafter=1 Laeven, L. y Valencia, F. (2018). \textit{Systemic banking crises revisited} (IMF Working Paper n.º 18/206). Fondo Monetario Internacional. \href{https://doi.org/10.5089/9781484376379.001}{\nolinkurl{https://doi.org/10.5089/9781484376379.001}}

\noindent\hangindent=1.5em\hangafter=1 Lago Peñas, S. (2024). \textit{La dinámica de la deuda pública en España: pasado, presente y futuro} (Nota técnica). Funcas. \href{https://www.funcas.es/documentos_trabajo/la-dinamica-de-la-deuda-publica-en-espana-presente-pasado-y-futuro/}{\nolinkurl{https://www.funcas.es/documentos_trabajo/la-dinamica-de-la-deuda-publica-en-espana-presente-pasado-y-futuro/}}

\noindent\hangindent=1.5em\hangafter=1 Lane, P. R. (2012). The European sovereign debt crisis. \textit{Journal of Economic Perspectives, 26}(3), 49--68. \href{https://doi.org/10.1257/jep.26.3.49}{\nolinkurl{https://doi.org/10.1257/jep.26.3.49}}

\noindent\hangindent=1.5em\hangafter=1 Lucas, R. E., Jr. y Stokey, N. L. (1983). Optimal fiscal and monetary policy in an economy without capital. \textit{Journal of Monetary Economics, 12}(1), 55--93. \href{https://doi.org/10.1016/0304-3932(83)90049-1}{\nolinkurl{https://doi.org/10.1016/0304-3932(83)90049-1}}

\noindent\hangindent=1.5em\hangafter=1 Malvey, P. F. y Archibald, C. M. (1998). \textit{Uniform-price auctions: Update of the Treasury experience}. U.S. Department of the Treasury. \href{https://home.treasury.gov/system/files/276/update-of-the-treasury-experience-october-1998.pdf}{\nolinkurl{https://home.treasury.gov/system/files/276/update-of-the-treasury-experience-october-1998.pdf}}

\noindent\hangindent=1.5em\hangafter=1 Martí, F. y Pérez, J. J. (2015). Spanish public finances through the financial crisis. \textit{Fiscal Studies, 36}(4), 527--554. \href{https://doi.org/10.1111/j.1475-5890.2015.12079}{\nolinkurl{https://doi.org/10.1111/j.1475-5890.2015.12079}}

\noindent\hangindent=1.5em\hangafter=1 Milesi-Ferretti, G. M. (2004). Good, bad or ugly? On the effects of fiscal rules with creative accounting. \textit{Journal of Public Economics, 88}(1--2), 377--394. \href{https://doi.org/10.1016/S0047-2727(02)00076-2}{\nolinkurl{https://doi.org/10.1016/S0047-2727(02)00076-2}}

\noindent\hangindent=1.5em\hangafter=1 Missale, A. y Blanchard, O. J. (1994). The debt burden and debt maturity. \textit{The American Economic Review, 84}(1), 309--319. \href{https://ideas.repec.org/a/aea/aecrev/v84y1994i1p309-19.html}{\nolinkurl{https://ideas.repec.org/a/aea/aecrev/v84y1994i1p309-19.html}}

\noindent\hangindent=1.5em\hangafter=1 Myerson, R. B. (1981). Optimal auction design. \textit{Mathematics of Operations Research, 6}(1), 58--73. \href{https://doi.org/10.1287/moor.6.1.58}{\nolinkurl{https://doi.org/10.1287/moor.6.1.58}}

\noindent\hangindent=1.5em\hangafter=1 Perotti, R. (2011). \textit{The ``austerity myth'': Gain without pain?} (NBER Working Paper n.º 17571). National Bureau of Economic Research. \href{https://doi.org/10.3386/w17571}{\nolinkurl{https://doi.org/10.3386/w17571}}

\noindent\hangindent=1.5em\hangafter=1 Persson, T. y Svensson, L. E. O. (1989). Why a stubborn conservative would run a deficit: Policy with time-inconsistent preferences. \textit{The Quarterly Journal of Economics, 104}(2), 325--345. \href{https://doi.org/10.2307/2937850}{\nolinkurl{https://doi.org/10.2307/2937850}}

\noindent\hangindent=1.5em\hangafter=1 Prats Albentosa, M. A. y Rocamora Martí, A. M. (2016). Análisis de la sostenibilidad de la deuda pública en España. \textit{Revista de Ciencias Sociales, 22}(2), 10--23. \href{https://www.redalyc.org/pdf/280/28049145002.pdf}{\nolinkurl{https://www.redalyc.org/pdf/280/28049145002.pdf}}

\noindent\hangindent=1.5em\hangafter=1 Reinhart, C. M. y Rogoff, K. S. (2010). Growth in a time of debt. \textit{American Economic Review, 100}(2), 573--578. \href{https://doi.org/10.1257/aer.100.2.573}{\nolinkurl{https://doi.org/10.1257/aer.100.2.573}}

\noindent\hangindent=1.5em\hangafter=1 Reinhart, C. M. y Sbrancia, M. B. (2015). The liquidation of government debt. \textit{Economic Policy, 30}(82), 291--333. \href{https://doi.org/10.1093/epolic/eiv003}{\nolinkurl{https://doi.org/10.1093/epolic/eiv003}}

\noindent\hangindent=1.5em\hangafter=1 Roubini, N. y Sachs, J. D. (1989). Political and economic determinants of budget deficits in the industrial democracies. \textit{European Economic Review, 33}(5), 903--933. \href{https://doi.org/10.1016/0014-2921(89)90002-0}{\nolinkurl{https://doi.org/10.1016/0014-2921(89)90002-0}}

\noindent\hangindent=1.5em\hangafter=1 Sargent, T. J. y Wallace, N. (1981). Some unpleasant monetarist arithmetic. \textit{Federal Reserve Bank of Minneapolis Quarterly Review, 5}(3), 1--17. \href{https://doi.org/10.21034/qr.531}{\nolinkurl{https://doi.org/10.21034/qr.531}}

\noindent\hangindent=1.5em\hangafter=1 Sudrià, C. (2014). Las crisis bancarias en España: una perspectiva histórica. \textit{Estudios de Economía Aplicada, 32}(2), 473--496. \href{https://www.redalyc.org/pdf/301/30130732001.pdf}{\nolinkurl{https://www.redalyc.org/pdf/301/30130732001.pdf}}

\noindent\hangindent=1.5em\hangafter=1 Tanzi, V. y Schuknecht, L. (2000). \textit{Public spending in the 20th century: A global perspective}. Cambridge University Press. \href{https://doi.org/10.1017/CBO9780511625800}{\nolinkurl{https://doi.org/10.1017/CBO9780511625800}}

\noindent\hangindent=1.5em\hangafter=1 Tribunal de Cuentas (2025). \textit{Actualización del Informe de fiscalización sobre el proceso de reestructuración bancaria} (coste a 1 de enero de 2022). \href{https://www.tcu.es/es/sala-de-prensa/noticias/El-TCu-actualiza-hasta-los-71.833-millones-de-euros-el-coste-de-la-reestructuracion-bancaria/}{\nolinkurl{https://www.tcu.es/es/sala-de-prensa/noticias/El-TCu-actualiza-hasta-los-71.833-millones-de-euros-el-coste-de-la-reestructuracion-bancaria/}}

\noindent\hangindent=1.5em\hangafter=1 Umlauf, S. R. (1993). An empirical study of the Mexican Treasury bill auction. \textit{Journal of Financial Economics, 33}(3), 313--340. \href{https://doi.org/10.1016/0304-405X(93)90010-9}{\nolinkurl{https://doi.org/10.1016/0304-405X(93)90010-9}}

\noindent\hangindent=1.5em\hangafter=1 Vickrey, W. (1961). Counterspeculation, auctions, and competitive sealed tenders. \textit{The Journal of Finance, 16}(1), 8--37. \href{https://doi.org/10.1111/j.1540-6261.1961.tb02789.x}{\nolinkurl{https://doi.org/10.1111/j.1540-6261.1961.tb02789.x}}

\noindent\hangindent=1.5em\hangafter=1 von Hagen, J. (1992). \textit{Budgeting procedures and fiscal performance in the European Communities} (Economic Papers n.º 96). Comisión de las Comunidades Europeas, Dirección General de Asuntos Económicos y Financieros.

\noindent\hangindent=1.5em\hangafter=1 von Hagen, J. y Wolff, G. B. (2006). What do deficits tell us about debt? Empirical evidence on creative accounting with fiscal rules in the EU. \textit{Journal of Banking \& Finance, 30}(12), 3259--3279. \href{https://doi.org/10.1016/j.jbankfin.2006.05.011}{\nolinkurl{https://doi.org/10.1016/j.jbankfin.2006.05.011}}

\bigskip
{\bfseries\large Complementario\par}
\medskip

\noindent\hangindent=1.5em\hangafter=1 Banco Central Europeo (2026). \textit{Decisiones de política monetaria} (nota de prensa, 10 de septiembre). \href{https://www.ecb.europa.eu/press/pr/date/2026/html/ecb.mp260910~314e508016.es.html}{\nolinkurl{https://www.ecb.europa.eu/press/pr/date/2026/html/ecb.mp260910~314e508016.es.html}}

\noindent\hangindent=1.5em\hangafter=1 Banco de España (2024). \textit{Nota informativa del BCE sobre la remuneración de los depósitos de las administraciones públicas}. \href{https://www.bde.es/f/webbe/GAP/Secciones/SalaPrensa/ComunicadosBCE/NotasInformativasBCE/24/presbce2024-42.pdf}{\nolinkurl{https://www.bde.es/f/webbe/GAP/Secciones/SalaPrensa/ComunicadosBCE/NotasInformativasBCE/24/presbce2024-42.pdf}}

\noindent\hangindent=1.5em\hangafter=1 Banco de España (2026a). \textit{El BCE sube los tipos de interés oficiales en 25 puntos básicos} (junio). \href{https://www.bde.es/wbe/en/noticias-eventos/actualidad-bce/decisiones-politica-monetaria/bce-tipos-junio26.html}{\nolinkurl{https://www.bde.es/wbe/en/noticias-eventos/actualidad-bce/decisiones-politica-monetaria/bce-tipos-junio26.html}}

\noindent\hangindent=1.5em\hangafter=1 Banco de España (2026b). \textit{La deuda de las Administraciones Públicas se situó en el 100,7\% del PIB al cierre de 2025} (nota de prensa, 31 de marzo). \href{https://www.bde.es}{\nolinkurl{https://www.bde.es}}

\noindent\hangindent=1.5em\hangafter=1 Bankinter (2014). \textit{Primera emisión de un bono indexado a la inflación}. \href{https://www.bankinter.com/blog/lo-ultimo/primera-emision-de-un-bono-indexado-a-la-inflacion-deuda-bajo-presion}{\nolinkurl{https://www.bankinter.com/blog/lo-ultimo/primera-emision-de-un-bono-indexado-a-la-inflacion-deuda-bajo-presion}}

\noindent\hangindent=1.5em\hangafter=1 Bolsamanía (2026). \textit{El Banco de España gana 234 millones en 2025 tras liberar 541 millones en provisiones}. \href{https://www.bolsamania.com/noticias/mercados/economiafinanzas--banco-de-espana-gana-234-millones-en-2025-tras-liberar-541-millones-en-provisiones--22209834.html}{\nolinkurl{https://www.bolsamania.com/noticias/mercados/economiafinanzas--banco-de-espana-gana-234-millones-en-2025-tras-liberar-541-millones-en-provisiones--22209834.html}}

\noindent\hangindent=1.5em\hangafter=1 Conde-Ruiz, J. I. (s. f.). Lecciones del pasado: el fondo de reserva que pudo ser y no fue. \textit{Nada es Gratis}. \href{https://nadaesgratis.es/j-ignacio-conde-ruiz/lecciones-del-pasado-el-fondo-de-reserva-que-pudo-ser-y-no-fue}{\nolinkurl{https://nadaesgratis.es/j-ignacio-conde-ruiz/lecciones-del-pasado-el-fondo-de-reserva-que-pudo-ser-y-no-fue}}

\noindent\hangindent=1.5em\hangafter=1 countryeconomy.com (s. f.). \textit{España: déficit público} [conjunto de datos, a partir de Eurostat]. \href{https://countryeconomy.com/deficit/spain}{\nolinkurl{https://countryeconomy.com/deficit/spain}}

\noindent\hangindent=1.5em\hangafter=1 De Grauwe, P. y Ji, Y. (2023, 13 de marzo). Monetary policies with fewer subsidies for banks: A two-tier system of minimum reserve requirements. \textit{VoxEU}. \href{https://cepr.org/voxeu/columns/monetary-policies-fewer-subsidies-banks-two-tier-system-minimum-reserve-requirements}{\nolinkurl{https://cepr.org/voxeu/columns/monetary-policies-fewer-subsidies-banks-two-tier-system-minimum-reserve-requirements}}

\noindent\hangindent=1.5em\hangafter=1 EFPA España (2024). \textit{2023, el año en que los inversores particulares se fijaron en las Letras del Tesoro}. \href{https://www.asesoresfinancierosefpa.es/actualidad/2023-el-ano-en-que-los-inversores-particulares-se-fijaron-en-las-letras-del-tesoro/}{\nolinkurl{https://www.asesoresfinancierosefpa.es/actualidad/2023-el-ano-en-que-los-inversores-particulares-se-fijaron-en-las-letras-del-tesoro/}}

\noindent\hangindent=1.5em\hangafter=1 El Blog Salmón (s. f.). \textit{Este es el coste del rescate al sector financiero en la UE}. \href{https://www.elblogsalmon.com/economia/este-es-el-coste-del-rescate-al-sector-financiero-en-la-ue}{\nolinkurl{https://www.elblogsalmon.com/economia/este-es-el-coste-del-rescate-al-sector-financiero-en-la-ue}}

\noindent\hangindent=1.5em\hangafter=1 El Boletín (2026). \textit{Italia lidera la caída de las primas de riesgo en la eurozona en septiembre y España recorta por debajo de 60 pb}. \href{https://www.elboletin.com/italia-lidera-la-caida-de-las-primas-de-riesgo-en-la-eurozona-en-septiembre-y-espana-recorta-por-debajo-de-60-pb/}{\nolinkurl{https://www.elboletin.com/italia-lidera-la-caida-de-las-primas-de-riesgo-en-la-eurozona-en-septiembre-y-espana-recorta-por-debajo-de-60-pb/}}

\noindent\hangindent=1.5em\hangafter=1 El Diario de Madrid (2026, 17 de junio). \textit{España gastó 4,5 euros en intereses de la deuda por cada euro invertido en 2025}. \href{https://www.eldiariodemadrid.es/articulo/economia/espana-gasto-45-euros-intereses-deuda-cada-euro-invertido-2025/20260617080650134214.html}{\nolinkurl{https://www.eldiariodemadrid.es/articulo/economia/espana-gasto-45-euros-intereses-deuda-cada-euro-invertido-2025/20260617080650134214.html}}

\noindent\hangindent=1.5em\hangafter=1 El Economista (2024, enero). \textit{Cuerpo anuncia una emisión neta del Tesoro de 55.000 millones en 2024}. \href{https://www.eleconomista.es/economia/noticias/12615361/01/24/cuerpo-anuncia-una-emision-neta-del-tesoro-55000-millones-en-2024-10000-millones-menos.html}{\nolinkurl{https://www.eleconomista.es/economia/noticias/12615361/01/24/cuerpo-anuncia-una-emision-neta-del-tesoro-55000-millones-en-2024-10000-millones-menos.html}}

\noindent\hangindent=1.5em\hangafter=1 El Economista (2026a, marzo). \textit{La recaudación fiscal en 2025 hace historia al superar los 325.000 millones}. \href{https://www.eleconomista.es/economia/noticias/13850918/03/26/la-recaudacion-fiscal-en-2025-hace-historia-al-superar-los-325000-millones-y-subir-un-104.html}{\nolinkurl{https://www.eleconomista.es/economia/noticias/13850918/03/26/la-recaudacion-fiscal-en-2025-hace-historia-al-superar-los-325000-millones-y-subir-un-104.html}}

\noindent\hangindent=1.5em\hangafter=1 El Economista (2026b, julio). \textit{Francia paga la mayor prima de riesgo de su historia frente a España por su debilidad fiscal}. \href{https://www.eleconomista.es/mercados-cotizaciones/noticias/13999958/07/26/francia-paga-la-mayor-prima-de-riesgo-de-su-historia-frente-a-espana-por-su-debilidad-fiscal-con-el-bono-a-un-paso-del-4.html}{\nolinkurl{https://www.eleconomista.es/mercados-cotizaciones/noticias/13999958/07/26/francia-paga-la-mayor-prima-de-riesgo-de-su-historia-frente-a-espana-por-su-debilidad-fiscal-con-el-bono-a-un-paso-del-4.html}}

\noindent\hangindent=1.5em\hangafter=1 El Economista (2026c, septiembre). \textit{El diferencial de la deuda francesa con la española se sitúa en máximos históricos}. \href{https://www.eleconomista.es/mercados-cotizaciones/noticias/14011436/09/26/el-diferencial-de-la-deuda-francesa-con-la-espanola-se-situa-en-maximos-historicos.html}{\nolinkurl{https://www.eleconomista.es/mercados-cotizaciones/noticias/14011436/09/26/el-diferencial-de-la-deuda-francesa-con-la-espanola-se-situa-en-maximos-historicos.html}}

\noindent\hangindent=1.5em\hangafter=1 El Español (2021, 23 de marzo). \textit{La Sareb sumará 35.000 millones a la deuda pública}. \href{https://www.elespanol.com/invertia/empresas/banca/20210323/sareb-sumara-millones-publica-gobierno-puerta-extender/568194202_0.html}{\nolinkurl{https://www.elespanol.com/invertia/empresas/banca/20210323/sareb-sumara-millones-publica-gobierno-puerta-extender/568194202_0.html}}

\noindent\hangindent=1.5em\hangafter=1 El Español (2025, 6 de noviembre). \textit{España pagará más de 40.000 millones al año en intereses de la deuda}. \href{https://www.elespanol.com/invertia/economia/20251106/espana-pagara-millones-intereses-deuda-ano-equivalente-recaudacion-sociedades/1003744001419_0.html}{\nolinkurl{https://www.elespanol.com/invertia/economia/20251106/espana-pagara-millones-intereses-deuda-ano-equivalente-recaudacion-sociedades/1003744001419_0.html}}

\noindent\hangindent=1.5em\hangafter=1 El Independiente (2021, 23 de marzo). \textit{Europa obliga al Gobierno a que incluya a la Sareb en el cómputo de deuda pública}. \href{https://www.elindependiente.com/economia/2021/03/23/europa-obliga-al-gobierno-a-que-incluya-a-la-sareb-en-el-computo-de-deuda-publica-que-sube-al-120-en-2020/}{\nolinkurl{https://www.elindependiente.com/economia/2021/03/23/europa-obliga-al-gobierno-a-que-incluya-a-la-sareb-en-el-computo-de-deuda-publica-que-sube-al-120-en-2020/}}

\noindent\hangindent=1.5em\hangafter=1 elDiario.es (2025, diciembre). \textit{El Tesoro mantendrá en 2026 una estrategia similar, con una emisión neta de 55.000 millones}. \href{https://www.eldiario.es/economia/tesoro-mantendra-2026-estrategia-similar-emision-neta-55-000-millones_1_12882825.html}{\nolinkurl{https://www.eldiario.es/economia/tesoro-mantendra-2026-estrategia-similar-emision-neta-55-000-millones_1_12882825.html}}

\noindent\hangindent=1.5em\hangafter=1 Esteve, V. y Prats, M. A. (2024, 12 de febrero). Evolución histórica de las variables macroeconómicas fiscales de la economía española, 1964-2022, y proyecciones 2023-2026. \textit{(bAg) Blog de Economía de la Aldea Global}, ALdE. \href{https://alde.es/blog/evolucion-historica-de-las-variables-macroeconomicas-fiscales-de-la-economia-espanola/}{\nolinkurl{https://alde.es/blog/evolucion-historica-de-las-variables-macroeconomicas-fiscales-de-la-economia-espanola/}}

\noindent\hangindent=1.5em\hangafter=1 Eurostat (2011, 21 de octubre). \textit{Euro area and EU27 government deficit at 6.2\% and 6.6\% of GDP respectively} (Euro indicators). \href{https://ec.europa.eu/eurostat/documents/2995521/5037534/2-21102011-AP-EN.PDF/3eb0463f-c019-44de-ac16-fc360b4ef4f9?version=1.0}{\nolinkurl{https://ec.europa.eu/eurostat/documents/2995521/5037534/2-21102011-AP-EN.PDF/3eb0463f-c019-44de-ac16-fc360b4ef4f9?version=1.0}}

\noindent\hangindent=1.5em\hangafter=1 Eurostat (2026a, 22 de abril). \textit{Euro area government deficit at 2.9\% and EU at 3.1\% of GDP} (Euro indicators). \href{https://ec.europa.eu/eurostat/web/products-euro-indicators/w/2-22042026-ap}{\nolinkurl{https://ec.europa.eu/eurostat/web/products-euro-indicators/w/2-22042026-ap}}

\noindent\hangindent=1.5em\hangafter=1 Eurostat (2026b, 22 de abril). \textit{Government debt at 87.8\% of GDP in euro area} (Euro indicators). \href{https://ec.europa.eu/eurostat/web/products-euro-indicators/w/2-22042026-bp}{\nolinkurl{https://ec.europa.eu/eurostat/web/products-euro-indicators/w/2-22042026-bp}}

\noindent\hangindent=1.5em\hangafter=1 Eurostat (2026c). \textit{Government finance statistics} (Statistics Explained). \href{https://ec.europa.eu/eurostat/statistics-explained/index.php?title=Government_finance_statistics}{\nolinkurl{https://ec.europa.eu/eurostat/statistics-explained/index.php?title=Government_finance_statistics}}

\noindent\hangindent=1.5em\hangafter=1 Fondo Monetario Internacional (s. f.). \textit{Government gross debt (\% of GDP)} [conjunto de datos, DataMapper, Historical Public Debt Database]. \href{https://www.imf.org/external/datamapper/api/v1/GG_DEBT_GDP/ESP}{\nolinkurl{https://www.imf.org/external/datamapper/api/v1/GG_DEBT_GDP/ESP}}

\noindent\hangindent=1.5em\hangafter=1 Forbes España (2025). \textit{La emisión neta del Tesoro se mantendrá en 55.000 millones en 2026, pero la bruta subirá un 4,2\%}. \href{https://forbes.es/economia/842646/la-emision-neta-del-tesoro-se-mantendra-en-55-000-millones-en-2026-pero-la-bruta-subira-un-42/}{\nolinkurl{https://forbes.es/economia/842646/la-emision-neta-del-tesoro-se-mantendra-en-55-000-millones-en-2026-pero-la-bruta-subira-un-42/}}

\noindent\hangindent=1.5em\hangafter=1 Forbes España (2026a). \textit{El rating de España buscará el sobresaliente en 2026 tras las subidas cosechadas en 2025}. \href{https://forbes.es/economia/850407/el-rating-de-espana-buscara-el-sobresaliente-en-2026-tras-las-subidas-cosechadas-en-2025/}{\nolinkurl{https://forbes.es/economia/850407/el-rating-de-espana-buscara-el-sobresaliente-en-2026-tras-las-subidas-cosechadas-en-2025/}}

\noindent\hangindent=1.5em\hangafter=1 Forbes España (2026b). \textit{El Tesoro emite 4.000 millones en una emisión sindicada a 20 años}. \href{https://forbes.es/ultima-hora/1013933/el-tesoro-emite-4-000-millones-en-una-emision-sindicada-a-20-anos-con-demanda-inicial-de-55-000-millones/}{\nolinkurl{https://forbes.es/ultima-hora/1013933/el-tesoro-emite-4-000-millones-en-una-emision-sindicada-a-20-anos-con-demanda-inicial-de-55-000-millones/}}

\noindent\hangindent=1.5em\hangafter=1 Hoy Aragón (2025, 16 de marzo). \textit{El Tribunal de Cuentas y el coste del rescate bancario}. \href{https://www.hoyaragon.es/articulo/noticias-espana/tribunal-cuentas-rescate-bancario/20250316192541090616.html}{\nolinkurl{https://www.hoyaragon.es/articulo/noticias-espana/tribunal-cuentas-rescate-bancario/20250316192541090616.html}}

\noindent\hangindent=1.5em\hangafter=1 Infobae (2025, 22 de mayo). \textit{El Tesoro vende 1.646 millones en una nueva reapertura del bono verde y reduce su interés}. \href{https://www.infobae.com/espana/agencias/2025/05/22/el-tesoro-vende-1646-millones-en-una-nueva-reapertura-del-bono-verde-y-reduce-su-interes/}{\nolinkurl{https://www.infobae.com/espana/agencias/2025/05/22/el-tesoro-vende-1646-millones-en-una-nueva-reapertura-del-bono-verde-y-reduce-su-interes/}}

\noindent\hangindent=1.5em\hangafter=1 Infobae (2026a, 31 de marzo). \textit{El déficit público cierra 2025 en el 2,18\% del PIB}. \href{https://www.infobae.com/america/agencias/2026/03/31/el-deficit-publico-cierra-2025-en-el-218-del-pib-con-36780-millones-y-mejora-el-objetivo-del-gobierno/}{\nolinkurl{https://www.infobae.com/america/agencias/2026/03/31/el-deficit-publico-cierra-2025-en-el-218-del-pib-con-36780-millones-y-mejora-el-objetivo-del-gobierno/}}

\noindent\hangindent=1.5em\hangafter=1 Infobae (2026b, 31 de marzo). \textit{La deuda pública cerró 2025 en el 100,7\% del PIB, el nivel más bajo desde la pandemia}. \href{https://www.infobae.com/espana/agencias/2026/03/31/la-deuda-publica-cerro-2025-en-el-1007-del-pib-el-nivel-mas-bajo-desde-la-pandemia/}{\nolinkurl{https://www.infobae.com/espana/agencias/2026/03/31/la-deuda-publica-cerro-2025-en-el-1007-del-pib-el-nivel-mas-bajo-desde-la-pandemia/}}

\noindent\hangindent=1.5em\hangafter=1 Infobae (2026c, 31 de agosto). \textit{España encara el final de los fondos europeos con 25.862 millones aún pendientes}. \href{https://www.infobae.com/espana/2026/08/31/espana-encara-el-final-de-los-fondos-europeos-con-25862-millones-aun-pendientes-tras-haberse-embolsado-78000-millones/}{\nolinkurl{https://www.infobae.com/espana/2026/08/31/espana-encara-el-final-de-los-fondos-europeos-con-25862-millones-aun-pendientes-tras-haberse-embolsado-78000-millones/}}

\noindent\hangindent=1.5em\hangafter=1 Infobae (2026d, 29 de septiembre). \textit{El Congreso reactiva la ley para condonar deuda a las comunidades autónomas}. \href{https://www.infobae.com/america/agencias/2026/09/29/el-congreso-reactiva-la-ley-para-condonar-deuda-a-las-comunidades-autonomas-tras-casi-un-ano-guardada-en-un-cajon/}{\nolinkurl{https://www.infobae.com/america/agencias/2026/09/29/el-congreso-reactiva-la-ley-para-condonar-deuda-a-las-comunidades-autonomas-tras-casi-un-ano-guardada-en-un-cajon/}}

\noindent\hangindent=1.5em\hangafter=1 La Crónica del Henares (2026a, septiembre). \textit{Déficit público de España en 2026 e intereses de la deuda}. \href{https://www.cronicadelhenares.com/2026/09/deficit-publico-espana-2026-intereses-deuda-comunidades.html}{\nolinkurl{https://www.cronicadelhenares.com/2026/09/deficit-publico-espana-2026-intereses-deuda-comunidades.html}}

\noindent\hangindent=1.5em\hangafter=1 La Crónica del Henares (2026b, septiembre). \textit{Scope eleva la calificación de España a A+}. \href{https://www.cronicadelhenares.com/2026/09/scope-rating-comunidad-madrid-a-plus-techo-espana.html}{\nolinkurl{https://www.cronicadelhenares.com/2026/09/scope-rating-comunidad-madrid-a-plus-techo-espana.html}}

\noindent\hangindent=1.5em\hangafter=1 La Moncloa (2023a, 12 de enero). \textit{El Tesoro prevé una emisión neta de 70.000 millones este año y cierra 2022 con un sólido acceso al mercado} (nota de prensa). \href{https://www.lamoncloa.gob.es/serviciosdeprensa/notasprensa/asuntos-economicos/Paginas/2023/120123-tesoro.aspx}{\nolinkurl{https://www.lamoncloa.gob.es/serviciosdeprensa/notasprensa/asuntos-economicos/Paginas/2023/120123-tesoro.aspx}}

\noindent\hangindent=1.5em\hangafter=1 La Moncloa (2023b, 13 de abril). \textit{Primera reapertura del bono verde soberano español en 2023} (nota de prensa). \href{https://www.lamoncloa.gob.es/serviciosdeprensa/notasprensa/asuntos-economicos/Paginas/2023/130423-bono-verde-soberano-espanol.aspx}{\nolinkurl{https://www.lamoncloa.gob.es/serviciosdeprensa/notasprensa/asuntos-economicos/Paginas/2023/130423-bono-verde-soberano-espanol.aspx}}

\noindent\hangindent=1.5em\hangafter=1 La Moncloa (2025, 30 de septiembre). \textit{El Tesoro recorta la emisión neta de 2025 en 5.000 millones} (nota de prensa). \href{https://www.lamoncloa.gob.es/serviciosdeprensa/notasprensa/economia-comercio-empresa/Paginas/2025/300925-emision-deuda.aspx}{\nolinkurl{https://www.lamoncloa.gob.es/serviciosdeprensa/notasprensa/economia-comercio-empresa/Paginas/2025/300925-emision-deuda.aspx}}

\noindent\hangindent=1.5em\hangafter=1 La Moncloa (2026a, 20 de enero). \textit{El Tesoro emite 15.000 millones de un nuevo bono a 10 años, con la mayor demanda de su historia} (nota de prensa). \href{https://www.lamoncloa.gob.es/serviciosdeprensa/notasprensa/economia-comercio-empresa/Paginas/2026/200126-bono-tesoro.aspx}{\nolinkurl{https://www.lamoncloa.gob.es/serviciosdeprensa/notasprensa/economia-comercio-empresa/Paginas/2026/200126-bono-tesoro.aspx}}

\noindent\hangindent=1.5em\hangafter=1 La Moncloa (2026b, 27 de mayo). \textit{Emisión sindicada de una nueva obligación del Estado a 10 años} (nota de prensa). \href{https://www.lamoncloa.gob.es/serviciosdeprensa/notasprensa/economia-comercio-empresa/paginas/2026/270526-emision-bonos-tesoro-publico.aspx}{\nolinkurl{https://www.lamoncloa.gob.es/serviciosdeprensa/notasprensa/economia-comercio-empresa/paginas/2026/270526-emision-bonos-tesoro-publico.aspx}}

\noindent\hangindent=1.5em\hangafter=1 La Moncloa (2026c, 31 de agosto). \textit{Plan de respuesta a la guerra en Oriente Próximo} (nota de prensa). \href{https://www.lamoncloa.gob.es/serviciosdeprensa/notasprensa/presidencia/paginas/2026/plan-respuesta-guerra-oriente.aspx}{\nolinkurl{https://www.lamoncloa.gob.es/serviciosdeprensa/notasprensa/presidencia/paginas/2026/plan-respuesta-guerra-oriente.aspx}}

\noindent\hangindent=1.5em\hangafter=1 Merca2 (2026, 9 de septiembre). \textit{El Tesoro coloca 4.000 millones en bonos verdes ante una demanda récord}. \href{https://www.merca2.es/2026/09/09/bonos-verdes-espana-tesoro-emision-2450986/}{\nolinkurl{https://www.merca2.es/2026/09/09/bonos-verdes-espana-tesoro-emision-2450986/}}

\noindent\hangindent=1.5em\hangafter=1 Ministerio de Economía y Hacienda (2008, 21 de febrero). \textit{Presentación del cierre presupuestario de 2007}. \href{https://www.hacienda.gob.es/GabineteMinistro/presentaciones/21-02-08\%20presentacion\%20cierre\%202007.pdf}{\nolinkurl{https://www.hacienda.gob.es/GabineteMinistro/presentaciones/21-02-08\%20presentacion\%20cierre\%202007.pdf}}

\noindent\hangindent=1.5em\hangafter=1 Ministerio de Hacienda (2026a, 31 de marzo). \textit{España cierra 2025 con un déficit público del 2,2\% del PIB} (nota de prensa). \href{https://www.hacienda.gob.es}{\nolinkurl{https://www.hacienda.gob.es}}

\noindent\hangindent=1.5em\hangafter=1 Ministerio de Hacienda (2026b, 11 de agosto). \textit{España recibe 6.234 millones del sexto desembolso del Plan de Recuperación} (nota de prensa). \href{https://www.hacienda.gob.es/es-es/prensa/noticias/paginas/2026/20260811-np-pago-sexto-desembolso-prtr.aspx}{\nolinkurl{https://www.hacienda.gob.es/es-es/prensa/noticias/paginas/2026/20260811-np-pago-sexto-desembolso-prtr.aspx}}

\noindent\hangindent=1.5em\hangafter=1 Ministerio de Inclusión, Seguridad Social y Migraciones (s. f.). \textit{El Fondo de Reserva de la Seguridad Social alcanzó 57.223,17 millones de euros en 2008}. \href{https://www.inclusion.gob.es/w/el-fondo-de-reserva-de-la-seguridad-social-alcanzo-57.223-17-millones-de-euros-en-2008}{\nolinkurl{https://www.inclusion.gob.es/w/el-fondo-de-reserva-de-la-seguridad-social-alcanzo-57.223-17-millones-de-euros-en-2008}}

\noindent\hangindent=1.5em\hangafter=1 Plan de Recuperación, Transformación y Resiliencia (2025). \textit{El Tesoro Público presenta su Estrategia de Financiación para 2026}. \href{https://planderecuperacion.gob.es/noticias/tesoro-publico-presenta-estrategia-financiacion-2026-prtr}{\nolinkurl{https://planderecuperacion.gob.es/noticias/tesoro-publico-presenta-estrategia-financiacion-2026-prtr}}

\noindent\hangindent=1.5em\hangafter=1 Público (2025). \textit{El Tesoro prevé emitir 55.000 millones de deuda neta en 2026}. \href{https://www.publico.es/economia/tesoro-preve-emitir-55-000-millones-2026-deuda-neta-2026-ano.html}{\nolinkurl{https://www.publico.es/economia/tesoro-preve-emitir-55-000-millones-2026-deuda-neta-2026-ano.html}}

\noindent\hangindent=1.5em\hangafter=1 Qué.es (2026, 15 de septiembre). \textit{El IPC de agosto de 2026 sube al 4,3\%, máximo en tres años}. \href{https://www.que.es/2026/09/15/ipc-agosto-2026-maximo-tres-anos/}{\nolinkurl{https://www.que.es/2026/09/15/ipc-agosto-2026-maximo-tres-anos/}}

\noindent\hangindent=1.5em\hangafter=1 Rappler (2023, febrero). \textit{European Union national fiscal policy responses to the energy crisis (Bruegel)}. \href{https://www.rappler.com/business/european-union-national-fiscal-policy-responses-energy-crisis-bruegel-february-2023/}{\nolinkurl{https://www.rappler.com/business/european-union-national-fiscal-policy-responses-energy-crisis-bruegel-february-2023/}}

\noindent\hangindent=1.5em\hangafter=1 Secretaría General de Presupuestos y Gastos (2026). \textit{Presupuestos Generales del Estado prorrogados}. \href{https://www.sepg.pap.hacienda.gob.es/sitios/sepg/es-ES/Presupuestos/PGE/PGE2025Prorroga/Paginas/PGE2025Prorroga.aspx}{\nolinkurl{https://www.sepg.pap.hacienda.gob.es/sitios/sepg/es-ES/Presupuestos/PGE/PGE2025Prorroga/Paginas/PGE2025Prorroga.aspx}}

\noindent\hangindent=1.5em\hangafter=1 Sgaravatti, G., Tagliapietra, S., Trasi, C. y Zachmann, G. (2023). \textit{National fiscal policy responses to the energy crisis} [conjunto de datos]. Bruegel. \href{https://www.bruegel.org/dataset/national-policies-shield-consumers-rising-energy-prices}{\nolinkurl{https://www.bruegel.org/dataset/national-policies-shield-consumers-rising-energy-prices}}

\noindent\hangindent=1.5em\hangafter=1 Tesoro Público (2025). \textit{Guía práctica de compra de deuda pública por inversores particulares 2025}. \href{https://www.tesoropublico.gob.es/sites/default/files/Gu\%C3\%ADa_Pr\%C3\%A1ctica_TP_2025_Interactivo.pdf}{\nolinkurl{https://www.tesoropublico.gob.es/sites/default/files/Gu\%C3\%ADa_Pr\%C3\%A1ctica_TP_2025_Interactivo.pdf}}

\noindent\hangindent=1.5em\hangafter=1 The Objective (2026, 15 de julio). \textit{La AIReF reclama al Gobierno un ajuste del 0,6\% del PIB para cumplir las normas fiscales}. \href{https://theobjective.com/economia/2026-07-15/airef-gobierno-ajuste-pib-normas-fiscales/}{\nolinkurl{https://theobjective.com/economia/2026-07-15/airef-gobierno-ajuste-pib-normas-fiscales/}}

\noindent\hangindent=1.5em\hangafter=1 Vozpópuli (2026). \textit{Bruselas ve alto riesgo en la deuda pública de España a medio plazo}. \href{https://www.vozpopuli.com/economia/macroeconomia/bruselas-alto-riesgo-deuda-publica-espana.html}{\nolinkurl{https://www.vozpopuli.com/economia/macroeconomia/bruselas-alto-riesgo-deuda-publica-espana.html}}


% ANEXOS
\clearpage
\end{document}

